\chapter{Problems Related to Volume VIII}
\label{ch:companion-viii}

\paragraph{Main-text correspondence.}
This Part accompanies \emph{Volume VIII: Algebraic Topology}, whose canonical chapter range is \texttt{VIII/01--VIII/35}. The problems are grouped by mathematical theme and provide worked examples across homotopy, coverings, homology, cohomology, and bundle theory.


\section{Homotopy}

\begin{problem}[CP-VIII-0001 --- The antipodal map on odd-dimensional spheres]
Let
\[
a:S^n\to S^n,\qquad a(x)=-x,
\]
be the antipodal map. Show that \(a\) is homotopic to the identity map when \(n\) is odd.
\end{problem}

\begin{solution}
We must show that if \(n\) is odd, then there exists a continuous map
\[
H:S^n\times [0,1]\to S^n
\]
such that
\[
H(x,0)=a(x)=-x
\qquad\text{and}\qquad
H(x,1)=x.
\]


\medskip
\noindent\textbf{Step 1: The case \(n=1\)}\par
\smallskip


First consider
\[
S^1=\{z\in \mathbb C:|z|=1\}.
\]
The antipodal map is
\[
a(z)=-z.
\]
On the circle, multiplication by \(-1\) is just rotation by angle \(\pi\). So we can connect \(-1\) to \(1\) by a continuous path in the unit circle.

Define
\[
H(z,t)=e^{i\pi(1-t)}z.
\]
Then \(H\) is continuous, and since \(|e^{i\pi(1-t)}|=1\), we have \(H(z,t)\in S^1\) for all \(z\in S^1\), \(t\in[0,1]\).

Moreover,
\[
H(z,0)=e^{i\pi}z=-z=a(z),
\]
and
\[
H(z,1)=e^0 z=z.
\]
Thus \(a\) is homotopic to the identity on \(S^1\).


\medskip
\noindent\textbf{Step 2: The general odd-dimensional case}\par
\smallskip


Now suppose \(n\) is odd. Then we may write
\[
n=2m+1
\]
for some integer \(m\ge 0\). Realize \(S^{2m+1}\) as the unit sphere in \(\mathbb C^{m+1}\):
\[
S^{2m+1}
=
\left\{
(z_0,\dots,z_m)\in \mathbb C^{m+1}
:
\sum_{j=0}^m |z_j|^2=1
\right\}.
\]

Define
\[
H\big((z_0,\dots,z_m),t\big)
=
\big(e^{i\pi(1-t)}z_0,\dots,e^{i\pi(1-t)}z_m\big).
\]

We check that this is well-defined. Since multiplication by a complex number of modulus \(1\) preserves norms,
\[
\sum_{j=0}^m \left|e^{i\pi(1-t)}z_j\right|^2
=
\sum_{j=0}^m |z_j|^2
=
1.
\]
Hence \(H((z_0,\dots,z_m),t)\in S^{2m+1}\) for all \(t\in[0,1]\).

The map \(H\) is continuous, and at the endpoints we obtain:
\[
H((z_0,\dots,z_m),0)
=
(-z_0,\dots,-z_m)
=
-(z_0,\dots,z_m)
=
a(z),
\]
while
\[
H((z_0,\dots,z_m),1)
=
(z_0,\dots,z_m).
\]

Therefore \(H\) is a homotopy from the antipodal map \(a\) to the identity map on \(S^{2m+1}\).

Thus, whenever \(n\) is odd, the antipodal map
\[
a(x)=-x
\]
is homotopic to the identity on \(S^n\).


\medskip
\noindent\textbf{Conclusion}\par
\smallskip


If \(n\) is odd, then \(S^n=S^{2m+1}\) can be viewed as the unit sphere in \(\mathbb C^{m+1}\), and the homotopy
\[
H(z,t)=e^{i\pi(1-t)}z
\]
provides a continuous deformation from the antipodal map to the identity. Hence
\[
a\simeq \mathrm{id}_{S^n}
\qquad\text{for odd }n.
\]


\medskip
\noindent\textbf{Remark}\par
\smallskip


For even \(n\), the antipodal map is not homotopic to the identity. Indeed,
\[
\deg(a)=\deg(-\mathrm{id}_{S^n})=(-1)^{n+1}.
\]
So if \(n\) is even, then \(\deg(a)=-1\), whereas
\[
\deg(\mathrm{id}_{S^n})=1.
\]
Since homotopic maps have the same degree, \(a\) cannot be homotopic to the identity when \(n\) is even.
\end{solution}

\begin{problem}[CP-VIII-0002 --- A map of \(S^1\) not homotopic to the identity]
Let
\[
f:S^1\to S^1
\]
be a map which is not homotopic to the identity map. Show that there exists an
\[
x\in S^1
\]
such that
\[
f(x)=x,
\]
and a
\[
y\in S^1
\]
such that
\[
f(y)=-y.
\]
\end{problem}

\begin{solution}
We will prove separately that:

\begin{enumerate}
	\item there exists \(y\in S^1\) such that \(f(y)=-y\),
	\item there exists \(x\in S^1\) such that \(f(x)=x\).
\end{enumerate}

The key observation is that if \(f\) avoids antipodal points, then \(f\) is homotopic to the identity, and if \(f\) avoids fixed points, then \(f\) is homotopic to the antipodal map. Since on \(S^1\) the antipodal map is homotopic to the identity, either situation contradicts the assumption that \(f\) is not homotopic to the identity.


\medskip
\noindent\textbf{Step 1: Existence of a point \(y\) such that \(f(y)=-y\)}\par
\smallskip


Assume, for contradiction, that
\[
f(z)\neq -z
\qquad\text{for all }z\in S^1.
\]

We identify \(S^1\) with the unit circle in \(\mathbb C\). Then for each \(z\in S^1\),
\[
\frac{f(z)}{z}\in S^1
\]
and the assumption \(f(z)\neq -z\) means precisely that
\[
\frac{f(z)}{z}\neq -1.
\]
Hence the map
\[
z\longmapsto \frac{f(z)}{z}
\]
takes values in
\[
S^1\setminus\{-1\},
\]
which is an open arc. Therefore there exists a continuous function
\[
\theta:S^1\to (-\pi,\pi)
\]
such that
\[
\frac{f(z)}{z}=e^{i\theta(z)}.
\]

Now define
\[
H:S^1\times[0,1]\to S^1,
\qquad
H(z,t)=e^{it\theta(z)}z.
\]

This map is continuous, and for all \(z\in S^1\),
\[
H(z,0)=z,
\]
while
\[
H(z,1)=e^{i\theta(z)}z=f(z).
\]
Thus \(H\) is a homotopy from \(\mathrm{id}_{S^1}\) to \(f\). Hence
\[
f\simeq \mathrm{id}_{S^1},
\]
contrary to the assumption.

Therefore our assumption was false. So there exists \(y\in S^1\) such that
\[
f(y)=-y.
\]


\medskip
\noindent\textbf{Step 2: Existence of a point \(x\) such that \(f(x)=x\)}\par
\smallskip


Assume, for contradiction, that
\[
f(z)\neq z
\qquad\text{for all }z\in S^1.
\]

Let
\[
a:S^1\to S^1,\qquad a(z)=-z
\]
be the antipodal map. We claim that \(f\) is homotopic to \(a\).

Indeed, for each \(z\in S^1\),
\[
\frac{f(z)}{-z}\in S^1.
\]
The assumption \(f(z)\neq z\) implies
\[
\frac{f(z)}{-z}\neq -1,
\]
because
\[
\frac{f(z)}{-z}=-1
\quad\Longleftrightarrow\quad
f(z)=z.
\]
Thus the map
\[
z\longmapsto \frac{f(z)}{-z}
\]
also takes values in \(S^1\setminus\{-1\}\). Therefore there exists a continuous function
\[
\phi:S^1\to (-\pi,\pi)
\]
such that
\[
\frac{f(z)}{-z}=e^{i\phi(z)}.
\]

Define
\[
K:S^1\times[0,1]\to S^1,
\qquad
K(z,t)=e^{it\phi(z)}(-z).
\]

Then
\[
K(z,0)=-z=a(z),
\]
and
\[
K(z,1)=e^{i\phi(z)}(-z)=f(z).
\]
So \(K\) is a homotopy from \(a\) to \(f\). Therefore
\[
f\simeq a.
\]

But on \(S^1\), the antipodal map is homotopic to the identity. For example,
\[
L(z,t)=e^{i\pi t}z
\]
defines a homotopy from \(z\) to \(-z\). Hence
\[
a\simeq \mathrm{id}_{S^1}.
\]
It follows that
\[
f\simeq a\simeq \mathrm{id}_{S^1},
\]
again contradicting the assumption.

Therefore our assumption was false. So there exists \(x\in S^1\) such that
\[
f(x)=x.
\]


\medskip
\noindent\textbf{Conclusion}\par
\smallskip


We have proved that if
\[
f:S^1\to S^1
\]
is not homotopic to the identity, then:

\begin{itemize}
	\item there exists \(x\in S^1\) such that
	\[
	f(x)=x,
	\]
	\item there exists \(y\in S^1\) such that
	\[
	f(y)=-y.
	\]
\end{itemize}

This is exactly the required statement.
\end{solution}

\begin{problem}[CP-VIII-0003 --- From one-sided homotopy inverses to a homotopy equivalence]
Suppose
\[
f:X\to Y
\]
is a map for which there exist maps
\[
g,h:Y\to X
\]
such that
\[
g\circ f \simeq \operatorname{Id}_X
\qquad\text{and}\qquad
f\circ h \simeq \operatorname{Id}_Y.
\]
Show that \(f\), \(g\), and \(h\) are homotopy equivalences.
\end{problem}

\begin{solution}

\medskip
\noindent\textbf{Idea}\par
\smallskip


We are given a left homotopy inverse \(g\) of \(f\) and a right homotopy inverse \(h\) of \(f\).  
The key point is to show that
\[
g\simeq h.
\]
Then \(g\) becomes both a left and a right homotopy inverse of \(f\), so \(f\) is a homotopy equivalence.


\medskip
\noindent\textbf{Proof}\par
\smallskip


We have
\[
h \simeq \operatorname{Id}_X\circ h.
\]
Since
\[
g\circ f \simeq \operatorname{Id}_X,
\]
it follows that
\[
\operatorname{Id}_X\circ h \simeq (g\circ f)\circ h.
\]
By associativity,
\[
(g\circ f)\circ h = g\circ(f\circ h).
\]
Since
\[
f\circ h \simeq \operatorname{Id}_Y,
\]
we obtain
\[
g\circ(f\circ h)\simeq g\circ\operatorname{Id}_Y = g.
\]
Therefore
\[
h\simeq g.
\]

Now we already know
\[
g\circ f \simeq \operatorname{Id}_X.
\]
Also,
\[
f\circ g \simeq f\circ h \simeq \operatorname{Id}_Y.
\]
Hence \(g\) is a homotopy inverse of \(f\), so \(f\) is a homotopy equivalence.

Then \(g\) is also a homotopy equivalence, with homotopy inverse \(f\).

Finally, since \(h\simeq g\), we get
\[
h\circ f \simeq g\circ f \simeq \operatorname{Id}_X
\]
and by assumption
\[
f\circ h \simeq \operatorname{Id}_Y.
\]
So \(h\) is also a homotopy inverse of \(f\), and thus \(h\) is a homotopy equivalence.

Therefore \(f\), \(g\), and \(h\) are all homotopy equivalences.


\medskip
\noindent\textbf{Interpretation}\par
\smallskip


This shows that if a map has a left inverse up to homotopy and a right inverse up to homotopy, then these inverses agree up to homotopy, and the map is a genuine homotopy equivalence.
\end{solution}

\begin{problem}[CP-VIII-0004 --- A retract of a contractible space is contractible]
Show that a retract of a contractible space is contractible.
\end{problem}

\begin{solution}

\medskip
\noindent\textbf{Definitions}\par
\smallskip


A space \(A\) is called a retract of a space \(X\) if there exist maps
\[
i:A\to X,
\qquad
r:X\to A
\]
such that
\[
r\circ i = \operatorname{Id}_A.
\]

A space \(X\) is contractible if
\[
\operatorname{Id}_X \simeq c
\]
for some constant map \(c:X\to X\).


\medskip
\noindent\textbf{Idea}\par
\smallskip


If \(X\) contracts to a point, then a retract \(A\) of \(X\) also contracts: perform the contraction in \(X\), then apply the retraction back to \(A\).


\medskip
\noindent\textbf{Proof}\par
\smallskip


Let \(A\) be a retract of \(X\). Then there exist maps
\[
i:A\to X,
\qquad
r:X\to A
\]
such that
\[
r\circ i = \operatorname{Id}_A.
\]

Assume that \(X\) is contractible. Then there exists a point \(x_0\in X\) such that
\[
\operatorname{Id}_X \simeq c_{x_0},
\]
where \(c_{x_0}:X\to X\) is the constant map with value \(x_0\).

Compose this homotopy on the left with \(r\) and on the right with \(i\). Since homotopy is preserved under composition, we get
\[
r\circ \operatorname{Id}_X \circ i \simeq r\circ c_{x_0}\circ i.
\]
But
\[
r\circ \operatorname{Id}_X \circ i = r\circ i = \operatorname{Id}_A.
\]
Hence
\[
\operatorname{Id}_A \simeq r\circ c_{x_0}\circ i.
\]
Since \(c_{x_0}\) is constant, the composite \(r\circ c_{x_0}\circ i\) is also constant, with value \(r(x_0)\in A\).

Therefore \(\operatorname{Id}_A\) is homotopic to a constant map, so \(A\) is contractible.


\medskip
\noindent\textbf{Remark}\par
\smallskip


This result is often used as an obstruction: if a space \(A\) is not contractible, then it cannot be a retract of any contractible space.
\end{solution}

\begin{problem}[CP-VIII-0005 --- A space retracting onto both an annulus and a M\"obius band]
Construct a space which contains both the annulus \(S^1\times I\) and the Möbius band as deformation retracts.

\par

\textbf{Idea.}
Both the annulus and the Möbius band deformation retract onto a circle, namely their core circle. So we glue an annulus and a Möbius band along their core circles. Then collapsing one of them onto the common core circle while keeping the other fixed gives a deformation retraction onto the other.

\par

\textbf{Construction of the space.}
Let \(A\) be an annulus and \(M\) a Möbius band. Let
\[
C_A\subset A
\qquad\text{and}\qquad
C_M\subset M
\]
be their core circles.

Choose a homeomorphism
\[
\varphi:C_A\to C_M.
\]
Now form the quotient space
\[
Y=(A\sqcup M)/\sim,
\]
where \(c\sim \varphi(c)\) for all \(c\in C_A\).

Equivalently, we identify the core circle of the annulus with the core circle of the Möbius band. Denote the common image of these circles by
\[
C\subset Y.
\]
Then \(Y\) contains a copy of the annulus \(A\) and a copy of the Möbius band \(M\), intersecting exactly in the common circle \(C\).

\par

\textbf{Why \(Y\) deformation retracts onto the annulus.}
The Möbius band deformation retracts onto its core circle. Let
\[
r_M:M\times I\to M
\]
be such a deformation retraction. Thus
\[
r_M(m,0)=m,
\qquad
r_M(m,1)\in C,
\qquad
r_M(c,t)=c
\]
for all \(m\in M\), \(c\in C\), and \(t\in I\).

Define
\[
H:Y\times I\to Y
\]
by
\[
H(y,t)=
\begin{cases}
	y,& y\in A,\\[4pt]
	r_M(y,t),& y\in M.
\end{cases}
\]
This is well-defined and continuous because on the overlap \(A\cap M=C\), the two formulas agree:
\[
r_M(c,t)=c.
\]

At time \(t=0\), we have
\[
H(y,0)=y
\]
for all \(y\in Y\), so \(H(-,0)=\mathrm{id}_Y\).

At time \(t=1\), every point of \(M\) has been moved into the common circle \(C\subset A\), while every point of \(A\) has remained fixed. Hence
\[
H(Y,1)=A.
\]
Thus \(H\) is a deformation retraction of \(Y\) onto the annulus \(A\).

\par

\textbf{Why \(Y\) deformation retracts onto the Möbius band.}
Exactly similarly, the annulus deformation retracts onto its core circle. Let
\[
r_A:A\times I\to A
\]
be a deformation retraction onto \(C\), fixing \(C\) pointwise.

Define
\[
K:Y\times I\to Y
\]
by
\[
K(y,t)=
\begin{cases}
	r_A(y,t),& y\in A,\\[4pt]
	y,& y\in M.
\end{cases}
\]
Again this is continuous because on the common circle \(C\),
\[
r_A(c,t)=c.
\]

At time \(t=0\), \(K\) is the identity on \(Y\). At time \(t=1\), all of \(A\) has collapsed into \(C\subset M\), while \(M\) has remained fixed. Therefore
\[
K(Y,1)=M.
\]
So \(K\) is a deformation retraction of \(Y\) onto the Möbius band \(M\).

\par

\textbf{Conclusion.}
The space obtained by gluing an annulus and a Möbius band along their core circles contains both of them as deformation retracts.
\end{problem}

\begin{solution}

\medskip
\noindent\textbf{Explicit coordinate models for the two bands}\par
\smallskip


It is often useful to write the annulus and the Möbius band explicitly as quotients.

The annulus may be written as
\[
A=[0,1]\times [-1,1]\big/\bigl((0,s)\sim (1,s)\bigr),
\]
and its core circle is
\[
C_A=[0,1]\times \{0\}\big/\bigl((0,0)\sim (1,0)\bigr).
\]
A deformation retraction of \(A\) onto \(C_A\) is given by
\[
r_A([u,s],t)=[u,(1-t)s].
\]

The Möbius band may be written as
\[
M=[0,1]\times [-1,1]\big/\bigl((0,s)\sim (1,-s)\bigr),
\]
and its core circle is
\[
C_M=[0,1]\times \{0\}\big/\bigl((0,0)\sim (1,0)\bigr).
\]
A deformation retraction of \(M\) onto \(C_M\) is given by
\[
r_M([u,s],t)=[u,(1-t)s].
\]

Since both core circles are naturally homeomorphic to \(S^1\), we may identify them and form the space \(Y\) above.
\end{solution}

\section{CW Complexes}

\begin{problem}[CP-VIII-0006 --- The dunce cap is contractible]
The \emph{dunce cap} is obtained from a solid triangle by identifying its three boundary edges according to the standard directed edge identifications. Show that the resulting space is contractible.
\end{problem}

\begin{solution}
Let \(D\) denote the dunce cap. By the edge identifications, the quotient of the boundary of the triangle is a circle:
\[
\partial \Delta^2/\!\sim \;\cong S^1.
\]
Hence \(D\) can be written as
\[
D \cong S^1\cup_g D^2
\]
for some attaching map
\[
g:S^1\to S^1.
\]

Reading the boundary according to the arrows shows that the total winding number of \(g\) around the unique circle is
\[
\deg(g)=\pm1.
\]
Therefore \(g\) is homotopic to a homeomorphism
\[
h:S^1\to S^1.
\]

By Problem 9,
\[
S^1\cup_g D^2 \simeq S^1\cup_h D^2.
\]
Since \(h\) is a homeomorphism, attaching \(D^2\) along \(h\) yields a disk:
\[
S^1\cup_h D^2 \cong D^2.
\]
Thus
\[
D \simeq D^2.
\]
Because \(D^2\) is contractible, the dunce cap is contractible.


\medskip
\noindent\textbf{Conclusion.}\par
\smallskip

The dunce cap is contractible.


\medskip
\noindent\textbf{Remark}\par
\smallskip


The dunce cap is a famous example because it is contractible but not collapsible. Thus it is topologically trivial from the homotopy point of view, but combinatorially quite subtle.
\end{solution}

\begin{problem}[CP-VIII-0007 --- Homotopic attaching maps give the same homotopy type]

\medskip
\noindent\textbf{Statement.}\par
\smallskip

If \(f_0,f_1:S^{n-1}\to X\) and \(f_0\simeq f_1\), show that
\[
X\cup_{f_0} D^n \simeq X\cup_{f_1} D^n .
\]
\end{problem}

\begin{solution}
This is a standard fact about attaching cells: attaching an \(n\)-cell along two homotopic maps produces homotopy equivalent spaces.

Write
\[
Y_i:=X\cup_{f_i}D^n \qquad (i=0,1).
\]
Each \(Y_i\) is the pushout of the diagram


\[
\begin{array}{ccc}
S^{n-1} & \xrightarrow{\ f_i\ } & X\\
\big\downarrow{\scriptstyle \mathrm{incl.}} && \big\downarrow\\
D^n & \longrightarrow & Y_i
\end{array}
\]



Now let
\[
H:S^{n-1}\times I \to X
\]
be a homotopy from \(f_0\) to \(f_1\), so
\[
H(-,0)=f_0,\qquad H(-,1)=f_1.
\]

The key point is that the inclusion
\[
S^{n-1}\hookrightarrow D^n
\]
is a \emph{cofibration}. For cofibrations, ordinary pushouts agree with homotopy pushouts. Therefore the homotopy type of the adjunction space depends only on the homotopy class of the attaching map.

Equivalently, the two pushout diagrams


\[
\begin{array}{ccc@{\qquad\qquad}ccc}
S^{n-1} & \xrightarrow{\ f_0\ } & X
&
S^{n-1} & \xrightarrow{\ f_1\ } & X\\
\big\downarrow{\scriptstyle \mathrm{incl.}} && \big\downarrow
&
\big\downarrow{\scriptstyle \mathrm{incl.}} && \big\downarrow\\
D^n & \longrightarrow & X\cup_{f_0}D^n
&
D^n & \longrightarrow & X\cup_{f_1}D^n
\end{array}
\]


have the same homotopy type because \(f_0\) and \(f_1\) are homotopic.

Hence
\[
X\cup_{f_0}D^n \simeq X\cup_{f_1}D^n.
\]


\medskip
\noindent\textbf{Remark.}\par
\smallskip

This is one of the basic principles behind CW-complexes: the homotopy type of a CW-complex depends only on the homotopy classes of its attaching maps, not on the precise representatives.

\par
\end{solution}

\begin{problem}[CP-VIII-0008 --- Cellular homology of the lens space \texorpdfstring{$L_p^{2k-1}$}{L\_p\string^{2k-1}}]
Write
\[
S^{2k-1}=\left\{(z_1,\dots,z_k)\in \mathbb C^k \;\middle|\; \sum_{i=1}^k |z_i|^2=1\right\}.
\]
The group \(\mathbb Z/p\) acts on \(S^{2k-1}\) by
\[
a\cdot z=\lambda^a z,\qquad \lambda=e^{2\pi i/p}.
\]
The lens space \(L_p^{2k-1}\) is the quotient
\[
L_p^{2k-1}=S^{2k-1}/(\mathbb Z/p).
\]
Show that \(L_p^{2k-1}\) has a cell decomposition with one cell in each dimension, find \(C_*^{\mathrm{cell}}(L_p^{2k-1})\), and compute its homology.
\end{problem}

\begin{solution}

\medskip
\noindent\textbf{Remark.}\par
\smallskip

The standard CW decomposition of the lens space \(L_p^{2k-1}\) has one cell in each dimension
\[
0,1,2,\dots,2k-1.
\]
This is the decomposition used below.

A classical CW structure on the lens space \(L_p^{2k-1}\) has exactly one cell \(e^q\) in every dimension \(q=0,1,\dots,2k-1\). Therefore
\[
C_q^{\mathrm{cell}}(L_p^{2k-1})\cong \mathbb Z
\qquad \text{for } 0\le q\le 2k-1,
\]
and \(C_q^{\mathrm{cell}}(L_p^{2k-1})=0\) otherwise.

Hence the cellular chain complex has the form
\[
0 \to \mathbb Z \xrightarrow{d_{2k-1}} \mathbb Z \xrightarrow{d_{2k-2}} \cdots
\xrightarrow{d_2} \mathbb Z \xrightarrow{d_1} \mathbb Z \to 0.
\]

For the standard cell structure on lens spaces, the differentials alternate between \(0\) and multiplication by \(p\):
\[
d_q=
\begin{cases}
	0, & q \text{ odd},\\[4pt]
	\times p, & q \text{ even}.
\end{cases}
\]
Thus the cellular chain complex is
\[
0 \to \mathbb Z
\xrightarrow{0}
\mathbb Z
\xrightarrow{\times p}
\mathbb Z
\xrightarrow{0}
\mathbb Z
\xrightarrow{\times p}
\cdots
\xrightarrow{0}
\mathbb Z
\xrightarrow{\times p}
\mathbb Z
\xrightarrow{0}
\mathbb Z
\to 0.
\]

We now compute homology.

\par
\noindent
For degree \(0\),
\[
H_0(L_p^{2k-1})\cong \mathbb Z
\]
because the space is path connected.

\par
\noindent
Let \(q=2r-1\) be an odd degree with \(1\le 2r-1\le 2k-3\). Since \(d_{2r-1}=0\),
\[
\ker(d_{2r-1})=\mathbb Z.
\]
Since \(d_{2r}=\times p\),
\[
\operatorname{im}(d_{2r})=p\mathbb Z.
\]
Therefore
\[
H_{2r-1}(L_p^{2k-1})
=
\ker(d_{2r-1})/\operatorname{im}(d_{2r})
\cong
\mathbb Z/p\mathbb Z.
\]

\par
\noindent
If \(q=2r\) is even with \(2\le 2r\le 2k-2\), then \(d_{2r}=\times p\), so
\[
\ker(d_{2r})=0.
\]
Hence
\[
H_{2r}(L_p^{2k-1})=0.
\]

\par
\noindent
Finally, in top degree \(2k-1\), since there is no cell in dimension \(2k\), we have
\[
H_{2k-1}(L_p^{2k-1})=\ker(d_{2k-1})=\mathbb Z
\]
because \(d_{2k-1}=0\).

Therefore the homology groups are
\[
H_q(L_p^{2k-1})=
\begin{cases}
	\mathbb Z, & q=0,\\[4pt]
	\mathbb Z/p\mathbb Z, & q \text{ odd and }1\le q\le 2k-3,\\[4pt]
	0, & q \text{ even and }0<q<2k-1,\\[4pt]
	\mathbb Z, & q=2k-1.
\end{cases}
\]

So the cellular chain complex is completely determined by
\[
C_q^{\mathrm{cell}}(L_p^{2k-1})\cong \mathbb Z \quad (0\le q\le 2k-1),
\qquad
d_q=
\begin{cases}
	0, & q\text{ odd},\\
	\times p, & q\text{ even}.
\end{cases}
\]

\par
\end{solution}

\begin{problem}[CP-VIII-0009 --- Realizing finitely generated abelian groups as homology groups of a finite CW complex]
Suppose \(G_i\) (\(1\le i\le n\)) are finitely generated abelian groups. Show that there is a finite cell complex \(X\) with
\[
H_i(X)=G_i \qquad \text{for } 1\le i\le n,
\]
and
\[
H_i(X)=0 \qquad \text{for } i>n.
\]
\end{problem}

\begin{solution}
We use the structure theorem for finitely generated abelian groups. For each \(i\),
\[
G_i \cong \mathbb Z^{r_i}\oplus \bigoplus_{j=1}^{s_i}\mathbb Z/m_{ij},
\]
for some integers \(r_i\ge 0\), \(s_i\ge 0\), and \(m_{ij}\ge 2\).

Thus each \(G_i\) is a direct sum of a free part \(\mathbb Z^{r_i}\) and finitely many cyclic torsion groups.

\par
\noindent
We first realize the free part. The sphere \(S^i\) has reduced homology
\[
\widetilde H_q(S^i)=
\begin{cases}
	\mathbb Z, & q=i,\\
	0, & q\ne i.
\end{cases}
\]
Hence the wedge
\[
\bigvee^{r_i} S^i
\]
has homology \(\mathbb Z^{r_i}\) in degree \(i\) and no other reduced homology.

\par
\noindent
Next, we realize a cyclic torsion summand \(\mathbb Z/m\) by a Moore space
\[
M(\mathbb Z/m,i)=S^i\cup_f e^{i+1},
\]
where the attaching map
\[
f:S^i\to S^i
\]
has degree \(m\).

The cellular chain complex of this space is
\[
0\to \mathbb Z \xrightarrow{\times m} \mathbb Z \to 0,
\]
with the left \(\mathbb Z\) in degree \(i+1\) and the right one in degree \(i\). Therefore
\[
H_i(M(\mathbb Z/m,i))\cong \mathbb Z/m,
\]
and all homology in degrees \(>i\) vanishes.

So each cyclic torsion summand \(\mathbb Z/m_{ij}\) can be realized by one Moore space \(M(\mathbb Z/m_{ij},i)\).

\par
\noindent
For each \(i\), define
\[
X_i
=
\left(\bigvee^{r_i} S^i\right)
\vee
\left(\bigvee_{j=1}^{s_i} M(\mathbb Z/m_{ij},i)\right).
\]
Then set
\[
X=\bigvee_{i=1}^n X_i.
\]

This is a finite wedge of finite CW complexes, hence \(X\) is a finite CW complex.

Now reduced homology of a wedge sum is the direct sum of the reduced homologies:
\[
\widetilde H_q\!\left(\bigvee_\alpha Y_\alpha\right)
\cong
\bigoplus_\alpha \widetilde H_q(Y_\alpha).
\]
Therefore in degree \(i\) we obtain
\[
H_i(X)
\cong
\mathbb Z^{r_i}\oplus \bigoplus_{j=1}^{s_i}\mathbb Z/m_{ij}
\cong
G_i.
\]

Moreover, all the pieces used to build \(X\) for indices \(1,\dots,n\) have no homology above degree \(n\). Hence
\[
H_i(X)=0 \qquad \text{for } i>n.
\]

This proves the claim.


\medskip
\noindent\textbf{Conclusion}\par
\smallskip

There exists a finite CW complex
\[
X=\bigvee_{i=1}^n
\left[
\left(\bigvee^{r_i} S^i\right)
\vee
\left(\bigvee_{j=1}^{s_i} M(\mathbb Z/m_{ij},i)\right)
\right]
\]
such that
\[
H_i(X)\cong G_i \qquad (1\le i\le n),
\]
and
\[
H_i(X)=0 \qquad (i>n).
\]
Thus every finite list of finitely generated abelian groups can be realized as the homology groups of a finite cell complex.
\end{solution}

\section{Fundamental Groups and Coverings}

\begin{problem}[CP-VIII-0010 --- Strong Deformation Retractions and a Space Containing Both an Annulus and a Möbius Band as Deformation Retracts]
Show that if a space \(X\) strongly deformation retracts to a point \(x_0 \in X\), then for every open neighbourhood \(x_0 \in U\) there exists a smaller open neighbourhood \(x_0 \in V \subset U\) such that the inclusion
\[
(V,x_0)\hookrightarrow (U,x_0)
\]
is based homotopic to the constant map.

\par

\textbf{Intuition.}
A strong deformation retraction gives a homotopy
\[
H:X\times I\to X
\]
from the identity to the constant map at \(x_0\), and throughout the homotopy the point \(x_0\) stays fixed:
\[
H(x_0,t)=x_0.
\]
Since \(x_0\) stays inside \(U\) for all \(t\), continuity implies that points sufficiently close to \(x_0\) also stay inside \(U\) for all \(t\). That smaller neighbourhood is the required \(V\).

\par
\end{problem}

\begin{solution}
Let
\[
H:X\times I\to X
\]
be a strong deformation retraction of \(X\) onto the point \(x_0\). Thus
\[
H(x,0)=x,
\qquad
H(x,1)=x_0,
\qquad
H(x_0,t)=x_0
\]
for all \(x\in X\) and all \(t\in I\).

Now let \(U\) be any open neighbourhood of \(x_0\).

For each \(t\in I\), we have
\[
H(x_0,t)=x_0\in U.
\]
Since \(H\) is continuous at the point \((x_0,t)\), there exist an open neighbourhood \(W_t\) of \(x_0\) in \(X\) and an open interval \(J_t\subset I\) containing \(t\) such that
\[
H(W_t\times J_t)\subset U.
\]

The family \(\{J_t\}_{t\in I}\) is an open cover of the compact interval \(I=[0,1]\). Hence there exist finitely many points
\[
t_1,\dots,t_n\in I
\]
such that
\[
I=J_{t_1}\cup\cdots\cup J_{t_n}.
\]

Define
\[
V=\Bigl(\bigcap_{i=1}^n W_{t_i}\Bigr)\cap U.
\]
Then \(V\) is an open neighbourhood of \(x_0\), and clearly
\[
V\subset U.
\]

We now claim that
\[
H(V\times I)\subset U.
\]
Indeed, let \(v\in V\) and \(s\in I\). Since the intervals \(J_{t_1},\dots,J_{t_n}\) cover \(I\), there exists some index \(i\) such that
\[
s\in J_{t_i}.
\]
Also,
\[
v\in V\subset W_{t_i}.
\]
Therefore
\[
H(v,s)\in H(W_{t_i}\times J_{t_i})\subset U.
\]
This proves the claim.

Hence the restriction
\[
K=H|_{V\times I}:V\times I\to U
\]
is well-defined. Moreover, for all \(v\in V\) and \(t\in I\),
\[
K(v,0)=H(v,0)=v,
\]
\[
K(v,1)=H(v,1)=x_0,
\]
and
\[
K(x_0,t)=H(x_0,t)=x_0.
\]
Thus \(K\) is a based homotopy in \((U,x_0)\) from the inclusion
\[
(V,x_0)\hookrightarrow (U,x_0)
\]
to the constant map at \(x_0\).

Therefore the inclusion
\[
(V,x_0)\hookrightarrow (U,x_0)
\]
is based homotopic to the constant map.
\end{solution}

\begin{problem}[CP-VIII-0011 --- The M\"obius band does not retract onto its boundary]
Show that the Möbius band does not retract onto its boundary.
\end{problem}

\begin{solution}
Let
\[
M=\bigl([0,1]\times[-1,1]\bigr)/\sim,
\qquad
(0,t)\sim(1,-t).
\]
This is the Möbius band. Its boundary is homeomorphic to a circle:
\[
\partial M \cong S^1.
\]

Also, the Möbius band deformation retracts onto its central circle
\[
C=\{[(s,0)] : s\in[0,1]\}\cong S^1,
\]
so
\[
\pi_1(M)\cong \pi_1(C)\cong \mathbb Z.
\]

Now consider the inclusion
\[
i:\partial M\hookrightarrow M.
\]
Geometrically, when one goes once around the boundary circle of the Möbius band, one winds twice around the central circle. Therefore, under the identifications
\[
\pi_1(\partial M)\cong \mathbb Z,
\qquad
\pi_1(M)\cong \mathbb Z,
\]
the induced homomorphism
\[
i_*:\pi_1(\partial M)\to \pi_1(M)
\]
is multiplication by \(2\):
\[
i_*(n)=2n.
\]

Suppose, for contradiction, that there exists a retraction
\[
r:M\to \partial M
\]
such that
\[
r\circ i=\operatorname{id}_{\partial M}.
\]
Applying the fundamental group functor, we get
\[
r_*\circ i_*=\operatorname{id}_{\pi_1(\partial M)}.
\]
Thus \(i_*\) would have to admit a left inverse.

But \(i_*\) is the map
\[
\mathbb Z \to \mathbb Z,
\qquad
n\mapsto 2n.
\]
There is no homomorphism
\[
\varphi:\mathbb Z\to\mathbb Z
\]
such that
\[
\varphi(2n)=n
\qquad\text{for all }n\in\mathbb Z,
\]
because every homomorphism \(\mathbb Z\to\mathbb Z\) has the form \(m\mapsto km\) for some \(k\in\mathbb Z\), and then
\[
k(2n)=n
\quad\Longrightarrow\quad
2k=1,
\]
which is impossible in \(\mathbb Z\).

This contradiction shows that no such retraction can exist.

Hence the Möbius band does not retract onto its boundary.
\end{solution}

\begin{problem}[CP-VIII-0012 --- Fundamental group of a product]
For based spaces \((X,x_0)\) and \((Y,y_0)\), show that
\[
\pi_1(X\times Y,(x_0,y_0))
\cong
\pi_1(X,x_0)\times \pi_1(Y,y_0).
\]
\end{problem}

\begin{solution}
Let
\[
p_X:X\times Y\to X,
\qquad
p_Y:X\times Y\to Y
\]
be the projection maps.

We define
\[
\Phi:\pi_1(X\times Y,(x_0,y_0))
\longrightarrow
\pi_1(X,x_0)\times \pi_1(Y,y_0)
\]
by
\[
\Phi([\gamma])
=
\bigl([p_X\circ \gamma],[p_Y\circ \gamma]\bigr).
\]
This is well defined, because homotopic loops in \(X\times Y\) project to homotopic loops in \(X\) and \(Y\).

Conversely, define
\[
\Psi:\pi_1(X,x_0)\times \pi_1(Y,y_0)
\longrightarrow
\pi_1(X\times Y,(x_0,y_0))
\]
by
\[
\Psi([\alpha],[\beta])=[(\alpha,\beta)],
\]
where
\[
(\alpha,\beta)(t)=(\alpha(t),\beta(t)).
\]
Again this is well defined: if \(\alpha\simeq \alpha'\) rel endpoints and \(\beta\simeq \beta'\) rel endpoints, then \((\alpha,\beta)\simeq(\alpha',\beta')\) rel endpoints.

Now we check that \(\Phi\) and \(\Psi\) are inverse maps.

First,
\[
\Phi(\Psi([\alpha],[\beta]))
=
\Phi([(\alpha,\beta)])
=
\bigl([p_X\circ(\alpha,\beta)],[p_Y\circ(\alpha,\beta)]\bigr)
=
([\alpha],[\beta]).
\]
So
\[
\Phi\circ\Psi=\operatorname{id}.
\]

Second, if \(\gamma:I\to X\times Y\) is a loop, write
\[
\gamma(t)=(\gamma_X(t),\gamma_Y(t)).
\]
Then
\[
\Psi(\Phi([\gamma]))
=
\Psi\bigl([\gamma_X],[\gamma_Y]\bigr)
=
[(\gamma_X,\gamma_Y)]
=
[\gamma].
\]
So
\[
\Psi\circ\Phi=\operatorname{id}.
\]

Thus \(\Phi\) is a bijection with inverse \(\Psi\).

It remains to check that \(\Phi\) is a group homomorphism. If \(\gamma,\delta\) are loops in \(X\times Y\) based at \((x_0,y_0)\), then
\[
\Phi([\gamma]*[\delta])
=
\Phi([\gamma * \delta])
=
\bigl([p_X\circ(\gamma * \delta)],[p_Y\circ(\gamma * \delta)]\bigr).
\]
Since projection commutes with concatenation of paths,
\[
p_X\circ(\gamma * \delta)=(p_X\circ\gamma)*(p_X\circ\delta),
\]
and similarly for \(p_Y\). Hence
\[
\Phi([\gamma]*[\delta])
=
\bigl([p_X\circ\gamma]*[p_X\circ\delta],\,
[p_Y\circ\gamma]*[p_Y\circ\delta]\bigr)
=
\Phi([\gamma])\cdot\Phi([\delta]).
\]
So \(\Phi\) is an isomorphism of groups.

Therefore,
\[
\pi_1(X\times Y,(x_0,y_0))
\cong
\pi_1(X,x_0)\times \pi_1(Y,y_0).
\]
\end{solution}

\begin{problem}[CP-VIII-0013 --- Universal covering of the Klein bottle]
Construct a covering map \(\pi:\mathbb R^2\to K\) of the Klein bottle, and hence show that \(\pi_1(K,k_0)\) is isomorphic to the group \(G\) with elements \((m,n)\in\mathbb Z^2\) and group operation
\[
(m,n)*(p,q)=\bigl(m+(-1)^n p,\; n+q\bigr).
\]

Show that \(K\) has a covering space homeomorphic to the torus \(S^1\times S^1\), but that the torus does not have a covering space homeomorphic to \(K\).
\end{problem}

\begin{solution}
We begin with the standard square model of the Klein bottle. Take the unit square \([0,1]\times[0,1]\), identify the vertical sides in the same direction, and identify the horizontal sides in opposite directions. The quotient is the Klein bottle \(K\).

We now realize this quotient as a quotient of \(\mathbb R^2\) by a discrete group of homeomorphisms.

Define two homeomorphisms of \(\mathbb R^2\):
\[
a(x,y)=(x+1,y),
\qquad
b(x,y)=(-x,y+1).
\]
The map \(a\) is a horizontal translation, while \(b\) shifts upward by \(1\) and reflects in the \(y\)-axis.

Let \(\Gamma\) be the group generated by \(a\) and \(b\). Every element of \(\Gamma\) has the form
\[
g_{m,n}(x,y)=\bigl((-1)^n x+m,\; y+n\bigr),
\qquad (m,n)\in\mathbb Z^2.
\]
Indeed, \(n\) records how many times we apply the vertical gluing, and each odd vertical step reverses the horizontal direction.

The action of \(\Gamma\) on \(\mathbb R^2\) is free and properly discontinuous, so the quotient
\[
\mathbb R^2/\Gamma
\]
is a surface, and the quotient map
\[
\pi:\mathbb R^2\to \mathbb R^2/\Gamma
\]
is a covering map.

Now the identifications produced by \(\Gamma\) are exactly
\[
(x,y)\sim(x+1,y),
\qquad
(x,y)\sim(-x,y+1),
\]
which are precisely the standard edge identifications for the Klein bottle. Hence
\[
\mathbb R^2/\Gamma \cong K.
\]
Therefore we obtain a covering map
\[
\pi:\mathbb R^2\to K.
\]

Since \(\mathbb R^2\) is simply connected, this is the universal covering of the Klein bottle. It follows that the deck transformation group of this covering is isomorphic to the fundamental group:
\[
\pi_1(K,k_0)\cong \Gamma.
\]

We now compute the group law in terms of \(\mathbb Z^2\). Under the correspondence
\[
(m,n)\longleftrightarrow g_{m,n},
\qquad
g_{m,n}(x,y)=\bigl((-1)^n x+m,\; y+n\bigr),
\]
we have
\[
g_{m,n}\circ g_{p,q}(x,y)
=
g_{m,n}\bigl((-1)^q x+p,\; y+q\bigr).
\]
Applying \(g_{m,n}\), we get
\[
g_{m,n}\circ g_{p,q}(x,y)
=
\bigl((-1)^n((-1)^q x+p)+m,\; y+q+n\bigr).
\]
This simplifies to
\[
g_{m,n}\circ g_{p,q}(x,y)
=
\bigl((-1)^{n+q}x + m + (-1)^n p,\; y+n+q\bigr).
\]
Hence
\[
g_{m,n}\circ g_{p,q}=g_{\,m+(-1)^n p,\; n+q}.
\]
So the group law on \(\mathbb Z^2\) corresponding to composition is
\[
(m,n)*(p,q)=\bigl(m+(-1)^n p,\; n+q\bigr).
\]
Thus
\[
\pi_1(K,k_0)\cong G,
\]
where \(G=\mathbb Z^2\) with the above multiplication.

We next show that the Klein bottle has a covering space homeomorphic to the torus. Consider the subgroup
\[
H=\langle a,b^2\rangle\subset \Gamma.
\]
Now
\[
a(x,y)=(x+1,y),
\qquad
b^2(x,y)=b(-x,y+1)=(x,y+2).
\]
So \(H\) consists only of translations:
\[
H=\{(x,y)\mapsto(x+m,y+2n):m,n\in\mathbb Z\}.
\]
Therefore
\[
\mathbb R^2/H
\]
is obtained by quotienting \(\mathbb R^2\) by two independent translations, so
\[
\mathbb R^2/H\cong S^1\times S^1.
\]
That is, \(\mathbb R^2/H\) is a torus.

Since \(H\subset \Gamma\), there is an induced covering map
\[
\mathbb R^2/H \to \mathbb R^2/\Gamma.
\]
In other words,
\[
T^2\to K
\]
is a covering map. Thus the Klein bottle has a covering space homeomorphic to the torus.

Finally, we show that the torus does not have a covering space homeomorphic to the Klein bottle.

Assume, for contradiction, that there exists a covering map
\[
p:K\to T^2.
\]
Then the induced homomorphism on fundamental groups is injective:
\[
p_*:\pi_1(K)\hookrightarrow \pi_1(T^2).
\]
But
\[
\pi_1(T^2)\cong \mathbb Z^2,
\]
which is abelian. Therefore every subgroup of \(\mathbb Z^2\) is abelian.

On the other hand, \(\pi_1(K)\cong G\) is not abelian. For example,
\[
(1,1)*(1,0)=\bigl(1+(-1)^1\cdot 1,\;1+0\bigr)=(0,1),
\]
while
\[
(1,0)*(1,1)=\bigl(1+(-1)^0\cdot 1,\;0+1\bigr)=(2,1).
\]
These are not equal, so
\[
(1,1)*(1,0)\neq (1,0)*(1,1).
\]
Hence \(G\) is nonabelian.

This is impossible if \(\pi_1(K)\) injects into \(\mathbb Z^2\). Therefore no such covering \(K\to T^2\) can exist.

So the Klein bottle has a covering space homeomorphic to the torus, but the torus does not have a covering space homeomorphic to the Klein bottle.
\end{solution}

\begin{problem}[CP-VIII-0014 --- Covering groups and the universal cover of $SO(3)$]
Let \(G\) be a path-connected, locally path connected topological group, and let
\[
p:\widehat{G}\to G
\]
be a path-connected covering space. Let \(\epsilon\in p^{-1}(e)\) be a point in the fibre over the identity \(e\in G\).

\begin{enumerate}
	\item[(i)] Show that \(\widehat{G}\) has a unique structure of a topological group with unit \(\epsilon\) such that \(p\) is a homomorphism.
	\item[(ii)] Show that \(\ker(p)\subset \widehat{G}\) lies in the centre of \(\widehat{G}\).
	\item[(iii)] Show that \(SO(3)\), the group of rotations of \(\mathbb R^3\) (equivalently the group of orthogonal \(3\times 3\) matrices of determinant \(1\)), is homeomorphic to the projective space \(\mathbb{RP}^3\).
	\item[(iv)] Together, (i) and (iii) give a group \(\widetilde{SO(3)}\) homeomorphic to \(S^3\). Identify this group with a well-known matrix group.
\end{enumerate}
\end{problem}

\begin{solution}
We treat the four parts in order.

\textbf{(i) Constructing the group structure on \(\widehat{G}\).}

Let
\[
H=p_*\bigl(\pi_1(\widehat{G},\epsilon)\bigr)\subset \pi_1(G,e).
\]
Since \(p\) is a covering map, \(p_*\) is injective, so \(H\) is a subgroup of \(\pi_1(G,e)\).

A standard fact about topological groups is that \(\pi_1(G,e)\) is abelian. Therefore \(H\) is automatically a normal subgroup, and it is closed under multiplication and inversion.

Now let
\[
m:G\times G\to G,
\qquad
m(g,h)=gh
\]
be the multiplication map. Since \(\widehat{G}\times \widehat{G}\) is path-connected and locally path connected, and
\[
p\times p:\widehat{G}\times \widehat{G}\to G\times G
\]
is again a covering map, we can ask whether the map
\[
m\circ (p\times p):\widehat{G}\times \widehat{G}\to G
\]
lifts to \(\widehat{G}\).

Under the standard isomorphism
\[
\pi_1(G\times G,(e,e))\cong \pi_1(G,e)\times \pi_1(G,e),
\]
the induced homomorphism
\[
m_*:\pi_1(G,e)\times \pi_1(G,e)\to \pi_1(G,e)
\]
is given by
\[
m_*([\alpha],[\beta])=[\,t\mapsto \alpha(t)\beta(t)\,].
\]
Since \(\pi_1(G,e)\) is abelian, this corresponds exactly to multiplication in \(\pi_1(G,e)\). Hence if \([\alpha],[\beta]\in H\), then
\[
m_*([\alpha],[\beta])\in H.
\]
So by the lifting criterion, the map \(m\circ (p\times p)\) lifts to a continuous map
\[
\widehat{m}:\widehat{G}\times \widehat{G}\to \widehat{G}
\]
such that
\[
\widehat{m}(\epsilon,\epsilon)=\epsilon.
\]

We now define multiplication on \(\widehat{G}\) by
\[
x\cdot y:=\widehat{m}(x,y).
\]
By construction,
\[
p(x\cdot y)=p(x)p(y).
\]

Next consider inversion
\[
\iota:G\to G,
\qquad
\iota(g)=g^{-1}.
\]
The induced map on \(\pi_1(G,e)\) sends a loop class to its inverse, so
\[
\iota_*(H)\subset H.
\]
Therefore \(\iota\circ p:\widehat{G}\to G\) lifts uniquely to a continuous map
\[
\widehat{\iota}:\widehat{G}\to \widehat{G}
\]
with
\[
\widehat{\iota}(\epsilon)=\epsilon.
\]
Define
\[
x^{-1}:=\widehat{\iota}(x).
\]
Then
\[
p(x^{-1})=p(x)^{-1}.
\]

We now verify the group axioms.

First, \(p(\epsilon)=e\), and both maps
\[
x\mapsto \epsilon\cdot x
\qquad\text{and}\qquad
x\mapsto x
\]
are lifts of the same map \(x\mapsto p(x)\) from \(\widehat{G}\) to \(G\). They also agree at \(x=\epsilon\), since
\[
\epsilon\cdot \epsilon=\epsilon.
\]
By uniqueness of lifts, they are equal. Hence
\[
\epsilon\cdot x=x.
\]
Similarly,
\[
x\cdot \epsilon=x.
\]
So \(\epsilon\) is the identity element.

Next, both maps
\[
(x,y,z)\mapsto (x\cdot y)\cdot z
\qquad\text{and}\qquad
(x,y,z)\mapsto x\cdot (y\cdot z)
\]
from \(\widehat{G}^3\) to \(\widehat{G}\) are lifts of the same map
\[
(x,y,z)\mapsto p(x)p(y)p(z)
\]
from \(\widehat{G}^3\) to \(G\), because multiplication in \(G\) is associative. They agree at
\[
(\epsilon,\epsilon,\epsilon),
\]
so by uniqueness of lifts they are equal everywhere. Thus multiplication on \(\widehat{G}\) is associative.

Finally, both maps
\[
x\mapsto x\cdot x^{-1}
\qquad\text{and}\qquad
x\mapsto \epsilon
\]
are lifts of the constant map \(x\mapsto e\), and they agree at \(x=\epsilon\). Hence
\[
x\cdot x^{-1}=\epsilon.
\]
Similarly,
\[
x^{-1}\cdot x=\epsilon.
\]

Therefore \(\widehat{G}\) is a topological group with unit \(\epsilon\), and by construction \(p\) is a continuous group homomorphism.

To prove uniqueness, note that the multiplication on \(\widehat{G}\) must be a lift of
\[
m\circ (p\times p)
\]
sending \((\epsilon,\epsilon)\) to \(\epsilon\), and such a lift is unique. The same is true for inversion. Hence the group structure is unique.

\textbf{(ii) The kernel lies in the centre.}

Let
\[
\ker(p)=p^{-1}(e).
\]
Take \(a\in \ker(p)\). Consider left multiplication and right multiplication by \(a\):
\[
L_a:\widehat{G}\to \widehat{G},
\qquad
L_a(x)=ax,
\]
\[
R_a:\widehat{G}\to \widehat{G},
\qquad
R_a(x)=xa.
\]
For every \(x\in \widehat{G}\),
\[
p(L_a(x))=p(a)p(x)=ep(x)=p(x),
\]
and similarly
\[
p(R_a(x))=p(x)p(a)=p(x)e=p(x).
\]
So both \(L_a\) and \(R_a\) are lifts of the identity map
\[
\operatorname{id}_G:G\to G
\]
through the covering \(p\).

Moreover,
\[
L_a(\epsilon)=a\epsilon=a,
\qquad
R_a(\epsilon)=\epsilon a=a.
\]
Thus \(L_a\) and \(R_a\) are two lifts of the same map that agree at the point \(\epsilon\). By uniqueness of lifts, they must be equal:
\[
L_a=R_a.
\]
Hence for every \(x\in \widehat{G}\),
\[
ax=xa.
\]
Therefore every element of \(\ker(p)\) commutes with every element of \(\widehat{G}\), so
\[
\ker(p)\subset Z(\widehat{G}),
\]
the centre of \(\widehat{G}\).

\textbf{(iii) Showing that \(SO(3)\cong \mathbb{RP}^3\) as spaces.}

We use the axis-angle description of rotations.

Every rotation \(R\in SO(3)\) is determined by:

\begin{itemize}
	\item an axis \(u\in S^2\),
	\item an angle \(\theta\in [0,\pi]\),
\end{itemize}
with the following convention:

\begin{itemize}
	\item if \(0\le \theta<\pi\), then the axis \(u\) and angle \(\theta\) are uniquely determined;
	\item if \(\theta=\pi\), then the axis is determined only up to sign, because rotation by \(\pi\) about \(u\) is the same as rotation by \(\pi\) about \(-u\).
\end{itemize}

Now let
\[
B^3=\{x\in \mathbb R^3:\|x\|\le \pi\}.
\]
Every nonzero point \(x\in B^3\) may be written uniquely as
\[
x=\theta u,
\qquad
0<\theta\le \pi,\quad u\in S^2.
\]
Define
\[
F:B^3\to SO(3)
\]
by sending \(x=\theta u\) to the rotation of \(\mathbb R^3\) through angle \(\theta\) about the axis \(u\), and let
\[
F(0)=I.
\]
This map is continuous and surjective.

If \(\|x\|=\pi\), say \(x=\pi u\), then
\[
F(\pi u)=F(-\pi u),
\]
because rotation by angle \(\pi\) about \(u\) is the same as rotation by angle \(\pi\) about \(-u\). Thus \(F\) identifies antipodal points on the boundary sphere \(\partial B^3=S^2\).

Therefore \(F\) factors through the quotient
\[
B^3/\!\sim,
\qquad
x\sim -x \quad \text{for } x\in \partial B^3.
\]
So we get a continuous map
\[
\overline{F}:B^3/\!\sim\;\to SO(3).
\]

Now the quotient \(B^3/\!\sim\) is homeomorphic to \(\mathbb{RP}^3\). Indeed, one may view \(B^3\) as a closed hemisphere in \(S^3\); when we pass to the quotient by the antipodal relation on the boundary, we obtain exactly
\[
S^3/\{\pm 1\}=\mathbb{RP}^3.
\]

It remains to show that \(\overline{F}\) is bijective.

\begin{itemize}
	\item It is surjective because every element of \(SO(3)\) has an axis-angle description.
	\item It is injective because:
	\begin{itemize}
		\item if \(0\le \theta<\pi\), the axis-angle description is unique;
		\item if \(\theta=\pi\), the only ambiguity is \(u\leftrightarrow -u\), and those two points have already been identified in the quotient.
	\end{itemize}
\end{itemize}

Thus \(\overline{F}\) is a continuous bijection
\[
\mathbb{RP}^3\to SO(3).
\]
Since \(\mathbb{RP}^3\) is compact and \(SO(3)\) is Hausdorff, it follows that \(\overline{F}\) is a homeomorphism.

Therefore
\[
SO(3)\cong \mathbb{RP}^3.
\]

\textbf{(iv) Identifying the universal covering group of \(SO(3)\).}

From part (iii), we know that
\[
SO(3)\cong \mathbb{RP}^3.
\]
The universal covering of \(\mathbb{RP}^3\) is
\[
S^3\to \mathbb{RP}^3,
\]
given by the antipodal quotient
\[
x\mapsto [x].
\]
Hence \(SO(3)\) has a universal covering space homeomorphic to \(S^3\).

By part (i), once we choose a point over the identity, this covering space \(S^3\) acquires a unique topological group structure making the covering map
\[
S^3\to SO(3)
\]
a continuous group homomorphism. This group is the universal covering group \(\widetilde{SO(3)}\).

We now identify it with a standard matrix group. The correct group is
\[
SU(2)=\{A\in M_2(\mathbb C):A^*A=I,\ \det A=1\}.
\]

Every element of \(SU(2)\) has the form
\[
\begin{pmatrix}
	\alpha & \beta\\
	-\overline{\beta} & \overline{\alpha}
\end{pmatrix},
\qquad
|\alpha|^2+|\beta|^2=1.
\]
Write
\[
\alpha=a+bi,
\qquad
\beta=c+di.
\]
Then the condition
\[
|\alpha|^2+|\beta|^2=1
\]
becomes
\[
a^2+b^2+c^2+d^2=1.
\]
So \(SU(2)\) is homeomorphic to
\[
S^3\subset \mathbb R^4.
\]

There is a standard surjective homomorphism
\[
SU(2)\to SO(3)
\]
whose kernel is
\[
\{\pm I\}.
\]
Equivalently, if one identifies \(SU(2)\) with the group of unit quaternions, then a unit quaternion \(q\) acts on the space of imaginary quaternions \(\operatorname{Im}\mathbb H\cong \mathbb R^3\) by
\[
v\mapsto qvq^{-1},
\]
and this action gives a homomorphism
\[
S^3\to SO(3)
\]
with kernel
\[
\{\pm 1\}.
\]
Thus
\[
SO(3)\cong SU(2)/\{\pm I\}.
\]

Since \(SU(2)\) is connected, simply connected, and homeomorphic to \(S^3\), it is the universal covering group of \(SO(3)\). Therefore the group \(\widetilde{SO(3)}\) obtained from parts (i) and (iii) is precisely
\[
\widetilde{SO(3)}\cong SU(2).
\]

So the well-known matrix group is
\[
\boxed{SU(2)}.
\]

Equivalently, one may identify it with the group of unit quaternions.
\end{solution}

\begin{problem}[CP-VIII-0015 --- Universal cover of \texorpdfstring{$S^1 \vee S^2$}{S1 vee S2}]
Recall that for based spaces \((X,x_0)\), \((Y,y_0)\), the wedge product \(X \vee Y\) is the quotient of the space
\[
X \sqcup Y
\]
by the relation generated by \(x_0 \sim y_0\).
What is the universal cover of \(S^1 \vee S^2\)? Draw a picture.
\end{problem}

\begin{solution}
Since \(S^2\) is simply connected, wedging it to \(S^1\) does not change the fundamental group:
\[
\pi_1(S^1 \vee S^2)\cong \pi_1(S^1)\cong \mathbb{Z}.
\]
So the universal cover should be a simply connected space on which \(\mathbb{Z}\) acts by deck transformations.

Start from the universal cover of \(S^1\), namely
\[
\mathbb{R}\to S^1.
\]
Now in \(S^1\vee S^2\), the sphere \(S^2\) is attached at the wedge point. In the universal cover, every lift of that wedge point gets its own copy of \(S^2\). The lifts of the wedge point are the integers in \(\mathbb{R}\). Therefore the universal cover is:

\[
\text{a copy of }\mathbb{R}\text{ with one }S^2\text{ attached at each integer point.}
\]

Equivalently,
\[
\widetilde{S^1\vee S^2}
\cong
\left(\mathbb{R}\right)\cup \bigcup_{n\in\mathbb{Z}} S^2_n,
\]
where each \(S^2_n\) is attached to \(\mathbb{R}\) at the point \(n\in\mathbb{R}\).

This space is simply connected: the line \(\mathbb{R}\) is simply connected, and each attached sphere is simply connected, so no new fundamental group appears.

A picture is:

\[
\cdots \;-\; \bullet \;-\; \bullet \;-\; \bullet \;-\; \bullet \;-\; \cdots
\]
with a \(2\)-sphere attached at each bullet. The bullets correspond to the integers in \(\mathbb{R}\).

So the universal cover of \(S^1\vee S^2\) is an infinite chain of copies of \(S^2\), indexed by \(\mathbb{Z}\), joined by a line.
\end{solution}

\begin{problem}[CP-VIII-0016 --- Quotients by free properly discontinuous actions and odd-dimensional lens spaces]
Let \(X\) be a Hausdorff space, and \(G\) a group acting on \(X\) by homeomorphisms, freely (i.e. if \(g\in G\) satisfies \(g\cdot x=x\) for some \(x\in X\), then \(g=e\)) and properly discontinuously (i.e. each \(x\in X\) has an open neighbourhood \(U\ni x\) such that
\[
\{g\in G\mid g(U)\cap U\neq \varnothing\}
\]
is finite).

\begin{enumerate}
	\item[(i)] Show that the quotient map \(X\to X/G\) is a covering map.
	\item[(ii)] Deduce that if \(X\) is simply connected and locally path-connected then for any point \([x]\in X/G\) we have an isomorphism of groups
	\[
	\pi_1(X/G,[x])\cong G.
	\]
	\item[(iii)] Hence show that for \(n\ge 2\) odd and any \(m\ge 2\) there is a space \(Y\) with fundamental group \(\mathbb{Z}/m\) and universal cover \(S^n\). \emph{[Hint: Consider \(S^n\) as the unit sphere in \(\mathbb{C}^k\).]}
\end{enumerate}
\end{problem}

\begin{solution}
\textbf{(i) The quotient map is a covering map.}
Let \(p:X\to X/G\) be the quotient map, and fix \(x\in X\).

By proper discontinuity, there exists an open neighbourhood \(U\ni x\) such that
\[
F:=\{g\in G\mid g(U)\cap U\neq \varnothing\}
\]
is finite.

For each \(g\in F\setminus\{e\}\), we have \(gx\neq x\) because the action is free. Since \(X\) is Hausdorff, we can choose an open neighbourhood \(U_g\subseteq U\) of \(x\) such that
\[
U_g\cap g(U_g)=\varnothing.
\]
Now set
\[
W:=\bigcap_{g\in F\setminus\{e\}} U_g.
\]
Then \(W\) is an open neighbourhood of \(x\), and for every \(g\neq e\),
\[
g(W)\cap W=\varnothing.
\]
Indeed, if \(g(W)\cap W\neq\varnothing\), then since \(W\subseteq U\) we get \(g(U)\cap U\neq\varnothing\), hence \(g\in F\), contradicting the construction of \(W\).

Now consider \(p(W)\subseteq X/G\). We claim that
\[
p^{-1}(p(W))=\bigsqcup_{g\in G} g(W).
\]
This is clear: a point \(y\in X\) lies in \(p^{-1}(p(W))\) exactly when \(p(y)=p(w)\) for some \(w\in W\), i.e. \(y=g(w)\) for some \(g\in G\), so \(y\in g(W)\).

Also the sets \(g(W)\) are pairwise disjoint. Indeed, if
\[
g_1(W)\cap g_2(W)\neq\varnothing,
\]
then
\[
W\cap (g_1^{-1}g_2)(W)\neq\varnothing,
\]
so \(g_1^{-1}g_2=e\), hence \(g_1=g_2\).

Finally, for each \(g\in G\), the restriction
\[
p|_{g(W)}:g(W)\to p(W)
\]
is a homeomorphism, because \(p\) is injective on \(g(W)\) and surjective onto \(p(W)\), and \(p\) is a quotient map.

Thus \(p^{-1}(p(W))\) is a disjoint union of open sets each mapped homeomorphically onto \(p(W)\). Therefore \(p:X\to X/G\) is a covering map.

\textbf{(ii) Computation of the fundamental group.}
Assume now that \(X\) is simply connected and locally path-connected.

By part (i), the quotient map
\[
p:X\to X/G
\]
is a covering map. Since \(X\) is simply connected, this covering is the universal covering of \(X/G\).

The group \(G\) acts on \(X\) by homeomorphisms, and each element \(g\in G\) preserves fibres of \(p\), so each \(g\) is a deck transformation of the universal cover. Conversely, if two points \(x,y\in X\) lie in the same fibre of \(p\), then \(y=gx\) for a unique \(g\in G\). Hence the deck transformation group is exactly \(G\).

For a universal cover of a connected, locally path-connected space, the deck transformation group is isomorphic to the fundamental group of the base. Therefore
\[
\pi_1(X/G,[x])\cong G.
\]

\textbf{(iii) A space with fundamental group \(\mathbb{Z}/m\) and universal cover \(S^n\).}
Let \(n\ge 2\) be odd. Then
\[
n=2k-1
\]
for some \(k\ge 2\), and we may view
\[
S^n=S^{2k-1}=\left\{(z_1,\dots,z_k)\in \mathbb{C}^k:\sum_{j=1}^k |z_j|^2=1\right\}.
\]

Let \(\mu_m\subset \mathbb{C}^\times\) be the group of \(m\)-th roots of unity:
\[
\mu_m=\{\zeta\in\mathbb{C}:\zeta^m=1\}\cong \mathbb{Z}/m.
\]
Define an action of \(\mu_m\) on \(S^{2k-1}\) by
\[
\zeta\cdot (z_1,\dots,z_k)=(\zeta z_1,\dots,\zeta z_k).
\]

This action is free. Indeed, if
\[
\zeta\cdot (z_1,\dots,z_k)=(z_1,\dots,z_k),
\]
then for every coordinate \(j\),
\[
\zeta z_j=z_j.
\]
Since the point lies on the sphere, not all \(z_j\) are zero, so for some \(j\) we have \(z_j\neq 0\), hence \(\zeta=1\).

Since \(\mu_m\) is finite, the action is automatically properly discontinuous.

Therefore, by part (ii), the quotient
\[
Y:=S^{2k-1}/\mu_m
\]
has
\[
\pi_1(Y)\cong \mu_m\cong \mathbb{Z}/m,
\]
and its universal cover is \(S^{2k-1}=S^n\).

This quotient is a standard lens space.
\end{solution}

\begin{problem}[CP-VIII-0017 --- Isomorphisms of groups given by presentations]
\noindent\textbf{(i)} Show that the group
\[
G=\langle a,b\mid a\rangle
\]
is infinite cyclic; equivalently,
\[
G\cong \langle t\rangle\cong\mathbb Z.
\]

\begin{enumerate}
	\item[(ii)] Show that the groups
	\[
	G=\langle a,b\mid ab^{-3},\,ba^{-2}\rangle
	\qquad\text{and}\qquad
	H=\langle t\mid t^5\rangle
	\]
	are isomorphic.

	\item[(iii)] Show that the groups
	\[
	G=\langle a,b\mid a^3b^{-2}\rangle
	\qquad\text{and}\qquad
	H=\langle x,y\mid xyxy^{-1}x^{-1}y^{-1}\rangle
	\]
	are isomorphic. Show that this group is non-abelian and infinite.
\end{enumerate}
\end{problem}

\begin{solution}
\textbf{(i)}
In
\[
G=\langle a,b\mid a\rangle,
\]
the relation \(a=1\) kills the generator \(a\). Hence \(G\) is generated only by \(b\), with no further relation:
\[
G\cong \langle b\mid \ \rangle \cong \mathbb{Z}.
\]
Also
\[
H=\langle t\mid \ \rangle \cong \mathbb{Z}.
\]
Therefore
\[
G\cong H.
\]

\textbf{(ii)}
The relations are
\[
ab^{-3}=1,\qquad ba^{-2}=1,
\]
so
\[
a=b^3,\qquad b=a^2.
\]
Substituting \(a=b^3\) into \(b=a^2\), we get
\[
b=(b^3)^2=b^6,
\]
hence
\[
b^5=1.
\]
Also \(a=b^3\), so every element of \(G\) is a power of \(b\). Thus
\[
G\cong \langle b\mid b^5\rangle \cong \mathbb{Z}/5.
\]
But
\[
H=\langle t\mid t^5\rangle \cong \mathbb{Z}/5.
\]
So
\[
G\cong H.
\]

An explicit isomorphism is given by
\[
\varphi:G\to H,\qquad \varphi(b)=t,\quad \varphi(a)=t^3.
\]

\textbf{(iii)}
First rewrite the relation in \(H\):
\[
xyxy^{-1}x^{-1}y^{-1}=1.
\]
This is equivalent to
\[
xyx=yxy.
\]
Indeed,
\[
xyxy^{-1}x^{-1}y^{-1}=1
\iff xyxy^{-1}x^{-1}=y
\iff xyxy^{-1}=yx
\iff xyx=yxy.
\]

So
\[
H=\langle x,y\mid xyx=yxy\rangle,
\]
the standard braid-group-on-three-strands presentation.

Now define
\[
\varphi:G\to H
\]
by
\[
\varphi(a)=xy,\qquad \varphi(b)=xyx.
\]
We check the defining relation of \(G\):
\[
\varphi(a)^3=(xy)^3=xyxyxy.
\]
Using \(xyx=yxy\),
\[
(xy)^3=(xyx)(yxy)=(xyx)^2=\varphi(b)^2.
\]
Hence
\[
\varphi(a^3b^{-2})=1,
\]
so \(\varphi\) is a well-defined homomorphism.

Now define
\[
\psi:H\to G
\]
by
\[
\psi(x)=a^{-1}b,\qquad \psi(y)=b^{-1}a^2.
\]
Then
\[
\psi(x)\psi(y)=a^{-1}bb^{-1}a^2=a,
\]
so
\[
\psi(x)\psi(y)\psi(x)=a(a^{-1}b)=b.
\]
Also
\[
\psi(y)\psi(x)\psi(y)=\psi(y)\,a
= b^{-1}a^2a
= b^{-1}a^3
= b^{-1}b^2
= b.
\]
Thus
\[
\psi(x)\psi(y)\psi(x)=\psi(y)\psi(x)\psi(y),
\]
so \(\psi\) respects the relation \(xyx=yxy\), hence is well-defined.

Now compute:
\[
\psi(\varphi(a))=\psi(xy)=a,\qquad
\psi(\varphi(b))=\psi(xyx)=b,
\]
and also
\[
\varphi(\psi(x))=\varphi(a^{-1}b)=(xy)^{-1}(xyx)=x,
\]
\[
\varphi(\psi(y))=\varphi(b^{-1}a^2)=(xyx)^{-1}(xy)^2=y.
\]
Therefore \(\varphi\) and \(\psi\) are inverse isomorphisms, so
\[
G\cong H.
\]

\textbf{The group is non-abelian.}
Consider the homomorphism
\[
G\to S_3
\]
defined by
\[
a\mapsto (123),\qquad b\mapsto (12).
\]
Then
\[
a^3=e,\qquad b^2=e,
\]
so the relation \(a^3=b^2\) is satisfied. Hence this is a well-defined homomorphism. Its image contains a \(3\)-cycle and a transposition, so the image is all of \(S_3\), which is non-abelian. Therefore \(G\), and hence \(H\), is non-abelian.

\textbf{The group is infinite.}
Define a homomorphism
\[
G\to \mathbb{Z}
\]
by
\[
a\mapsto 2,\qquad b\mapsto 3.
\]
This is well-defined because the relation \(a^3=b^2\) becomes
\[
3\cdot 2 = 2\cdot 3,
\]
which is true in \(\mathbb{Z}\). The image contains \(2\) and \(3\), hence contains \(\gcd(2,3)=1\), so it is all of \(\mathbb{Z}\). Thus \(G\) surjects onto an infinite group, and therefore \(G\) itself is infinite.

So the group in part (iii) is isomorphic to \(B_3\), non-abelian, and infinite.
\end{solution}

\begin{problem}[CP-VIII-0018 --- The fundamental group of \texorpdfstring{$S^1\vee S^1$}{S1 vee S1} and nonexistence of a retraction]
Use Seifert--van Kampen to show that the fundamental group of \(S^1\vee S^1\) is a free group on two generators. Deduce that the inclusion
\[
i:S^1\vee S^1=(S^1\times\{*\})\cup(\{*\}\times S^1)\hookrightarrow S^1\times S^1
\]
does not admit a retraction.
\end{problem}

\begin{solution}
Let
\[
X=S^1\vee S^1.
\]
We regard \(X\) as the union of two circles meeting in one point \(x_0\). Choose open neighbourhoods \(U,V\subseteq X\) such that

\[
U\simeq S^1,\qquad V\simeq S^1,\qquad U\cap V \text{ is contractible,}
\]
and
\[
X=U\cup V.
\]

Then Seifert--van Kampen gives
\[
\pi_1(X,x_0)\cong \pi_1(U,x_0)*\pi_1(V,x_0).
\]
Since
\[
\pi_1(U,x_0)\cong \mathbb Z,\qquad \pi_1(V,x_0)\cong \mathbb Z,
\]
we obtain
\[
\pi_1(S^1\vee S^1,x_0)\cong \mathbb Z * \mathbb Z.
\]
Thus \(\pi_1(S^1\vee S^1,x_0)\) is the free group on two generators. If \(a\) and \(b\) denote the two characteristic loops, we may write
\[
\pi_1(S^1\vee S^1,x_0)\cong F_2=\langle a,b\rangle.
\]

Now consider the inclusion
\[
i:S^1\vee S^1\hookrightarrow S^1\times S^1.
\]
The torus has fundamental group
\[
\pi_1(S^1\times S^1)\cong \mathbb Z^2,
\]
and the induced homomorphism
\[
i_*:\pi_1(S^1\vee S^1,x_0)\to \pi_1(S^1\times S^1,i(x_0))
\]
sends the two generators \(a,b\) to the two standard generators of \(\mathbb Z^2\). Hence
\[
i_*:F_2\to \mathbb Z^2
\]
is the abelianization map. In particular,
\[
i_*([a,b])=0,
\]
so \(i_*\) is not injective.

Assume for contradiction that there exists a retraction
\[
r:S^1\times S^1\to S^1\vee S^1
\]
such that
\[
r\circ i=\mathrm{id}_{S^1\vee S^1}.
\]
Then on fundamental groups we would have
\[
r_*\circ i_*=\mathrm{id}_{\pi_1(S^1\vee S^1,x_0)}.
\]
Therefore \(i_*\) would have to be injective, which is impossible.

Hence the inclusion
\[
i:S^1\vee S^1\hookrightarrow S^1\times S^1
\]
does not admit a retraction.

\par
\end{solution}

\begin{problem}[CP-VIII-0019 --- Connected degree \texorpdfstring{$2$}{2} and degree \texorpdfstring{$3$}{3} coverings of \texorpdfstring{$S^1\vee S^1$}{S1 vee S1}]
A covering space is called normal if it corresponds to a normal subgroup. Draw pictures of all the connected degree \(2\) covering spaces of \(S^1\vee S^1\). Show that they are all normal coverings. Now do the same thing for the connected degree \(3\) covering spaces of \(S^1\vee S^1\). Which of them are normal coverings?
\end{problem}

\begin{solution}
Let
\[
X=S^1\vee S^1,
\qquad
\pi_1(X)\cong F_2=\langle a,b\rangle.
\]
A connected degree \(n\) covering of \(X\) is described by a transitive permutation representation
\[
\rho:F_2\to S_n.
\]
Equivalently, it is determined by a pair of permutations
\[
\sigma_a=\rho(a),\qquad \sigma_b=\rho(b)
\]
in \(S_n\) such that the subgroup generated by \(\sigma_a,\sigma_b\) acts transitively on \(\{1,\dots,n\}\). Two such coverings are isomorphic if and only if the two pairs are simultaneously conjugate in \(S_n\).

So we classify connected coverings by classifying transitive pairs \((\sigma_a,\sigma_b)\) up to simultaneous conjugacy.

\par

\textbf{Degree \(2\).}
Here
\[
S_2=\{1,\tau\},
\qquad
\tau=(12).
\]
For the action to be transitive, at least one of \(\sigma_a,\sigma_b\) must be \(\tau\). Thus, up to isomorphism, there are exactly three connected degree \(2\) coverings.

\par

\textbf{(1) \(\sigma_a=\tau,\ \sigma_b=\mathrm{id}\).}
There are two vertices \(1,2\). The \(a\)-edges form a \(2\)-cycle joining the two vertices, while the \(b\)-edges are loops, one at each vertex.

In graph form:
\begin{itemize}
	\item \(a\): one edge from \(1\) to \(2\) and one from \(2\) to \(1\),
	\item \(b\): one loop at \(1\), one loop at \(2\).
\end{itemize}

\par

\textbf{(2) \(\sigma_a=\mathrm{id},\ \sigma_b=\tau\).}
This is the same picture with \(a\) and \(b\) interchanged:
\begin{itemize}
	\item \(a\): one loop at each vertex,
	\item \(b\): a \(2\)-cycle joining the two vertices.
\end{itemize}

\par

\textbf{(3) \(\sigma_a=\tau,\ \sigma_b=\tau\).}
Again there are two vertices \(1,2\), and both \(a\) and \(b\) swap the two sheets:
\begin{itemize}
	\item \(a\): a \(2\)-cycle joining \(1\) and \(2\),
	\item \(b\): another \(2\)-cycle joining \(1\) and \(2\).
\end{itemize}

\par

\textbf{Why all degree \(2\) connected coverings are normal.}
A connected degree \(2\) covering corresponds to a subgroup of index \(2\) in \(F_2\). Every subgroup of index \(2\) is normal. Therefore every connected degree \(2\) covering of \(S^1\vee S^1\) is normal.

So:
\[
\boxed{\text{There are exactly \(3\) connected degree \(2\) coverings, and all of them are normal.}}
\]

\par

\textbf{Degree \(3\).}
Now we classify transitive pairs in \(S_3\), up to simultaneous conjugacy. Let
\[
c=(123),\qquad t=(12).
\]
The transitive subgroups of \(S_3\) are \(A_3=\langle c\rangle\) and \(S_3\). Hence there are two kinds of connected degree \(3\) coverings: those whose monodromy image is \(A_3\), and those whose monodromy image is \(S_3\).

\par

\textbf{Case A: image \(A_3\).}
Then both \(\sigma_a,\sigma_b\) lie in
\[
\{1,c,c^{-1}\},
\]
and at least one of them must be a \(3\)-cycle. Up to simultaneous conjugacy, there are four possibilities.

\par

\textbf{(1) \((\sigma_a,\sigma_b)=(c,\mathrm{id})\).}
The \(a\)-edges form a \(3\)-cycle on the vertices \(1,2,3\), and the \(b\)-edges are loops at each vertex.

\par

\textbf{(2) \((\sigma_a,\sigma_b)=(\mathrm{id},c)\).}
This is the same with \(a\) and \(b\) interchanged: the \(b\)-edges form a \(3\)-cycle, and the \(a\)-edges are loops.

\par

\textbf{(3) \((\sigma_a,\sigma_b)=(c,c)\).}
Both \(a\) and \(b\) run around the same \(3\)-cycle in the same direction.

\par

\textbf{(4) \((\sigma_a,\sigma_b)=(c,c^{-1})\).}
The \(a\)-edges and \(b\)-edges both form the same \(3\)-cycle, but with opposite orientations.

\par

These four coverings are all normal, because in each case the image is
\[
A_3\cong \mathbb Z/3,
\]
so the corresponding subgroup has normal quotient of order \(3\).

\par

\textbf{Case B: image \(S_3\).}
Now the pair \((\sigma_a,\sigma_b)\) generates all of \(S_3\). Up to simultaneous conjugacy, there are three possibilities.

\par

\textbf{(5) \((\sigma_a,\sigma_b)=(t,c)\).}
Here
\[
\sigma_a=(12),\qquad \sigma_b=(123).
\]
So:
\begin{itemize}
	\item \(a\): one vertex carries an \(a\)-loop and the other two vertices are joined by an \(a\)-edge-pair,
	\item \(b\): the three vertices form a \(3\)-cycle.
\end{itemize}

\par

\textbf{(6) \((\sigma_a,\sigma_b)=(c,t)\).}
This is the same with \(a\) and \(b\) interchanged.

\par

\textbf{(7) \((\sigma_a,\sigma_b)=((12),(13))\).}
Here both \(a\) and \(b\) are transpositions, but different ones. One may describe the graph as follows:
\begin{itemize}
	\item \(a\): fixes one vertex and swaps the other two,
	\item \(b\): fixes a different vertex and swaps the other two.
\end{itemize}
Since \((12)\) and \((13)\) generate \(S_3\), the action is transitive, so the covering is connected.

\par

These three coverings are not normal. Indeed, if an index \(3\) subgroup were normal, the quotient would have order \(3\), so the monodromy image would be cyclic of order \(3\). Here the image is \(S_3\), so the subgroup is not normal.

Therefore:
\[
\boxed{\text{There are exactly \(7\) connected degree \(3\) coverings.}}
\]
Among them:
\[
\boxed{\text{the first \(4\) are normal, and the last \(3\) are not normal.}}
\]

\par
\end{solution}

\begin{problem}[CP-VIII-0020 --- Coverings associated to specific subgroups of \texorpdfstring{$\pi_1(S^1\vee S^1)$}{pi1(S1 vee S1)}]
Consider
\[
X=S^1\vee S^1
\]
with basepoint \(x_0\) the wedge point, and
\[
\pi_1(X,x_0)=\langle a,b\rangle.
\]
Describe covering spaces associated to
\begin{enumerate}
	\item[(i)] \(\langle\!\langle a\rangle\!\rangle\), the normal subgroup generated by \(a\),
	\item[(ii)] \(\langle a\rangle\), the subgroup generated by \(a\),
	\item[(iii)] the kernel of the homomorphism
	\[
	\phi:\langle a,b\rangle\to \mathbb Z/4
	\]
	given by
	\[
	\phi(a)=[1],\qquad \phi(b)=[3]=[-1].
	\]
\end{enumerate}
What are the fundamental groups of these covering spaces?
\end{problem}

\begin{solution}

\medskip
\noindent\textbf{(i) The covering corresponding to \texorpdfstring{$\langle\!\langle a\rangle\!\rangle$}{<<a>>}}\par
\smallskip


Let
\[
H=\langle\!\langle a\rangle\!\rangle,
\]
the normal closure of \(a\) in \(F_2=\langle a,b\rangle\). This is the kernel of the homomorphism
\[
F_2\to \mathbb Z,\qquad a\mapsto 0,\quad b\mapsto 1.
\]
Hence
\[
F_2/H\cong \mathbb Z.
\]
So the corresponding covering is a connected normal infinite-sheeted covering with deck transformation group \(\mathbb Z\).

A convenient picture is the following graph. The vertices are indexed by the integers
\[
\cdots,-2,-1,0,1,2,\cdots,
\]
the \(b\)-edges join \(n\) to \(n+1\), and at each vertex there is one \(a\)-loop. Thus the covering graph is an infinite line in the \(b\)-direction with an \(a\)-circle attached at every vertex.

Schematically:
\[
\cdots \text{---} \bullet \text{---} \bullet \text{---} \bullet \text{---} \bullet \text{---} \cdots
\]
with an \(a\)-loop at every bullet.

This agrees with the quotient description: \(a\) has been killed in the quotient \(F_2/H\), while \(b\) survives and acts by deck translations.

The fundamental group of the covering is precisely the subgroup \(H\). Therefore
\[
\pi_1(\widetilde X)\cong H=\langle\!\langle a\rangle\!\rangle.
\]
Topologically, this covering graph has infinitely many independent loops, so its fundamental group is a free group of countably infinite rank. A convenient generating family is
\[
\{\, b^n a b^{-n}: n\in \mathbb Z\,\}.
\]

Hence
\[
\pi_1(\widetilde X)\cong \text{free group of countably infinite rank}.
\]


\medskip
\noindent\textbf{(ii) The covering corresponding to \texorpdfstring{$\langle a\rangle$}{<a>}}\par
\smallskip


Now let
\[
H=\langle a\rangle\cong \mathbb Z.
\]
This subgroup is not normal in \(F_2\), so the corresponding covering is not normal.

A geometric way to describe this cover is to start with the universal cover of \(S^1\vee S^1\), which is a tree, and then quotient by the action of the subgroup \(\langle a\rangle\). The element \(a\) acts by translation along one axis in the tree. After quotienting, this axis becomes a single \(a\)-loop, while all the remaining branches of the tree descend to infinite trees attached to that loop.

So the covering has:
\begin{itemize}
	\item one \(a\)-cycle,
	\item infinitely many tree-like branches attached to it,
	\item no other essential cycles.
\end{itemize}

Equivalently, the core of the covering graph is a single loop labelled \(a\), and outside that core the graph is a forest.

Since the covering corresponds to \(H\), we have
\[
\pi_1(\widetilde X)\cong H\cong \mathbb Z.
\]
Thus the fundamental group of this covering is
\[
\pi_1(\widetilde X)\cong \mathbb Z.
\]


\medskip
\noindent\textbf{(iii) The covering corresponding to \texorpdfstring{$\ker\phi$}{ker phi}, where \texorpdfstring{$\phi(a)=1$}{phi(a)=1} and \texorpdfstring{$\phi(b)=3=-1$}{phi(b)=3=-1} in \texorpdfstring{$\mathbb Z/4$}{Z/4}}\par
\smallskip


Let
\[
H=\ker\phi.
\]
Since \(\phi\) is surjective onto \(\mathbb Z/4\), the subgroup \(H\) has index \(4\). Hence the corresponding covering is a connected \(4\)-sheeted normal covering.

Label the four sheets by the elements of \(\mathbb Z/4\):
\[
0,1,2,3.
\]
Then the monodromy action is:
\[
a:i\mapsto i+1,\qquad b:i\mapsto i-1,
\]
with indices taken modulo \(4\).

Therefore the covering graph has four vertices \(0,1,2,3\), and:
\begin{itemize}
	\item an \(a\)-edge from \(i\) to \(i+1\),
	\item a \(b\)-edge from \(i\) to \(i-1\).
\end{itemize}
So the \(a\)-edges form a \(4\)-cycle in one direction, while the \(b\)-edges form a \(4\)-cycle in the opposite direction. As an unlabeled graph, one may think of it as a \(4\)-cycle with parallel \(a\)- and \(b\)-edges joining consecutive vertices.

Since the covering corresponds to \(H=\ker\phi\), we have
\[
\pi_1(\widetilde X)\cong H.
\]

To compute its rank, use the Nielsen--Schreier formula: a subgroup of index \(d\) in a free group of rank \(r\) has rank
\[
1+d(r-1).
\]
Here \(d=4\) and \(r=2\), so
\[
\operatorname{rank}(H)=1+4(2-1)=5.
\]
Hence
\[
\pi_1(\widetilde X)\cong \text{free group of rank }5.
\]

\par

\textbf{Summary.}
For the three coverings in this problem we obtain:

\begin{enumerate}
	\item[(i)] For \(H=\langle\!\langle a\rangle\!\rangle\):
	\begin{itemize}
		\item the covering is an infinite line of \(b\)-edges with one \(a\)-loop at each vertex,
		\item it is normal,
		\item the deck group is \(\mathbb Z\),
		\item
		\[
		\pi_1(\widetilde X)\cong \text{free group of countably infinite rank}.
		\]
	\end{itemize}

	\item[(ii)] For \(H=\langle a\rangle\):
	\begin{itemize}
		\item the covering has one \(a\)-loop as core, with infinite trees attached,
		\item it is not normal,
		\item
		\[
		\pi_1(\widetilde X)\cong \mathbb Z.
		\]
	\end{itemize}

	\item[(iii)] For \(H=\ker\phi\):
	\begin{itemize}
		\item the covering is a \(4\)-sheeted normal cover with vertices \(0,1,2,3\),
		\item \(a\) acts by \(i\mapsto i+1\), \(b\) acts by \(i\mapsto i-1\),
		\item the deck group is \(\mathbb Z/4\),
		\item
		\[
		\pi_1(\widetilde X)\cong \text{free group of rank }5.
		\]
	\end{itemize}
\end{enumerate}
\end{solution}

\begin{problem}[CP-VIII-0021 --- Triangulations of the torus, Klein bottle, and projective plane]
Show that there are triangulations of the torus, Klein bottle, and projective plane as indicated in the figure.

How many vertices, edges, and faces does each triangulation have? What is the number
\[
\chi = \text{vertices} - \text{edges} + \text{faces}
\]
for each triangulation?
\end{problem}

\begin{solution}
We count the numbers of vertices, edges, and triangular faces in each quotient triangulation.

\par

\textbf{1. The torus and Klein bottle}

In the figure for the torus and Klein bottle, the square is divided into a \(3\times 3\) grid, and each small square is cut by a diagonal. Thus before taking the quotient, the square is triangulated as follows:

\begin{itemize}
	\item vertices: \(4\times 4=16\),
	\item small squares: \(3\times 3=9\),
	\item triangles: \(2\cdot 9=18\).
\end{itemize}

So the number of faces after quotienting is still
\[
F=18,
\]
since no two triangles are identified with one another.

\par

\textbf{Counting vertices.}

There are:
\begin{itemize}
	\item \(4\) corner vertices,
	\item \(8\) non-corner boundary vertices,
	\item \(4\) interior vertices.
\end{itemize}

For both the torus and the Klein bottle, the four corners are all identified to a single vertex in the quotient.

The \(8\) non-corner boundary vertices are identified in pairs under the side gluings, giving \(4\) boundary vertex classes.

The \(4\) interior vertices remain distinct.

Hence the number of vertices is
\[
V = 1 + 4 + 4 = 9.
\]

\par

\textbf{Counting edges.}

Before quotienting, the subdivided square has:
\begin{itemize}
	\item \(12\) horizontal edges,
	\item \(12\) vertical edges,
	\item \(9\) diagonal edges.
\end{itemize}
Thus in total there are
\[
12+12+9=33
\]
edges in the square triangulation.

Now pass to the quotient.

\begin{itemize}
	\item Among the horizontal edges, \(6\) lie on the top and bottom boundary and are identified in pairs, so they give \(3\) quotient edges. The remaining \(6\) horizontal edges are interior and remain distinct.
	\item Among the vertical edges, \(6\) lie on the left and right boundary and are identified in pairs, so they give \(3\) quotient edges. The remaining \(6\) vertical edges are interior and remain distinct.
	\item All \(9\) diagonal edges remain distinct.
\end{itemize}

Therefore
\[
E = 3+6+3+6+9 = 27.
\]

So for both the torus and the Klein bottle we obtain
\[
V=9,\qquad E=27,\qquad F=18.
\]
Hence
\[
\chi = V-E+F = 9-27+18 = 0.
\]

Thus:

\[
\boxed{\chi(T)=0}, \qquad \boxed{\chi(K)=0}.
\]

\par

\textbf{2. The projective plane}

Now consider the triangulation shown for \(\mathbb{RP}^2\).

We first count the triangles in the square. From the drawing, the square is divided into \(10\) triangles, so
\[
F=10.
\]

\par

\textbf{Counting vertices before quotient.}

In the square picture there are \(9\) vertices:

\begin{itemize}
	\item the \(4\) corner vertices,
	\item \(2\) additional boundary vertices: one on the left side and one on the right side,
	\item \(3\) interior vertices.
\end{itemize}

Under the projective-plane identifications:
\begin{itemize}
	\item the four corners split into two equivalence classes:
	\[
	\{\text{top-left},\text{bottom-right}\},\qquad
	\{\text{top-right},\text{bottom-left}\},
	\]
	\item the two extra boundary vertices are identified with each other,
	\item the \(3\) interior vertices remain distinct.
\end{itemize}

Therefore the number of vertices in the quotient is
\[
V = 2 + 1 + 3 = 6.
\]

\par

\textbf{Counting edges before quotient.}

In the square triangulation there are \(18\) edges altogether. This can be checked directly, or by using Euler's formula for the disk:
\[
V-E+F=1.
\]
Since \(V=9\) and \(F=10\), we get
\[
9-E+10=1,
\]
so
\[
E=18.
\]

Among these \(18\) edges, \(6\) lie on the boundary of the square, and the remaining \(12\) are interior.

Under the projective-plane identifications, the \(6\) boundary edges are identified in pairs, so they become \(3\) quotient edges. The \(12\) interior edges remain distinct.

Hence
\[
E = 12+3=15.
\]

Thus for the projective plane:
\[
V=6,\qquad E=15,\qquad F=10.
\]
Therefore
\[
\chi = V-E+F = 6-15+10 = 1.
\]

So
\[
\boxed{\chi(\mathbb{RP}^2)=1}.
\]


\medskip
\noindent\textbf{Final counts}\par
\smallskip


\[
\begin{array}{c|c|c|c|c}
	\text{Surface} & V & E & F & \chi=V-E+F\\
	\hline
	T & 9 & 27 & 18 & 0\\
	K & 9 & 27 & 18 & 0\\
	\mathbb{RP}^2 & 6 & 15 & 10 & 1
\end{array}
\]


\medskip
\noindent\textbf{Conclusion}\par
\smallskip


For the triangulations shown in the figure, we have:

\[
\boxed{
	\begin{aligned}
		T &: (V,E,F)=(9,27,18),\qquad \chi(T)=0,\\
		K &: (V,E,F)=(9,27,18),\qquad \chi(K)=0,\\
		\mathbb{RP}^2 &: (V,E,F)=(6,15,10),\qquad \chi(\mathbb{RP}^2)=1.
\end{aligned}}
\]

These values agree with the known Euler characteristics of the torus, Klein bottle, and projective plane.
\end{solution}

\begin{problem}[CP-VIII-0022 --- Free subgroups of \(F_2\)]
Show that the free group \(F_2\) contains subgroups isomorphic to the free group \(F_n\) for any \(n>1\).

\par
\end{problem}

\begin{solution}
Write
\[
F_2=\langle x,y\rangle.
\]
Fix an integer \(n>1\). Define a homomorphism
\[
\psi:F_2\longrightarrow \mathbb Z/(n-1)\mathbb Z
\]
by
\[
\psi(x)=1,\qquad \psi(y)=0.
\]
This is surjective, so its kernel
\[
H=\ker(\psi)
\]
is a subgroup of index \(n-1\) in \(F_2\).

By the Nielsen--Schreier theorem, every subgroup of a free group is itself free. Moreover, if a subgroup \(H\) has index \(d\) in a free group of rank \(r\), then
\[
\operatorname{rank}(H)=1+d(r-1).
\]
Here \(r=2\) and \(d=n-1\), so
\[
\operatorname{rank}(H)=1+(n-1)(2-1)=1+(n-1)=n.
\]
Hence \(H\) is a free group of rank \(n\). Therefore
\[
H\cong F_n.
\]

So \(F_2\) contains a subgroup isomorphic to \(F_n\) for every \(n>1\).

\par
\end{solution}

\begin{problem}[CP-VIII-0023 --- The wedge \(\mathbb{RP}^2\vee \mathbb{RP}^2\)]
Let \(Y=\mathbb{RP}^2\vee \mathbb{RP}^2\) and let \(*\) be the wedge point.

\begin{enumerate}
	\item[(i)] Show that
	\[
	\pi_1(Y,*)\cong \mathbb Z/2 * \mathbb Z/2 \cong \langle a,b \mid a^2,b^2\rangle.
	\]

	\item[(ii)] Describe the covering space of \(Y\) corresponding to the kernel of the homomorphism
	\[
	\phi:\langle a,b \mid a^2,b^2\rangle \to \mathbb Z/2
	\]
	given by \(\phi(a)=1\) and \(\phi(b)=0\). Hence show that \(\ker(\phi)\) is isomorphic to \(\langle a,b \mid a^2,b^2\rangle\).

	\item[(iii)] Draw a picture of the universal cover \(\widetilde{Y}\). Deduce that \(ab\) has infinite order in \(\langle a,b \mid a^2,b^2\rangle\).
\end{enumerate}

\par
\end{problem}

\begin{solution}
Since \(Y=\mathbb{RP}^2\vee \mathbb{RP}^2\) is a wedge of two path-connected spaces, van Kampen's theorem gives
\[
\pi_1(Y,*)\cong \pi_1(\mathbb{RP}^2)*\pi_1(\mathbb{RP}^2).
\]
Now
\[
\pi_1(\mathbb{RP}^2)\cong \mathbb Z/2,
\]
so
\[
\pi_1(Y,*)\cong \mathbb Z/2 * \mathbb Z/2.
\]
If \(a\) and \(b\) denote the nontrivial loops in the two copies of \(\mathbb{RP}^2\), then this free product has presentation
\[
\pi_1(Y,*)\cong \langle a,b \mid a^2=1,\; b^2=1\rangle.
\]
Equivalently,
\[
\pi_1(Y,*)\cong \langle a,b \mid a^2,b^2\rangle.
\]

\par


Let
\[
G=\langle a,b \mid a^2,b^2\rangle \cong \mathbb Z/2 * \mathbb Z/2,
\]
and consider the surjective homomorphism
\[
\phi:G\to \mathbb Z/2,\qquad \phi(a)=1,\quad \phi(b)=0.
\]
Its kernel \(H=\ker(\phi)\) has index \(2\), so the corresponding covering space is a connected \(2\)-sheeted cover of \(Y\).

We now describe this cover geometrically. Over the first copy of \(\mathbb{RP}^2\), corresponding to the generator \(a\), the map \(\phi(a)=1\) means that this summand is lifted by its nontrivial double cover
\[
S^2 \longrightarrow \mathbb{RP}^2.
\]
Let \(p,q\in S^2\) be the two points lying over the wedge point \(*\in \mathbb{RP}^2\).

Over the second copy of \(\mathbb{RP}^2\), corresponding to the generator \(b\), we have \(\phi(b)=0\). This means that this summand does not unwrap; instead, there are two separate copies of \(\mathbb{RP}^2\), one attached at \(p\) and one attached at \(q\).

Thus the covering space \(X\) corresponding to \(\ker(\phi)\) is obtained by taking:
\begin{itemize}
	\item one sphere \(S^2\),
	\item one copy of \(\mathbb{RP}^2\) attached to \(S^2\) at \(p\),
	\item another copy of \(\mathbb{RP}^2\) attached to \(S^2\) at \(q\).
\end{itemize}

The covering map \(X\to Y\) is:
\begin{itemize}
	\item the standard double cover \(S^2\to \mathbb{RP}^2\) on the first summand,
	\item a homeomorphism from each attached copy of \(\mathbb{RP}^2\) onto the second summand.
\end{itemize}

Now compute the fundamental group of \(X\). Since \(S^2\) is simply connected, it contributes nothing to the fundamental group. So
\[
\pi_1(X)\cong \pi_1(\mathbb{RP}^2)*\pi_1(\mathbb{RP}^2)
\cong \mathbb Z/2 * \mathbb Z/2
\cong \langle u,v \mid u^2,v^2\rangle.
\]
But \(\pi_1(X)\) is precisely the subgroup corresponding to the cover, namely \(\ker(\phi)\). Hence
\[
\ker(\phi)\cong \mathbb Z/2 * \mathbb Z/2
\cong \langle a,b \mid a^2,b^2\rangle.
\]

So \(\ker(\phi)\) is isomorphic to the same abstract group again.

\par


The universal cover \(\widetilde{Y}\) may be understood from the fact that
\[
\pi_1(Y)\cong \mathbb Z/2 * \mathbb Z/2.
\]
The Bass--Serre tree of this free product is an infinite line. Over each vertex sits a copy of the universal cover of \(\mathbb{RP}^2\), namely \(S^2\). Therefore \(\widetilde{Y}\) looks like an infinite chain of \(2\)-spheres, each touching the next one in a single point:
\[
\cdots - S^2 - S^2 - S^2 - S^2 - \cdots
\]
where adjacent spheres meet at one point, and the spheres alternate according to the two \(\mathbb{RP}^2\)-summands.

The generators \(a\) and \(b\) each act as involutions, while the product \(ab\) acts as a translation along this infinite chain. A nonzero power \((ab)^m\) therefore moves a chosen lift of the basepoint to a different point of the universal cover. Hence
\[
(ab)^m\neq 1 \qquad \text{for all } m>0.
\]
So \(ab\) has infinite order in
\[
\langle a,b \mid a^2,b^2\rangle.
\]

\par
\end{solution}

\begin{problem}[CP-VIII-0024 --- A cell structure for the Klein bottle]
Show that the Klein bottle has a cell structure with a single \(0\)-cell, two \(1\)-cells, and a single \(2\)-cell. Deduce that its fundamental group has a presentation
\[
\langle a,b \mid baba^{-1}\rangle,
\]
and show this is isomorphic to the group \(G\) in Q13 of Example Sheet 1.

\par
\end{problem}

\begin{solution}
Take the usual square model of the Klein bottle. In this model the boundary edges are identified in pairs, with one pair identified in the same direction and the other pair identified in opposite directions.

Because all four vertices of the square are identified together, there is exactly one \(0\)-cell.

The two edge-identification classes give two \(1\)-cells, which we denote by \(a\) and \(b\).

The interior of the square gives one \(2\)-cell.

Therefore the Klein bottle has a CW-structure consisting of:
\[
1 \text{ zero-cell},\qquad 2 \text{ one-cells},\qquad 1 \text{ two-cell}.
\]

Now compute the fundamental group from this cell structure. The \(1\)-skeleton is a wedge of two circles, so its fundamental group is the free group on two generators:
\[
\pi_1(K^{(1)})\cong \langle a,b\rangle.
\]
The attaching map of the \(2\)-cell is read off from the boundary of the square. This gives the relation
\[
baba^{-1}=1,
\]
up to cyclic permutation of the boundary word. Hence van Kampen gives
\[
\pi_1(K)\cong \langle a,b \mid baba^{-1}\rangle.
\]

This presentation is equivalent to the standard one for the Klein bottle group. Indeed,
\[
baba^{-1}=1
\]
implies
\[
bab=a,
\]
and rearranging gives the familiar relation
\[
aba^{-1}=b^{-1}.
\]
Equivalently,
\[
ab=b^{-1}a.
\]

So we may also write
\[
\pi_1(K)\cong \langle a,b \mid aba^{-1}=b^{-1}\rangle.
\]

This shows that the Klein bottle group is the semidirect product
\[
\mathbb Z \rtimes \mathbb Z,
\]
where one copy of \(\mathbb Z\) acts on the other by inversion. Explicitly,
\[
\pi_1(K)\cong \mathbb Z \rtimes_{(-1)} \mathbb Z.
\]

If the group \(G\) in Q13 of Example Sheet 1 is given in this semidirect-product form, then the isomorphism is obtained by sending
\[
a\mapsto (0,1),\qquad b\mapsto (1,0).
\]
Under this identification,
\[
aba^{-1}=b^{-1},
\]
which is exactly the Klein bottle relation. Hence
\[
\pi_1(K)\cong G.
\]
\end{solution}

\begin{problem}[CP-VIII-0025 --- Hopf link versus the two-component unlink in \(S^3\)]
Consider the following configurations of pairs of circles in \(S^3\) (we have drawn them in \(\mathbb{R}^3\); add a point at infinity). By computing the fundamental groups of the complements of the circles, show there is no homeomorphism of \(S^3\) taking one configuration to the other.

\par
\end{problem}

\begin{solution}
Let \(L_H\) denote the left-hand configuration and \(L_U\) the right-hand configuration.

The left-hand configuration is the \emph{Hopf link}: the two circles are linked once.
The right-hand configuration is the \emph{two-component unlink}: the two circles are disjoint and unlinked.

We compute the fundamental groups of their complements.


\medskip
\noindent\textbf{1. The complement of the unlink}\par
\smallskip

Let
\[
X_U = S^3 \setminus L_U.
\]
It is standard that removing the circles themselves or removing disjoint open tubular neighborhoods of the circles gives spaces with the same homotopy type. So we may work with the \emph{link exterior}
\[
E_U = S^3 \setminus \operatorname{int}(N(L_U)).
\]

Since \(L_U\) is the unlink of two components, \(N(L_U)\) is the disjoint union of two solid tori embedded in \(S^3\) in an unlinked way. The complement \(E_U\) is therefore a handlebody of genus \(2\). A genus \(2\) handlebody deformation retracts onto a wedge of two circles, so
\[
\pi_1(X_U)\cong \pi_1(E_U)\cong F_2,
\]
the free group on two generators.

Thus
\[
\pi_1(S^3\setminus L_U)\cong F_2.
\]


\medskip
\noindent\textbf{2. The complement of the Hopf link}\par
\smallskip

Now let
\[
X_H = S^3 \setminus L_H.
\]

Take one component \(K_1\) of the Hopf link. The complement of an unknot in \(S^3\) is a solid torus:
\[
S^3 \setminus K_1 \cong S^1 \times D^2.
\]
Inside this solid torus, the second component \(K_2\) of the Hopf link is isotopic to the core circle
\[
S^1 \times \{0\}.
\]
Hence
\[
X_H \cong (S^1\times D^2)\setminus (S^1\times \{0\})
\cong S^1 \times (D^2\setminus \{0\}).
\]
But
\[
D^2\setminus \{0\}\cong S^1\times (0,1),
\]
so
\[
X_H \cong S^1\times S^1\times (0,1)\cong T^2\times (0,1).
\]
Therefore
\[
\pi_1(X_H)\cong \pi_1(T^2)\cong \mathbb{Z}\oplus \mathbb{Z}.
\]

Thus
\[
\pi_1(S^3\setminus L_H)\cong \mathbb{Z}^2.
\]


\medskip
\noindent\textbf{3. Comparison}\par
\smallskip

We have shown:
\[
\pi_1(S^3\setminus L_H)\cong \mathbb{Z}^2,
\qquad
\pi_1(S^3\setminus L_U)\cong F_2.
\]
These groups are not isomorphic, since \(\mathbb{Z}^2\) is abelian while \(F_2\) is nonabelian.

If there were a homeomorphism of \(S^3\) taking one configuration to the other, it would induce a homeomorphism between the complements, hence an isomorphism of their fundamental groups. Since no such isomorphism exists, no such homeomorphism exists.

\par

Therefore the two configurations are not equivalent under any homeomorphism of \(S^3\).
\end{solution}

\begin{problem}[CP-VIII-0026 --- Trees, quotient graphs, coverings, and free groups]
A graph \(G\) is a space obtained by starting with a set \(E(G)\) of copies of the interval \(I\) and an equivalence relation \(\sim\) on \(E(G)\times \{0,1\}\), and forming the quotient space of \(E(G)\times I\) by the minimal equivalence relation containing \(\sim\). (That is, it is a space obtained from a set of copies of \(I\) by gluing their ends together.) The \emph{vertices} are the equivalence classes represented by the ends of the intervals.

\begin{enumerate}
	\item[(i)] A tree is a simply-connected graph. A star is a tree with a vertex \(x_0\) such that one end of each edge is attached to \(x_0\). A leaf of a tree is a vertex attached to only one edge. Prove that every tree is homotopy equivalent to a star, relative to its leaves.

	\item[(ii)] If \(T\subset G\) is a tree, show that the quotient map
	\(G\to G/T\) is a homotopy equivalence and that \(G/T\) is again a graph.
	Hence show that every connected graph is homotopy equivalent to a graph
	with a single vertex.

	\emph{You may assume that every graph has a maximal tree.}

	\item[(iii)] Show that the fundamental group of a graph with one vertex,
	based at that vertex, is a free group with one generator for each edge.
	Hence show that every free group occurs as the fundamental group of a graph.

	\emph{Graphs are not assumed to have finitely many edges.}

	\item[(iv)] Show that a covering space of a graph is again a graph, and deduce that a subgroup of a free group is again a free group.
\end{enumerate}

\par
\end{problem}

\begin{solution}

\medskip
\noindent\textbf{(i) Every tree is homotopy equivalent to a star, relative to its leaves}\par
\smallskip

Let \(T\) be a tree.

Consider the subgraph \(C\subset T\) consisting of all edges that are \emph{not} incident to a leaf, together with their endpoints.

If \(C=\varnothing\), then every edge of \(T\) is incident to a leaf. In that case \(T\) is already a star.

Assume \(C\neq \varnothing\). Then \(C\) is connected: if \(e_1\) and \(e_2\) are two edges of \(C\), the unique path in the tree joining them cannot pass through a leaf, since a leaf cannot lie in the interior of a path joining two other edges. Hence every edge on that path also lies in \(C\).

Thus \(C\) is a connected subtree, hence contractible.

Now collapse \(C\) to a point. Since \(C\) is contractible, the quotient map
\[
T \longrightarrow T/C
\]
is a homotopy equivalence. Moreover, each component attached to \(C\) outside \(C\) is just an edge ending in a leaf, so after collapsing \(C\) to one point, the quotient \(T/C\) is a star. The leaves remain distinct and are fixed during the collapse, so this is a homotopy equivalence \emph{relative to the leaves}.

Therefore every tree is homotopy equivalent to a star, relative to its leaves.


\medskip
\noindent\textbf{(ii) Collapsing a tree inside a graph}\par
\smallskip

Let \(T\subset G\) be a tree.

Since \(T\) is contractible, there is a homotopy in \(T\) from the identity map to a constant map \(T\to \{x_0\}\), where \(x_0\in T\). Because \(T\) is a subcomplex of the graph \(G\), this homotopy extends to a homotopy of all of \(G\). Hence collapsing \(T\) to a point does not change the homotopy type: the quotient map
\[
q:G\longrightarrow G/T
\]
is a homotopy equivalence.

Also, \(G/T\) is again a graph: all of the edges and vertices of \(T\) are collapsed to one vertex, while every edge of \(G\setminus T\) remains an interval whose endpoints are attached to quotient vertices.

Now let \(G\) be connected, and choose a maximal tree \(T\subset G\). Then \(G/T\) is a graph with exactly one vertex, since all vertices of \(G\) lie in the maximal tree \(T\). By the first part, \(G\) is homotopy equivalent to \(G/T\), which is a graph with a single vertex.

Therefore every connected graph is homotopy equivalent to a graph with one vertex.


\medskip
\noindent\textbf{(iii) Fundamental group of a one-vertex graph}\par
\smallskip

Let \(G\) be a graph with one vertex \(v\). Then every edge begins and ends at \(v\), so topologically
\[
G \cong \bigvee_{e\in E(G)} S^1_e,
\]
a wedge of circles, one for each edge.

If \(E(G)\) is finite, say \(E(G)=\{e_1,\dots,e_n\}\), then by van Kampen,
\[
\pi_1(G,v)\cong F_n,
\]
the free group on \(n\) generators, one generator for each loop \(e_i\).

If \(E(G)\) is infinite, any loop \(\alpha:S^1\to G\) has compact image, so it meets only finitely many open edges. Hence \(\alpha\) is contained in a finite subgraph
\[
G_F=\bigvee_{e\in F} S^1_e
\]
for some finite subset \(F\subset E(G)\). Therefore every element of \(\pi_1(G,v)\) lies in the image of \(\pi_1(G_F,v)\cong F(F)\), and passing to the direct limit over finite \(F\), we obtain
\[
\pi_1(G,v)\cong F(E(G)),
\]
the free group on the set of edges.

Thus the fundamental group of a graph with one vertex is a free group with one generator for each edge.

Conversely, given any free group \(F(S)\) on a set \(S\), take a graph with one vertex and one loop for each element of \(S\). Its fundamental group is exactly \(F(S)\).

Hence every free group occurs as the fundamental group of some graph.


\medskip
\noindent\textbf{(iv) Coverings of graphs and subgroups of free groups}\par
\smallskip

Let
\[
p:\widetilde{G}\longrightarrow G
\]
be a covering map, where \(G\) is a graph.

Each open edge \(e^\circ\subset G\) is homeomorphic to an open interval. Since coverings are local homeomorphisms, \(p^{-1}(e^\circ)\) is a disjoint union of open intervals. The preimage of the vertex set consists of points, and a neighborhood of each such point is mapped homeomorphically onto a star-neighborhood of a vertex in \(G\). Therefore \(\widetilde{G}\) is again obtained by gluing intervals along endpoints; that is, \(\widetilde{G}\) is again a graph.

Now let \(F\) be a free group, and let \(H\leq F\) be a subgroup. By part (iii), choose a connected graph \(G\) with
\[
\pi_1(G,v)\cong F.
\]
By covering space theory, the subgroup \(H\) corresponds to a connected covering space
\[
p:(\widetilde{G},\widetilde{v})\to (G,v)
\]
such that
\[
p_*\bigl(\pi_1(\widetilde{G},\widetilde{v})\bigr)=H.
\]
Since \(\widetilde{G}\) is again a graph, parts (ii) and (iii) imply that
\[
\pi_1(\widetilde{G},\widetilde{v})
\]
is a free group. Therefore \(H\) is free.

This proves that every subgroup of a free group is again a free group.
\end{solution}

\section{Homology}

\begin{problem}[CP-VIII-0027 --- Abstract simplicial complexes, geometric realizations, and realization in Euclidean space]
An abstract simplicial complex consists of a finite set \(V_X\) (called the vertices) and a collection \(X\) (called the simplices) of subsets of \(V_X\) such that if \(\sigma\in X\) and \(\tau\subseteq \sigma\), then \(\tau\in X\). A map
\[
f:(V_X,X)\to (V_Y,Y)
\]
of abstract simplicial complexes is a function
\[
f:V_X\to V_Y
\]
such that \(f(\sigma)\in Y\) for all \(\sigma\in X\).

\begin{enumerate}
	\item[(i)] For a simplicial complex \(K\) in \(\mathbb R^m\), show that the abstraction of \(K\),
	\[
	V_X=\{0\text{-simplices of }K\},
	\qquad
	X=\bigl\{\{a_0,\dots,a_n\}\subset V_X \mid \langle a_0,\dots,a_n\rangle\in K\bigr\},
	\]
	is an abstract simplicial complex. Show that any abstract simplicial complex arises in this way (up to isomorphism), and that if simplicial complexes \(K\) and \(L\) have isomorphic abstractions, then \(|K|\) and \(|L|\) are homeomorphic.

	\item[(ii)] Show that there exists an infinite sequence of points \((x_1,x_2,\dots)\in \mathbb R^m\) such that any \((m+1)\) of them are affinely independent. Hence show that if \((V_X,X)\) is an abstract simplicial complex with all simplices of dimension \(\le n\), then there is a simplicial complex \(K\) in \(\mathbb R^{2n+1}\) with abstraction isomorphic to \((V_X,X)\).
\end{enumerate}
\end{problem}

\begin{solution}
\textbf{Part (i): the abstraction of a geometric simplicial complex.}

Let \(K\) be a simplicial complex in \(\mathbb R^m\). Define
\[
V_X=\{\text{vertices of }K\},
\]
and
\[
X=\bigl\{\{a_0,\dots,a_n\}\subseteq V_X \mid \langle a_0,\dots,a_n\rangle\in K\bigr\}.
\]

We must show that \((V_X,X)\) is an abstract simplicial complex.

Take any simplex
\[
\sigma=\{a_0,\dots,a_n\}\in X.
\]
By definition, the geometric simplex
\[
\langle a_0,\dots,a_n\rangle
\]
belongs to \(K\). Now let \(\tau\subseteq \sigma\). Then \(\tau\) is obtained by choosing some subset of the vertices \(a_0,\dots,a_n\), so the geometric simplex spanned by \(\tau\) is a face of \(\langle a_0,\dots,a_n\rangle\). Since every face of a simplex in \(K\) again belongs to \(K\), it follows that
\[
\langle \tau\rangle\in K.
\]
Hence \(\tau\in X\). Therefore \(X\) is closed under taking subsets, so \((V_X,X)\) is indeed an abstract simplicial complex.

\par

\textbf{Every abstract simplicial complex arises this way.}

Now let \((V_X,X)\) be any abstract simplicial complex, where
\[
V_X=\{v_1,\dots,v_r\}.
\]
We construct a geometric simplicial complex whose abstraction is isomorphic to \((V_X,X)\).

Place the vertices at the standard basis vectors
\[
e_1,\dots,e_r\in \mathbb R^r.
\]
For each simplex
\[
\sigma=\{v_{i_0},\dots,v_{i_k}\}\in X,
\]
consider the geometric simplex
\[
\langle e_{i_0},\dots,e_{i_k}\rangle.
\]
Let \(K\) be the collection of all such simplices.

We claim that \(K\) is a simplicial complex.

First, if
\[
\langle e_{i_0},\dots,e_{i_k}\rangle\in K,
\]
then any face is of the form
\[
\langle e_{j_0},\dots,e_{j_\ell}\rangle
\]
for some subset \(\{v_{j_0},\dots,v_{j_\ell}\}\subseteq \{v_{i_0},\dots,v_{i_k}\}\). Since \(X\) is an abstract simplicial complex, every subset of a simplex is again a simplex, so
\[
\{v_{j_0},\dots,v_{j_\ell}\}\in X.
\]
Hence this face also belongs to \(K\).

Second, let
\[
\sigma=\{v_{i_0},\dots,v_{i_k}\},\qquad \tau=\{v_{j_0},\dots,v_{j_\ell}\}
\]
be simplices of \(X\). Then
\[
\langle e_{i_0},\dots,e_{i_k}\rangle\cap \langle e_{j_0},\dots,e_{j_\ell}\rangle
=
\langle e_s : v_s\in \sigma\cap\tau\rangle,
\]
which is a face of both simplices. Thus \(K\) satisfies the intersection condition for a simplicial complex.

So \(K\) is a simplicial complex in \(\mathbb R^r\), and by construction its abstraction is exactly \((V_X,X)\), up to the obvious identification \(v_i\leftrightarrow e_i\).

Therefore every abstract simplicial complex arises as the abstraction of a geometric simplicial complex, up to isomorphism.

\par

\textbf{Isomorphic abstractions imply homeomorphic realizations.}

Suppose simplicial complexes \(K\) and \(L\) have isomorphic abstractions. Then there is a bijection
\[
\phi:V(K)\to V(L)
\]
such that a finite set of vertices spans a simplex of \(K\) if and only if its image under \(\phi\) spans a simplex of \(L\).

We define a map
\[
|\phi|:|K|\to |L|
\]
simplexwise by barycentric coordinates. If
\[
x=\sum_{i=0}^k t_i a_i\in \langle a_0,\dots,a_k\rangle\subseteq |K|,
\qquad t_i\ge 0,\quad \sum_{i=0}^k t_i=1,
\]
then set
\[
|\phi|(x)=\sum_{i=0}^k t_i\,\phi(a_i)\in \langle \phi(a_0),\dots,\phi(a_k)\rangle\subseteq |L|.
\]

This is well-defined, because \(\phi\) sends simplices of \(K\) to simplices of \(L\), and on intersections of simplices the barycentric coordinate descriptions agree. On each simplex, \(|\phi|\) is an affine homeomorphism. Likewise, \(\phi^{-1}\) induces a map
\[
|\phi^{-1}|:|L|\to |K|
\]
which is inverse to \(|\phi|\).

Hence \(|\phi|:|K|\to |L|\) is a homeomorphism.

Therefore, if two simplicial complexes have isomorphic abstractions, then their realizations are homeomorphic:
\[
|K|\cong |L|.
\]

\par

\textbf{Part (ii): infinite affine-independent sequences and realization in \(\mathbb R^{2n+1}\).}

We first prove the existence of an infinite sequence of points in \(\mathbb R^m\) such that any \(m+1\) of them are affinely independent.

Consider the \emph{moment curve}
\[
\gamma(t)=(t,t^2,\dots,t^m)\in \mathbb R^m.
\]
Choose distinct real numbers
\[
t_1,t_2,t_3,\dots
\]
and define
\[
x_i=\gamma(t_i)=(t_i,t_i^2,\dots,t_i^m)\in \mathbb R^m.
\]

We claim that any \(m+1\) of these points are affinely independent.

Indeed, points \(x_{i_0},\dots,x_{i_m}\in \mathbb R^m\) are affinely independent if and only if the vectors
\[
(x_{i_0},1),\dots,(x_{i_m},1)\in \mathbb R^{m+1}
\]
are linearly independent. But these vectors form the rows of the matrix
\[
\begin{pmatrix}
	1 & t_{i_0} & t_{i_0}^2 & \cdots & t_{i_0}^m\\
	1 & t_{i_1} & t_{i_1}^2 & \cdots & t_{i_1}^m\\
	\vdots & \vdots & \vdots & & \vdots\\
	1 & t_{i_m} & t_{i_m}^2 & \cdots & t_{i_m}^m
\end{pmatrix}.
\]
This is a Vandermonde matrix, whose determinant is
\[
\prod_{p<q}(t_{i_q}-t_{i_p}).
\]
Since the \(t_{i_j}\) are distinct, this determinant is nonzero. Therefore the vectors are linearly independent, and the points \(x_{i_0},\dots,x_{i_m}\) are affinely independent.

Thus there exists an infinite sequence of points in \(\mathbb R^m\) such that any \(m+1\) of them are affinely independent.

\par

\textbf{Realization in \(\mathbb R^{2n+1}\).}

Now let \((V_X,X)\) be an abstract simplicial complex such that every simplex has dimension at most \(n\). Equivalently, every simplex has at most \(n+1\) vertices.

Write
\[
V_X=\{v_1,\dots,v_r\}.
\]
Choose points
\[
x_1,\dots,x_r\in \mathbb R^{2n+1}
\]
such that any \(2n+2\) of them are affinely independent. This is possible by the first part, applied with \(m=2n+1\), for example by choosing the \(x_i\) on the moment curve in \(\mathbb R^{2n+1}\).

For each simplex
\[
\sigma=\{v_{i_0},\dots,v_{i_k}\}\in X,
\qquad k\le n,
\]
let
\[
\langle x_{i_0},\dots,x_{i_k}\rangle
\]
be the geometric simplex spanned by the corresponding points. Let \(K\) be the collection of all such simplices.

We show that \(K\) is a simplicial complex in \(\mathbb R^{2n+1}\).

First, every face of a simplex of \(K\) again belongs to \(K\), because if
\[
\{v_{i_0},\dots,v_{i_k}\}\in X,
\]
then every subset is again in \(X\), since \(X\) is an abstract simplicial complex.

So only the intersection condition remains.

\par

\textbf{Lemma.}
If \(A\) and \(B\) are finite sets of points in \(\mathbb R^N\), and \(A\cup B\) is affinely independent, then
\[
\operatorname{conv}(A)\cap \operatorname{conv}(B)=\operatorname{conv}(A\cap B).
\]

\textbf{Proof.}
Certainly
\[
\operatorname{conv}(A\cap B)\subseteq \operatorname{conv}(A)\cap \operatorname{conv}(B).
\]
For the reverse inclusion, let
\[
x\in \operatorname{conv}(A)\cap \operatorname{conv}(B).
\]
Then
\[
x=\sum_{a\in A}\lambda_a a=\sum_{b\in B}\mu_b b,
\qquad
\lambda_a,\mu_b\ge 0,
\qquad
\sum_{a\in A}\lambda_a=\sum_{b\in B}\mu_b=1.
\]
Subtracting gives
\[
\sum_{a\in A\setminus B}\lambda_a a
+\sum_{c\in A\cap B}(\lambda_c-\mu_c)c
-\sum_{b\in B\setminus A}\mu_b b=0.
\]
Also the sum of the coefficients here is
\[
\sum_{a\in A\setminus B}\lambda_a
+\sum_{c\in A\cap B}(\lambda_c-\mu_c)
-\sum_{b\in B\setminus A}\mu_b
=
1-1=0.
\]
So this is an affine dependence relation among the points of \(A\cup B\). Since \(A\cup B\) is affinely independent, all coefficients must be zero. Hence
\[
\lambda_a=0 \text{ for } a\in A\setminus B,
\qquad
\mu_b=0 \text{ for } b\in B\setminus A.
\]
Therefore
\[
x\in \operatorname{conv}(A\cap B).
\]
So indeed
\[
\operatorname{conv}(A)\cap \operatorname{conv}(B)=\operatorname{conv}(A\cap B).
\]
This proves the lemma.

\par

Now apply the lemma to two simplices of \(K\), say corresponding to
\[
\sigma,\tau\in X.
\]
Since each simplex has at most \(n+1\) vertices,
\[
|\sigma\cup\tau|\le (n+1)+(n+1)=2n+2.
\]
By construction, any \(2n+2\) of the chosen points are affinely independent, so the set of vertices corresponding to \(\sigma\cup\tau\) is affinely independent. Therefore the lemma gives
\[
\langle \sigma\rangle \cap \langle \tau\rangle = \langle \sigma\cap\tau\rangle,
\]
that is, the intersection of any two simplices of \(K\) is a common face.

Hence \(K\) is a simplicial complex in \(\mathbb R^{2n+1}\).

Finally, by construction, the abstraction of \(K\) is isomorphic to \((V_X,X)\): the vertices of \(K\) are in bijection with the vertices \(v_1,\dots,v_r\), and a subset of vertices spans a simplex of \(K\) exactly when the corresponding subset belongs to \(X\).

Thus there exists a simplicial complex
\[
K\subseteq \mathbb R^{2n+1}
\]
whose abstraction is isomorphic to \((V_X,X)\).


\medskip
\noindent\textbf{Conclusion}\par
\smallskip


We have proved the following:

\begin{itemize}
	\item The abstraction of a geometric simplicial complex is an abstract simplicial complex.
	\item Every abstract simplicial complex can be realized geometrically, up to isomorphism.
	\item If two simplicial complexes have isomorphic abstractions, then their realizations are homeomorphic.
	\item If all simplices have dimension at most \(n\), then the abstract simplicial complex can be realized by a simplicial complex in \(\mathbb R^{2n+1}\).
\end{itemize}

So abstract simplicial complexes and geometric simplicial complexes encode the same combinatorial data, and bounded dimension allows realization in a controlled Euclidean dimension.
\end{solution}

\begin{problem}[CP-VIII-0028 --- Simplicial approximation: countability of homotopy classes, maps between spheres, and the \(2\)-skeleton]
Use the simplicial approximation theorem to show that

\begin{enumerate}
	\item[(i)] if \(K\) and \(L\) are simplicial complexes, there are at most countably many homotopy classes of continuous maps
	\[
	f:|K|\to |L|;
	\]

	\item[(ii)] if \(m<n\), then any continuous map
	\[
	S^m\to S^n
	\]
	is homotopic to a constant map;

	\item[(iii)] for any vertex \(v\) of a simplicial complex \(K\), the based map
	\[
	(|K^{(2)}|,v)\to (|K|,v)
	\]
	(i.e. the inclusion of the \(2\)-skeleton) induces an isomorphism on fundamental groups. Show how to prove the same result using the Seifert--van Kampen theorem.
\end{enumerate}
\end{problem}

\begin{solution}
\textbf{(i) There are at most countably many homotopy classes of maps \( |K|\to |L| \).}

Let
\[
f:|K|\to |L|
\]
be continuous. By the simplicial approximation theorem, there exists some integer \(r\ge 0\) such that after \(r\) barycentric subdivisions of \(K\), the map \(f\) is homotopic to the realization of a simplicial map
\[
g:\operatorname{sd}^r K \to L.
\]
Hence every homotopy class of continuous maps contains a simplicial representative from some subdivision \(\operatorname{sd}^r K\).

Now fix \(r\). A simplicial map
\[
g:\operatorname{sd}^r K \to L
\]
is completely determined by its effect on vertices, subject to the simplicial condition that the image of each simplex is again a simplex. If \(K\) and \(L\) are finite, then \(\operatorname{sd}^r K\) and \(L\) have only finitely many vertices, so for fixed \(r\) there are only finitely many such simplicial maps. More generally, if the complexes are countable, then for fixed \(r\) there are at most countably many simplicial maps.

Therefore, for each fixed \(r\), the set of simplicial maps
\[
\operatorname{sd}^r K \to L
\]
is at most countable. Since \(r\) ranges over the countable set \(\{0,1,2,\dots\}\), the total set of simplicial maps from all subdivisions
\[
\bigcup_{r\ge 0}\operatorname{Map}_{\mathrm{simp}}(\operatorname{sd}^r K,L)
\]
is at most countable.

Since every homotopy class of continuous maps \(|K|\to |L|\) contains one of these simplicial representatives, it follows that there are at most countably many homotopy classes of continuous maps
\[
|K|\to |L|.
\]

\par

\textbf{Conclusion for (i).}
The set
\[
[\,|K|,|L|\,]
\]
of homotopy classes of continuous maps is at most countable.

\par

\textbf{(ii) If \(m<n\), every map \(S^m\to S^n\) is null-homotopic.}

Let
\[
f:S^m\to S^n
\]
be continuous. Choose triangulations
\[
|K|\cong S^m,
\qquad
|L|\cong S^n.
\]
By simplicial approximation, after sufficiently many barycentric subdivisions of \(K\), the map \(f\) is homotopic to the realization of a simplicial map
\[
g:\operatorname{sd}^r K \to L.
\]

Now \(\operatorname{sd}^r K\) is still an \(m\)-dimensional simplicial complex. Therefore the image of each simplex of \(\operatorname{sd}^r K\) under \(g\) is a simplex of \(L\) of dimension at most \(m\). Hence the image of \(g\) is contained in the \(m\)-skeleton
\[
L^{(m)} \subseteq L.
\]

Since \(m<n\), the \(m\)-skeleton \(L^{(m)}\) is a proper subcomplex of \(L\), so it cannot fill all of \(|L|=S^n\). Therefore there exists some point
\[
y\in S^n
\]
which does not belong to the image of \(g\). Thus
\[
g:S^m \to S^n\setminus \{y\}.
\]

But
\[
S^n\setminus \{y\}\cong \mathbb R^n,
\]
and \(\mathbb R^n\) is contractible. Therefore \(g\) is homotopic to a constant map. Since \(f\simeq g\), it follows that \(f\) is also homotopic to a constant map.

\par

\textbf{Conclusion for (ii).}
If \(m<n\), every continuous map
\[
S^m\to S^n
\]
is null-homotopic.

\par

\textbf{(iii) The inclusion \( |K^{(2)}| \hookrightarrow |K| \) induces an isomorphism on fundamental groups.}

Let
\[
i:(|K^{(2)}|,v)\hookrightarrow (|K|,v)
\]
be the inclusion.

We show that
\[
i_*:\pi_1(|K^{(2)}|,v)\to \pi_1(|K|,v)
\]
is an isomorphism.

\par

\textbf{Surjectivity.}
Let
\[
\alpha:(S^1,*)\to (|K|,v)
\]
be a based loop. Triangulate \(S^1\) as a \(1\)-dimensional simplicial complex, with the basepoint as a vertex. By simplicial approximation, after subdividing \(S^1\) sufficiently many times, the loop \(\alpha\) is homotopic through based maps to the realization of a simplicial map
\[
\beta:\operatorname{sd}^r(S^1)\to K.
\]
Since the domain is \(1\)-dimensional, the image of \(\beta\) lies in the \(1\)-skeleton \(K^{(1)}\), hence certainly in \(K^{(2)}\). Thus every loop in \(|K|\) is homotopic to a loop in \(|K^{(2)}|\). Therefore
\[
i_*:\pi_1(|K^{(2)}|,v)\to \pi_1(|K|,v)
\]
is surjective.

\par

\textbf{Injectivity.}
Now suppose two based loops in \(|K^{(2)}|\) are homotopic in \(|K|\). Then there exists a based homotopy
\[
H:S^1\times I \to |K|
\]
between them. Triangulate the cylinder \(S^1\times I\) as a \(2\)-dimensional simplicial complex, with boundary a subcomplex. By the relative form of simplicial approximation, after subdividing sufficiently many times, the homotopy \(H\) is homotopic relative to the boundary to the realization of a simplicial map
\[
H':\operatorname{sd}^r(S^1\times I)\to K
\]
which agrees with the given loops on the boundary.

Since the domain is \(2\)-dimensional, the image of \(H'\) lies in the \(2\)-skeleton \(K^{(2)}\). Thus the original two loops are homotopic already inside \(|K^{(2)}|\). Therefore
\[
i_*:\pi_1(|K^{(2)}|,v)\to \pi_1(|K|,v)
\]
is injective.

Combining injectivity and surjectivity, we conclude that
\[
i_*:\pi_1(|K^{(2)}|,v)\xrightarrow{\cong}\pi_1(|K|,v)
\]
is an isomorphism.

\par

\textbf{The same result via the Seifert--van Kampen theorem.}

There is also a more geometric proof. Start with the \(2\)-skeleton \(K^{(2)}\), and obtain \(K\) by attaching simplices of dimensions \(3,4,\dots\).

Suppose \(\tau\) is a simplex of dimension \(q\ge 3\). Topologically, \(\tau\) is a \(q\)-disk \(D^q\), and it is attached along its boundary
\[
\partial \tau \cong S^{q-1}.
\]
Since \(q-1\ge 2\), the sphere \(S^{q-1}\) is simply connected. Also \(D^q\) is contractible, hence simply connected. When one attaches \(\tau\) to the space already constructed, the Seifert--van Kampen theorem shows that the fundamental group does not change, because one is gluing along a path-connected, simply connected subspace.

More explicitly, if
\[
X = Y \cup D^q
\]
with
\[
Y\cap D^q = S^{q-1},
\]
then van Kampen gives
\[
\pi_1(X)\cong \pi_1(Y) *_{\pi_1(S^{q-1})} \pi_1(D^q).
\]
But
\[
\pi_1(S^{q-1})=0
\qquad\text{and}\qquad
\pi_1(D^q)=0,
\]
so
\[
\pi_1(X)\cong \pi_1(Y).
\]
Thus attaching simplices of dimension at least \(3\) does not change the fundamental group.

Applying this repeatedly, we conclude that passing from \(K^{(2)}\) to \(K\) leaves the fundamental group unchanged:
\[
\pi_1(|K^{(2)}|,v)\cong \pi_1(|K|,v).
\]


\medskip
\noindent\textbf{Conclusion}\par
\smallskip


We have proved:

\begin{itemize}
	\item there are at most countably many homotopy classes of maps \(|K|\to |L|\);
	\item if \(m<n\), every map \(S^m\to S^n\) is null-homotopic;
	\item the inclusion
	\[
	|K^{(2)}|\hookrightarrow |K|
	\]
	induces an isomorphism on fundamental groups.
\end{itemize}

These are typical and important consequences of the simplicial approximation theorem: low-dimensional homotopy information is controlled by low-dimensional skeleta, and continuous maps may be studied combinatorially after subdivision.

\par
\end{solution}

\begin{problem}[CP-VIII-0029 --- Elementary collapse of a simplex pair and invariance of homology]
Let \(K\) be a simplicial complex, and suppose that \(\sigma\in K\) is not a proper face of any simplex, and that \(\tau\le \sigma\) is a face of one dimension lower which is not a face of any other simplex. Show that
\[
L=K\setminus \{\sigma,\tau\}
\]
is again a simplicial complex, and that the inclusion
\[
V_L\to V_K
\]
defines a simplicial map
\[
i:L\to K.
\]

\begin{enumerate}
	\item[(i)] By constructing a chain homotopy inverse to
	\[
	i_\#:C_\bullet(L)\to C_\bullet(K),
	\]
	show that
	\[
	i_*:H_j(L)\to H_j(K)
	\]
	is an isomorphism for all \(j\).

	\item[(ii)] Prove the same thing using the previous question twice instead.
\end{enumerate}
\end{problem}

\begin{solution}
Write
\[
\dim \sigma = n,
\qquad
\dim \tau = n-1.
\]
Thus \(\tau\) is a codimension-one face of \(\sigma\), \(\sigma\) is maximal, and \(\tau\) is a face of no simplex other than \(\sigma\).

\par

\textbf{Step 1: \(L=K\setminus\{\sigma,\tau\}\) is again a simplicial complex.}

We must show that \(L\) is closed under taking faces.

Let \(\eta\in L\), and let \(\rho\subseteq \eta\) be a face. Since \(K\) is a simplicial complex, \(\rho\in K\). We must show that \(\rho\neq \sigma\) and \(\rho\neq \tau\).

First, \(\rho\neq \sigma\), since \(\rho\subseteq \eta\) and \(\eta\neq \sigma\), while \(\sigma\) cannot be a face of a proper sub-simplex.

Next, suppose \(\rho=\tau\). Then \(\tau\subseteq \eta\). Since \(\eta\neq \tau\) and \(\tau\) is assumed to be a face of no simplex other than \(\sigma\), it would follow that \(\eta=\sigma\), contradicting \(\eta\in L\).

Thus \(\rho\notin \{\sigma,\tau\}\), so \(\rho\in L\). Therefore \(L\) is again a simplicial complex.

\par

\textbf{Step 2: the inclusion \(V_L\to V_K\) defines a simplicial map \(i:L\to K\).}

The vertex set of \(L\) is contained in the vertex set of \(K\), so there is a natural inclusion
\[
V_L\hookrightarrow V_K.
\]
If \(\eta\in L\), then \(\eta\) is certainly also a simplex of \(K\). Hence the image of every simplex of \(L\) is a simplex of \(K\). So the inclusion of vertex sets defines a simplicial map
\[
i:L\to K.
\]

\par

\textbf{(i) Constructing a chain homotopy inverse.}

We now show that the induced chain map
\[
i_\#:C_\bullet(L)\to C_\bullet(K)
\]
is a chain homotopy equivalence.

\par

\textbf{Step 3: choose orientations.}

Choose orientations of \(\sigma\) and \(\tau\) so that \(\tau\) appears in \(\partial \sigma\) with coefficient \(+1\). Then we can write
\[
\partial \sigma = \tau + \beta,
\]
where \(\beta\in C_{n-1}(L)\) is the sum of the other codimension-one faces of \(\sigma\), with the appropriate signs.

Indeed, all faces of \(\sigma\) other than \(\tau\) belong to \(L\), because only \(\sigma\) and \(\tau\) were removed.

\par

\textbf{Step 4: define a chain map \(r:C_\bullet(K)\to C_\bullet(L)\).}

We define \(r\) on oriented simplices and extend linearly.

\begin{itemize}
	\item If \(\eta\) is a simplex of \(L\) and \(\eta\neq \tau\), define
	\[
	r(\eta)=\eta.
	\]

	\item Define
	\[
	r(\sigma)=0.
	\]

	\item Define
	\[
	r(\tau)=-\beta.
	\]
\end{itemize}

Since \(\beta\in C_{n-1}(L)\), this is well-defined and lands in \(C_\bullet(L)\).

\par

\textbf{Step 5: verify that \(r\) is a chain map.}

We must show that
\[
r\partial = \partial r.
\]

\par

\emph{Case 1: \(\eta\in L\), \(\eta\neq \tau,\sigma\).}

Then \(r(\eta)=\eta\). Since \(L\) is a subcomplex, every face of \(\eta\) also lies in \(L\). Therefore
\[
r(\partial \eta)=\partial \eta=\partial(r(\eta)).
\]

\par

\emph{Case 2: \(\eta=\sigma\).}

We have
\[
r(\sigma)=0,
\qquad
\partial r(\sigma)=0.
\]
On the other hand,
\[
\partial \sigma = \tau+\beta,
\]
so
\[
r(\partial \sigma)=r(\tau)+r(\beta)=(-\beta)+\beta=0.
\]
Hence
\[
r(\partial \sigma)=\partial r(\sigma).
\]

\par

\emph{Case 3: \(\eta=\tau\).}

By definition,
\[
r(\tau)=-\beta,
\]
so
\[
\partial r(\tau)= -\partial \beta.
\]
Now from
\[
\partial \sigma = \tau+\beta
\]
and \(\partial^2=0\), we get
\[
0=\partial^2 \sigma = \partial(\tau+\beta)=\partial \tau+\partial \beta.
\]
Hence
\[
\partial \beta = -\partial \tau,
\]
and therefore
\[
\partial r(\tau)= -(-\partial \tau)=\partial \tau.
\]
Since every proper face of \(\tau\) lies in \(L\), we have \(r(\partial \tau)=\partial \tau\). Thus
\[
r(\partial \tau)=\partial r(\tau).
\]

So in all cases,
\[
r\partial=\partial r.
\]
Therefore \(r\) is a chain map.

\par

\textbf{Step 6: \(r\circ i_\#=\mathrm{id}_{C_\bullet(L)}\).}

If \(c\in C_\bullet(L)\), then \(i_\#(c)\) is simply the same chain viewed inside \(C_\bullet(K)\), and \(r\) fixes every simplex of \(L\). Therefore
\[
r\circ i_\# = \mathrm{id}_{C_\bullet(L)}.
\]

\par

\textbf{Step 7: construct a chain homotopy between \(i_\# r\) and \(\mathrm{id}_{C_\bullet(K)}\).}

Define
\[
h:C_j(K)\to C_{j+1}(K)
\]
on basis simplices by
\[
h(\tau)=\sigma,
\]
and
\[
h(\eta)=0
\qquad\text{for all simplices } \eta\neq \tau.
\]
Extend linearly.

We claim that
\[
\partial h + h\partial = \mathrm{id}_{C_\bullet(K)} - i_\# r.
\]

We verify this on basis simplices.

\par

\emph{Case 1: \(\eta\in L\), \(\eta\neq \tau\).}

Then \(h(\eta)=0\). Also, because \(\tau\) is not a face of any simplex of \(L\), the boundary \(\partial \eta\) contains no \(\tau\)-term. Hence
\[
h(\partial \eta)=0.
\]
So
\[
(\partial h+h\partial)(\eta)=0.
\]
On the other hand, \(r(\eta)=\eta\), so
\[
(\mathrm{id}-i_\# r)(\eta)=\eta-\eta=0.
\]

\par

\emph{Case 2: \(\eta=\sigma\).}

Since \(h(\sigma)=0\),
\[
\partial h(\sigma)=0.
\]
And because
\[
\partial \sigma=\tau+\beta,
\]
we get
\[
h(\partial \sigma)=h(\tau+\beta)=\sigma.
\]
Therefore
\[
(\partial h+h\partial)(\sigma)=\sigma.
\]
Also \(r(\sigma)=0\), so
\[
(\mathrm{id}-i_\# r)(\sigma)=\sigma.
\]

\par

\emph{Case 3: \(\eta=\tau\).}

We have
\[
h(\tau)=\sigma,
\]
so
\[
\partial h(\tau)=\partial \sigma=\tau+\beta.
\]
Since every face of \(\tau\) lies in \(L\), and \(h\) vanishes on simplices other than \(\tau\), we have
\[
h(\partial \tau)=0.
\]
Thus
\[
(\partial h+h\partial)(\tau)=\tau+\beta.
\]
On the other hand,
\[
r(\tau)=-\beta,
\]
so
\[
(\mathrm{id}-i_\# r)(\tau)=\tau-(-\beta)=\tau+\beta.
\]

Hence in all cases,
\[
\partial h+h\partial = \mathrm{id}_{C_\bullet(K)} - i_\# r.
\]
So \(i_\# r\) is chain homotopic to the identity on \(C_\bullet(K)\).

Together with
\[
r i_\# = \mathrm{id}_{C_\bullet(L)},
\]
this shows that \(i_\#\) and \(r\) are chain homotopy inverses.

\par

\textbf{Step 8: conclude on homology.}

Chain homotopy equivalent chain complexes have isomorphic homology. Therefore
\[
i_*:H_j(L)\xrightarrow{\cong} H_j(K)
\qquad\text{for all } j.
\]

This proves part (i).

\par

\textbf{(ii) Proof using the previous question twice.}

We now prove the same result using the previous problem, which described the effect of deleting a maximal simplex.

Let
\[
K_1 = K\setminus\{\sigma\}.
\]
Since \(\sigma\) is not a proper face of any simplex of \(K\), the previous problem applies to the pair
\[
K_1 \subset K.
\]
Thus there is an exact sequence
\[
0 \to H_n(K_1)\to H_n(K)\to \mathbb Z
\to H_{n-1}(K_1)\to H_{n-1}(K)\to 0,
\]
and
\[
H_j(K_1)\cong H_j(K)
\qquad\text{for } j\neq n,n-1.
\]

Now observe that in \(K_1\), the simplex \(\tau\) is no longer a proper face of any simplex, because its only coface in \(K\) was \(\sigma\), and \(\sigma\) has been removed. Therefore the previous problem applies again, this time to
\[
L = K_1\setminus\{\tau\}.
\]
Since \(\dim \tau=n-1\), we obtain a second exact sequence
\[
0 \to H_{n-1}(L)\to H_{n-1}(K_1)\to \mathbb Z
\to H_{n-2}(L)\to H_{n-2}(K_1)\to 0,
\]
and
\[
H_j(L)\cong H_j(K_1)
\qquad\text{for } j\neq n-1,n-2.
\]

Now the key point is that these two changes cancel each other, because \(\tau\) is a free face of \(\sigma\). Algebraically, removing \(\sigma\) creates a possible change controlled by the class of \(\partial \sigma\), and removing \(\tau\) immediately removes exactly that remaining contribution. Geometrically, the pair \((\sigma,\tau)\) is an elementary collapse pair, so deleting both should not change homology.

Thus the composite inclusion
\[
L \hookrightarrow K_1 \hookrightarrow K
\]
induces an isomorphism
\[
H_j(L)\xrightarrow{\cong} H_j(K)
\qquad\text{for all } j.
\]

This proves part (ii).


\medskip
\noindent\textbf{Conclusion}\par
\smallskip


We have shown the following.

\begin{itemize}
	\item If \(\sigma\) is a maximal simplex and \(\tau<\sigma\) is a codimension-one face which is not a face of any other simplex, then
	\[
	L=K\setminus\{\sigma,\tau\}
	\]
	is again a simplicial complex.

	\item The inclusion of vertex sets defines a simplicial map
	\[
	i:L\to K.
	\]

	\item The induced chain map
	\[
	i_\#:C_\bullet(L)\to C_\bullet(K)
	\]
	is a chain homotopy equivalence.

	\item Consequently,
	\[
	i_*:H_j(L)\xrightarrow{\cong} H_j(K)
	\qquad\text{for all } j.
	\]

	\item The same conclusion may also be derived by applying the previous deletion result twice: first delete \(\sigma\), then delete \(\tau\).
\end{itemize}

Thus deleting a free face together with the unique simplex containing it preserves homology in all degrees. This is the algebraic-topological content of an elementary collapse.
\end{solution}

\begin{problem}[CP-VIII-0030 --- A five-lemma type isomorphism criterion]
Consider a commutative diagram


\[
\begin{array}{ccccccccc}
A_1&\xrightarrow{f_1}&A_2&\xrightarrow{f_2}&A_3&\xrightarrow{f_3}&A_4&\xrightarrow{f_4}&A_5\\
\big\downarrow{\scriptstyle h_1}&&\big\downarrow{\scriptstyle h_2}&&
\big\downarrow{\scriptstyle h_3}&&\big\downarrow{\scriptstyle h_4}&&
\big\downarrow{\scriptstyle h_5}\\
B_1&\xrightarrow{g_1}&B_2&\xrightarrow{g_2}&B_3&\xrightarrow{g_3}&B_4&\xrightarrow{g_4}&B_5
\end{array}
\]


in which the rows are exact and each square commutes. If \(h_1,h_2,h_4,\) and \(h_5\) are isomorphisms, show that \(h_3\) is also an isomorphism.


This is a standard form of the five lemma. We prove that \(h_3\) is injective and surjective.

\par

\textbf{Injectivity of \(h_3\).}
Let \(a\in A_3\) satisfy
\[
h_3(a)=0.
\]
Apply \(g_3\) to this equality. Since the diagram commutes,
\[
g_3(h_3(a))=h_4(f_3(a)).
\]
The left-hand side is zero, so
\[
h_4(f_3(a))=0.
\]
Because \(h_4\) is an isomorphism, it is injective, hence
\[
f_3(a)=0.
\]
By exactness of the top row at \(A_3\),
\[
\ker(f_3)=\operatorname{im}(f_2),
\]
so there exists \(b\in A_2\) such that
\[
a=f_2(b).
\]

Now
\[
0=h_3(a)=h_3(f_2(b))=g_2(h_2(b)),
\]
again by commutativity. Hence
\[
h_2(b)\in \ker(g_2)=\operatorname{im}(g_1)
\]
by exactness of the bottom row. Therefore there exists \(c\in B_1\) such that
\[
g_1(c)=h_2(b).
\]
Since \(h_1\) is an isomorphism, there exists \(d\in A_1\) with
\[
h_1(d)=c.
\]
Then commutativity of the left square gives
\[
h_2(f_1(d))=g_1(h_1(d))=g_1(c)=h_2(b).
\]
Because \(h_2\) is injective, we conclude
\[
f_1(d)=b.
\]
Thus
\[
a=f_2(b)=f_2(f_1(d))=0,
\]
since consecutive maps in an exact sequence compose to zero. Therefore \(\ker(h_3)=0\), so \(h_3\) is injective.

\par

\textbf{Surjectivity of \(h_3\).}
Let \(y\in B_3\). We want to find \(x\in A_3\) such that
\[
h_3(x)=y.
\]

First consider \(g_3(y)\in B_4\). Since \(h_4\) is surjective, there exists \(a_4\in A_4\) with
\[
h_4(a_4)=g_3(y).
\]
Apply \(g_4\) to both sides:
\[
g_4(h_4(a_4))=g_4(g_3(y))=0.
\]
By commutativity,
\[
h_5(f_4(a_4))=0.
\]
Since \(h_5\) is injective, it follows that
\[
f_4(a_4)=0.
\]
By exactness of the top row at \(A_4\),
\[
a_4\in \ker(f_4)=\operatorname{im}(f_3),
\]
so there exists \(x_0\in A_3\) such that
\[
f_3(x_0)=a_4.
\]

Now compute:
\[
g_3(h_3(x_0))=h_4(f_3(x_0))=h_4(a_4)=g_3(y).
\]
Hence
\[
g_3\bigl(y-h_3(x_0)\bigr)=0.
\]
So
\[
y-h_3(x_0)\in \ker(g_3)=\operatorname{im}(g_2)
\]
by exactness of the bottom row. Therefore there exists \(b_2\in B_2\) such that
\[
g_2(b_2)=y-h_3(x_0).
\]

Since \(h_2\) is surjective, choose \(a_2\in A_2\) with
\[
h_2(a_2)=b_2.
\]
Set
\[
x=x_0+f_2(a_2)\in A_3.
\]
Then
\[
\begin{aligned}
h_3(x)
&=h_3(x_0)+h_3(f_2(a_2))\\
&=h_3(x_0)+g_2(h_2(a_2))\\
&=h_3(x_0)+g_2(b_2)\\
&=h_3(x_0)+\bigl(y-h_3(x_0)\bigr)=y.
\end{aligned}
\]
Thus every \(y\in B_3\) lies in the image of \(h_3\), so \(h_3\) is surjective.

\par

Since \(h_3\) is both injective and surjective, it is an isomorphism. Therefore
\[
\boxed{h_3 \text{ is an isomorphism}.}
\]
\end{problem}

\begin{solution}
\emph{The source solution requires editorial reconstruction in the next curation pass.}
\end{solution}

\begin{problem}[CP-VIII-0031 --- Suspension of a simplicial complex and induced maps on homology]
Let \(K\) be a simplicial complex in \(\mathbb R^m\), and consider this as lying inside \(\mathbb R^{m+1}\) as the vectors of the form
\[
(x_1,\dots,x_m,0).
\]
Let
\[
e_+=(0,\dots,0,1)\in \mathbb R^{m+1},
\qquad
e_-=(0,\dots,0,-1)\in \mathbb R^{m+1}.
\]
The \emph{suspension} of \(K\) is the simplicial complex in \(\mathbb R^{m+1}\)
\[
SK:=K\cup \{\langle v_0,\dots,v_n,e_+\rangle,\ \langle v_0,\dots,v_n,e_-\rangle
\mid \langle v_0,\dots,v_n\rangle\in K\}.
\]

\begin{enumerate}
	\item[(a)] Show that \(SK\) is a simplicial complex, and that if \(|K|\cong S^n\) then \(|SK|\cong S^{n+1}\).
	\item[(b)] Using the Mayer--Vietoris sequence, show that if \(K\) is connected then
	\[
	H_0(SK)\cong \mathbb Z,\qquad H_1(SK)=0,
	\qquad H_i(SK)\cong H_{i-1}(K)\ \text{if } i\ge 2.
	\]
	\item[(c)] If \(f:K\to K\) is a simplicial map, let \(Sf:SK\to SK\) be the unique simplicial map which agrees with \(f\) on the subcomplex \(K\) and fixes the points \(e_+\) and \(e_-\). Show that under the isomorphism in (b), the maps \(f_*\) and \((Sf)_*\) agree.
	\item[(d)] Deduce that for every \(n\ge 1\) and \(d\in \mathbb Z\) there is a map
	\[
	f:S^n\to S^n
	\]
	so that \(f_*\) induces multiplication by \(d\) on
	\[
	H_n(S^n)\cong \mathbb Z.
	\]
\end{enumerate}


\par

\textbf{(a) The suspension \(SK\) is a simplicial complex.}

Let
\[
\sigma=\langle v_0,\dots,v_n\rangle\in K.
\]
Then the added simplices are the cones
\[
\langle v_0,\dots,v_n,e_+\rangle,
\qquad
\langle v_0,\dots,v_n,e_-\rangle.
\]
To check that \(SK\) is a simplicial complex, we verify the two defining properties.

First, every face of a simplex in \(SK\) again lies in \(SK\). Indeed, if a simplex already lies in \(K\), then all its faces lie in \(K\) because \(K\) is a simplicial complex. If the simplex is of the form
\[
\langle v_0,\dots,v_n,e_+\rangle
\quad\text{or}\quad
\langle v_0,\dots,v_n,e_-\rangle,
\]
then any face is either:

\begin{itemize}
	\item a face \(\langle v_{i_0},\dots,v_{i_r}\rangle\) of \(\sigma\), hence in \(K\subset SK\), or
	\item a simplex of the form \(\langle v_{i_0},\dots,v_{i_r},e_+\rangle\) or \(\langle v_{i_0},\dots,v_{i_r},e_-\rangle\), again by construction in \(SK\).
\end{itemize}

Second, the intersection of two simplices of \(SK\) is either empty or a common face. This is clear if both lie in \(K\), since \(K\) is a simplicial complex. If both contain \(e_+\), then they are cones on simplices of \(K\), and their intersection is the cone on the intersection of the base simplices, hence a common face. The same argument applies to two simplices containing \(e_-\). If one simplex contains \(e_+\) and the other contains \(e_-\), then they cannot share either \(e_+\) or \(e_-\), so their intersection is exactly the intersection of the corresponding base simplices in \(K\), again a common face or empty.

Thus \(SK\) is a simplicial complex.

\par

Now suppose \(|K|\cong S^n\). Geometrically, \(|SK|\) is obtained by taking two cones on \(|K|\), one with apex \(e_+\) and one with apex \(e_-\), and gluing them along the common base \(|K|\). This is exactly the topological suspension of \(|K|\):
\[
|SK|\cong \Sigma |K|.
\]
Since the suspension of \(S^n\) is \(S^{n+1}\), we obtain
\[
|SK|\cong \Sigma S^n\cong S^{n+1}.
\]
Hence
\[
\boxed{|K|\cong S^n \implies |SK|\cong S^{n+1}.}
\]

\par

\textbf{(b) Homology of the suspension.}

Let
\[
U=\text{the cone on }K\text{ with apex }e_+,
\qquad
V=\text{the cone on }K\text{ with apex }e_-.
\]
Then
\[
SK=U\cup V,
\qquad
U\cap V=K.
\]
Both \(U\) and \(V\) are cones, hence contractible. Therefore
\[
H_i(U)=H_i(V)=0 \qquad (i\ge 1),
\]
and
\[
H_0(U)\cong \mathbb Z,\qquad H_0(V)\cong \mathbb Z.
\]

Apply the Mayer--Vietoris sequence to the decomposition \(SK=U\cup V\):
\[
\cdots \to H_i(K)\to H_i(U)\oplus H_i(V)\to H_i(SK)\xrightarrow{\delta} H_{i-1}(K)\to H_{i-1}(U)\oplus H_{i-1}(V)\to \cdots
\]

If \(i\ge 2\), then
\[
H_i(U)\oplus H_i(V)=0
\qquad\text{and}\qquad
H_{i-1}(U)\oplus H_{i-1}(V)=0,
\]
so the exact sequence becomes
\[
0\to H_i(SK)\xrightarrow{\delta} H_{i-1}(K)\to 0.
\]
Thus
\[
\boxed{H_i(SK)\cong H_{i-1}(K)\qquad (i\ge 2).}
\]

Now consider the low-dimensional part:
\[
0\to H_1(SK)\to H_0(K)\xrightarrow{\phi} H_0(U)\oplus H_0(V)\to H_0(SK)\to 0.
\]
Since \(K\) is connected,
\[
H_0(K)\cong \mathbb Z.
\]
Also
\[
H_0(U)\oplus H_0(V)\cong \mathbb Z\oplus \mathbb Z.
\]
The map \(\phi\) is induced by the inclusions \(K\hookrightarrow U\) and \(K\hookrightarrow V\), so on generators it is the diagonal map
\[
\phi(1)=(1,1).
\]
Hence \(\phi\) is injective, so
\[
H_1(SK)=0.
\]
Its cokernel is
\[
(\mathbb Z\oplus \mathbb Z)/\Delta(\mathbb Z)\cong \mathbb Z,
\]
so
\[
H_0(SK)\cong \mathbb Z.
\]

Therefore
\[
\boxed{H_0(SK)\cong \mathbb Z,\qquad H_1(SK)=0,\qquad H_i(SK)\cong H_{i-1}(K)\text{ for }i\ge 2.}
\]

\par

\textbf{(c) Compatibility of \(f_*\) and \((Sf)_*\).}

Let \(f:K\to K\) be simplicial, and let \(Sf:SK\to SK\) be its suspension, fixing \(e_+\) and \(e_-\). Since \(Sf\) fixes the two cone points, it preserves the subcomplexes \(U\) and \(V\), and on
\[
U\cap V=K
\]
it restricts exactly to \(f\). Thus we get a morphism between the Mayer--Vietoris sequences of the decomposition
\[
SK=U\cup V.
\]

For \(i\ge 2\), the isomorphism
\[
H_i(SK)\xrightarrow{\cong} H_{i-1}(K)
\]
from part (b) is precisely the connecting homomorphism \(\delta\) in Mayer--Vietoris. By naturality of the Mayer--Vietoris sequence, the diagram


\[
\begin{array}{ccc}
H_i(SK)&\xrightarrow{\ \delta\ }&H_{i-1}(K)\\
\big\downarrow{\scriptstyle (Sf)_*}&&\big\downarrow{\scriptstyle f_*}\\
H_i(SK)&\xrightarrow{\ \delta\ }&H_{i-1}(K)
\end{array}
\]


commutes.

Since \(\delta\) is an isomorphism, under the identification
\[
H_i(SK)\cong H_{i-1}(K),
\]
the map \((Sf)_*\) corresponds exactly to \(f_*\). That is,
\[
\boxed{\delta\circ (Sf)_*=f_*\circ \delta.}
\]

\par

\textbf{(d) Maps of arbitrary degree on spheres.}

We prove by induction on \(n\ge 1\) that for every integer \(d\) there exists a map
\[
f:S^n\to S^n
\]
whose induced map on \(H_n(S^n)\cong \mathbb Z\) is multiplication by \(d\).

For \(n=1\), consider the map
\[
f_d:S^1\to S^1,
\qquad
f_d(z)=z^d
\quad (z\in S^1\subset \mathbb C).
\]
This is a continuous map of degree \(d\), so
\[
(f_d)_*:H_1(S^1)\to H_1(S^1)
\]
is multiplication by \(d\).

Now assume the statement holds for \(S^n\). Then there is a map
\[
f:S^n\to S^n
\]
with \(f_*\) equal to multiplication by \(d\) on \(H_n(S^n)\). Suspend \(f\) to obtain
\[
Sf:S^{n+1}\to S^{n+1},
\]
using part (a), since \(S^{n+1}\cong \Sigma S^n\). By part (c), under the isomorphism
\[
H_{n+1}(S^{n+1})\cong H_n(S^n),
\]
the map \((Sf)_*\) agrees with \(f_*\). Hence \((Sf)_*\) is also multiplication by \(d\).

Therefore, by induction,
\[
\boxed{\text{for every }n\ge 1\text{ and }d\in\mathbb Z,\text{ there exists }f:S^n\to S^n
	\text{ with }f_*=\times d\text{ on }H_n(S^n)\cong \mathbb Z.}
\]
\end{problem}

\begin{solution}
\emph{The source solution requires editorial reconstruction in the next curation pass.}
\end{solution}

\begin{problem}[CP-VIII-0032 --- Rational homology detects only the free part]
If \(K\) is a simplicial complex with
\[
H_i(K)\cong \mathbb Z^r\oplus F,
\]
where \(F\) is a finite abelian group, show that
\[
H_i(K;\mathbb Q)\cong \mathbb Q^r.
\]

\emph{Hint.} Use
\[
C_\bullet(K;\mathbb Q)
=
C_\bullet(K)\otimes_{\mathbb Z}\mathbb Q
\]
and the flatness of \(\mathbb Q\) as a \(\mathbb Z\)-module.
\end{problem}

\begin{solution}
Since
\[
C_\bullet(K;\mathbb Q)
=
C_\bullet(K)\otimes_{\mathbb Z}\mathbb Q
\]
and \(\mathbb Q\) is flat over \(\mathbb Z\), tensoring with \(\mathbb Q\)
commutes with homology. Hence
\[
H_i(K;\mathbb Q)
\cong
H_i(K)\otimes_{\mathbb Z}\mathbb Q.
\]
Using
\[
H_i(K)\cong \mathbb Z^r\oplus F,
\]
we obtain
\[
H_i(K;\mathbb Q)
\cong
(\mathbb Z^r\otimes_{\mathbb Z}\mathbb Q)
\oplus
(F\otimes_{\mathbb Z}\mathbb Q).
\]
The free part gives
\[
\mathbb Z^r\otimes_{\mathbb Z}\mathbb Q\cong\mathbb Q^r.
\]
For the finite group \(F\), every element \(x\in F\) has finite order:
\(mx=0\) for some \(m>0\). Therefore, for \(q\in\mathbb Q\),
\[
x\otimes q
=
mx\otimes\frac qm
=
0,
\]
so
\[
F\otimes_{\mathbb Z}\mathbb Q=0.
\]
Consequently,
\[
\boxed{H_i(K;\mathbb Q)\cong\mathbb Q^r.}
\]

Thus rational homology preserves the free rank and kills finite torsion.
For example,
\[
\mathbb Z^2\oplus\mathbb Z/6
\quad\longmapsto\quad
\mathbb Q^2,
\qquad
\mathbb Z/4\oplus\mathbb Z/9
\quad\longmapsto\quad
0.
\]
\end{solution}

\begin{problem}[CP-VIII-0033 --- Homology of the triangulated torus, Klein bottle, and projective plane]
Using the two previous questions, compute the homology groups of the simplicial complexes described in Q2, and describe generators for each homology group.
\end{problem}

\begin{solution}
The simplicial complexes from Q2 triangulate the following spaces:

\[
T,\qquad K,\qquad \mathbb{RP}^2.
\]

By repeated elementary collapses and deletions of maximal simplices, as described in the previous two questions, each triangulation may be reduced to the usual quotient-square model of the corresponding surface. Therefore the homology groups are the standard homology groups of the torus, Klein bottle, and projective plane.

We compute each in turn.

\par

\textbf{1. The torus.}

The torus has the standard CW structure with:

\begin{itemize}
	\item one \(0\)-cell,
	\item two \(1\)-cells \(a,b\),
	\item one \(2\)-cell attached along the loop
	\[
	aba^{-1}b^{-1}.
	\]
\end{itemize}

Hence the cellular chain complex is
\[
0 \longrightarrow C_2(T)\cong \mathbb Z
\xrightarrow{d_2}
C_1(T)\cong \mathbb Z^2
\xrightarrow{d_1}
C_0(T)\cong \mathbb Z
\longrightarrow 0.
\]

Since there is only one \(0\)-cell, we have
\[
d_1=0.
\]
Also, in the attaching word \(aba^{-1}b^{-1}\), the total exponent of both \(a\) and \(b\) is \(0\), so
\[
d_2=0.
\]

Therefore
\[
H_2(T)\cong \ker d_2 \cong \mathbb Z,
\]
\[
H_1(T)\cong \ker d_1/\operatorname{im} d_2 \cong \mathbb Z^2,
\]
and
\[
H_0(T)\cong \mathbb Z.
\]
All higher homology groups vanish:
\[
H_j(T)=0 \qquad (j\ge 3).
\]

\par

\textbf{Generators for the torus.}

\begin{itemize}
	\item \(H_0(T)\cong \mathbb Z\) is generated by the class of any vertex.
	\item \(H_1(T)\cong \mathbb Z\oplus \mathbb Z\) is generated by the two basic loops:
	\[
	[a],\qquad [b],
	\]
	where \(a\) is the horizontal loop and \(b\) is the vertical loop in the quotient-square model.
	\item \(H_2(T)\cong \mathbb Z\) is generated by the fundamental class
	\[
	[T],
	\]
	which may be represented simplicially by the sum of all oriented \(2\)-simplices in the triangulation, with compatible orientations.
\end{itemize}

Thus
\[
H_0(T)\cong \mathbb Z,\qquad
H_1(T)\cong \mathbb Z^2,\qquad
H_2(T)\cong \mathbb Z.
\]

\par

\textbf{2. The Klein bottle.}

The Klein bottle has the standard CW structure with:

\begin{itemize}
	\item one \(0\)-cell,
	\item two \(1\)-cells \(a,b\),
	\item one \(2\)-cell attached along the loop
	\[
	aba\,b^{-1},
	\]
	equivalently \(aba^{-1}b\), depending on conventions.
\end{itemize}

With the usual convention, the boundary map in degree \(2\) is
\[
d_2(1)=2a,
\]
up to a choice of basis in \(C_1(K)\cong \mathbb Z^2\). Hence the cellular chain complex is
\[
0 \longrightarrow \mathbb Z
\xrightarrow{d_2}
\mathbb Z^2
\xrightarrow{0}
\mathbb Z
\longrightarrow 0,
\]
with
\[
d_2(1)=(2,0)
\]
after choosing a suitable basis of \(C_1\).

Therefore
\[
H_2(K)=\ker d_2=0,
\]
and
\[
H_1(K)\cong \mathbb Z^2/\langle (2,0)\rangle \cong \mathbb Z\oplus \mathbb Z/2.
\]
Also
\[
H_0(K)\cong \mathbb Z,
\]
and
\[
H_j(K)=0 \qquad (j\ge 3).
\]

\par

\textbf{Generators for the Klein bottle.}

\begin{itemize}
	\item \(H_0(K)\cong \mathbb Z\) is generated by the class of any vertex.
	\item \(H_1(K)\cong \mathbb Z\oplus \mathbb Z/2\) is generated by:
	\begin{itemize}
		\item one class \([b]\) of infinite order, represented by the untwisted direction in the square model;
		\item one class \([a]\) of order \(2\), represented by the twisted direction.
	\end{itemize}
	Thus
	\[
	2[a]=0,
	\qquad
	[b]\text{ has infinite order}.
	\]
	\item Since the Klein bottle is nonorientable,
	\[
	H_2(K)=0.
	\]
	There is no global compatible choice of orientations on all \(2\)-simplices producing a nonzero \(2\)-cycle.
\end{itemize}

So
\[
H_0(K)\cong \mathbb Z,\qquad
H_1(K)\cong \mathbb Z\oplus \mathbb Z/2,\qquad
H_2(K)=0.
\]

\par

\textbf{3. The projective plane.}

The projective plane has the standard CW structure with:

\begin{itemize}
	\item one \(0\)-cell,
	\item one \(1\)-cell \(a\),
	\item one \(2\)-cell attached by the map
	\[
	a^2.
	\]
\end{itemize}

So the cellular chain complex is
\[
0 \longrightarrow \mathbb Z
\xrightarrow{2}
\mathbb Z
\xrightarrow{0}
\mathbb Z
\longrightarrow 0.
\]

Therefore
\[
H_2(\mathbb{RP}^2)=0,
\]
\[
H_1(\mathbb{RP}^2)\cong \mathbb Z/2,
\]
and
\[
H_0(\mathbb{RP}^2)\cong \mathbb Z.
\]
Also
\[
H_j(\mathbb{RP}^2)=0 \qquad (j\ge 3).
\]

\par

\textbf{Generators for the projective plane.}

\begin{itemize}
	\item \(H_0(\mathbb{RP}^2)\cong \mathbb Z\) is generated by the class of any vertex.
	\item \(H_1(\mathbb{RP}^2)\cong \mathbb Z/2\) is generated by the class \([a]\) of the basic loop in the quotient-square model. It satisfies
	\[
	2[a]=0.
	\]
	\item Since \(\mathbb{RP}^2\) is nonorientable,
	\[
	H_2(\mathbb{RP}^2)=0.
	\]
\end{itemize}

So
\[
H_0(\mathbb{RP}^2)\cong \mathbb Z,\qquad
H_1(\mathbb{RP}^2)\cong \mathbb Z/2,\qquad
H_2(\mathbb{RP}^2)=0.
\]


\medskip
\noindent\textbf{Final answer}\par
\smallskip


We conclude that
\[
\boxed{
	\begin{aligned}
		H_0(T)&\cong \mathbb Z, &
		H_1(T)&\cong \mathbb Z^2, &
		H_2(T)&\cong \mathbb Z,\\[0.5ex]
		H_0(K)&\cong \mathbb Z, &
		H_1(K)&\cong \mathbb Z\oplus \mathbb Z/2, &
		H_2(K)&=0,\\[0.5ex]
		H_0(\mathbb{RP}^2)&\cong \mathbb Z, &
		H_1(\mathbb{RP}^2)&\cong \mathbb Z/2, &
		H_2(\mathbb{RP}^2)&=0,
\end{aligned}}
\]
and all higher homology groups vanish.

Generators may be chosen as follows:

\begin{itemize}
	\item \(H_0\): the class of any vertex;
	\item \(H_1(T)\): the two basic loops \(a\) and \(b\);
	\item \(H_2(T)\): the fundamental class, represented by the sum of all oriented \(2\)-simplices;
	\item \(H_1(K)\): one untwisted loop of infinite order and one twisted loop of order \(2\);
	\item \(H_1(\mathbb{RP}^2)\): the basic projective loop of order \(2\).
\end{itemize}

\par
\end{solution}

\begin{problem}[CP-VIII-0034 --- Top homology of an \(n\)-dimensional pseudomanifold-type simplicial complex]
Let \(K\) be an \(n\)-dimensional simplicial complex such that

\begin{enumerate}
	\item[(i)] every \((n-1)\)-simplex is a face of precisely two \(n\)-simplices, and
	\item[(ii)] every pair of \(n\)-simplices can be connected by a sequence of \(n\)-simplices such that adjacent terms share an \((n-1)\)-dimensional face.
\end{enumerate}

Show that
\[
H_n(K)
\]
is either \(\mathbb Z\) or trivial. In the first case show that \(H_n(K)\) is generated by a cycle which is a sum of all \(n\)-simplices with suitable orientations.
\end{problem}

\begin{solution}
Let the \(n\)-simplices of \(K\) be
\[
\sigma_1,\dots,\sigma_r.
\]
Then
\[
C_n(K)\cong \bigoplus_{i=1}^r \mathbb Z\,\sigma_i.
\]
Since \(K\) has dimension \(n\), there are no \((n+1)\)-simplices. Therefore
\[
B_n(K)=0,
\]
so
\[
H_n(K)=Z_n(K)=\ker(\partial_n).
\]
Thus it is enough to determine all \(n\)-cycles.

\par

\textbf{Step 1: a general \(n\)-cycle.}

Take an arbitrary \(n\)-chain
\[
z=\sum_{i=1}^r a_i\sigma_i
\qquad (a_i\in \mathbb Z).
\]
Assume that \(z\) is a cycle:
\[
\partial z=0.
\]

Let \(\tau\) be any \((n-1)\)-simplex of \(K\). By hypothesis (i), \(\tau\) is a face of exactly two \(n\)-simplices, say \(\sigma_i\) and \(\sigma_j\). In the boundary \(\partial z\), the coefficient of \(\tau\) is
\[
a_i[\tau:\sigma_i] + a_j[\tau:\sigma_j],
\]
where each incidence number is \(\pm 1\). Since \(\partial z=0\), this coefficient must vanish:
\[
a_i[\tau:\sigma_i] + a_j[\tau:\sigma_j]=0.
\]
Hence
\[
a_j=\pm a_i.
\]

So whenever two \(n\)-simplices share an \((n-1)\)-face, the absolute values of their coefficients in a cycle are equal.

\par

\textbf{Step 2: connectedness forces all absolute values to agree.}

By hypothesis (ii), any two \(n\)-simplices can be connected by a chain of adjacent \(n\)-simplices
\[
\sigma_{i_0},\sigma_{i_1},\dots,\sigma_{i_m}
\]
such that consecutive simplices share an \((n-1)\)-face. Repeatedly applying the conclusion of Step 1 along this chain, we obtain
\[
|a_{i_0}|=|a_{i_1}|=\cdots=|a_{i_m}|.
\]
Since any pair of \(n\)-simplices can be connected in this way, it follows that all coefficients of the cycle satisfy
\[
|a_1|=|a_2|=\cdots=|a_r|.
\]

Thus every \(n\)-cycle is of the form
\[
z=m\sum_{i=1}^r \varepsilon_i \sigma_i,
\]
where
\[
m\in \mathbb Z,
\qquad
\varepsilon_i\in \{\pm 1\}.
\]
Therefore the group of \(n\)-cycles has rank at most \(1\). Since
\[
H_n(K)=Z_n(K),
\]
it follows that \(H_n(K)\) is either trivial or isomorphic to \(\mathbb Z\).

This proves the first assertion.

\par

\textbf{Step 3: if \(H_n(K)\neq 0\), construct a generator.}

Assume now that \(H_n(K)\neq 0\). Then there exists a nonzero \(n\)-cycle
\[
z=\sum_{i=1}^r a_i\sigma_i.
\]
By the previous step, all \(|a_i|\) are equal. Since \(z\neq 0\), let
\[
|a_i|=m>0
\qquad\text{for all } i.
\]
Then
\[
z=m\sum_{i=1}^r \varepsilon_i\sigma_i
\qquad (\varepsilon_i=\pm 1).
\]

Now orient each simplex \(\sigma_i\) so that its sign agrees with \(\varepsilon_i\). With these chosen orientations, the cycle takes the form
\[
z=m(\sigma_1+\cdots+\sigma_r).
\]
Therefore the chain
\[
c=\sigma_1+\cdots+\sigma_r
\]
(with suitable orientations) is itself an \(n\)-cycle.

Since every \(n\)-cycle is an integer multiple of this one, the cycle \(c\) generates \(Z_n(K)\), and hence generates \(H_n(K)\cong \mathbb Z\).

Thus, in the nontrivial case, \(H_n(K)\) is generated by the sum of all \(n\)-simplices with suitably chosen orientations.

\par

\textbf{Step 4: geometric interpretation.}

The two possibilities correspond to orientability behavior.

\begin{itemize}
	\item If one can choose orientations on the \(n\)-simplices so that along each shared \((n-1)\)-face the contributions cancel, then the sum of all \(n\)-simplices is a nonzero cycle. In that case
	\[
	H_n(K)\cong \mathbb Z.
	\]

	\item If no such global consistent choice exists, then there is no nonzero \(n\)-cycle at all, and
	\[
	H_n(K)=0.
	\]
\end{itemize}

So this theorem is the simplicial version of the familiar statement that the top homology of a connected closed \(n\)-dimensional pseudomanifold is either \(\mathbb Z\) or \(0\), depending on orientability.


\medskip
\noindent\textbf{Conclusion}\par
\smallskip


We have proved that
\[
\boxed{H_n(K)\text{ is either }0\text{ or }\mathbb Z.}
\]
Moreover, if
\[
H_n(K)\cong \mathbb Z,
\]
then a generator is represented by
\[
\sum_{\sigma\in K^{(n)}} \sigma,
\]
where the sum runs over all \(n\)-simplices and the orientations are chosen so that adjacent simplices cancel along their common \((n-1)\)-faces.

Thus the top homology is completely controlled by whether the top-dimensional simplices admit a globally compatible orientation.
\end{solution}

\begin{problem}[CP-VIII-0035 --- Products of polyhedra are polyhedra of simplicial complexes]
For simplicial complexes \(K\) and \(L\) inside \(\mathbb R^m\) and \(\mathbb R^n\) respectively, show that
\[
|K|\times |L| \subseteq \mathbb R^{m+n}=\mathbb R^m\times \mathbb R^n
\]
is the polyhedron of a simplicial complex.

\medskip
\noindent\emph{Hint.} Prove the result first when both \(K\) and \(L\) consist of a single simplex together with all of its faces.
\smallskip
\end{problem}

\begin{solution}
We first treat the case of two simplices.

\par

\textbf{Step 1: product of two simplices.}

Suppose
\[
K=\{\text{all faces of a simplex }\sigma\},
\qquad
L=\{\text{all faces of a simplex }\tau\},
\]
where
\[
\sigma=\langle v_0,\dots,v_p\rangle \subseteq \mathbb R^m,
\qquad
\tau=\langle w_0,\dots,w_q\rangle \subseteq \mathbb R^n.
\]
Then
\[
|\sigma|\times |\tau| \subseteq \mathbb R^{m+n}
\]
has vertices
\[
(v_i,w_j)\in \mathbb R^m\times \mathbb R^n
\qquad
(0\le i\le p,\ 0\le j\le q).
\]

We triangulate \(\sigma\times \tau\) using these vertices.

\par

\textbf{Step 2: monotone lattice paths and staircase simplices.}

Consider the rectangular grid of index pairs
\[
(i,j),
\qquad
0\le i\le p,\ 0\le j\le q.
\]
A \emph{monotone lattice path} from \((0,0)\) to \((p,q)\) is a path that moves only:
\begin{itemize}
	\item one step to the right:
	\[
	(i,j)\to (i+1,j),
	\]
	\item or one step upward:
	\[
	(i,j)\to (i,j+1).
	\]
\end{itemize}

Such a path has exactly \(p+q+1\) vertices:
\[
(i_0,j_0),(i_1,j_1),\dots,(i_{p+q},j_{p+q}).
\]
To such a path \(P\), associate the simplex
\[
\Delta_P=
\langle
(v_{i_0},w_{j_0}),
(v_{i_1},w_{j_1}),
\dots,
(v_{i_{p+q}},w_{j_{p+q}})
\rangle
\subseteq \mathbb R^{m+n}.
\]
Each \(\Delta_P\) is a simplex of dimension \(p+q\).

The collection of all such simplices \(\Delta_P\), as \(P\) ranges over all monotone lattice paths from \((0,0)\) to \((p,q)\), is called the \emph{staircase triangulation} of \(\sigma\times\tau\).

\par

\textbf{Step 3: why this triangulates the product.}

The staircase simplices cover the whole product \(\sigma\times\tau\), and any two of them intersect in a common face. This is a standard fact about the staircase construction.

Geometrically, one may think of \(\sigma\times\tau\) as subdivided according to the relative order of the barycentric coordinates in the two factors. The monotone lattice paths record these possible orderings, and each such ordering determines one simplex.

Thus \(\sigma\times\tau\) is the polyhedron of a simplicial complex.

\par

\textbf{Example 1: \(\Delta^1\times \Delta^1\).}

The product of two edges is a square. The staircase triangulation divides it into two triangles, corresponding to the two monotone lattice paths from \((0,0)\) to \((1,1)\):
\[
(0,0)\to (1,0)\to (1,1),
\qquad
(0,0)\to (0,1)\to (1,1).
\]
So the square is triangulated by a diagonal.

\par

\textbf{Example 2: \(\Delta^1\times \Delta^2\).}

The product of an edge and a triangle is a triangular prism. The staircase triangulation subdivides it into \(3\) tetrahedra, corresponding to the \(3\) monotone paths from \((0,0)\) to \((1,2)\).

\par

\textbf{Step 4: pass to general simplicial complexes.}

Now let \(K\subseteq \mathbb R^m\) and \(L\subseteq \mathbb R^n\) be arbitrary simplicial complexes. Then
\[
|K|\times |L|
=
\bigcup_{\sigma\in K,\ \tau\in L} (\sigma\times\tau).
\]
For each pair \((\sigma,\tau)\), triangulate \(\sigma\times\tau\) by the staircase construction above.

We claim that these triangulations fit together to form a simplicial complex on the whole product.

Indeed, if
\[
\sigma'\le \sigma,
\qquad
\tau'\le \tau,
\]
then
\[
\sigma'\times\tau' \subseteq \sigma\times\tau,
\]
and the staircase triangulation of \(\sigma'\times\tau'\) is exactly the restriction of the staircase triangulation of \(\sigma\times\tau\). Therefore the triangulations agree on common faces.

Thus the simplices coming from all products \(\sigma\times\tau\) together form a simplicial complex \(M\subseteq \mathbb R^{m+n}\).

Its polyhedron is exactly
\[
|M|=|K|\times |L|.
\]

\par

\textbf{Example 3: \(S^1\times I\).}

If \(K\) triangulates a circle and \(L\) triangulates an interval, then
\[
|K|\times |L| \cong S^1\times I
\]
is a cylinder. The theorem shows that this cylinder is the polyhedron of a simplicial complex.

\par

\textbf{Example 4: \(S^1\times S^1\).}

If \(K\) and \(L\) both triangulate \(S^1\), then
\[
|K|\times |L| \cong S^1\times S^1=T^2.
\]
So the torus admits a triangulation obtained from a product triangulation.


\medskip
\noindent\textbf{Conclusion}\par
\smallskip


We have shown that the product of two simplices admits a staircase triangulation, and that these triangulations are compatible on faces. Therefore for arbitrary simplicial complexes \(K\subseteq \mathbb R^m\) and \(L\subseteq \mathbb R^n\), the product
\[
|K|\times |L| \subseteq \mathbb R^{m+n}
\]
is the polyhedron of a simplicial complex.

Hence
\[
\boxed{|K|\times |L|\text{ is the polyhedron of a simplicial complex.}}
\]
\end{solution}

\begin{problem}[CP-VIII-0036 --- Torus endomorphism induced by an integer matrix]
Let \(A\) be a \(2\times 2\) matrix with entries in \(\mathbb Z\). Show that the linear map
\[
A:\mathbb R^2\to \mathbb R^2
\]
preserves the equivalence relation
\[
(a,b)\sim (a',b')
\quad\Longleftrightarrow\quad
(a-a',\,b-b')\in \mathbb Z^2,
\]
and therefore induces a continuous map
\[
f_A:T\to T
\]
on the torus
\[
T=\mathbb R^2/\mathbb Z^2.
\]
Compute the induced homomorphisms
\[
(f_A)_*:H_k(T)\to H_k(T)
\]
on the homology of \(T\).

\par
\end{problem}

\begin{solution}
Recall that the torus is the quotient
\[
T=\mathbb R^2/\mathbb Z^2,
\]
where two points of \(\mathbb R^2\) are identified if they differ by an element of the lattice \(\mathbb Z^2\).

We first show that \(A\) respects this equivalence relation. Suppose
\[
(a,b)\sim (a',b').
\]
Then
\[
(a-a',\,b-b')\in \mathbb Z^2.
\]
Writing \(x=(a,b)\) and \(y=(a',b')\), this means \(x-y\in \mathbb Z^2\). Since \(A\) has integer entries, it sends integer vectors to integer vectors:
\[
A(\mathbb Z^2)\subseteq \mathbb Z^2.
\]
Hence
\[
A(x-y)=Ax-Ay\in \mathbb Z^2,
\]
so \(Ax\sim Ay\). Therefore \(A\) preserves the quotient relation and induces a well-defined map
\[
f_A:T\to T,
\qquad
f_A([x])=[Ax].
\]
Because the quotient map \(\pi:\mathbb R^2\to T\) is continuous and \(A\) is continuous, the induced map \(f_A\) is continuous.

We now compute the induced map on homology.

For the torus,
\[
H_0(T)\cong \mathbb Z,\qquad
H_1(T)\cong \mathbb Z^2,\qquad
H_2(T)\cong \mathbb Z,
\]
and \(H_k(T)=0\) for all \(k\ge 3\).


\medskip
\noindent\textbf{Action on \(H_0(T)\)}\par
\smallskip

Since \(T\) is connected, any continuous self-map \(T\to T\) induces the identity on \(H_0\). Thus
\[
(f_A)_0=\mathrm{id}_{\mathbb Z}.
\]


\medskip
\noindent\textbf{Action on \(H_1(T)\)}\par
\smallskip

Choose the standard generators of \(H_1(T)\cong \mathbb Z^2\) represented by the two coordinate circles:
\[
\alpha(t)=[(t,0)],
\qquad
\beta(t)=[(0,t)],
\qquad t\in \mathbb R/\mathbb Z.
\]
Write
\[
A=
\begin{pmatrix}
	a & b\\
	c & d
\end{pmatrix}.
\]
Then
\[
A(t,0)=(at,ct),
\qquad
A(0,t)=(bt,dt).
\]
Therefore the induced loops satisfy
\[
f_A\circ \alpha \sim a\alpha+c\beta,
\qquad
f_A\circ \beta \sim b\alpha+d\beta.
\]
Hence, with respect to the basis \([\alpha],[\beta]\) of \(H_1(T)\),
\[
(f_A)_1=
\begin{pmatrix}
	a & b\\
	c & d
\end{pmatrix}
=A.
\]


\medskip
\noindent\textbf{Action on \(H_2(T)\)}\par
\smallskip

Since \(H_2(T)\cong \mathbb Z\), the induced map on \(H_2\) must be multiplication by the degree of \(f_A\). For a torus map induced by the linear map \(A\), the degree is
\[
\deg(f_A)=\det(A).
\]
Thus
\[
(f_A)_2=\times \det(A).
\]

So the induced maps on homology are:
\[
(f_A)_*:H_0(T)\to H_0(T)
\quad\text{is}\quad
\mathrm{id}_{\mathbb Z},
\]
\[
(f_A)_*:H_1(T)\to H_1(T)
\quad\text{is}\quad
A,
\]
\[
(f_A)_*:H_2(T)\to H_2(T)
\quad\text{is multiplication by}\quad
\det(A),
\]
and for \(k\ge 3\),
\[
H_k(T)=0.
\]

Thus the complete answer is
\[
(f_A)_0=\mathrm{id},\qquad
(f_A)_1=A,\qquad
(f_A)_2=\times \det(A).
\]
\end{solution}

\begin{problem}[CP-VIII-0037 --- Connected sum of an orientable and a nonorientable surface]
For triangulated surfaces \(X\) and \(Y\), let \(X\#Y\) be the surface obtained by removing a \(2\)-simplex from each and gluing the two boundary circles formed by \(\partial\Delta^2\). Let \(\Sigma_g\) denote the closed orientable surface of genus \(g\), and let \(S_n\) denote the closed nonorientable surface of genus \(n\). Use the Mayer--Vietoris sequence to compute the homology of
\[
\Sigma_g\# S_n,
\]
and deduce that
\[
\Sigma_g\# S_n\cong S_{n+2g}.
\]

\par
\end{problem}

\begin{solution}
Let
\[
M=\Sigma_g\# S_n.
\]
By construction, we may write
\[
M=U\cup V,
\]
where
\[
U=\Sigma_g\setminus \operatorname{int}(D^2),
\qquad
V=S_n\setminus \operatorname{int}(D^2),
\]
and the intersection \(U\cap V\) is a collar neighborhood of the glued boundary circle. Therefore
\[
U\cap V\simeq S^1.
\]

We now describe the homology of the pieces.


\medskip
\noindent\textbf{Step 1: homology of \(U\), \(V\), and \(U\cap V\)}\par
\smallskip

The space \(U\) is an orientable surface of genus \(g\) with one boundary component. It deformation retracts onto a wedge of \(2g\) circles, so
\[
H_0(U)\cong \mathbb Z,
\qquad
H_1(U)\cong \mathbb Z^{2g},
\qquad
H_i(U)=0\ \text{for }i\ge 2.
\]

Similarly, \(V\) is a nonorientable surface of genus \(n\) with one boundary component. It deformation retracts onto a wedge of \(n\) circles, so
\[
H_0(V)\cong \mathbb Z,
\qquad
H_1(V)\cong \mathbb Z^{n},
\qquad
H_i(V)=0\ \text{for }i\ge 2.
\]

Finally, since \(U\cap V\simeq S^1\),
\[
H_0(U\cap V)\cong \mathbb Z,
\qquad
H_1(U\cap V)\cong \mathbb Z,
\qquad
H_i(U\cap V)=0\ \text{for }i\ge 2.
\]


\medskip
\noindent\textbf{Step 2: the Mayer--Vietoris sequence}\par
\smallskip

The relevant part of the Mayer--Vietoris sequence for \(M=U\cup V\) is
\[
\begin{aligned}
0&\to H_2(M)\to H_1(U\cap V)
\xrightarrow{\alpha}H_1(U)\oplus H_1(V)\\
&\to H_1(M)\to H_0(U\cap V)
\to H_0(U)\oplus H_0(V)\\
&\to H_0(M)\to 0.
\end{aligned}
\]
Substituting the homology groups above gives
\[
0\to H_2(M)\to \mathbb Z
\overset{\alpha}{\longrightarrow}
\mathbb Z^{2g}\oplus \mathbb Z^n
\longrightarrow H_1(M)
\longrightarrow \mathbb Z
\longrightarrow \mathbb Z\oplus \mathbb Z
\longrightarrow H_0(M)\to 0.
\]

Since \(U\), \(V\), and \(U\cap V\) are all connected, the map
\[
H_0(U\cap V)\to H_0(U)\oplus H_0(V)
\]
is the diagonal map \(1\mapsto (1,1)\), hence injective. It follows that
\[
H_0(M)\cong \mathbb Z,
\]
and exactness reduces the computation of \(H_1(M)\) and \(H_2(M)\) to the short exact sequence
\[
0\to H_2(M)\to \mathbb Z
\overset{\alpha}{\longrightarrow}
\mathbb Z^{2g}\oplus \mathbb Z^n
\longrightarrow H_1(M)\to 0.
\]

Thus the key point is to determine the map \(\alpha\).


\medskip
\noindent\textbf{Step 3: identify the boundary class}\par
\smallskip

The generator of \(H_1(U\cap V)\cong \mathbb Z\) is represented by the common boundary circle.


\medskip
\noindent\textbf{Its image in \(H_1(U)\).}\par
\smallskip

For the orientable surface \(U=\Sigma_g\setminus \operatorname{int}(D^2)\), the boundary loop is homologous to zero. One way to see this is that
\[
\pi_1(U)=
\left\langle
a_1,b_1,\dots,a_g,b_g
\right\rangle
\]
with the boundary represented by the word
\[
[a_1,b_1]\cdots [a_g,b_g].
\]
After abelianization, every commutator becomes trivial, so the boundary class maps to \(0\in H_1(U)\).


\medskip
\noindent\textbf{Its image in \(H_1(V)\).}\par
\smallskip

For the nonorientable surface \(V=S_n\setminus \operatorname{int}(D^2)\), one has a standard presentation
\[
\pi_1(V)=\langle x_1,\dots,x_n\rangle,
\]
and the boundary loop is represented by
\[
x_1^2x_2^2\cdots x_n^2.
\]
After abelianization, this becomes
\[
2x_1+2x_2+\cdots+2x_n
\in H_1(V)\cong \mathbb Z^n.
\]

Hence the map
\[
\alpha:H_1(U\cap V)\to H_1(U)\oplus H_1(V)
\]
is given by
\[
\alpha(1)=\bigl(0,\;2x_1+2x_2+\cdots+2x_n\bigr).
\]

Now \(x_1+\cdots+x_n\) is a primitive element of \(\mathbb Z^n\), so after a change of basis in \(\mathbb Z^n\), the map \(\alpha\) becomes
\[
\mathbb Z\to \mathbb Z^{2g}\oplus \mathbb Z^n,
\qquad
1\mapsto (0,2,0,\dots,0).
\]
This map is injective. Therefore
\[
H_2(M)=0.
\]

Moreover,
\[
H_1(M)\cong
\bigl(\mathbb Z^{2g}\oplus \mathbb Z^n\bigr)\big/\langle (0,2,0,\dots,0)\rangle.
\]
Thus
\[
H_1(M)\cong \mathbb Z^{2g+n-1}\oplus \mathbb Z/2.
\]

So we have proved
\[
H_0(\Sigma_g\# S_n)\cong \mathbb Z,
\]
\[
H_1(\Sigma_g\# S_n)\cong \mathbb Z^{2g+n-1}\oplus \mathbb Z/2,
\]
\[
H_2(\Sigma_g\# S_n)=0,
\]
and all higher homology groups vanish.


\medskip
\noindent\textbf{Step 4: identify the surface}\par
\smallskip

For the closed nonorientable surface \(S_m\) of genus \(m\), the homology groups are
\[
H_0(S_m)\cong \mathbb Z,
\qquad
H_1(S_m)\cong \mathbb Z^{m-1}\oplus \mathbb Z/2,
\qquad
H_2(S_m)=0.
\]
If we put
\[
m=n+2g,
\]
then
\[
H_1(S_{n+2g})\cong \mathbb Z^{n+2g-1}\oplus \mathbb Z/2,
\]
which agrees exactly with the homology computed above for \(\Sigma_g\# S_n\).

Thus \(\Sigma_g\# S_n\) has the same homology as \(S_{n+2g}\). By the classification theorem for closed connected surfaces, the connected sum of an orientable genus-\(g\) surface with a nonorientable genus-\(n\) surface is the nonorientable genus-\((n+2g)\) surface. Therefore
\[
\Sigma_g\# S_n\cong S_{n+2g}.
\]

So the final result is
\[
H_0(\Sigma_g\# S_n)\cong \mathbb Z,
\qquad
H_1(\Sigma_g\# S_n)\cong \mathbb Z^{n+2g-1}\oplus \mathbb Z/2,
\qquad
H_2(\Sigma_g\# S_n)=0,
\]
and hence
\[
\boxed{\Sigma_g\# S_n\cong S_{n+2g}.}
\]
\end{solution}

\begin{problem}[CP-VIII-0038 --- Finite-sheeted coverings, triangulations, Euler characteristic, and coverings of surfaces]
Let
\[
p:\widetilde X \to X
\]
be a finite-sheeted covering space, and let
\[
h:|K|\to X
\]
be a triangulation. Show that there is an \(r\ge 1\) and a triangulation
\[
g:|L|\to \widetilde X
\]
such that the composition
\[
h^{-1}\circ p\circ g:|L|\to |K^{(r)}|
\]
is a simplicial map.

If \(p\) has \(n\) sheets, show that
\[
\chi(\widetilde X)=n\,\chi(X).
\]
Hence show that \(\Sigma_g\) is a covering space of \(\Sigma_h\) if and only if
\[
\frac{1-g}{1-h}
\]
is an integer.

\par
\end{problem}

\begin{solution}
We proceed in three steps.

\par

\textbf{Step 1: after subdivision, the covering becomes simplicial.}

Because \(p:\widetilde X\to X\) is a covering map, each point of \(X\) has an evenly covered neighborhood. Since \(h:|K|\to X\) is a triangulation and \(|K|\) is compact on each simplex, after sufficiently many barycentric subdivisions there exists \(r\ge 1\) such that for every simplex \(\sigma\) of \(K^{(r)}\), the set \(h(|\sigma|)\) lies inside some evenly covered open set \(U_\sigma\subset X\).

For each such simplex \(\sigma\), the preimage \(p^{-1}(U_\sigma)\) is a disjoint union of open sets
\[
p^{-1}(U_\sigma)=\bigsqcup_\alpha V_{\sigma,\alpha},
\]
and each restriction
\[
p|_{V_{\sigma,\alpha}}:V_{\sigma,\alpha}\to U_\sigma
\]
is a homeomorphism.

Since \(|\sigma|\) is contractible, each lift of \(h|_{|\sigma|}:|\sigma|\to U_\sigma\) to \(\widetilde X\) lands in one sheet \(V_{\sigma,\alpha}\), and under the inverse of
\[
p|_{V_{\sigma,\alpha}},
\]
it identifies with a copy of \(|\sigma|\). Doing this for every simplex of \(K^{(r)}\), and gluing along common faces, produces a simplicial complex \(L\) together with a triangulation
\[
g:|L|\to \widetilde X
\]
such that
\[
h^{-1}\circ p\circ g:|L|\to |K^{(r)}|
\]
is simplicial.

Thus the covering map may be made simplicial after sufficiently fine subdivision of the triangulation downstairs.

\par

\textbf{Step 2: Euler characteristic multiplies by the number of sheets.}

Assume now that \(p\) has exactly \(n\) sheets. Since every simplex of \(|K^{(r)}|\) lies in an evenly covered set, each simplex of \(K^{(r)}\) has exactly \(n\) distinct lifts to simplices of \(L\). Therefore, for each dimension \(q\),
\[
f_q(L)=n\,f_q\bigl(K^{(r)}\bigr),
\]
where \(f_q(\cdot)\) denotes the number of \(q\)-simplices.

Hence
\[
\chi(\widetilde X)
=\sum_{q\ge 0}(-1)^q f_q(L)
=\sum_{q\ge 0}(-1)^q n\,f_q\bigl(K^{(r)}\bigr)
=n\sum_{q\ge 0}(-1)^q f_q\bigl(K^{(r)}\bigr).
\]
Since barycentric subdivision does not change Euler characteristic,
\[
\chi\bigl(|K^{(r)}|\bigr)=\chi(|K|)=\chi(X).
\]
Therefore
\[
\boxed{\chi(\widetilde X)=n\,\chi(X).}
\]

\par

\textbf{Step 3: application to orientable closed surfaces.}

Recall that for the closed orientable surface \(\Sigma_m\) of genus \(m\),
\[
\chi(\Sigma_m)=2-2m.
\]
If \(\Sigma_g\to \Sigma_h\) is an \(n\)-sheeted covering, then by the formula just proved,
\[
2-2g=n(2-2h).
\]
Dividing by \(2\),
\[
1-g=n(1-h),
\]
so
\[
\boxed{\frac{1-g}{1-h}=n\in \mathbb Z_{>0}.}
\]
Thus integrality of \((1-g)/(1-h)\) is necessary.

Conversely, assume
\[
1-g=n(1-h)
\]
for some positive integer \(n\), equivalently
\[
g=1+n(h-1).
\]
We explain a standard construction of an \(n\)-sheeted covering
\[
\Sigma_g\to \Sigma_h.
\]

Start with a torus \(T^2\), viewed as \(\mathbb R^2/\mathbb Z^2\). Let \(\mathbb Z/n\) act freely on \(T^2\) by translation in one direction:
\[
[j]\cdot [x,y]=\left[x+\frac{j}{n},\,y\right].
\]
This action is free, and the quotient is again a torus.

Choose \(n\) disjoint open disks in \(T^2\) forming one orbit under this action. Remove them. The resulting surface is a torus with \(n\) boundary components. Now attach to each boundary circle a copy of the surface \(\Sigma_{h-1,1}\), meaning the orientable surface of genus \(h-1\) with one boundary component. Let \(\mathbb Z/n\) cyclically permute these \(n\) attached copies.

The resulting closed orientable surface has genus
\[
1+n(h-1)=g,
\]
and the quotient by the \(\mathbb Z/n\)-action has genus
\[
1+(h-1)=h.
\]
Hence the quotient map gives an \(n\)-sheeted covering
\[
\Sigma_g\to \Sigma_h.
\]

So for \(h\neq 1\),
\[
\boxed{\Sigma_g \text{ covers } \Sigma_h \iff \frac{1-g}{1-h}\in \mathbb Z.}
\]

\par

\textbf{Remark.}
When \(h=1\), the denominator \(1-h\) is \(0\), so the displayed ratio is not defined. In that case one uses Euler characteristic directly:
\[
\chi(\Sigma_g)=n\chi(\Sigma_1)=n\cdot 0=0,
\]
hence \(g=1\). Thus only the torus covers the torus.
\end{solution}

\begin{problem}[CP-VIII-0039 --- Coverings of spaces whose universal cover is an even-dimensional sphere]
Let
\[
p:S^{2k}\to X
\]
be a covering map, let
\[
G=\pi_1(X,x_0),
\]
and recall that \(G\) then acts freely on \(S^{2k}\). Show that for any
\[
g\in G\setminus\{e\},
\]
the map
\[
g_*:H_{2k}(S^{2k})\to H_{2k}(S^{2k})
\]
is multiplication by \(-1\). Deduce that \(G\) is either trivial or \(\mathbb Z/2\), and that \(\mathbb{RP}^{2k}\) is not a proper covering space of any other space.

\par
\end{problem}

\begin{solution}
Since \(S^{2k}\) is simply connected for \(k\ge 1\), the covering
\[
p:S^{2k}\to X
\]
is the universal covering of \(X\). Therefore the fundamental group
\[
G=\pi_1(X,x_0)
\]
identifies with the group of deck transformations of \(p\), acting freely on \(S^{2k}\).

Let \(g\in G\setminus\{e\}\). Then \(g:S^{2k}\to S^{2k}\) is a fixed-point-free homeomorphism.

\par

\textbf{Step 1: compute the induced map on top homology.}

For any continuous map
\[
f:S^{2k}\to S^{2k},
\]
the Lefschetz number is
\[
L(f)=\sum_i (-1)^i \operatorname{tr}\bigl(f_*:H_i(S^{2k})\to H_i(S^{2k})\bigr).
\]
The only nonzero homology groups of \(S^{2k}\) are
\[
H_0(S^{2k})\cong \mathbb Z,
\qquad
H_{2k}(S^{2k})\cong \mathbb Z.
\]
Hence
\[
L(f)=\operatorname{tr}(f_*|_{H_0})+\operatorname{tr}(f_*|_{H_{2k}}).
\]
Since \(S^{2k}\) is connected, every self-map induces the identity on \(H_0\), so
\[
\operatorname{tr}(f_*|_{H_0})=1.
\]
On top homology, the induced map is multiplication by the degree:
\[
f_*|_{H_{2k}}=\deg(f).
\]
Therefore
\[
L(f)=1+\deg(f).
\]

Apply this to the deck transformation \(g\). Since \(g\) has no fixed point, the Lefschetz fixed point theorem implies
\[
L(g)=0.
\]
Thus
\[
0=1+\deg(g),
\]
so
\[
\deg(g)=-1.
\]
Equivalently,
\[
\boxed{g_*:H_{2k}(S^{2k})\to H_{2k}(S^{2k}) \text{ is multiplication by } -1.}
\]

\par

\textbf{Step 2: deduce that \(G\) is trivial or \(\mathbb Z/2\).}

We now have a homomorphism
\[
\rho:G\to \operatorname{Aut}\bigl(H_{2k}(S^{2k})\bigr)\cong \{\pm 1\},
\]
and every nontrivial element of \(G\) is sent to \(-1\).

Suppose \(g,h\in G\setminus\{e\}\). Then
\[
\rho(gh)=\rho(g)\rho(h)=(-1)(-1)=1.
\]
But the only element of \(G\) that can act by \(+1\) is the identity, because every nontrivial element acts by \(-1\). Hence
\[
gh=e.
\]
In particular, taking \(h=g\), we get
\[
g^2=e.
\]
So every nontrivial element has order \(2\), and there can be at most one nontrivial element. Therefore
\[
\boxed{G\cong \{e\}\ \text{or}\ \mathbb Z/2.}
\]

\par

\textbf{Step 3: \(\mathbb{RP}^{2k}\) is not a proper covering of another space.}

Assume for contradiction that there exists a proper covering
\[
q:\mathbb{RP}^{2k}\to Y.
\]
Compose it with the universal double covering
\[
\pi:S^{2k}\to \mathbb{RP}^{2k}.
\]
Then
\[
q\circ \pi:S^{2k}\to Y
\]
is a covering map. Since \(S^{2k}\) is simply connected, it is the universal covering of \(Y\), and thus
\[
\pi_1(Y)
\]
must be either trivial or \(\mathbb Z/2\), by the result of Step 2.

On the other hand,
\[
\pi_1(\mathbb{RP}^{2k})\cong \mathbb Z/2.
\]
For connected coverings, intermediate coverings correspond to subgroups of the fundamental group of the base. Since
\[
q:\mathbb{RP}^{2k}\to Y
\]
is assumed to be a proper covering, the subgroup corresponding to \(\mathbb{RP}^{2k}\) must be a proper subgroup of \(\pi_1(Y)\). But this subgroup is isomorphic to \(\mathbb Z/2\), and neither of the groups
\[
\{e\},\qquad \mathbb Z/2
\]
has a proper subgroup isomorphic to \(\mathbb Z/2\).

This is impossible. Therefore no such proper covering exists. Hence
\[
\boxed{\mathbb{RP}^{2k}\text{ is not a proper covering space of any other space.}}
\]
\end{solution}

\begin{problem}[CP-VIII-0040 --- A simplicial automorphism, its fixed set, and the Lefschetz number]
If
\[
f:K\to K
\]
is a simplicial isomorphism, let
\[
X\subset |K|
\]
be the fixed set of \(|f|\), that is,
\[
X=\{x\in |K| \mid |f|(x)=x\}.
\]
Show that the Lefschetz number \(L(f)\) is equal to \(\chi(X)\).

\[
\textit{Hint: barycentrically subdivide \(K\) so that \(X\) is the polyhedron of a sub-simplicial complex.}
\]

\par
\end{problem}

\begin{solution}
Let \(K'=\operatorname{Sd}(K)\) be the barycentric subdivision of \(K\). The simplicial isomorphism
\[
f:K\to K
\]
induces a simplicial isomorphism
\[
f':K'\to K'.
\]
Since barycentric subdivision does not change the underlying polyhedron and does not change the induced map on homology up to canonical identification, we have
\[
L(f')=L(f).
\]
Thus it suffices to prove that
\[
L(f')=\chi(X).
\]

\par

\textbf{Step 1: describe the fixed set as a subcomplex of the barycentric subdivision.}

Recall that the vertices of \(K'\) are the barycenters \(b_\sigma\) of simplices \(\sigma\) of \(K\). The induced simplicial map \(f'\) acts on vertices by
\[
f'(b_\sigma)=b_{f(\sigma)}.
\]
Therefore a vertex \(b_\sigma\) of \(K'\) is fixed by \(f'\) if and only if
\[
f(\sigma)=\sigma.
\]

A simplex of \(K'\) is a chain
\[
\sigma_0<\sigma_1<\cdots<\sigma_q
\]
of simplices of \(K\). This simplex is fixed pointwise by \(f'\) precisely when every \(\sigma_i\) is fixed setwise by \(f\), because then each vertex \(b_{\sigma_i}\) is fixed.

Hence the collection of all simplices of \(K'\) whose vertices are fixed by \(f'\) forms a subcomplex \(L\subset K'\). Its polyhedron is exactly the fixed-point set:
\[
|L|=X.
\]
Indeed, if a point of \(|K|\) is fixed by \(|f|\), then in the barycentric subdivision it lies in a simplex whose vertices correspond to \(f\)-invariant simplices, and conversely every point of such a simplex is fixed.

Thus
\[
X=|L|
\]
for a subcomplex \(L\subset K'\).

\par

\textbf{Step 2: compute the Lefschetz number on the chain level.}

Let \(C_q(K')\) be the simplicial chain group in dimension \(q\), with basis the oriented \(q\)-simplices of \(K'\). The map \(f'\) induces a chain map
\[
f'_q:C_q(K')\to C_q(K').
\]

The trace of \(f'_q\) is determined as follows.

If an oriented \(q\)-simplex \(\tau\) is not carried to itself by \(f'\), then \(\tau\) is sent to a different basis element, so its contribution to the trace is \(0\).

If \(\tau\) is fixed by \(f'\), then every vertex of \(\tau\) is fixed. Hence \(f'\) fixes \(\tau\) pointwise, not merely setwise, and therefore preserves its orientation. So the contribution of such a simplex to the trace is \(+1\).

Thus
\[
\operatorname{tr}(f'_q)=\#\{\text{\(q\)-simplices of }L\}.
\]

Now the Lefschetz number is
\[
L(f')=\sum_{q\ge 0}(-1)^q \operatorname{tr}(f'_q)
=\sum_{q\ge 0}(-1)^q \#\{\text{\(q\)-simplices of }L\}.
\]
But this alternating sum is exactly the Euler characteristic of \(L\):
\[
L(f')=\chi(L).
\]
Since \(X=|L|\), Euler characteristic is preserved under passing from a finite simplicial complex to its polyhedron, so
\[
\chi(L)=\chi(X).
\]
Therefore
\[
L(f)=L(f')=\chi(X).
\]

Hence
\[
\boxed{L(f)=\chi(X).}
\]
\end{solution}

\begin{problem}[CP-VIII-0041 --- A singular cycle in \(\mathbb{R}\) that is a boundary]
Let
\[
\sigma_1,\sigma_2:[0,1]\to \mathbb{R}
\]
be given by
\[
\sigma_1(x)=x,
\qquad
\sigma_2(x)=1-x.
\]
Identifying \([0,1]\) with the standard \(1\)-simplex \(\Delta^1\), the singular \(1\)-chain
\[
\sigma_1+\sigma_2\in C_1(\mathbb{R})
\]
is a cycle. Find an element \(x\in C_2(\mathbb{R})\) such that
\[
\partial x=\sigma_1+\sigma_2.
\]
\end{problem}

\begin{solution}

\medskip
\noindent\textbf{Idea of the construction}\par
\smallskip


\textbf{The chain \(\sigma_1+\sigma_2\) goes from \(0\) to \(1\) and then immediately back from \(1\) to \(0\), so geometrically it should bound a degenerate \(2\)-simplex.}

Since the target space is just \(\mathbb{R}\), any singular \(2\)-simplex \(\Delta^2\to \mathbb{R}\) must collapse the triangle onto an interval. That is enough for singular homology: the simplex may be degenerate, and its boundary may still produce a nontrivial \(1\)-chain.

We will exhibit explicit singular \(2\)-simplices whose boundaries combine to \(\sigma_1+\sigma_2\).


\medskip
\noindent\textbf{The singular \(2\)-simplices}\par
\smallskip


\textbf{We define two singular \(2\)-simplices on the standard triangle \(\Delta^2\).}

Let
\[
\Delta^2=\{(t_0,t_1,t_2)\in \mathbb{R}^3 \mid t_0,t_1,t_2\ge 0,\ t_0+t_1+t_2=1\}.
\]
Define maps
\[
\tau_1,\tau_2:\Delta^2\to \mathbb{R}
\]
by
\[
\tau_1(t_0,t_1,t_2)=t_1,
\qquad
\tau_2(t_0,t_1,t_2)=t_2.
\]

These are singular \(2\)-simplices in \(\mathbb{R}\). In vertex notation one may write
\[
\tau_1=[0,1,0],
\qquad
\tau_2=[0,0,1].
\]
This means:

\begin{itemize}
	\item \(\tau_1\) sends the three vertices of \(\Delta^2\) to \(0,1,0\),
	\item \(\tau_2\) sends the three vertices of \(\Delta^2\) to \(0,0,1\).
\end{itemize}


\medskip
\noindent\textbf{Boundary formula}\par
\smallskip


\textbf{For a singular \(2\)-simplex \([a,b,c]\), the boundary is}
\[
\partial[a,b,c]=[b,c]-[a,c]+[a,b].
\]

We now apply this formula to \(\tau_1\) and \(\tau_2\).


\medskip
\noindent\textbf{Computing \(\partial \tau_1\)}\par
\smallskip


\textbf{First compute the boundary of the simplex \([0,1,0]\).}

We have
\[
\partial\tau_1
=
\partial[0,1,0]
=
[1,0]-[0,0]+[0,1].
\]

Now interpret these edges:

\begin{itemize}
	\item \([0,1]\) is the path \(\sigma_1\),
	\item \([1,0]\) is the path \(\sigma_2\),
	\item \([0,0]\) is the constant singular \(1\)-simplex at \(0\), denote it by \(c_0\).
\end{itemize}

Therefore
\[
\partial\tau_1=\sigma_2-c_0+\sigma_1.
\]


\medskip
\noindent\textbf{Computing \(\partial \tau_2\)}\par
\smallskip


\textbf{Now compute the boundary of the simplex \([0,0,1]\).}

We get
\[
\partial\tau_2
=
\partial[0,0,1]
=
[0,1]-[0,1]+[0,0].
\]
Hence
\[
\partial\tau_2=c_0.
\]


\medskip
\noindent\textbf{Choosing the required \(2\)-chain}\par
\smallskip


\textbf{Now we add the two \(2\)-simplices so that the constant edge cancels.}

Set
\[
x=\tau_1+\tau_2.
\]
Then
\[
\partial x
=
\partial\tau_1+\partial\tau_2
=
(\sigma_2-c_0+\sigma_1)+c_0
=
\sigma_1+\sigma_2.
\]

Thus \(x\in C_2(\mathbb{R})\) satisfies exactly the required condition.


\medskip
\noindent\textbf{Final answer}\par
\smallskip


\textbf{A valid singular \(2\)-chain is}
\[
\boxed{x=\tau_1+\tau_2}
\]
where
\[
\tau_1(t_0,t_1,t_2)=t_1,
\qquad
\tau_2(t_0,t_1,t_2)=t_2.
\]
Its boundary is
\[
\boxed{\partial x=\sigma_1+\sigma_2.}
\]

Equivalently, in vertex notation,
\[
\boxed{x=[0,1,0]+[0,0,1].}
\]


\medskip
\noindent\textbf{Geometric comment}\par
\smallskip


\textbf{This example shows well that in singular homology a simplex need not be nondegenerate.}

The \(2\)-simplices above are collapsed triangles in \(\mathbb{R}\), but they are still legitimate singular simplices, and their boundaries can produce ordinary paths. This is one reason singular homology is flexible: it allows arbitrary continuous maps from simplices, not only embedded geometric simplices.
\end{solution}

\begin{problem}[CP-VIII-0042 --- Cone contractibility and homology of the suspension]
Let \(X\) be a topological space. The \emph{cone} on \(X\) is defined by
\[
CX := X\times I /(X\times \{0\}).
\]
Show that \(CX\) is contractible.

The \emph{suspension} of \(X\) is defined by
\[
\Sigma X := CX /(X\times \{1\}).
\]
Express \(\widetilde H_*(\Sigma X)\) in terms of \(\widetilde H_*(X)\).
\end{problem}

\begin{solution}
\noindent
We first prove that the cone \(CX\) is contractible.

By definition, all points of \(X\times \{0\}\) are identified to a single point, called the cone point. Denote the equivalence class of \((x,t)\in X\times I\) by \([x,t]\). Consider the homotopy
\[
F:CX\times I \longrightarrow CX,
\qquad
F([x,t],s)=[x,(1-s)t].
\]
This is well-defined, because if \(t=0\), then \([x,0]\) is already the cone point, independent of \(x\).

Now:
\[
F([x,t],0)=[x,t],
\]
so \(F(-,0)=\mathrm{id}_{CX}\), and
\[
F([x,t],1)=[x,0],
\]
which is the cone point for every \([x,t]\in CX\). Therefore \(F\) is a contraction of \(CX\) to a point. Hence
\[
CX \simeq *.
\]
So \(CX\) is contractible.

Next, let
\[
A:=X\times \{1\}\subset CX.
\]
This subspace is naturally homeomorphic to \(X\). By the definition of suspension,
\[
\Sigma X = CX/A.
\]
For reduced homology of a quotient, we have
\[
\widetilde H_n(\Sigma X)\cong H_n(CX,A).
\]
Now consider the long exact sequence of the pair \((CX,A)\):
\[
\cdots \to H_n(A)\to H_n(CX)\to H_n(CX,A)\to H_{n-1}(A)\to H_{n-1}(CX)\to \cdots
\]
Since \(CX\) is contractible,
\[
\widetilde H_n(CX)=0
\qquad \text{for all } n,
\]
so the exact sequence reduces to isomorphisms
\[
H_n(CX,A)\cong \widetilde H_{n-1}(A).
\]
Because \(A\cong X\), this becomes
\[
H_n(CX,A)\cong \widetilde H_{n-1}(X).
\]
Combining this with \(\widetilde H_n(\Sigma X)\cong H_n(CX,A)\), we obtain
\[
\widetilde H_n(\Sigma X)\cong \widetilde H_{n-1}(X)
\qquad \text{for all } n.
\]

Thus the reduced homology of the suspension is the reduced homology of \(X\), shifted up by one degree.


\medskip
\noindent\textbf{Conclusion}\par
\smallskip


We have shown that
\[
CX \text{ is contractible}
\]
and
\[
\boxed{\widetilde H_n(\Sigma X)\cong \widetilde H_{n-1}(X)\quad \text{for all } n.}
\]
\end{solution}

\begin{problem}[CP-VIII-0043 --- Decomposition of \(S^{n+m+1}\) and the homology of \(S^n \times S^m\)]
Show that \(S^{n+m+1}\) can be decomposed as the union of
\[
S^n \times D^{m+1}
\quad\text{and}\quad
D^{n+1} \times S^m
\]
along their common boundary
\[
S^n \times S^m.
\]
Compute
\[
H_*(S^n \times S^m)
\quad\text{and}\quad
H_*(D^{n+1} \times S^m,\; S^n \times S^m).
\]
\end{problem}

\begin{solution}
A very clean way to see the required decomposition is to start with the product
\[
D^{n+1} \times D^{m+1}.
\]
Since \(D^{n+1} \times D^{m+1}\) is homeomorphic to \(D^{n+m+2}\), its boundary is homeomorphic to \(S^{n+m+1}\). On the other hand,
\[
\partial\bigl(D^{n+1} \times D^{m+1}\bigr)
=
\bigl(\partial D^{n+1} \times D^{m+1}\bigr)
\cup
\bigl(D^{n+1} \times \partial D^{m+1}\bigr).
\]
Because \(\partial D^{n+1} = S^n\) and \(\partial D^{m+1} = S^m\), this becomes
\[
\partial\bigl(D^{n+1} \times D^{m+1}\bigr)
=
\bigl(S^n \times D^{m+1}\bigr)
\cup
\bigl(D^{n+1} \times S^m\bigr).
\]
Their intersection is
\[
\bigl(S^n \times D^{m+1}\bigr)\cap\bigl(D^{n+1} \times S^m\bigr)
=
S^n \times S^m.
\]
Hence
\[
S^{n+m+1}
\cong
\bigl(S^n \times D^{m+1}\bigr)\cup_{\,S^n \times S^m}\bigl(D^{n+1} \times S^m\bigr).
\]


\medskip
\noindent\textbf{Computation of \(H_*(S^n \times S^m)\)}\par
\smallskip


Since the homology groups of spheres are free abelian, the K\"unneth formula has no \(\operatorname{Tor}\)-term, so
\[
H_k(S^n \times S^m)
\cong
\bigoplus_{p+q=k} H_p(S^n)\otimes H_q(S^m).
\]
Now
\[
H_p(S^n)\cong
\begin{cases}
	\mathbb Z, & p=0,n,\\
	0, & \text{otherwise},
\end{cases}
\qquad
H_q(S^m)\cong
\begin{cases}
	\mathbb Z, & q=0,m,\\
	0, & \text{otherwise}.
\end{cases}
\]
Therefore the only nonzero contributions occur for
\[
(p,q)=(0,0),\ (n,0),\ (0,m),\ (n,m).
\]
Thus
\[
H_k(S^n \times S^m)\cong
\begin{cases}
	\mathbb Z, & k=0,\\[1mm]
	\mathbb Z, & k=n \text{ and } n\neq m,\\[1mm]
	\mathbb Z, & k=m \text{ and } n\neq m,\\[1mm]
	\mathbb Z\oplus \mathbb Z, & k=n=m,\\[1mm]
	\mathbb Z, & k=n+m,\\[1mm]
	0, & \text{otherwise}.
\end{cases}
\]
Equivalently, one may simply say that \(H_*(S^n \times S^m)\) is generated by:
\begin{itemize}
	\item the class of a point in degree \(0\),
	\item the copy of \(S^n \times \{\ast\}\) in degree \(n\),
	\item the copy of \(\{\ast\} \times S^m\) in degree \(m\),
	\item the fundamental class of \(S^n \times S^m\) in degree \(n+m\).
\end{itemize}


\medskip
\noindent\textbf{Computation of \(H_*(D^{n+1} \times S^m,\; S^n \times S^m)\)}\par
\smallskip


Observe that
\[
(D^{n+1} \times S^m,\; S^n \times S^m)
=
(D^{n+1},S^n)\times S^m
\]
as a relative product situation. Again by the relative K\"unneth formula,
\[
H_k(D^{n+1} \times S^m,\; S^n \times S^m)
\cong
\bigoplus_{p+q=k} H_p(D^{n+1},S^n)\otimes H_q(S^m).
\]
Now the relative homology of the pair \((D^{n+1},S^n)\) is
\[
H_p(D^{n+1},S^n)\cong
\begin{cases}
	\mathbb Z, & p=n+1,\\
	0, & \text{otherwise}.
\end{cases}
\]
Hence
\[
H_k(D^{n+1} \times S^m,\; S^n \times S^m)
\cong
H_{k-(n+1)}(S^m).
\]
Therefore
\[
H_k(D^{n+1} \times S^m,\; S^n \times S^m)
\cong
\begin{cases}
	\mathbb Z, & k=n+1,\\[1mm]
	\mathbb Z, & k=n+m+1,\\[1mm]
	0, & \text{otherwise}.
\end{cases}
\]


\medskip
\noindent\textbf{Conclusion}\par
\smallskip


We have shown that
\[
S^{n+m+1}
\cong
\bigl(S^n \times D^{m+1}\bigr)\cup_{\,S^n \times S^m}\bigl(D^{n+1} \times S^m\bigr),
\]
and the homology groups are
\[
H_k(S^n \times S^m)\cong
\begin{cases}
	\mathbb Z, & k=0,n,m,n+m
	\text{ (with the degree \(n=m\) case giving } \mathbb Z\oplus\mathbb Z\text{)},\\
	0, & \text{otherwise},
\end{cases}
\]
and
\[
H_k(D^{n+1} \times S^m,\; S^n \times S^m)
\cong
\begin{cases}
	\mathbb Z, & k=n+1,\ n+m+1,\\
	0, & \text{otherwise}.
\end{cases}
\]
\end{solution}

\begin{problem}[CP-VIII-0044 --- The pair \((S^3,K)\) and the complement \(S^3 - K\)]
Let \(V \subset S^3\) be an open subset of \(S^3\) homeomorphic to
\[
S^1 \times \operatorname{int}(D^2),
\]
and let
\[
K=S^1 \times 0 \subset S^1 \times \operatorname{int}(D^2)\cong V.
\]
Compute
\[
H_*(S^3,K)
\quad\text{and}\quad
H_*(S^3-K).
\]
\end{problem}

\begin{solution}
The space \(K\) is homeomorphic to \(S^1\), so
\[
H_i(K)\cong
\begin{cases}
	\mathbb Z, & i=0,1,\\
	0, & \text{otherwise}.
\end{cases}
\]
Also
\[
H_i(S^3)\cong
\begin{cases}
	\mathbb Z, & i=0,3,\\
	0, & \text{otherwise}.
\end{cases}
\]


\medskip
\noindent\textbf{Computation of \(H_*(S^3,K)\)}\par
\smallskip


Use the long exact sequence of the pair \((S^3,K)\):
\[
\cdots \to H_i(K)\to H_i(S^3)\to H_i(S^3,K)\to H_{i-1}(K)\to H_{i-1}(S^3)\to \cdots
\]

We now read off the groups degree by degree.

For \(i=3\), we have
\[
H_3(K)\to H_3(S^3)\to H_3(S^3,K)\to H_2(K),
\]
that is,
\[
0\to \mathbb Z \to H_3(S^3,K)\to 0.
\]
Hence
\[
H_3(S^3,K)\cong \mathbb Z.
\]

For \(i=2\), we have
\[
H_2(K)\to H_2(S^3)\to H_2(S^3,K)\to H_1(K)\to H_1(S^3),
\]
that is,
\[
0\to 0 \to H_2(S^3,K)\to \mathbb Z \to 0.
\]
Hence
\[
H_2(S^3,K)\cong \mathbb Z.
\]

For \(i=1\), we get
\[
H_1(S^3)\to H_1(S^3,K)\to H_0(K)\to H_0(S^3),
\]
that is,
\[
0\to H_1(S^3,K)\to \mathbb Z \xrightarrow{\cong} \mathbb Z.
\]
The map \(H_0(K)\to H_0(S^3)\) is an isomorphism because both spaces are connected and \(K\subset S^3\) is nonempty. Therefore
\[
H_1(S^3,K)=0.
\]

Finally,
\[
H_0(S^3,K)=0
\]
for every pair with nonempty subspace in a connected space.

Thus
\[
H_i(S^3,K)\cong
\begin{cases}
	\mathbb Z, & i=3,\\[1mm]
	\mathbb Z, & i=2,\\[1mm]
	0, & \text{otherwise}.
\end{cases}
\]


\medskip
\noindent\textbf{Computation of \(H_*(S^3-K)\)}\par
\smallskip


To compute the homology of the complement, Alexander duality is the most natural tool. For a compact locally contractible subspace \(K\subset S^3\),
\[
\widetilde H_i(S^3-K)\cong \widetilde H^{\,2-i}(K).
\]
Since \(K\cong S^1\), its reduced cohomology is
\[
\widetilde H^j(K)\cong
\begin{cases}
	\mathbb Z, & j=1,\\
	0, & \text{otherwise}.
\end{cases}
\]
Hence
\[
\widetilde H_i(S^3-K)\cong
\begin{cases}
	\mathbb Z, & 2-i=1,\\
	0, & \text{otherwise},
\end{cases}
\]
so the only nonzero reduced homology group is
\[
\widetilde H_1(S^3-K)\cong \mathbb Z.
\]
Also
\[
\widetilde H_0(S^3-K)\cong \widetilde H^2(K)=0,
\]
so \(S^3-K\) is connected. Therefore
\[
H_i(S^3-K)\cong
\begin{cases}
	\mathbb Z, & i=0,\\[1mm]
	\mathbb Z, & i=1,\\[1mm]
	0, & \text{otherwise}.
\end{cases}
\]


\medskip
\noindent\textbf{Conclusion}\par
\smallskip


The homology groups are
\[
H_i(S^3,K)\cong
\begin{cases}
	\mathbb Z, & i=3,2,\\
	0, & \text{otherwise},
\end{cases}
\]
and
\[
H_i(S^3-K)\cong
\begin{cases}
	\mathbb Z, & i=0,1,\\
	0, & \text{otherwise}.
\end{cases}
\]
So the complement \(S^3-K\) has the homology of a circle.
\end{solution}

\begin{problem}[CP-VIII-0045 --- A nonvanishing vector field on \(S^n\) and the antipodal map]
A vector field on \(S^n\) is a continuous map
\[
v:S^n\to \mathbb R^{n+1}
\]
such that
\[
\langle x,v(x)\rangle =0
\qquad\text{for all }x\in S^n.
\]
If \(v(x)\neq 0\) for all \(x\in S^n\), show that \(\operatorname{id}_{S^n}\) is homotopic to the antipodal map
\[
a(x)=-x.
\]
Deduce that if \(S^n\) admits a nonvanishing vector field, then \(n\) is odd.
\end{problem}

\begin{solution}
Assume \(v(x)\neq 0\) for every \(x\in S^n\). Define a normalized vector field
\[
u(x)=\frac{v(x)}{\|v(x)\|}.
\]
Then \(u:S^n\to S^n\) is continuous and still tangent to the sphere, so
\[
\langle x,u(x)\rangle =0
\qquad\text{for all }x\in S^n.
\]
Thus for each \(x\in S^n\), the vectors \(x\) and \(u(x)\) are orthogonal unit vectors in \(\mathbb R^{n+1}\).

Now define
\[
H:S^n\times [0,1]\to S^n
\]
by
\[
H(x,t)=\cos(\pi t)\,x+\sin(\pi t)\,u(x).
\]
We must check that \(H(x,t)\) really lies in \(S^n\). Since \(\|x\|=1\), \(\|u(x)\|=1\), and \(\langle x,u(x)\rangle =0\), we have
\[
\|H(x,t)\|^2
=
\cos^2(\pi t)\|x\|^2
+
\sin^2(\pi t)\|u(x)\|^2
+
2\cos(\pi t)\sin(\pi t)\langle x,u(x)\rangle.
\]
Hence
\[
\|H(x,t)\|^2
=
\cos^2(\pi t)+\sin^2(\pi t)=1.
\]
So \(H(x,t)\in S^n\) for all \((x,t)\).

At the endpoints,
\[
H(x,0)=\cos(0)\,x+\sin(0)\,u(x)=x,
\]
and
\[
H(x,1)=\cos(\pi)\,x+\sin(\pi)\,u(x)=-x.
\]
Therefore \(H\) is a homotopy from the identity map \(\operatorname{id}_{S^n}\) to the antipodal map \(a(x)=-x\).


\medskip
\noindent\textbf{Deduction using the degree}\par
\smallskip


Homotopic maps \(S^n\to S^n\) have the same degree. Since \(\operatorname{id}_{S^n}\) is homotopic to the antipodal map \(a\), we obtain
\[
\deg(\operatorname{id}_{S^n})=\deg(a).
\]
But
\[
\deg(\operatorname{id}_{S^n})=1,
\]
while the antipodal map has degree
\[
\deg(a)=(-1)^{n+1}.
\]
Thus
\[
1=(-1)^{n+1}.
\]
This happens exactly when \(n+1\) is even, that is, when \(n\) is odd.

Therefore, if \(S^n\) admits a nonvanishing continuous tangent vector field, then \(n\) must be odd.


\medskip
\noindent\textbf{Conclusion}\par
\smallskip


A nowhere-zero tangent vector field on \(S^n\) gives a homotopy
\[
\operatorname{id}_{S^n}\simeq -\operatorname{id}_{S^n}.
\]
Comparing degrees then forces
\[
(-1)^{n+1}=1,
\]
so \(n\) must be odd.
\end{solution}

\begin{problem}[CP-VIII-0046 --- Homeomorphisms of the torus and the group \(GL(2,\mathbb Z)\)]
Suppose \(f:T^2\to T^2\) is a homeomorphism. Show that
\[
f_*:H_1(T^2)\to H_1(T^2)
\]
defines an element of \(GL(2,\mathbb Z)\), and that any element of \(GL(2,\mathbb Z)\) can be realized by a homeomorphism of \(T^2\).

\par
\end{problem}

\begin{solution}
Recall that
\[
T^2=S^1\times S^1=\mathbb R^2/\mathbb Z^2,
\qquad
H_1(T^2)\cong H_1(S^1)\oplus H_1(S^1)\cong \mathbb Z\oplus \mathbb Z.
\]
Hence \(f_*\) is a homomorphism
\[
f_*:\mathbb Z^2\to \mathbb Z^2.
\]

Since \(f\) is a homeomorphism, it has an inverse \(f^{-1}:T^2\to T^2\). By functoriality of homology,
\[
(f^{-1})_*\circ f_*=(f^{-1}\circ f)_*=(\mathrm{id}_{T^2})_*=\mathrm{id},
\]
and similarly
\[
f_*\circ (f^{-1})_*=\mathrm{id}.
\]
Therefore \(f_*\) is an automorphism of \(\mathbb Z^2\). But
\[
\mathrm{Aut}(\mathbb Z^2)=GL(2,\mathbb Z),
\]
so \(f_*\) defines an element of \(GL(2,\mathbb Z)\).

We now prove the converse. Let
\[
A=\begin{pmatrix}a&b\\ c&d\end{pmatrix}\in GL(2,\mathbb Z).
\]
Then \(A\) is an invertible linear map \(\mathbb R^2\to \mathbb R^2\) with integer entries and determinant \(\pm 1\). Since the entries are integers, we have
\[
A(\mathbb Z^2)\subseteq \mathbb Z^2.
\]
Since \(A\) is invertible over \(\mathbb Z\), in fact
\[
A(\mathbb Z^2)=\mathbb Z^2.
\]
Therefore \(A\) descends to a map on the quotient
\[
\bar A:\mathbb R^2/\mathbb Z^2\to \mathbb R^2/\mathbb Z^2,
\qquad
[x]\mapsto [Ax].
\]
This is well defined: if \(x-y\in \mathbb Z^2\), then
\[
A(x-y)\in \mathbb Z^2,
\]
so \(Ax\) and \(Ay\) represent the same point of \(\mathbb R^2/\mathbb Z^2\).

Since \(A^{-1}\in GL(2,\mathbb Z)\) as well, the inverse map descends too, and \(\bar A\) is a homeomorphism of \(T^2\).

Finally, the standard generators of \(H_1(T^2)\cong \mathbb Z^2\) are represented by the two coordinate circles:
\[
[S^1\times \{*\}],\qquad [\{*\}\times S^1].
\]
The induced map \((\bar A)_*\) on \(H_1(T^2)\) is exactly the matrix \(A\). Thus every element of \(GL(2,\mathbb Z)\) is realized by a homeomorphism of \(T^2\).

\par

Therefore every homeomorphism of \(T^2\) induces an element of
\(GL(2,\mathbb Z)\), and conversely every element of \(GL(2,\mathbb Z)\)
is realized by a homeomorphism of \(T^2\).
\end{solution}

\begin{problem}[CP-VIII-0047 --- Homology of \(X\times S^1\) and of the torus \(T^n\)]
If \(H_*(X)\) is a free abelian group, show that
\[
H_*(X\times S^1)\cong H_*(X)\oplus H_{*-1}(X).
\]
(In fact, this is true even if \(H_*(X)\) is not free.) Compute \(H_*(T^n)\).

\par
\end{problem}

\begin{solution}
We use the K\"unneth theorem.

The homology of the circle is
\[
H_0(S^1)\cong \mathbb Z,\qquad
H_1(S^1)\cong \mathbb Z,\qquad
H_j(S^1)=0 \quad \text{for }j\ge 2.
\]
Since these groups are free, the \(\operatorname{Tor}\)-terms in the K\"unneth formula vanish. Therefore
\[
H_k(X\times S^1)\cong \bigoplus_{i+j=k} H_i(X)\otimes H_j(S^1).
\]
But only the terms \(j=0\) and \(j=1\) survive, so
\[
H_k(X\times S^1)\cong H_k(X)\otimes H_0(S^1)\;\oplus\; H_{k-1}(X)\otimes H_1(S^1).
\]
Since \(H_0(S^1)\cong H_1(S^1)\cong \mathbb Z\), we conclude
\[
H_k(X\times S^1)\cong H_k(X)\oplus H_{k-1}(X).
\]
Equivalently, as graded groups,
\[
H_*(X\times S^1)\cong H_*(X)\oplus H_{*-1}(X).
\]

\par

Now compute \(H_*(T^n)\), where
\[
T^n=(S^1)^n.
\]
We proceed inductively. Since
\[
T^n=T^{n-1}\times S^1,
\]
the above formula gives
\[
H_k(T^n)\cong H_k(T^{n-1})\oplus H_{k-1}(T^{n-1}).
\]
Starting from
\[
T^1=S^1,\qquad H_0(S^1)\cong \mathbb Z,\qquad H_1(S^1)\cong \mathbb Z,
\]
one sees that the ranks satisfy Pascal's recursion
\[
\binom{n}{k}=\binom{n-1}{k}+\binom{n-1}{k-1}.
\]
Thus
\[
H_k(T^n)\cong \mathbb Z^{\binom{n}{k}}
\qquad (0\le k\le n),
\]
and
\[
H_k(T^n)=0\qquad (k>n).
\]

So explicitly,
\[
H_0(T^n)\cong \mathbb Z,
\]
\[
H_1(T^n)\cong \mathbb Z^n,
\]
\[
H_2(T^n)\cong \mathbb Z^{\binom{n}{2}},
\]
\[
\cdots
\]
\[
H_n(T^n)\cong \mathbb Z.
\]

Hence the full answer is
\[
\boxed{H_k(T^n)\cong \mathbb Z^{\binom{n}{k}}.}
\]

A useful conceptual reformulation is that
\[
H_*(T^n)\cong \Lambda^*(\mathbb Z^n),
\]
the exterior algebra on \(n\) generators in degree \(1\).
\end{solution}

\begin{problem}[CP-VIII-0048 --- Degree of the product map \(g(x,y)=(f(x),y)\)]
Given
\[
f:(D^n,S^{n-1})\to (D^n,S^{n-1}),
\]
define
\[
g:(D^n\times D^m,\partial(D^n\times D^m))\to (D^n\times D^m,\partial(D^n\times D^m))
\]
by
\[
g(x,y)=(f(x),y).
\]
Show that
\[
\deg g=\deg f.
\]

\par
\end{problem}

\begin{solution}
Recall first that
\[
H_n(D^n,S^{n-1})\cong \mathbb Z.
\]
If \(u\) is a generator of this group, then by definition
\[
f_*(u)=(\deg f)\,u.
\]

Similarly,
\[
H_{n+m}(D^n\times D^m,\partial(D^n\times D^m))\cong \mathbb Z.
\]
We use the relative K\"unneth theorem. Since
\[
(D^n\times D^m,\partial(D^n\times D^m))
\]
is the product of pairs
\[
(D^n,S^{n-1})\times (D^m,S^{m-1}),
\]
there is a natural isomorphism
\[
H_{n+m}(D^n\times D^m,\partial(D^n\times D^m))
\cong
H_n(D^n,S^{n-1})\otimes H_m(D^m,S^{m-1}).
\]
Let
\[
u\in H_n(D^n,S^{n-1}),\qquad v\in H_m(D^m,S^{m-1})
\]
be generators. Then \(u\times v\) is a generator of
\[
H_{n+m}(D^n\times D^m,\partial(D^n\times D^m)).
\]

Now the map \(g\) is just
\[
g=f\times \mathrm{id}_{D^m}.
\]
Hence on relative homology,
\[
g_*(u\times v)=f_*(u)\times (\mathrm{id}_{D^m})_*(v).
\]
Since \((\mathrm{id}_{D^m})_*(v)=v\), we obtain
\[
g_*(u\times v)=((\deg f)u)\times v=(\deg f)(u\times v).
\]
Thus \(g_*\) acts on the top relative homology group by multiplication by \(\deg f\). By definition, this means
\[
\deg g=\deg f.
\]

Therefore
\[
\boxed{\deg(f\times \mathrm{id}_{D^m})=\deg f.}
\]
\end{solution}

\begin{problem}[CP-VIII-0049 --- Relative homology of the genus \(2\) surface]
Let \(X\) be the closed orientable surface of genus \(2\), decomposed as in the figure into three pieces:
\[
X=A\cup B\cup C,
\]
where \(A\) and \(C\) are the two once-punctured tori on the right and left, and \(B\) is the middle annulus separating them. Compute \(H_*(X,A)\) and use this to compute \(H_*(X)\). What are \(H_*(X,B)\) and \(H_*(X,C)\)?

\par
\end{problem}

\begin{solution}
We interpret the figure as follows:

\begin{itemize}
	\item \(A\) is the right-hand once-punctured torus,
	\item \(C\) is the left-hand once-punctured torus,
	\item \(B\cong S^1\times I\) is the annulus between them.
\end{itemize}

Recall first that a once-punctured torus \(T^\circ\) has homology
\[
H_0(T^\circ)\cong \mathbb Z,\qquad
H_1(T^\circ)\cong \mathbb Z^2,\qquad
H_n(T^\circ)=0\quad(n\ge 2).
\]
Also, if \(\partial T^\circ\cong S^1\) denotes its boundary circle, then the map
\[
H_1(\partial T^\circ)\to H_1(T^\circ)
\]
is the zero map, because the boundary circle bounds the punctured torus.


\medskip
\noindent\textbf{Step 1: compute \(H_*(X,A)\)}\par
\smallskip


Let
\[
Y=\overline{X\setminus \operatorname{int}(A)}.
\]
Then \(Y\) is again a once-punctured torus, and
\[
Y\cap A=\partial A=\partial Y\cong S^1.
\]
By excision,
\[
H_*(X,A)\cong H_*(Y,\partial Y).
\]
So it remains to compute the relative homology of a once-punctured torus modulo its boundary.

Consider the long exact sequence of the pair \((Y,\partial Y)\):
\[
\begin{aligned}
\cdots&\to H_2(Y)\to H_2(Y,\partial Y)
\to H_1(\partial Y)\to H_1(Y)\\
&\to H_1(Y,\partial Y)\to H_0(\partial Y)
\to H_0(Y)\to H_0(Y,\partial Y)\to 0.
\end{aligned}
\]
Since \(Y\) is a once-punctured torus,
\[
H_2(Y)=0,\qquad
H_1(\partial Y)\cong \mathbb Z,\qquad
H_1(Y)\cong \mathbb Z^2,\qquad
H_0(\partial Y)\cong \mathbb Z,\qquad
H_0(Y)\cong \mathbb Z.
\]
The map
\[
H_1(\partial Y)\to H_1(Y)
\]
is zero, and the map
\[
H_0(\partial Y)\to H_0(Y)
\]
is an isomorphism because both spaces are connected. Hence exactness gives
\[
H_2(Y,\partial Y)\cong \mathbb Z,
\qquad
H_1(Y,\partial Y)\cong \mathbb Z^2,
\qquad
H_0(Y,\partial Y)=0.
\]
Therefore
\[
H_n(X,A)\cong
\begin{cases}
	\mathbb Z, & n=2,\\[1mm]
	\mathbb Z^2, & n=1,\\[1mm]
	0, & n=0\text{ and }n\ge 3.
\end{cases}
\]
So
\[
\boxed{
	H_2(X,A)\cong \mathbb Z,\qquad
	H_1(X,A)\cong \mathbb Z^2,\qquad
	H_0(X,A)=0.
}
\]


\medskip
\noindent\textbf{Step 2: use \(H_*(X,A)\) to compute \(H_*(X)\)}\par
\smallskip


Now consider the long exact sequence of the pair \((X,A)\):
\[
\begin{aligned}
\cdots&\to H_2(A)\to H_2(X)\to H_2(X,A)
\to H_1(A)\to H_1(X)\\
&\to H_1(X,A)\to H_0(A)\to H_0(X)\to 0.
\end{aligned}
\]
Since \(A\) is a once-punctured torus,
\[
H_2(A)=0,\qquad
H_1(A)\cong \mathbb Z^2,\qquad
H_0(A)\cong \mathbb Z.
\]
From Step 1 we already know
\[
H_2(X,A)\cong \mathbb Z,\qquad
H_1(X,A)\cong \mathbb Z^2,\qquad
H_0(X,A)=0.
\]
So the relevant part of the sequence becomes
\[
0\to H_2(X)\to \mathbb Z\to \mathbb Z^2\to H_1(X)\to \mathbb Z^2\to \mathbb Z\to H_0(X)\to 0.
\]
The map
\[
\mathbb Z\to \mathbb Z^2
\]
is zero, because it is induced by the boundary circle of \(A\), and that class is trivial in \(H_1(A)\). Also the map
\[
H_0(A)\to H_0(X)
\]
is an isomorphism since \(A\) and \(X\) are connected. Thus
\[
H_2(X)\cong \mathbb Z,
\]
and
\[
0\to \mathbb Z^2\to H_1(X)\to \mathbb Z^2\to 0.
\]
Since all groups involved are free abelian, this short exact sequence splits, so
\[
H_1(X)\cong \mathbb Z^4.
\]
Finally,
\[
H_0(X)\cong \mathbb Z.
\]
Therefore
\[
\boxed{
	H_n(X)\cong
	\begin{cases}
		\mathbb Z, & n=0,\\[1mm]
		\mathbb Z^4, & n=1,\\[1mm]
		\mathbb Z, & n=2,\\[1mm]
		0, & n\ge 3.
	\end{cases}
}
\]


\medskip
\noindent\textbf{Step 3: compute \(H_*(X,C)\)}\par
\smallskip


By symmetry, \(C\) plays exactly the same role as \(A\). Therefore
\[
H_n(X,C)\cong H_n(X,A)
\]
for every \(n\), and hence
\[
\boxed{
	H_n(X,C)\cong
	\begin{cases}
		\mathbb Z, & n=2,\\[1mm]
		\mathbb Z^2, & n=1,\\[1mm]
		0, & n=0\text{ and }n\ge 3.
	\end{cases}
}
\]


\medskip
\noindent\textbf{Step 4: compute \(H_*(X,B)\)}\par
\smallskip


Now \(B\cong S^1\times I\) is the middle annulus. If we remove the interior of \(B\), the remaining space is the disjoint union of the two once-punctured tori \(A\sqcup C\). By excision,
\[
H_*(X,B)\cong H_*(A\sqcup C,\partial A\sqcup \partial C).
\]
Since relative homology of disjoint unions splits as a direct sum,
\[
H_*(X,B)\cong H_*(A,\partial A)\oplus H_*(C,\partial C).
\]
But each pair \((A,\partial A)\) and \((C,\partial C)\) has relative homology
\[
H_2\cong \mathbb Z,\qquad
H_1\cong \mathbb Z^2,\qquad
H_0=0.
\]
Therefore
\[
\boxed{
	H_n(X,B)\cong
	\begin{cases}
		\mathbb Z^2, & n=2,\\[1mm]
		\mathbb Z^4, & n=1,\\[1mm]
		0, & n=0\text{ and }n\ge 3.
	\end{cases}
}
\]


\medskip
\noindent\textbf{Final answers}\par
\smallskip


We have proved:
\[
\boxed{
	H_n(X,A)\cong
	\begin{cases}
		\mathbb Z, & n=2,\\
		\mathbb Z^2, & n=1,\\
		0, & \text{otherwise,}
	\end{cases}
}
\]
\[
\boxed{
	H_n(X)\cong
	\begin{cases}
		\mathbb Z, & n=0,\\
		\mathbb Z^4, & n=1,\\
		\mathbb Z, & n=2,\\
		0, & \text{otherwise,}
	\end{cases}
}
\]
\[
\boxed{
	H_n(X,C)\cong
	\begin{cases}
		\mathbb Z, & n=2,\\
		\mathbb Z^2, & n=1,\\
		0, & \text{otherwise,}
	\end{cases}
}
\]
and
\[
\boxed{
	H_n(X,B)\cong
	\begin{cases}
		\mathbb Z^2, & n=2,\\
		\mathbb Z^4, & n=1,\\
		0, & \text{otherwise.}
	\end{cases}
}
\]



\[
\boxed{
\begin{array}{ccccc}
\text{once-punctured torus }C
&\;\cup\;&
\underbrace{S^1\times I}_{B}
&\;\cup\;&
\text{once-punctured torus }A\\[1mm]
&&\partial_-B\qquad\partial_+B&&
\end{array}}
\]
\[
X=A\cup B\cup C,\qquad
A\cong C\cong T^2\setminus\mathring D^2,\qquad
B\cong S^1\times I.
\]


% duplicate source sketch intentionally omitted after canonical schematic


% duplicate source sketch intentionally omitted after canonical schematic



The region \(B\) is the strip between the two indicated boundary curves
\(\partial_-B\) and \(\partial_+B\). Since it has exactly two boundary
components, each homeomorphic to \(S^1\), and is topologically a cylinder,
we have
\[
B \cong S^1\times I,
\]
so \(B\) is an annulus.

The complementary regions \(A\) and \(C\) each contain one handle of the
genus \(2\) surface, and each has exactly one boundary circle, namely one of
the boundary components of \(B\). Therefore each of \(A\) and \(C\) is
homeomorphic to a torus with an open disk removed, i.e. a once-punctured
torus:
\[
A \cong T^2\setminus \mathring{D}^2,
\qquad
C \cong T^2\setminus \mathring{D}^2.
\]
Hence \(X\) is decomposed as two once-punctured tori joined by an annulus.
\end{solution}

\begin{problem}[CP-VIII-0050 --- Realizing a homology class by a finite cell complex]
Suppose \(x\in H_n(X)\), where \(X\) is an arbitrary topological space. Show that there is a finite cell complex \(A\) and a map
\[
f:A\to X
\]
such that
\[
x\in \operatorname{im}\bigl(f_*:H_n(A)\to H_n(X)\bigr).
\]

\par
\end{problem}

\begin{solution}
Choose a singular cycle \(z\in Z_n(X)\) representing the class \(x\). Since singular chains are finite sums, we may write
\[
z=\sum_{i=1}^r a_i\sigma_i,
\]
where each \(\sigma_i:\Delta^n\to X\) is a singular \(n\)-simplex and each \(a_i\in \mathbb Z\).

We now replace each coefficient \(a_i\) by repeated copies of the corresponding simplex, reversing orientation when \(a_i<0\). Thus we may think of \(z\) as a finite sum
\[
z=\sum_{j=1}^N \varepsilon_j \sigma_j,
\qquad \varepsilon_j\in\{+1,-1\},
\]
of finitely many oriented singular \(n\)-simplices.

For each \(\sigma_j\), take a copy \(\Delta_j^n\) of the standard \(n\)-simplex. Since \(\partial z=0\), each oriented \((n-1)\)-face appearing in the boundary occurs with total coefficient zero. In particular, the faces can be paired so that whenever two faces are paired, the corresponding singular \((n-1)\)-simplices in \(X\) agree, but with opposite induced orientations.

We now glue the copies \(\Delta_j^n\) along these paired boundary faces. Because only finitely many simplices are involved, the resulting space \(A\) is a finite \(\Delta\)-complex, hence in particular a finite cell complex.

The singular maps \(\sigma_j:\Delta_j^n\to X\) fit together under these identifications, precisely because paired faces map to the same singular \((n-1)\)-simplex in \(X\). Therefore they descend to a continuous map
\[
f:A\to X.
\]

Let \(c_A\) denote the sum of the top-dimensional simplices of \(A\) with the induced orientations. By construction, the boundary faces cancel in pairs, so
\[
\partial c_A=0.
\]
Hence \(c_A\) defines a homology class
\[
[c_A]\in H_n(A).
\]
Moreover, by construction,
\[
f_*([c_A])=[z]=x.
\]
Therefore
\[
x\in \operatorname{im}(f_*).
\]

\par

Thus every homology class in \(H_n(X)\) is realized by a map from a finite cell complex.

\par
\end{solution}

\begin{problem}[CP-VIII-0051 --- Killing the \(n\)-skeleton and consequences for homology]
Let
\[
f:X\to Y
\]
be a map of finite cell complexes, and suppose that
\[
\pi_i(Y,*)=0
\qquad
\text{for all }0<i\leq n.
\]
Using Problem 10, show that there is a map
\[
\overline f:X/X_n\to Y
\]
such that
\[
\overline f\circ \pi\simeq f,
\]
where
\[
\pi:X\to X/X_n
\]
is the quotient map.

Deduce the following consequences:

\[
\text{(a)}\qquad H_i(Y)=0
\quad
\text{for }0<i\leq n.
\]

\[
\text{(b)}\qquad
\psi:\pi_{n+1}(Y,*)\to H_{n+1}(Y)
\]
is surjective.

\[
\text{(c)}\qquad
\text{If }\pi_i(Y,*)=0\text{ for all }i>0,\text{ then }Y\text{ is contractible.}
\]
\end{problem}

\begin{solution}
We assume that \(Y\) is path-connected, or equivalently that all maps under consideration land in the component of the chosen basepoint \(*\).

Let
\[
X_n\subset X
\]
be the \(n\)-skeleton of \(X\). Since
\[
\pi_i(Y,*)=0
\qquad
\text{for }0<i\leq n,
\]
Problem 10 implies that the restriction
\[
f|_{X_n}:X_n\to Y
\]
is homotopic to a constant map.

Thus, after replacing \(f\) by a homotopic map, we may assume that
\[
f(X_n)=\{*\}.
\]

Because \(f\) is constant on \(X_n\), it factors through the quotient
\[
X/X_n.
\]
Therefore there exists a map
\[
\overline f:X/X_n\to Y
\]
such that
\[
\overline f\circ \pi=f
\]
for the modified representative of the homotopy class. Hence for the original \(f\), we get
\[
\boxed{
	\overline f\circ \pi\simeq f.
}
\]

This proves the first claim.


\medskip
\noindent\textbf{(a) Vanishing of \(H_i(Y)\) for \(0<i\leq n\)}\par
\smallskip


Apply the first part to the identity map
\[
\operatorname{id}_Y:Y\to Y.
\]
Then there exists a map
\[
g:Y/Y_n\to Y
\]
such that
\[
g\circ \pi\simeq \operatorname{id}_Y,
\]
where
\[
\pi:Y\to Y/Y_n
\]
collapses the \(n\)-skeleton \(Y_n\) to a point.

Passing to homology gives
\[
g_*\circ \pi_*=(\operatorname{id}_Y)_*.
\]

Now observe that \(Y/Y_n\) has no cells in dimensions
\[
1,2,\ldots,n.
\]
Therefore
\[
H_i(Y/Y_n)=0
\qquad
\text{for }0<i\leq n.
\]

Hence, for \(0<i\leq n\), the map
\[
\pi_*:H_i(Y)\to H_i(Y/Y_n)
\]
has target zero. Thus
\[
\pi_*=0.
\]
Therefore
\[
(\operatorname{id}_Y)_*
=
g_*\circ \pi_*
=
0.
\]

But the identity map on \(H_i(Y)\) can be zero only if
\[
H_i(Y)=0.
\]

Hence
\[
\boxed{
	H_i(Y)=0
	\qquad
	\text{for }0<i\leq n.
}
\]


\medskip
\noindent\textbf{(b) Surjectivity of the Hurewicz homomorphism}\par
\smallskip


Again use the factorization of the identity map:
\[
Y\xrightarrow{\pi}Y/Y_n\xrightarrow{g}Y,
\]
with
\[
g\circ \pi\simeq \operatorname{id}_Y.
\]

Passing to homology in degree \(n+1\), every class
\[
\alpha\in H_{n+1}(Y)
\]
satisfies
\[
\alpha
=
(g\circ \pi)_*(\alpha)
=
g_*(\pi_*(\alpha)).
\]

So every class in \(H_{n+1}(Y)\) is the image under \(g_*\) of some class in
\[
H_{n+1}(Y/Y_n).
\]

Now \(Y/Y_n\) has a single \(0\)-cell and no cells in dimensions
\[
1,\ldots,n.
\]
Its \((n+1)\)-cells therefore contribute spheres
\[
S^{n+1}\to Y/Y_n.
\]

Consequently,
\[
H_{n+1}(Y/Y_n)
\]
is generated by homology classes represented by maps
\[
S^{n+1}\to Y/Y_n.
\]

Composing such maps with
\[
g:Y/Y_n\to Y
\]
gives maps
\[
S^{n+1}\to Y.
\]
Their homology classes lie in the image of the Hurewicz homomorphism
\[
\psi:\pi_{n+1}(Y,*)\to H_{n+1}(Y).
\]

Since every class of \(H_{n+1}(Y)\) is obtained in this way, the Hurewicz homomorphism is surjective:

\[
\boxed{
	\psi:\pi_{n+1}(Y,*)\to H_{n+1}(Y)
	\text{ is surjective.}
}
\]


\medskip
\noindent\textbf{(c) If all homotopy groups vanish, then \(Y\) is contractible}\par
\smallskip


Assume now that
\[
\pi_i(Y,*)=0
\qquad
\text{for all }i>0.
\]

Since \(Y\) is a finite cell complex, it has some finite dimension, say
\[
\dim Y=N.
\]

Apply the first part with
\[
n=N
\]
and with
\[
f=\operatorname{id}_Y.
\]

Then there exists a map
\[
g:Y/Y_N\to Y
\]
such that
\[
g\circ \pi\simeq \operatorname{id}_Y.
\]

But
\[
Y_N=Y,
\]
because \(N=\dim Y\). Hence
\[
Y/Y_N=Y/Y
\]
is a single point.

Therefore the identity map
\[
\operatorname{id}_Y
\]
is homotopic to a constant map:
\[
\operatorname{id}_Y\simeq c_*.
\]

This means precisely that \(Y\) is contractible.

Thus
\[
\boxed{
	\pi_i(Y,*)=0\text{ for all }i>0
	\quad\Longrightarrow\quad
	Y\text{ is contractible.}
}
\]

This is a finite CW-complex version of Whitehead's theorem.

\par
\end{solution}

\section{Homological Machinery}

\begin{problem}[CP-VIII-0052 --- Deleting a maximal simplex and the resulting homology exact sequence]
Let \(K\) be a simplicial complex, and suppose that \(\sigma\in K\) is not a proper face of any simplex. Show that
\[
L=K\setminus\{\sigma\}
\]
is again a simplicial complex, and that the inclusion \(V_L\to V_K\) defines a simplicial map
\[
i:L\to K.
\]

If \(\sigma\) has dimension \(n\), note that \(d_n(\sigma)\) is an \((n-1)\)-cycle and consists of simplices of \(L\), so represents a class
\[
[d_n(\sigma)]\in H_{n-1}(L);
\]
this defines a homomorphism
\[
\varphi:\mathbb Z\to H_{n-1}(L)
\]
by
\[
1\mapsto [d_n(\sigma)].
\]
Construct a homomorphism
\[
\phi:H_n(K)\to \mathbb Z
\]
such that
\[
0\longrightarrow H_n(L)\xrightarrow{i_*} H_n(K)\xrightarrow{\phi} \mathbb Z
\xrightarrow{\varphi} H_{n-1}(L)\xrightarrow{i_*} H_{n-1}(K)\longrightarrow 0
\]
is exact, and show that
\[
i_*:H_j(L)\to H_j(K)
\]
is an isomorphism for \(j\neq n-1,n\).
\end{problem}

\begin{solution}
\textbf{Step 1: \(L=K\setminus\{\sigma\}\) is again a simplicial complex.}

We must show that \(L\) is closed under taking faces.

Let \(\tau\in L\). Then \(\tau\in K\) and \(\tau\neq \sigma\). Let \(\eta\subseteq \tau\). Since \(K\) is a simplicial complex, we already know that \(\eta\in K\). We must show that \(\eta\neq \sigma\), so that \(\eta\in L\).

If \(\eta=\sigma\), then
\[
\sigma\subseteq \tau.
\]
Since \(\tau\neq \sigma\), this would mean that \(\sigma\) is a proper face of the simplex \(\tau\), contradicting the hypothesis that \(\sigma\) is not a proper face of any simplex of \(K\).

Therefore \(\eta\neq \sigma\), so \(\eta\in L\). Hence \(L\) is a simplicial complex.

\par

\textbf{Step 2: the inclusion \(V_L\to V_K\) defines a simplicial map \(i:L\to K\).}

The vertex set of \(L\) is contained in the vertex set of \(K\), so there is a natural inclusion
\[
V_L\hookrightarrow V_K.
\]
If \(\tau\in L\), then \(\tau\) is by definition also a simplex of \(K\). Therefore the image of every simplex of \(L\) is a simplex of \(K\). So the inclusion of vertex sets defines a simplicial map
\[
i:L\to K.
\]

\par

\textbf{Step 3: the relative chain complex \(C_*(K,L)\).}

Let \(\dim \sigma = n\).

Because \(\sigma\) is not a proper face of any simplex of \(K\), the only simplex present in \(K\) and not in \(L\) is \(\sigma\) itself. Therefore the simplicial chain groups satisfy:
\[
C_j(K)=C_j(L)\qquad \text{for } j\neq n,
\]
while in degree \(n\),
\[
C_n(K)=C_n(L)\oplus \mathbb Z\langle \sigma\rangle.
\]
Hence the relative chain groups are
\[
C_j(K,L)\cong
\begin{cases}
	\mathbb Z, & j=n,\\
	0, & j\neq n.
\end{cases}
\]
Moreover, since \(\sigma\) is not a face of any \((n+1)\)-simplex, the relative boundary map is zero. Thus
\[
H_j(K,L)\cong
\begin{cases}
	\mathbb Z, & j=n,\\
	0, & j\neq n.
\end{cases}
\]

So the pair \((K,L)\) has relative homology concentrated in degree \(n\), where it is generated by the class of the simplex \(\sigma\).

\par

\textbf{Step 4: the homomorphism \(\varphi:\mathbb Z\to H_{n-1}(L)\).}

The simplicial boundary
\[
d_n(\sigma)=\partial \sigma
\]
is a sum of the \((n-1)\)-faces of \(\sigma\). Since only the simplex \(\sigma\) itself was removed, all of its proper faces remain in \(L\). Therefore \(d_n(\sigma)\) is a simplicial \((n-1)\)-chain in \(L\).

Since
\[
d_{n-1}d_n(\sigma)=0,
\]
the chain \(d_n(\sigma)\) is an \((n-1)\)-cycle in \(L\), and so it defines a homology class
\[
[d_n(\sigma)]\in H_{n-1}(L).
\]
This gives a homomorphism
\[
\varphi:\mathbb Z\to H_{n-1}(L),
\qquad
1\mapsto [d_n(\sigma)].
\]

Under the identification
\[
H_n(K,L)\cong \mathbb Z
\]
with generator represented by the class of \(\sigma\), this map \(\varphi\) is exactly the connecting homomorphism in the long exact sequence of the pair \((K,L)\).

\par

\textbf{Step 5: construction of \(\phi:H_n(K)\to \mathbb Z\).}

The long exact sequence of the pair \((K,L)\) contains the map
\[
H_n(K)\to H_n(K,L).
\]
Since
\[
H_n(K,L)\cong \mathbb Z,
\]
we define
\[
\phi:H_n(K)\to \mathbb Z
\]
to be this map, followed by the chosen identification of \(H_n(K,L)\) with \(\mathbb Z\).

Concretely, if a class \([z]\in H_n(K)\) is represented by an \(n\)-cycle of the form
\[
z = a\,\sigma + c,
\qquad c\in C_n(L),
\]
then
\[
\phi([z])=a.
\]
In other words, \(\phi\) records the coefficient of the deleted simplex \(\sigma\).

This is well-defined because no \((n+1)\)-simplex has \(\sigma\) as a face, so boundaries in \(K\) cannot change the coefficient of \(\sigma\).

\par

\textbf{Step 6: the exact sequence.}

Now apply the long exact sequence of the pair \((K,L)\):
\[
\cdots \to H_n(L)\xrightarrow{i_*} H_n(K)\xrightarrow{\phi} H_n(K,L)
\xrightarrow{\partial} H_{n-1}(L)\xrightarrow{i_*} H_{n-1}(K)\to \cdots
\]
Using the identification
\[
H_n(K,L)\cong \mathbb Z
\]
and the fact that
\[
H_j(K,L)=0\qquad\text{for } j\neq n,
\]
this becomes
\[
0\longrightarrow H_n(L)\xrightarrow{i_*} H_n(K)\xrightarrow{\phi} \mathbb Z
\xrightarrow{\varphi} H_{n-1}(L)\xrightarrow{i_*} H_{n-1}(K)\longrightarrow 0.
\]

This is the desired exact sequence.

\par

\textbf{Interpretation of exactness.}

At \(H_n(K)\), exactness says that
\[
\ker \phi = \operatorname{im}(i_*).
\]
So an \(n\)-homology class in \(K\) comes from \(L\) exactly when its representative has zero coefficient on \(\sigma\).

At \(\mathbb Z\), exactness says
\[
\ker \varphi = \operatorname{im}(\phi).
\]
Thus an integer \(a\) lies in the image of \(\phi\) exactly when the cycle
\[
a\,d_n(\sigma)
\]
is homologous to zero in \(L\).

At \(H_{n-1}(L)\), exactness says
\[
\ker(i_*)=\operatorname{im}(\varphi).
\]
So the only \((n-1)\)-homology classes in \(L\) that vanish in \(K\) are those generated by the boundary of the deleted simplex.

\par

\textbf{Step 7: isomorphisms in all other degrees.}

Because
\[
H_j(K,L)=0\qquad \text{for } j\neq n,
\]
the long exact sequence of the pair shows that for every \(j\neq n,n-1\), the map
\[
i_*:H_j(L)\to H_j(K)
\]
is an isomorphism. Indeed, the relevant segment of the long exact sequence takes the form
\[
0\to H_j(L)\xrightarrow{i_*} H_j(K)\to 0.
\]

Therefore,
\[
i_*:H_j(L)\xrightarrow{\cong} H_j(K)
\qquad\text{for all } j\neq n-1,n.
\]


\medskip
\noindent\textbf{Conclusion}\par
\smallskip


We have proved the following:

\begin{itemize}
	\item If \(\sigma\in K\) is not a proper face of any simplex, then
	\[
	L=K\setminus\{\sigma\}
	\]
	is again a simplicial complex.

	\item The inclusion of vertex sets defines a simplicial map
	\[
	i:L\to K.
	\]

	\item If \(\dim \sigma=n\), then there is an exact sequence
	\[
	0\longrightarrow H_n(L)\xrightarrow{i_*} H_n(K)\xrightarrow{\phi} \mathbb Z
	\xrightarrow{\varphi} H_{n-1}(L)\xrightarrow{i_*} H_{n-1}(K)\longrightarrow 0,
	\]
	where
	\[
	\varphi(1)=[d_n(\sigma)].
	\]

	\item For every degree \(j\neq n-1,n\), the map
	\[
	i_*:H_j(L)\to H_j(K)
	\]
	is an isomorphism.
\end{itemize}

So removing a maximal simplex affects homology only in the two adjacent degrees \(n\) and \(n-1\), and the effect is described completely by the boundary of the removed simplex.
\end{solution}

\begin{problem}[CP-VIII-0053 --- Exact sequences of abelian groups: determining \(A\) and \(\alpha\)]
For each of the following exact sequences of abelian groups and homomorphisms, determine as much as possible about the unknown group \(A\) and the homomorphism \(\alpha\):
\[
\text{(i)}\qquad 0 \longrightarrow \mathbb Z/2 \longrightarrow A \longrightarrow \mathbb Z \longrightarrow 0,
\]
\[
\text{(ii)}\qquad 0 \longrightarrow \mathbb Z/2 \longrightarrow A \longrightarrow \mathbb Z/2 \longrightarrow 0,
\]
\[
\text{(iii)}\qquad 0 \longrightarrow A \longrightarrow \mathbb Z \xrightarrow{\alpha} \mathbb Z \longrightarrow \mathbb Z/2 \longrightarrow 0,
\]
\[
\text{(iv)}\qquad 0 \longrightarrow \mathbb Z/3 \longrightarrow A \longrightarrow \mathbb Z/2 \longrightarrow \mathbb Z \xrightarrow{\alpha} \mathbb Z \longrightarrow 0.
\]


We analyze the four cases separately.

\par

\textbf{(i)} Consider
\[
0 \longrightarrow \mathbb Z/2 \longrightarrow A \longrightarrow \mathbb Z \longrightarrow 0.
\]
This is a short exact sequence
\[
0 \longrightarrow \mathbb Z/2 \longrightarrow A \longrightarrow \mathbb Z \longrightarrow 0.
\]
Since \(\mathbb Z\) is a free abelian group, every short exact sequence ending in \(\mathbb Z\) splits. Hence
\[
A \cong \mathbb Z \oplus \mathbb Z/2.
\]
So in this case
\[
\boxed{A \cong \mathbb Z \oplus \mathbb Z/2.}
\]

\par

\textbf{(ii)} Now consider
\[
0 \longrightarrow \mathbb Z/2 \longrightarrow A \longrightarrow \mathbb Z/2 \longrightarrow 0.
\]
Here \(A\) is a finite abelian group of order \(4\). Up to isomorphism, there are exactly two abelian groups of order \(4\):
\[
\mathbb Z/4 \qquad\text{and}\qquad \mathbb Z/2 \oplus \mathbb Z/2.
\]
Both can occur:

\begin{itemize}
	\item the split extension gives \(A \cong \mathbb Z/2 \oplus \mathbb Z/2\),
	\item a non-split extension gives \(A \cong \mathbb Z/4\).
\end{itemize}

Therefore the exact sequence does not determine \(A\) uniquely, but it does determine that
\[
\boxed{A \cong \mathbb Z/4 \quad\text{or}\quad A \cong \mathbb Z/2 \oplus \mathbb Z/2.}
\]

\par

\textbf{(iii)} Consider
\[
0 \longrightarrow A \longrightarrow \mathbb Z \xrightarrow{\alpha} \mathbb Z \longrightarrow \mathbb Z/2 \longrightarrow 0.
\]
Any homomorphism \(\alpha:\mathbb Z \to \mathbb Z\) is multiplication by some integer \(n\):
\[
\alpha(k)=nk.
\]
Its cokernel is
\[
\operatorname{coker}(\alpha)\cong \mathbb Z/n\mathbb Z.
\]
By exactness, this cokernel must be isomorphic to \(\mathbb Z/2\), so necessarily
\[
|n|=2.
\]
Thus
\[
\alpha=\times 2 \qquad\text{or}\qquad \alpha=\times (-2).
\]

Now the kernel of multiplication by \(\pm 2\) on \(\mathbb Z\) is zero, hence
\[
A=\ker(\alpha)=0.
\]
Therefore
\[
\boxed{A=0,\qquad \alpha=\times 2 \text{ or } \times(-2).}
\]

\par

\textbf{(iv)} Finally consider
\[
0 \longrightarrow \mathbb Z/3 \longrightarrow A \longrightarrow \mathbb Z/2 \longrightarrow \mathbb Z \xrightarrow{\alpha} \mathbb Z \longrightarrow 0.
\]
First note that any homomorphism
\[
\mathbb Z/2 \longrightarrow \mathbb Z
\]
must be the zero homomorphism, because \(\mathbb Z\) has no nonzero element of order \(2\). Therefore the map
\[
\mathbb Z/2 \longrightarrow \mathbb Z
\]
is zero.

Exactness at the copy of \(\mathbb Z\) before \(\alpha\) now gives
\[
\ker(\alpha)=\operatorname{im}(\mathbb Z/2 \to \mathbb Z)=0,
\]
so \(\alpha\) is injective. Exactness at the last \(\mathbb Z\) says \(\alpha\) is surjective. Hence \(\alpha\) is an automorphism of \(\mathbb Z\), so
\[
\alpha=\pm \operatorname{id}_{\mathbb Z}.
\]

Since the map \(\mathbb Z/2 \to \mathbb Z\) is zero, exactness also gives that
\[
A \longrightarrow \mathbb Z/2
\]
is surjective, and its kernel is the image of \(\mathbb Z/3\). Thus we get a short exact sequence
\[
0 \longrightarrow \mathbb Z/3 \longrightarrow A \longrightarrow \mathbb Z/2 \longrightarrow 0.
\]
So \(A\) is an abelian group of order \(6\). Every abelian group of order \(6\) is cyclic, hence
\[
A \cong \mathbb Z/6.
\]
Therefore
\[
\boxed{A \cong \mathbb Z/6,\qquad \alpha=\pm \operatorname{id}_{\mathbb Z}.}
\]

\par

Collecting the answers:
\[
\boxed{\text{(i)}\ A\cong \mathbb Z\oplus \mathbb Z/2,}
\]
\[
\boxed{\text{(ii)}\ A\cong \mathbb Z/4 \text{ or } \mathbb Z/2\oplus \mathbb Z/2,}
\]
\[
\boxed{\text{(iii)}\ A=0,\ \alpha=\times 2 \text{ or } \times(-2),}
\]
\[
\boxed{\text{(iv)}\ A\cong \mathbb Z/6,\ \alpha=\pm \operatorname{id}_{\mathbb Z}.}
\]
\end{problem}

\begin{solution}

\medskip
\noindent\textbf{Concrete models for the groups in (i)--(iv)}\par
\smallskip


It is useful to describe the groups and maps in the four exact sequences by explicit concrete models.

\par

\textbf{(i) The case \(A \cong \mathbb Z \oplus \mathbb Z/2\).}

The symbol
\[
\mathbb Z \oplus \mathbb Z/2
\]
means the direct sum of the two abelian groups. Concretely,
\[
A=\mathbb Z \oplus \mathbb Z/2
=\{(n,\bar\varepsilon)\mid n\in\mathbb Z,\ \bar\varepsilon\in\mathbb Z/2\},
\]
with addition defined componentwise:
\[
(n,\bar\varepsilon)+(m,\bar\delta)
=(n+m,\overline{\varepsilon+\delta}).
\]
Thus elements of \(A\) look like
\[
(0,\bar 0),\quad (1,\bar 0),\quad (5,\bar 1),\quad (-3,\bar 1),\dots
\]
The first coordinate is any integer, while the second coordinate is either \(\bar 0\) or \(\bar 1\).

A concrete exact sequence is
\[
0\longrightarrow \mathbb Z/2
\xrightarrow{\ \iota\ }
\mathbb Z\oplus \mathbb Z/2
\xrightarrow{\ \pi\ }
\mathbb Z
\longrightarrow 0,
\]
where
\[
\iota(\bar\varepsilon)=(0,\bar\varepsilon),
\qquad
\pi(n,\bar\varepsilon)=n.
\]
Then
\[
\operatorname{im}(\iota)=\{(0,\bar0),(0,\bar1)\},
\]
and
\[
\ker(\pi)=\{(0,\bar0),(0,\bar1)\}.
\]
So indeed
\[
\ker(\pi)=\operatorname{im}(\iota),
\]
and this realizes the exact sequence concretely.

\par

\textbf{(ii) The two possibilities for \(A\).}

In
\[
0\longrightarrow \mathbb Z/2 \longrightarrow A \longrightarrow \mathbb Z/2 \longrightarrow 0,
\]
there are two possible isomorphism types for \(A\).

\par

\textbf{First possibility: \(A\cong \mathbb Z/2\oplus \mathbb Z/2\).}

Here
\[
A=\mathbb Z/2\oplus \mathbb Z/2
=\{(\bar a,\bar b)\mid \bar a,\bar b\in \mathbb Z/2\},
\]
whose elements are
\[
(\bar0,\bar0),\quad (\bar1,\bar0),\quad (\bar0,\bar1),\quad (\bar1,\bar1).
\]
Define
\[
\iota(\bar a)=(\bar a,\bar0),
\qquad
\pi(\bar a,\bar b)=\bar b.
\]
Then
\[
\operatorname{im}(\iota)=\{(\bar0,\bar0),(\bar1,\bar0)\},
\]
and
\[
\ker(\pi)=\{(\bar0,\bar0),(\bar1,\bar0)\}.
\]
So this gives the split exact sequence
\[
0\longrightarrow \mathbb Z/2
\longrightarrow \mathbb Z/2\oplus \mathbb Z/2
\longrightarrow \mathbb Z/2
\longrightarrow 0.
\]

\par

\textbf{Second possibility: \(A\cong \mathbb Z/4\).}

Take
\[
0\longrightarrow \mathbb Z/2
\xrightarrow{\ \iota\ }
\mathbb Z/4
\xrightarrow{\ \pi\ }
\mathbb Z/2
\longrightarrow 0,
\]
where
\[
\iota(\bar a)=\overline{2a}\in \mathbb Z/4,
\qquad
\pi(\bar x \bmod 4)=\bar x \bmod 2.
\]
Concretely,
\[
\mathbb Z/4=\{\bar0,\bar1,\bar2,\bar3\}.
\]
The image of \(\iota\) is
\[
\operatorname{im}(\iota)=\{\bar0,\bar2\},
\]
while the kernel of \(\pi\) is also
\[
\ker(\pi)=\{\bar0,\bar2\}.
\]
Hence
\[
\ker(\pi)=\operatorname{im}(\iota),
\]
so the sequence is exact.

Thus case (ii) admits both
\[
A\cong \mathbb Z/2\oplus \mathbb Z/2
\qquad\text{and}\qquad
A\cong \mathbb Z/4.
\]

\par

\textbf{(iii) The case \(A=0\) and \(\alpha=\times 2\) or \(\times(-2)\).}

We found that
\[
0\longrightarrow A \longrightarrow \mathbb Z \xrightarrow{\alpha} \mathbb Z \longrightarrow \mathbb Z/2 \longrightarrow 0
\]
forces
\[
A=0,\qquad \alpha(n)=2n \ \text{or}\ \alpha(n)=-2n.
\]
A concrete realization is
\[
0\longrightarrow 0
\longrightarrow \mathbb Z
\xrightarrow{\ \times 2\ }
\mathbb Z
\longrightarrow \mathbb Z/2
\longrightarrow 0.
\]
Here
\[
(\times 2)(n)=2n.
\]
Its image is the subgroup of even integers:
\[
2\mathbb Z=\{\dots,-4,-2,0,2,4,\dots\}.
\]
The last map is reduction modulo \(2\), and therefore
\[
\mathbb Z/(2\mathbb Z)\cong \mathbb Z/2.
\]
Since multiplication by \(2\) on \(\mathbb Z\) has trivial kernel,
\[
\ker(\times 2)=0,
\]
the group \(A\) is the trivial group:
\[
A=0.
\]

\par

\textbf{(iv) The case \(A\cong \mathbb Z/6\) and \(\alpha=\pm\operatorname{id}_{\mathbb Z}\).}

We found that
\[
0\longrightarrow \mathbb Z/3 \longrightarrow A \longrightarrow \mathbb Z/2
\longrightarrow \mathbb Z \xrightarrow{\alpha} \mathbb Z \longrightarrow 0
\]
forces
\[
A\cong \mathbb Z/6,
\qquad
\alpha=\pm\operatorname{id}_{\mathbb Z}.
\]

A concrete model for the beginning is
\[
0\longrightarrow \mathbb Z/3
\xrightarrow{\ \iota\ }
\mathbb Z/6
\xrightarrow{\ \pi\ }
\mathbb Z/2
\longrightarrow 0,
\]
where
\[
\iota(\bar k \bmod 3)=\overline{2k}\bmod 6,
\qquad
\pi(\bar x \bmod 6)=\bar x \bmod 2.
\]
The image of \(\iota\) is
\[
\operatorname{im}(\iota)=\{\bar0,\bar2,\bar4\}\subset \mathbb Z/6.
\]
The kernel of \(\pi\) is precisely the even residue classes modulo \(6\):
\[
\ker(\pi)=\{\bar0,\bar2,\bar4\}.
\]
Hence
\[
\ker(\pi)=\operatorname{im}(\iota),
\]
so this part is exact.

After that, the map
\[
\mathbb Z/2\longrightarrow \mathbb Z
\]
must be the zero map, since \(\mathbb Z\) has no nonzero element of order \(2\). Therefore the end of the sequence is
\[
\mathbb Z/2 \longrightarrow \mathbb Z
\xrightarrow{\ \alpha\ }
\mathbb Z \longrightarrow 0,
\]
with \(\alpha\) an isomorphism. Thus
\[
\alpha=\operatorname{id}_{\mathbb Z}
\qquad\text{or}\qquad
\alpha=-\operatorname{id}_{\mathbb Z}.
\]

So a concrete exact sequence is
\[
0\longrightarrow \mathbb Z/3
\longrightarrow \mathbb Z/6
\longrightarrow \mathbb Z/2
\longrightarrow \mathbb Z
\xrightarrow{\ \operatorname{id}\ }
\mathbb Z
\longrightarrow 0,
\]
where the map \(\mathbb Z/2\to \mathbb Z\) is the zero homomorphism.

\par

\textbf{Summary of the concrete models.}

\[
\text{(i)}\qquad A\cong \mathbb Z\oplus \mathbb Z/2,
\]
with elements \((n,\bar\varepsilon)\).

\[
\text{(ii)}\qquad A\cong \mathbb Z/4
\quad\text{or}\quad
A\cong \mathbb Z/2\oplus \mathbb Z/2.
\]

\[
\text{(iii)}\qquad A=0,\qquad \alpha(n)=\pm 2n.
\]

\[
\text{(iv)}\qquad A\cong \mathbb Z/6,\qquad \alpha=\pm \operatorname{id}_{\mathbb Z}.
\]

These explicit realizations show what the abstract conclusions mean in concrete algebraic terms.
\end{solution}

\begin{problem}[CP-VIII-0054 --- Acyclic chain complexes of free abelian groups are chain contractible]
Suppose that \((C_\bullet,d)\) is a chain complex of finitely-generated free abelian groups with
\[
H_*(C_\bullet,d)=0.
\]
Show that
\[
\mathrm{id}:C_\bullet\to C_\bullet
\]
is chain homotopic to the zero map.
\end{problem}

\begin{solution}
We must show that there exist homomorphisms
\[
s_n:C_n\to C_{n+1}
\]
such that
\[
d_{n+1}s_n+s_{n-1}d_n=\mathrm{id}_{C_n}
\qquad\text{for all }n.
\]

\par

\textbf{Step 1: use acyclicity.}

Since the homology vanishes, for every \(n\) we have
\[
H_n(C_\bullet)=0.
\]
Thus
\[
Z_n=\ker d_n = \operatorname{im} d_{n+1}=B_n.
\]
So every cycle is a boundary.

Now consider the sequence
\[
0\longrightarrow B_n \longrightarrow C_n \xrightarrow{d_n} B_{n-1}\longrightarrow 0.
\]
This sequence is exact, because:
\begin{itemize}
	\item \(\ker d_n=Z_n=B_n\),
	\item \(\operatorname{im} d_n=B_{n-1}\).
\end{itemize}

Since \(C_n\) is finitely-generated free abelian, the subgroup \(B_n\subseteq C_n\) is also free abelian. Hence \(B_{n-1}\) is free abelian as well, and the exact sequence splits. Therefore for each \(n\) there exists a subgroup \(D_n\subseteq C_n\) such that
\[
C_n=B_n\oplus D_n,
\]
and the restriction
\[
d_n|_{D_n}:D_n\to B_{n-1}
\]
is an isomorphism.

\par

\textbf{Step 2: define the contracting homotopy.}

We now define
\[
s_{n-1}:C_{n-1}\to C_n
\]
for each \(n\).

Since
\[
C_{n-1}=B_{n-1}\oplus D_{n-1},
\]
every element \(x\in C_{n-1}\) can be written uniquely as
\[
x=b+d,
\qquad
b\in B_{n-1},\ d\in D_{n-1}.
\]
Define
\[
s_{n-1}(x)=\bigl(d_n|_{D_n}\bigr)^{-1}(b).
\]
In other words:
\begin{itemize}
	\item on the boundary part \(B_{n-1}\), the map \(s_{n-1}\) is the inverse of \(d_n|_{D_n}\);
	\item on the complementary part \(D_{n-1}\), the map \(s_{n-1}\) is zero.
\end{itemize}

So explicitly,
\[
s_{n-1}|_{B_{n-1}}=(d_n|_{D_n})^{-1},
\qquad
s_{n-1}|_{D_{n-1}}=0.
\]

\par

\textbf{Step 3: verify the homotopy identity.}

Take \(x\in C_n\). Write
\[
x=b+c,
\qquad
b\in B_n,\ c\in D_n.
\]

We compute
\[
(d_{n+1}s_n+s_{n-1}d_n)(x).
\]

First,
\[
s_n(x)=s_n(b+c)=s_n(b),
\]
because \(s_n\) vanishes on \(D_n\). Since \(b\in B_n\), the element \(s_n(b)\in D_{n+1}\) is chosen so that
\[
d_{n+1}s_n(b)=b.
\]
Hence
\[
d_{n+1}s_n(x)=b.
\]

Next,
\[
d_n(x)=d_n(b+c)=d_n(c),
\]
because \(b\in B_n=\ker d_n\). Since \(c\in D_n\), we have
\[
d_n(c)\in B_{n-1},
\]
and by definition of \(s_{n-1}\),
\[
s_{n-1}d_n(x)=s_{n-1}(d_n(c))=c.
\]

Therefore
\[
(d_{n+1}s_n+s_{n-1}d_n)(x)=b+c=x.
\]
Since this holds for every \(x\in C_n\), we get
\[
d_{n+1}s_n+s_{n-1}d_n=\mathrm{id}_{C_n}.
\]

Thus
\[
\mathrm{id}_{C_\bullet}\simeq 0.
\]


\medskip
\noindent\textbf{Conclusion}\par
\smallskip


We have constructed homomorphisms
\[
s_n:C_n\to C_{n+1}
\]
such that
\[
d_{n+1}s_n+s_{n-1}d_n=\mathrm{id}_{C_n}
\qquad\text{for all }n.
\]
Hence the identity chain map is chain homotopic to the zero map.

So every acyclic chain complex of finitely-generated free abelian groups is \emph{chain contractible}.

\par
\end{solution}

\begin{problem}[CP-VIII-0055 --- Contractible chain complexes and the reduced chain complex of a simplex]
Suppose \((C_*,d)\) is a chain complex over a ring \(R\). We say that \(C\) is \emph{contractible} if there is an \(R\)-linear map
\[
H_k : C_k \to C_{k+1}
\]
for all \(k\), such that
\[
dH + Hd = \operatorname{id}_C.
\]
Show that if \(C\) is contractible, then
\[
H_*(C)=0.
\]
Show moreover that \(\widetilde{S}_*(\Delta^n)\) is contractible for all \(n \ge 0\).
\end{problem}

\begin{solution}
We prove the two statements separately.


\medskip
\noindent\textbf{If \(C\) is contractible, then \(H_*(C)=0\)}\par
\smallskip


Recall that
\[
H_k(C)=\frac{Z_k(C)}{B_k(C)},
\qquad
Z_k(C)=\ker(d_k),
\qquad
B_k(C)=\operatorname{im}(d_{k+1}).
\]
So it is enough to show that every cycle is a boundary.

Let \(z\in C_k\) be a cycle. Then
\[
dz=0.
\]
Since \(C\) is contractible, we have
\[
dH+Hd=\operatorname{id}_C.
\]
Applying this identity to \(z\), we obtain
\[
z=(dH+Hd)(z)=dH(z)+H(dz)=dH(z)+H(0)=dH(z).
\]
Thus \(z\in \operatorname{im}(d)=B_k(C)\). Therefore every cycle is a boundary, so
\[
Z_k(C)\subseteq B_k(C).
\]
On the other hand, for every chain complex one always has
\[
B_k(C)\subseteq Z_k(C),
\]
because \(d^2=0\). Hence
\[
Z_k(C)=B_k(C),
\]
and therefore
\[
H_k(C)=\frac{Z_k(C)}{B_k(C)}=0
\]
for every \(k\).

Thus a contractible chain complex has trivial homology:
\[
H_*(C)=0.
\]


\medskip
\noindent\textbf{The reduced simplicial chain complex \(\widetilde{S}_*(\Delta^n)\) is contractible}\par
\smallskip


Let
\[
\Delta^n=[v_0,v_1,\dots,v_n]
\]
be the standard \(n\)-simplex. We will define an explicit contracting homotopy
\[
H_q : \widetilde{S}_q(\Delta^n)\to \widetilde{S}_{q+1}(\Delta^n)
\]
by coning simplices to the vertex \(v_0\).


\medskip
\noindent\textbf{Step 1: Definition of \(H\) in degrees \(q\ge 0\).}\par
\smallskip


Let
\[
\sigma=[v_{i_0},v_{i_1},\dots,v_{i_q}]
\]
be an oriented simplex of \(\Delta^n\), written with increasing indices \(i_0<i_1<\cdots<i_q\). Define
\[
H_q(\sigma)=
\begin{cases}
	[v_0,v_{i_0},v_{i_1},\dots,v_{i_q}], & \text{if } i_0>0,\\[1mm]
	0, & \text{if } i_0=0.
\end{cases}
\]
Then extend \(R\)-linearly to all chains.

Since we are working with the reduced chain complex, we must also define the map in degree \(-1\):
\[
H_{-1}: \widetilde{S}_{-1}(\Delta^n)=R \longrightarrow \widetilde{S}_0(\Delta^n),
\qquad
H_{-1}(1)=[v_0].
\]

We will show that
\[
\partial H + H\partial=\operatorname{id}_{\widetilde{S}_*(\Delta^n)}.
\]


\medskip
\noindent\textbf{Step 2: Verification in degrees \(q\ge 1\).}\par
\smallskip


Take a simplex
\[
\sigma=[v_{i_0},\dots,v_{i_q}].
\]

\par

\noindent
\emph{Case 1: \(i_0>0\).}

Then
\[
H(\sigma)=[v_0,v_{i_0},\dots,v_{i_q}].
\]
Its boundary is
\[
\partial H(\sigma)
=
\partial [v_0,v_{i_0},\dots,v_{i_q}]
=
[v_{i_0},\dots,v_{i_q}]
+
\sum_{j=1}^{q+1}(-1)^j
[v_0,v_{i_0},\dots,\widehat{v_{i_{j-1}}},\dots,v_{i_q}].
\]
Reindexing the sum gives
\[
\partial H(\sigma)
=
\sigma
-
\sum_{j=0}^{q}(-1)^j
[v_0,v_{i_0},\dots,\widehat{v_{i_j}},\dots,v_{i_q}].
\]
But the second term is exactly \(H(\partial \sigma)\), because every face
\[
[v_{i_0},\dots,\widehat{v_{i_j}},\dots,v_{i_q}]
\]
still has smallest vertex \(>0\). Therefore
\[
\partial H(\sigma)=\sigma-H(\partial \sigma),
\]
hence
\[
(\partial H+H\partial)(\sigma)=\sigma.
\]

\par

\noindent
\emph{Case 2: \(i_0=0\).}

Then
\[
\sigma=[v_0,v_{i_1},\dots,v_{i_q}]
\]
and by definition
\[
H(\sigma)=0.
\]
Now
\[
\partial \sigma
=
[v_{i_1},\dots,v_{i_q}]
+
\sum_{j=1}^{q}(-1)^j
[v_0,v_{i_1},\dots,\widehat{v_{i_j}},\dots,v_{i_q}].
\]
Applying \(H\), all terms in the sum with first vertex \(v_0\) are sent to \(0\), while the first term becomes
\[
H([v_{i_1},\dots,v_{i_q}])=[v_0,v_{i_1},\dots,v_{i_q}]=\sigma.
\]
Thus
\[
H(\partial \sigma)=\sigma.
\]
Since \(H(\sigma)=0\), we obtain
\[
(\partial H+H\partial)(\sigma)=0+\sigma=\sigma.
\]

So for every simplex in degree \(q\ge 1\),
\[
(\partial H+H\partial)(\sigma)=\sigma.
\]


\medskip
\noindent\textbf{Step 3: Verification in degree \(0\).}\par
\smallskip


Now let us check the identity on vertices.

Take a \(0\)-simplex \([v_i]\).

\par

\noindent
\emph{If \(i=0\):} then \(H([v_0])=0\), and in the reduced complex
\[
\partial([v_0])=1\in \widetilde{S}_{-1}(\Delta^n)=R.
\]
Hence
\[
H\partial([v_0])=H_{-1}(1)=[v_0].
\]
Therefore
\[
(\partial H+H\partial)([v_0])=0+[v_0]=[v_0].
\]

\par

\noindent
\emph{If \(i>0\):} then
\[
H([v_i])=[v_0,v_i].
\]
So
\[
\partial H([v_i])=\partial [v_0,v_i]=[v_i]-[v_0].
\]
Also
\[
\partial([v_i])=1,
\qquad
H\partial([v_i])=H_{-1}(1)=[v_0].
\]
Hence
\[
(\partial H+H\partial)([v_i])
=
([v_i]-[v_0])+[v_0]
=
[v_i].
\]

Thus the identity also holds in degree \(0\).


\medskip
\noindent\textbf{Step 4: Verification in degree \(-1\).}\par
\smallskip


Finally, for the generator \(1\in R=\widetilde{S}_{-1}(\Delta^n)\), we have
\[
\partial H_{-1}(1)=\partial([v_0])=1.
\]
So
\[
(\partial H+H\partial)(1)=1.
\]

Therefore
\[
\partial H + H\partial=\operatorname{id}_{\widetilde{S}_*(\Delta^n)}.
\]
This proves that \(\widetilde{S}_*(\Delta^n)\) is contractible.


\medskip
\noindent\textbf{Conclusion}\par
\smallskip


We have shown:

\begin{enumerate}
	\item If a chain complex \(C\) is contractible, then all of its homology groups vanish:
	\[
	H_*(C)=0.
	\]
	\item For every \(n\ge 0\), the reduced simplicial chain complex of the simplex \(\Delta^n\) is contractible:
	\[
	\widetilde{S}_*(\Delta^n)\ \text{is contractible}.
	\]
\end{enumerate}

In particular,
\[
\widetilde{H}_*(\Delta^n)=0
\qquad \text{for all } n\ge 0.
\]
\end{solution}

\begin{problem}[CP-VIII-0056 --- Splitting over a free quotient and classification of two exact sequences]
Suppose
\[
0\longrightarrow H \longrightarrow G \longrightarrow \mathbb Z^n \longrightarrow 0
\]
is an exact sequence of abelian groups. Show that
\[
G\cong H\oplus \mathbb Z^n.
\]

Determine the possible isomorphism types of \(G\) in the exact sequences
\[
0\longrightarrow \mathbb Z/4 \longrightarrow G \longrightarrow \mathbb Z/4 \longrightarrow 0
\]
and
\[
0\longrightarrow \mathbb Z \longrightarrow G \longrightarrow \mathbb Z/4 \longrightarrow 0.
\]
\end{problem}

\begin{solution}

\medskip
\noindent\textbf{Splitting of the sequence with quotient \(\mathbb Z^n\)}\par
\smallskip


\noindent
Let
\[
\pi:G\to \mathbb Z^n
\]
be the surjective homomorphism in the exact sequence
\[
0\to H\to G\overset{\pi}{\to}\mathbb Z^n\to 0.
\]
Let \(e_1,\dots,e_n\) be the standard basis of \(\mathbb Z^n\). Since \(\pi\) is surjective, for each \(i\) there exists \(g_i\in G\) such that
\[
\pi(g_i)=e_i.
\]
Define
\[
s:\mathbb Z^n\to G,
\qquad
s(a_1,\dots,a_n)=a_1g_1+\cdots+a_ng_n.
\]
Then
\[
\pi(s(a_1,\dots,a_n))
=
a_1e_1+\cdots+a_ne_n
=
(a_1,\dots,a_n),
\]
so
\[
\pi\circ s=\mathrm{id}_{\mathbb Z^n}.
\]
Thus \(s\) is a right inverse of \(\pi\), so the sequence splits. Therefore
\[
G\cong H\oplus \mathbb Z^n.
\]

Hence every short exact sequence of abelian groups with quotient \(\mathbb Z^n\) is split.


\medskip
\noindent\textbf{Classification of}\par
\smallskip

\[
0\longrightarrow \mathbb Z/4 \longrightarrow G \longrightarrow \mathbb Z/4 \longrightarrow 0
\]

\noindent
Since both kernel and quotient have order \(4\), the group \(G\) has order
\[
|G|=4\cdot 4=16.
\]
By the structure theorem for finite abelian groups, the abelian groups of order \(16\) are
\[
\mathbb Z/16,\qquad
\mathbb Z/8\oplus \mathbb Z/2,\qquad
\mathbb Z/4\oplus \mathbb Z/4,\qquad
\mathbb Z/4\oplus \mathbb Z/2\oplus \mathbb Z/2,\qquad
(\mathbb Z/2)^4.
\]
We must determine which of these admit a subgroup isomorphic to \(\mathbb Z/4\) whose quotient is also isomorphic to \(\mathbb Z/4\).

\par

\noindent
\textbf{1. The group \(\mathbb Z/16\).}
It contains the subgroup \(4\mathbb Z/16\cong \mathbb Z/4\), and
\[
(\mathbb Z/16)/(4\mathbb Z/16)\cong \mathbb Z/4.
\]
So this case occurs.

\par

\noindent
\textbf{2. The group \(\mathbb Z/8\oplus \mathbb Z/2\).}
Take the subgroup generated by \((2,1)\). This element has order \(4\), so
\[
\langle (2,1)\rangle \cong \mathbb Z/4.
\]
The quotient has order \(4\), and one checks that it is cyclic, hence
\[
(\mathbb Z/8\oplus \mathbb Z/2)/\langle (2,1)\rangle \cong \mathbb Z/4.
\]
So this case also occurs.

\par

\noindent
\textbf{3. The group \(\mathbb Z/4\oplus \mathbb Z/4\).}
Take the first factor as kernel:
\[
\mathbb Z/4\oplus 0 \cong \mathbb Z/4.
\]
Then
\[
(\mathbb Z/4\oplus \mathbb Z/4)/(\mathbb Z/4\oplus 0)\cong \mathbb Z/4.
\]
So this case occurs as well.

\par

\noindent
\textbf{4. The groups \(\mathbb Z/4\oplus \mathbb Z/2\oplus \mathbb Z/2\) and \((\mathbb Z/2)^4\).}
These do not occur. Indeed, if \(G/N\cong \mathbb Z/4\) is cyclic, then \(G\) is generated by \(N\) together with one lift of a generator of \(G/N\). Since \(N\cong \mathbb Z/4\) is generated by one element, this would imply that \(G\) is generated by at most two elements. But both
\[
\mathbb Z/4\oplus \mathbb Z/2\oplus \mathbb Z/2
\quad \text{and} \quad
(\mathbb Z/2)^4
\]
require at least three generators, so they are impossible.

Therefore the only possible isomorphism types are
\[
\boxed{
	G\cong \mathbb Z/16,
	\qquad
	G\cong \mathbb Z/8\oplus \mathbb Z/2,
	\qquad
	G\cong \mathbb Z/4\oplus \mathbb Z/4.
}
\]


\medskip
\noindent\textbf{Classification of}\par
\smallskip

\[
0\longrightarrow \mathbb Z \longrightarrow G \longrightarrow \mathbb Z/4 \longrightarrow 0
\]

\noindent
Now \(G\) is a finitely generated abelian group whose quotient by a subgroup isomorphic to \(\mathbb Z\) is finite. Hence \(G\) has rank \(1\). By the structure theorem for finitely generated abelian groups,
\[
G\cong \mathbb Z\oplus T,
\]
where \(T\) is a finite abelian group.

Because \(G/\mathbb Z\cong \mathbb Z/4\), the quotient has order \(4\). Hence the subgroup \(\mathbb Z\subset G\) has index \(4\). This severely restricts the possibilities. The only possible isomorphism types are
\[
\mathbb Z,\qquad \mathbb Z\oplus \mathbb Z/2,\qquad \mathbb Z\oplus \mathbb Z/4.
\]
We now show that all three actually occur.

\par

\noindent
\textbf{1. The group \(\mathbb Z\).}
Consider
\[
0\longrightarrow \mathbb Z \overset{\iota}{\longrightarrow} \mathbb Z \longrightarrow \mathbb Z/4 \longrightarrow 0,
\]
where \(\iota(n)=4n\), and the map \(\mathbb Z\to \mathbb Z/4\) is reduction modulo \(4\). Then the kernel is \(4\mathbb Z\cong \mathbb Z\), so this gives an exact sequence of the required form. Thus \(G\cong \mathbb Z\) occurs.

\par

\noindent
\textbf{2. The group \(\mathbb Z\oplus \mathbb Z/2\).}
Let \(G=\mathbb Z\oplus \mathbb Z/2\), and define a homomorphism
\[
\varphi:\mathbb Z\oplus \mathbb Z/2\to \mathbb Z/4
\]
by
\[
\varphi(n,\overline a)=\overline{n+2a}.
\]
This is surjective. Its kernel is generated by \((2,\overline 1)\), hence
\[
\ker(\varphi)\cong \mathbb Z.
\]
Therefore we obtain an exact sequence
\[
0\longrightarrow \mathbb Z \longrightarrow \mathbb Z\oplus \mathbb Z/2 \longrightarrow \mathbb Z/4 \longrightarrow 0.
\]
So \(G\cong \mathbb Z\oplus \mathbb Z/2\) occurs.

\par

\noindent
\textbf{3. The group \(\mathbb Z\oplus \mathbb Z/4\).}
Let \(G=\mathbb Z\oplus \mathbb Z/4\), and define
\[
\psi:\mathbb Z\oplus \mathbb Z/4\to \mathbb Z/4
\]
by
\[
\psi(n,\overline a)=\overline{n+a}.
\]
This map is surjective. Its kernel is generated by \((1,\overline{-1})\), so
\[
\ker(\psi)\cong \mathbb Z.
\]
Hence we get an exact sequence
\[
0\longrightarrow \mathbb Z \longrightarrow \mathbb Z\oplus \mathbb Z/4 \longrightarrow \mathbb Z/4 \longrightarrow 0.
\]
So \(G\cong \mathbb Z\oplus \mathbb Z/4\) also occurs.

Therefore the possible isomorphism types are exactly
\[
\boxed{
	G\cong \mathbb Z,
	\qquad
	G\cong \mathbb Z\oplus \mathbb Z/2,
	\qquad
	G\cong \mathbb Z\oplus \mathbb Z/4.
}
\]


\medskip
\noindent\textbf{Conclusion}\par
\smallskip


For the split exact sequence with free quotient, we have
\[
\boxed{G\cong H\oplus \mathbb Z^n.}
\]

For
\[
0\to \mathbb Z/4\to G\to \mathbb Z/4\to 0,
\]
the possible isomorphism types are
\[
\boxed{
	\mathbb Z/16,\quad
	\mathbb Z/8\oplus \mathbb Z/2,\quad
	\mathbb Z/4\oplus \mathbb Z/4.
}
\]

For
\[
0\to \mathbb Z\to G\to \mathbb Z/4\to 0,
\]
the possible isomorphism types are
\[
\boxed{
	\mathbb Z,\quad
	\mathbb Z\oplus \mathbb Z/2,\quad
	\mathbb Z\oplus \mathbb Z/4.
}
\]
\end{solution}

\begin{problem}[CP-VIII-0057 --- The mapping torus exact sequence and computations]
Let \(f:X\to X\) be a homeomorphism, and let \(Y\) be the quotient of \(X\times [0,1]\) obtained by identifying \((x,0)\) and \((f(x),1)\). Show there is a long exact sequence
\[
\begin{aligned}
\cdots&\longrightarrow H_{n+1}(Y)\longrightarrow H_n(X)
\xrightarrow{\,1-f_*\,}H_n(X)\\
&\longrightarrow H_n(Y)\longrightarrow H_{n-1}(X)
\xrightarrow{\,1-f_*\,}H_{n-1}(X)\longrightarrow\cdots .
\end{aligned}
\]
Compute \(H_*(Y)\) in the following cases:
\begin{enumerate}
	\item \(X=S^n\) and \(f\) is the antipodal map;
	\item \(X=T^2=\mathbb R^2/\mathbb Z^2\) and \(f\) is multiplication by
	\[
	A=\begin{pmatrix}3&4\\2&3\end{pmatrix}.
	\]
\end{enumerate}

\par
\end{problem}

\begin{solution}
The space \(Y\) is the \emph{mapping torus} of \(f\), usually denoted \(T_f\). It is obtained from \(X\times [0,1]\) by gluing the two ends using \(f\):
\[
(x,0)\sim (f(x),1).
\]
There is a standard long exact sequence in homology for the mapping torus, often called the \emph{Wang exact sequence}:
\[
\begin{aligned}
\cdots&\longrightarrow H_{n+1}(Y)\longrightarrow H_n(X)
\xrightarrow{\,1-f_*\,}H_n(X)\\
&\longrightarrow H_n(Y)\longrightarrow H_{n-1}(X)
\xrightarrow{\,1-f_*\,}H_{n-1}(X)\longrightarrow\cdots .
\end{aligned}
\]
This is exactly the desired sequence.

A useful reformulation is that for each \(n\) there is a short exact sequence
\[
\begin{aligned}
0&\longrightarrow
\operatorname{coker}\!\bigl(1-f_*:H_n(X)\to H_n(X)\bigr)
\longrightarrow H_n(Y)\\
&\longrightarrow
\ker\!\bigl(1-f_*:H_{n-1}(X)\to H_{n-1}(X)\bigr)
\longrightarrow 0.
\end{aligned}
\]
We now compute the homology in the two cases.


\medskip
\noindent\textbf{Case 1: \(X=S^n\), \(f\) the antipodal map}\par
\smallskip


Recall that
\[
H_k(S^n)=
\begin{cases}
	\mathbb Z, & k=0,n,\\
	0, & \text{otherwise}.
\end{cases}
\]
Let \(a:S^n\to S^n\) denote the antipodal map \(a(x)=-x\). Its degree is
\[
\deg(a)=(-1)^{n+1}.
\]
Therefore:
\begin{itemize}
	\item on \(H_0(S^n)\cong \mathbb Z\), we have \(a_*=\mathrm{id}\), since \(S^n\) is connected;
	\item on \(H_n(S^n)\cong \mathbb Z\), we have \(a_*\) equal to multiplication by \((-1)^{n+1}\).
\end{itemize}

Thus on top homology,
\[
1-a_*=1-(-1)^{n+1}.
\]


\medskip
\noindent\textbf{If \(n\) is odd.}\par
\smallskip

Then \((-1)^{n+1}=1\), so
\[
1-a_*=0
\]
on \(H_n(S^n)\). Hence
\[
\ker(1-a_*)=\mathbb Z,\qquad \operatorname{coker}(1-a_*)=\mathbb Z.
\]
From the short exact sequence above we obtain:
\[
H_{n+1}(Y)\cong \mathbb Z,
\qquad
H_n(Y)\cong \mathbb Z.
\]
Also \(1-a_*=0\) on \(H_0(S^n)\), so
\[
H_1(Y)\cong \mathbb Z,
\qquad
H_0(Y)\cong \mathbb Z.
\]
All other homology groups vanish, provided \(n>1\). Therefore for odd \(n>1\),
\[
H_k(Y)=
\begin{cases}
	\mathbb Z, & k=0,1,n,n+1,\\
	0, & \text{otherwise}.
\end{cases}
\]

For \(n=1\), the antipodal map on \(S^1\) is rotation by \(\pi\), which is homotopic to the identity, and its mapping torus is a torus. Hence
\[
H_0(Y)\cong \mathbb Z,\qquad
H_1(Y)\cong \mathbb Z^2,\qquad
H_2(Y)\cong \mathbb Z.
\]


\medskip
\noindent\textbf{If \(n\) is even.}\par
\smallskip

Then \((-1)^{n+1}=-1\), so
\[
1-a_*=1-(-1)=2
\]
on \(H_n(S^n)\cong \mathbb Z\). Hence
\[
\ker(1-a_*)=0,\qquad \operatorname{coker}(1-a_*)\cong \mathbb Z/2.
\]
Therefore
\[
H_{n+1}(Y)=0,\qquad H_n(Y)\cong \mathbb Z/2.
\]
Again, on \(H_0\) we have \(1-a_*=0\), so
\[
H_1(Y)\cong \mathbb Z,\qquad H_0(Y)\cong \mathbb Z.
\]
All other groups vanish. Thus for even \(n\),
\[
H_k(Y)=
\begin{cases}
	\mathbb Z, & k=0,1,\\
	\mathbb Z/2, & k=n,\\
	0, & \text{otherwise}.
\end{cases}
\]


\medskip
\noindent\textbf{Conclusion for Case 1.}\par
\smallskip

For the mapping torus of the antipodal map on \(S^n\):
\[
H_k(Y)=
\begin{cases}
	\mathbb Z, & \text{if $k=0,1,n,n+1$ and $n$ odd and $n>1$,}\\
	\mathbb Z/2, & \text{if $k=n$ and $n$ even,}\\
	\mathbb Z, & \text{if $k=0,1$ and $n$ even,}\\
	0, & \text{otherwise.}
\end{cases}
\]
For \(n=1\), one gets the homology of \(T^2\).


\medskip
\noindent\textbf{Case 2: \(X=T^2\), \(f\) induced by \(A=\begin{pmatrix}3&4\\2&3\end{pmatrix}\)}\par
\smallskip


We know
\[
H_0(T^2)\cong \mathbb Z,\qquad
H_1(T^2)\cong \mathbb Z^2,\qquad
H_2(T^2)\cong \mathbb Z.
\]
The induced map on \(H_1(T^2)\) is the matrix \(A\):
\[
f_*|_{H_1}=A.
\]
Also
\[
\det(A)=3\cdot 3-4\cdot 2=1,
\]
so \(f\) preserves orientation. Hence on \(H_2(T^2)\cong \mathbb Z\),
\[
f_*=\mathrm{id}.
\]
On \(H_0(T^2)\cong \mathbb Z\), we also have \(f_*=\mathrm{id}\), since \(T^2\) is connected.

Therefore
\[
1-f_*=
\begin{cases}
	0 & \text{on }H_0(T^2),\\[1ex]
	I-A & \text{on }H_1(T^2),\\[1ex]
	0 & \text{on }H_2(T^2).
\end{cases}
\]
Now
\[
I-A=
\begin{pmatrix}
	1&0\\
	0&1
\end{pmatrix}
-
\begin{pmatrix}
	3&4\\
	2&3
\end{pmatrix}
=
\begin{pmatrix}
	-2&-4\\
	-2&-2
\end{pmatrix}.
\]
Its determinant is
\[
(-2)(-2)-(-4)(-2)=4-8=-4\neq 0,
\]
so \(I-A\) is injective as a homomorphism \(\mathbb Z^2\to \mathbb Z^2\). Hence
\[
\ker(I-A)=0.
\]

To compute the cokernel, we reduce to Smith normal form. Ignoring the harmless sign,
\[
\begin{pmatrix}
	2&4\\
	2&2
\end{pmatrix}
\sim
\begin{pmatrix}
	2&0\\
	0&2
\end{pmatrix},
\]
so
\[
\operatorname{coker}(I-A)\cong \mathbb Z/2\oplus \mathbb Z/2.
\]

We now compute \(H_k(Y)\).


\medskip
\noindent\textbf{Computation of \(H_3(Y)\).}\par
\smallskip

From the Wang sequence,
\[
H_3(Y)\cong \ker(1-f_*:H_2(T^2)\to H_2(T^2)).
\]
But \(1-f_*=0\) on \(H_2(T^2)\cong \mathbb Z\), hence
\[
H_3(Y)\cong \mathbb Z.
\]


\medskip
\noindent\textbf{Computation of \(H_2(Y)\).}\par
\smallskip

Using the short exact sequence,
\[
0\to \operatorname{coker}(1-f_*:H_2\to H_2)\to H_2(Y)\to \ker(1-f_*:H_1\to H_1)\to 0.
\]
Since \(1-f_*=0\) on \(H_2\), we have
\[
\operatorname{coker}(1-f_*:H_2\to H_2)\cong \mathbb Z,
\]
and since \(\ker(I-A)=0\), we get
\[
0\to \mathbb Z\to H_2(Y)\to 0\to 0.
\]
Therefore
\[
H_2(Y)\cong \mathbb Z.
\]


\medskip
\noindent\textbf{Computation of \(H_1(Y)\).}\par
\smallskip

Again,
\[
0\to \operatorname{coker}(I-A)\to H_1(Y)\to \ker(1-f_*:H_0\to H_0)\to 0.
\]
Since \(1-f_*=0\) on \(H_0\cong \mathbb Z\), we have
\[
\ker(1-f_*:H_0\to H_0)\cong \mathbb Z.
\]
Thus
\[
0\to \mathbb Z/2\oplus \mathbb Z/2\to H_1(Y)\to \mathbb Z\to 0.
\]
Because \(\mathbb Z\) is free, this extension splits, so
\[
H_1(Y)\cong \mathbb Z\oplus \mathbb Z/2\oplus \mathbb Z/2.
\]


\medskip
\noindent\textbf{Computation of \(H_0(Y)\).}\par
\smallskip

The mapping torus \(Y\) is connected, hence
\[
H_0(Y)\cong \mathbb Z.
\]


\medskip
\noindent\textbf{Conclusion for Case 2.}\par
\smallskip

We obtain
\[
H_k(Y)=
\begin{cases}
	\mathbb Z, & k=0,2,3,\\
	\mathbb Z\oplus \mathbb Z/2\oplus \mathbb Z/2, & k=1,\\
	0, & \text{otherwise}.
\end{cases}
\]
\end{solution}

\begin{problem}[CP-VIII-0058 --- Cancellation in a chain complex]
Suppose \((C,d)\) is a chain complex such that
\[
C_n=C_n'\oplus A,\qquad C_{n-1}=C_{n-1}'\oplus A,
\]
and suppose that the component of
\[
d_n:C_n'\oplus A\longrightarrow C_{n-1}'\oplus A
\]
mapping \(A\) to \(A\) is the identity map. Show that \((C,d)\) is chain homotopy equivalent to a chain complex \((C',d')\), where
\[
C_i'=C_i \quad \text{for } i\neq n,n-1,
\]
\[
C_n'=C_n',\qquad C_{n-1}'=C_{n-1}',
\]
\[
d_i'=d_i \quad \text{for } i\neq n-1,n,n+1,
\]
and \(d'_{n+1}\) is the composition of \(d_{n+1}\) with the projection onto \(C_n'\).

(Hint: use \(d^2=0\) to determine \(d_n'\).)

\par
\end{problem}

\begin{solution}
This is the standard \emph{cancellation lemma}. The hypothesis says that inside the complex there is a summand of the form
\[
A \xrightarrow{\;\mathrm{id}\;} A,
\]
and such a piece is contractible. The point is to remove it and adjust the neighboring differential accordingly.

Write
\[
d_n:C_n'\oplus A\longrightarrow C_{n-1}'\oplus A
\]
in block form:
\[
d_n=
\begin{pmatrix}
	u & v\\
	w & 1
\end{pmatrix},
\]
where
\[
u:C_n'\to C_{n-1}',\qquad
v:A\to C_{n-1}',\qquad
w:C_n'\to A,
\]
and \(1:A\to A\) is the identity.

We now determine the new differential \(d_n':C_n'\to C_{n-1}'\). Let
\[
d_{n-1}=(p,q):C_{n-1}'\oplus A\to C_{n-2},
\]
and let
\[
d_{n+1}=\binom{r}{s}:C_{n+1}\to C_n'\oplus A.
\]
Since \(d^2=0\), we have
\[
d_{n-1}d_n=0.
\]
Thus
\[
0=(p,q)
\begin{pmatrix}
	u & v\\
	w & 1
\end{pmatrix}
=
\bigl(pu+qw,\;pv+q\bigr).
\]
Hence
\[
q=-pv,
\]
and therefore
\[
pu-pvw=0,
\]
that is,
\[
p(u-vw)=0.
\]

Similarly, from
\[
d_n d_{n+1}=0
\]
we obtain
\[
0=
\begin{pmatrix}
	u & v\\
	w & 1
\end{pmatrix}
\binom{r}{s}
=
\binom{ur+vs}{wr+s}.
\]
So
\[
s=-wr,
\]
and hence
\[
ur-vwr=(u-vw)r=0.
\]

These identities show that the correct new differential is
\[
d_n'=u-vw.
\]
Also, by the statement of the problem,
\[
d'_{n+1}=r,
\]
that is, \(d_{n+1}\) followed by projection onto the \(C_n'\)-summand. The new complex is therefore
\[
\cdots \to C_{n+1}\xrightarrow{\,r\,}C_n'\xrightarrow{\,u-vw\,}C_{n-1}'\xrightarrow{\,p\,}C_{n-2}\to\cdots
\]
and agrees with \(C\) in every other degree.

We now construct explicit chain maps
\[
f:C\to C',
\qquad
g:C'\to C.
\]

For \(i\neq n,n-1\), let \(f_i=g_i=\mathrm{id}\).

In degree \(n\), define
\[
f_n:C_n'\oplus A\to C_n',
\qquad
f_n(x,a)=x,
\]
and
\[
g_n:C_n'\to C_n'\oplus A,
\qquad
g_n(x)=(x,-wx).
\]

In degree \(n-1\), define
\[
f_{n-1}:C_{n-1}'\oplus A\to C_{n-1}',
\qquad
f_{n-1}(y,b)=y-vb,
\]
and
\[
g_{n-1}:C_{n-1}'\to C_{n-1}'\oplus A,
\qquad
g_{n-1}(y)=(y,0).
\]

A direct computation shows that \(f\) and \(g\) are chain maps. For example, in degree \(n\),
\[
f_{n-1}d_n(x,a)
=
f_{n-1}(ux+va,\,wx+a)
=
ux+va-v(wx+a)
=
(u-vw)x
=
d_n' f_n(x,a),
\]
which is exactly the chain map condition.

Next we check that
\[
f\circ g=\mathrm{id}_{C'}.
\]
In degree \(n\),
\[
f_n(g_n(x))=f_n(x,-wx)=x,
\]
and in degree \(n-1\),
\[
f_{n-1}(g_{n-1}(y))=f_{n-1}(y,0)=y.
\]
So indeed \(fg=\mathrm{id}_{C'}\).

It remains to prove that
\[
g\circ f \simeq \mathrm{id}_C.
\]
Define a homotopy \(h:C\to C[1]\) by
\[
h_i=0 \quad \text{for } i\neq n-1,
\]
and
\[
h_{n-1}:C_{n-1}'\oplus A\to C_n'\oplus A,
\qquad
h_{n-1}(y,b)=(0,b).
\]
We now compute \(dh+hd\).

In degree \(n-1\),
\[
(dh+hd)_{n-1}(y,b)=d_n(0,b)=(vb,b).
\]
On the other hand,
\[
(\mathrm{id}-gf)_{n-1}(y,b)
=
(y,b)-g_{n-1}(y-vb)
=
(y,b)-(y-vb,0)
=
(vb,b).
\]

In degree \(n\),
\[
(dh+hd)_n(x,a)=h_{n-1}(d_n(x,a)).
\]
Since
\[
d_n(x,a)=(ux+va,\;wx+a),
\]
we get
\[
h_{n-1}(d_n(x,a))=(0,wx+a).
\]
Also,
\[
(\mathrm{id}-gf)_n(x,a)
=
(x,a)-g_n(f_n(x,a))
=
(x,a)-g_n(x)
=
(x,a)-(x,-wx)
=
(0,a+wx).
\]
These are equal.

In every other degree, both \(dh+hd\) and \(\mathrm{id}-gf\) vanish. Thus
\[
dh+hd=\mathrm{id}_C-gf.
\]
Therefore
\[
gf\simeq \mathrm{id}_C.
\]

We conclude that \(f\) and \(g\) are chain homotopy inverses, so
\[
\boxed{(C,d)\simeq (C',d').}
\]

\par

The essential geometric meaning is that a summand
\[
A\xrightarrow{\mathrm{id}}A
\]
is contractible and may be canceled, at the cost of replacing the remaining differential \(u\) by the corrected differential
\[
u-vw.
\]
\end{solution}

\begin{problem}[CP-VIII-0059 --- Acyclic free chain complexes, quasi-isomorphisms, and the Five Lemma]
Suppose \(C\) and \(C'\) are free, finitely generated chain complexes over \(\mathbb Z\).

\begin{enumerate}
	\item[(a)] Use Problem 14 to prove that if \(H_*(C)=0\), then \(C\) is contractible.
	\item[(b)] Suppose \(f:C\to C'\) is a chain map. Using Problem 5, show that \(f\) is a chain homotopy equivalence if and only if
	\[
	f_*:H_*(C)\to H_*(C')
	\]
	is an isomorphism.
	\item[(c)] Use Problems 5 and 14 to give a proof of the Five Lemma.
\end{enumerate}

\par
\end{problem}

\begin{solution}

\medskip
\noindent\textbf{(a) If \(H_*(C)=0\), then \(C\) is contractible}\par
\smallskip


We argue by induction on the total rank
\[
N=\sum_i \operatorname{rank}(C_i).
\]
If \(N=0\), then \(C=0\), so the statement is trivial.

Assume \(N>0\). Let \(n\) be the largest index such that
\[
C_n\neq 0.
\]
Then \(C_{n+1}=0\), so
\[
H_n(C)=\ker(d_n)/\operatorname{im}(d_{n+1})=\ker(d_n).
\]
Since \(H_n(C)=0\), it follows that
\[
\ker(d_n)=0,
\]
hence
\[
d_n:C_n\to C_{n-1}
\]
is injective.

Let
\[
B_{n-1}=\operatorname{im}(d_n)\subseteq C_{n-1}.
\]
Because \(C_n\) is free and \(d_n\) is injective, \(B_{n-1}\cong C_n\) is free. Since \(H_{n-1}(C)=0\), we have
\[
Z_{n-1}=B_{n-1}.
\]
Also,
\[
C_{n-1}/B_{n-1}\cong B_{n-2},
\]
and \(B_{n-2}\) is a subgroup of the free abelian group \(C_{n-2}\), hence is free. Therefore the short exact sequence
\[
0\to B_{n-1}\to C_{n-1}\to B_{n-2}\to 0
\]
splits. So there is a decomposition
\[
C_{n-1}\cong C_{n-1}'\oplus B_{n-1}.
\]

Since \(B_{n-1}\cong C_n\), we may identify both with a free abelian group \(A\). Under these identifications,
\[
C_n\cong A,\qquad C_{n-1}\cong C_{n-1}'\oplus A,
\]
and the differential \(d_n\) has \(A\to A\) component equal to the identity.

Thus Problem 14 applies: we may cancel this pair and obtain a smaller chain complex \(D\) such that
\[
C\simeq D.
\]
Since chain homotopy equivalent complexes have isomorphic homology,
\[
H_*(D)\cong H_*(C)=0.
\]
Also the total rank of \(D\) is strictly smaller than that of \(C\). By the induction hypothesis, \(D\) is contractible. Hence \(C\) is contractible as well.

Therefore
\[
\boxed{H_*(C)=0 \implies C \text{ is contractible}.}
\]


\medskip
\noindent\textbf{(b) A chain map is a chain homotopy equivalence if and only if it induces homology isomorphisms}\par
\smallskip


Let
\[
f:C\to C'
\]
be a chain map.


\medskip
\noindent\textbf{Forward implication.}\par
\smallskip

Assume first that \(f\) is a chain homotopy equivalence. Then there exists a chain map
\[
g:C'\to C
\]
such that
\[
gf\simeq \mathrm{id}_C,\qquad fg\simeq \mathrm{id}_{C'}.
\]
Chain-homotopic maps induce the same map on homology, so
\[
(gf)_*=\mathrm{id}_{H_*(C)},\qquad (fg)_*=\mathrm{id}_{H_*(C')}.
\]
Hence
\[
g_*f_*=\mathrm{id},\qquad f_*g_*=\mathrm{id},
\]
which shows that
\[
f_*:H_*(C)\to H_*(C')
\]
is an isomorphism.


\medskip
\noindent\textbf{Reverse implication.}\par
\smallskip

Now assume that
\[
f_*:H_*(C)\to H_*(C')
\]
is an isomorphism.

By Problem 5, the mapping cone \(\operatorname{Cone}(f)\) fits into a long exact sequence in homology, and from that sequence one obtains the standard criterion:
\[
f_* \text{ is an isomorphism } \Longleftrightarrow H_*(\operatorname{Cone}(f))=0.
\]
Thus
\[
H_*(\operatorname{Cone}(f))=0.
\]

Since \(C\) and \(C'\) are free finitely generated chain complexes, the cone \(\operatorname{Cone}(f)\) is also a free finitely generated chain complex. By part (a), every acyclic free finitely generated chain complex is contractible. Therefore
\[
\operatorname{Cone}(f)\text{ is contractible}.
\]

Again by Problem 5, a chain map has contractible mapping cone if and only if it is a chain homotopy equivalence. Hence \(f\) is a chain homotopy equivalence.

So we have proved
\[
\boxed{
	f \text{ is a chain homotopy equivalence }
	\Longleftrightarrow
	f_*:H_*(C)\to H_*(C') \text{ is an isomorphism}.
}
\]

This is the algebraic analogue of the principle that, for sufficiently good objects, a homology equivalence is already a stronger equivalence.


\medskip
\noindent\textbf{(c) Proof of the Five Lemma}\par
\smallskip


We now use Problems 5 and 14 to give a proof of the Five Lemma.

Suppose we have a commutative diagram of abelian groups with exact rows:
\[
\begin{array}{ccccccccc}
	A_1 &\longrightarrow& A_2 &\longrightarrow& A_3 &\longrightarrow& A_4 &\longrightarrow& A_5\\
	\downarrow \alpha_1 && \downarrow \alpha_2 && \downarrow \alpha_3 && \downarrow \alpha_4 && \downarrow \alpha_5\\
	B_1 &\longrightarrow& B_2 &\longrightarrow& B_3 &\longrightarrow& B_4 &\longrightarrow& B_5
\end{array}
\]
Assume that the rows are exact and that
\[
\alpha_1,\alpha_2,\alpha_4,\alpha_5
\]
are isomorphisms. We must prove that \(\alpha_3\) is also an isomorphism.

Treat the two rows as chain complexes concentrated in five consecutive degrees. Since the rows are exact, both chain complexes are acyclic. The vertical maps define a chain map
\[
\alpha:A_\bullet\to B_\bullet.
\]
Because both \(A_\bullet\) and \(B_\bullet\) are acyclic, the induced map on homology
\[
\alpha_*:H_*(A_\bullet)\to H_*(B_\bullet)
\]
is automatically an isomorphism, since both sides are zero. By part (b), it follows that \(\alpha\) is a chain homotopy equivalence. Equivalently, its mapping cone
\[
\operatorname{Cone}(\alpha)
\]
is contractible.

Now consider the structure of \(\operatorname{Cone}(\alpha)\). In the degrees corresponding to the first, second, fourth, and fifth columns, the maps
\[
\alpha_1,\alpha_2,\alpha_4,\alpha_5
\]
are isomorphisms. Hence within the mapping cone there appear four pieces of the form
\[
X\xrightarrow{\cong} X.
\]
By repeated application of the cancellation lemma from Problem 14, these pieces may be removed without changing the chain homotopy type.

After carrying out these cancellations, the cone reduces to a two-term complex concentrated in two consecutive degrees:
\[
0\longrightarrow A_3 \xrightarrow{\alpha_3} B_3 \longrightarrow 0.
\]
Because cancellation preserves chain homotopy type, this reduced complex is still contractible, hence acyclic.

But a two-term complex
\[
0\to A_3\xrightarrow{\alpha_3}B_3\to 0
\]
is acyclic if and only if \(\alpha_3\) is both injective and surjective. Indeed,
\[
H=\ker(\alpha_3)
\]
in the left degree, and
\[
H=\operatorname{coker}(\alpha_3)
\]
in the right degree. So acyclicity means
\[
\ker(\alpha_3)=0,\qquad \operatorname{coker}(\alpha_3)=0.
\]
Therefore \(\alpha_3\) is an isomorphism.

This proves the Five Lemma.

\par

So the three parts together give the following picture:

\begin{itemize}
	\item free finitely generated acyclic complexes are contractible;
	\item for such complexes, quasi-isomorphisms are the same as chain homotopy equivalences;
	\item the Five Lemma can be proved by passing to a mapping cone and then canceling the obvious isomorphic pieces.
\end{itemize}

The key philosophy is that exactness, cones, and cancellation turn homological information into a strong structural statement about the chain maps themselves.
\end{solution}

\begin{problem}[CP-VIII-0060 --- Mapping cone of a chain map]
Let
\[
f:(C,d)\to (C',d')
\]
be a chain map. The \emph{mapping cone} of \(f\) is the graded abelian group
\[
M(f)_n=C_{n-1}\oplus C'_n,
\]
with differential
\[
(d_f)_n=
\begin{pmatrix}
	d_{n-1} & 0\\
	(-1)^n f_{n-1} & d'_n
\end{pmatrix}.
\]
Show that \((M(f),d_f)\) is a chain complex, and that if \(f\sim g\), then \(M(f)\sim M(g)\). Show that
\[
H_*(M(f))=0
\quad\Longleftrightarrow\quad
f_*:H_*(C)\to H_*(C')
\]
is an isomorphism.

\par
\end{problem}

\begin{solution}
We proceed in three steps.


\medskip
\noindent\textbf{Step 1: \(M(f)\) is a chain complex}\par
\smallskip


We must prove that
\[
(d_f)_{n-1}(d_f)_n=0
\qquad\text{for all }n.
\]
Compute:
\[
(d_f)_{n-1}(d_f)_n
=
\begin{pmatrix}
	d_{n-2} & 0\\
	(-1)^{n-1}f_{n-2} & d'_{n-1}
\end{pmatrix}
\begin{pmatrix}
	d_{n-1} & 0\\
	(-1)^n f_{n-1} & d'_n
\end{pmatrix}.
\]
Multiplying the matrices gives
\[
(d_f)_{n-1}(d_f)_n
=
\begin{pmatrix}
	d_{n-2}d_{n-1} & 0\\[1mm]
	(-1)^{n-1}f_{n-2}d_{n-1}+(-1)^n d'_{n-1}f_{n-1} & d'_{n-1}d'_n
\end{pmatrix}.
\]
Now:

\[
d_{n-2}d_{n-1}=0
\qquad\text{and}\qquad
d'_{n-1}d'_n=0
\]
because \(C\) and \(C'\) are chain complexes.

Also, since \(f\) is a chain map,
\[
d'_{n-1}f_{n-1}=f_{n-2}d_{n-1}.
\]
Hence the lower-left entry becomes
\[
(-1)^{n-1}f_{n-2}d_{n-1}+(-1)^n d'_{n-1}f_{n-1}
=
(-1)^{n-1}\bigl(f_{n-2}d_{n-1}-d'_{n-1}f_{n-1}\bigr)=0.
\]
Therefore
\[
(d_f)_{n-1}(d_f)_n=0,
\]
so \((M(f),d_f)\) is indeed a chain complex.


\medskip
\noindent\textbf{Step 2: if \(f\sim g\), then \(M(f)\sim M(g)\)}\par
\smallskip


Assume \(f\sim g\). Thus there exist homomorphisms
\[
s_n:C_n\to C'_{n+1}
\]
such that
\[
f_n-g_n=d'_{n+1}s_n+s_{n-1}d_n
\qquad\text{for all }n.
\]
We define maps
\[
\Phi_n:M(f)_n\to M(g)_n
\]
by
\[
\Phi_n(x,y)=\bigl(x,\; y+(-1)^n s_{n-1}(x)\bigr),
\qquad x\in C_{n-1},\ y\in C'_n.
\]
In matrix form,
\[
\Phi_n=
\begin{pmatrix}
	1 & 0\\
	(-1)^n s_{n-1} & 1
\end{pmatrix}.
\]
This is an isomorphism, with inverse
\[
\Phi_n^{-1}=
\begin{pmatrix}
	1 & 0\\
	(-1)^{n+1}s_{n-1} & 1
\end{pmatrix}.
\]

We check that \(\Phi\) is a chain map. For \((x,y)\in M(f)_n\),
\[
(d_f)_n(x,y)=\bigl(d_{n-1}x,\; (-1)^n f_{n-1}x+d'_n y\bigr).
\]
Applying \(\Phi_{n-1}\), we get
\[
\Phi_{n-1}(d_f)_n(x,y)
=
\Bigl(
d_{n-1}x,\;
(-1)^n f_{n-1}x+d'_n y+(-1)^{n-1}s_{n-2}d_{n-1}x
\Bigr).
\]
On the other hand,
\[
\Phi_n(x,y)=\bigl(x,\; y+(-1)^n s_{n-1}x\bigr),
\]
so
\[
(d_g)_n\Phi_n(x,y)
=
\Bigl(
d_{n-1}x,\;
(-1)^n g_{n-1}x+d'_n y+(-1)^n d'_n s_{n-1}x
\Bigr).
\]
These second components agree precisely because
\[
f_{n-1}-g_{n-1}=d'_n s_{n-1}+s_{n-2}d_{n-1}.
\]
Thus
\[
(d_g)_n\Phi_n=\Phi_{n-1}(d_f)_n,
\]
so \(\Phi:M(f)\to M(g)\) is a chain isomorphism. In particular,
\[
M(f)\sim M(g).
\]


\medskip
\noindent\textbf{Step 3: \(H_*(M(f))=0\) if and only if \(f_*\) is an isomorphism}\par
\smallskip


There is a short exact sequence of chain complexes
\[
0\longrightarrow C' \xrightarrow{i} M(f)\xrightarrow{p} \Sigma C \longrightarrow 0,
\]
where \(\Sigma C\) is the suspension of \(C\), defined by
\[
(\Sigma C)_n=C_{n-1}.
\]
The maps are
\[
i(y)=(0,y),\qquad p(x,y)=x.
\]
This short exact sequence yields the long exact sequence in homology:
\[
\cdots
\to H_n(C')
\to H_n(M(f))
\to H_n(\Sigma C)
\xrightarrow{\partial}
H_{n-1}(C')
\to H_{n-1}(M(f))
\to \cdots
\]
Using the natural identification
\[
H_n(\Sigma C)\cong H_{n-1}(C),
\]
this becomes
\[
\cdots
\to H_n(C')
\to H_n(M(f))
\to H_{n-1}(C)
\xrightarrow{\partial}
H_{n-1}(C')
\to H_{n-1}(M(f))
\to \cdots
\]
and the connecting morphism \(\partial\) is exactly the map induced by \(f\), up to the standard sign convention. Hence the long exact sequence may be written as
\[
\cdots
\to H_n(C')
\to H_n(M(f))
\to H_{n-1}(C)
\xrightarrow{f_*}
H_{n-1}(C')
\to H_{n-1}(M(f))
\to \cdots
\]

Now suppose first that
\[
H_n(M(f))=0
\qquad\text{for all }n.
\]
Then exactness gives
\[
0\to H_{n-1}(C)\xrightarrow{f_*} H_{n-1}(C')\to 0,
\]
so \(f_*\) is an isomorphism in every degree.

Conversely, suppose
\[
f_*:H_*(C)\to H_*(C')
\]
is an isomorphism. Then by exactness, the group between \(H_n(C')\) and \(H_{n-1}(C)\) must vanish, i.e.
\[
H_n(M(f))=0
\qquad\text{for all }n.
\]

Therefore,
\[
H_*(M(f))=0
\quad\Longleftrightarrow\quad
f_*:H_*(C)\to H_*(C')
\text{ is an isomorphism.}
\]

\par
\end{solution}

\begin{problem}[CP-VIII-0061 --- The quotient \texorpdfstring{$(D^{n+1}\times S^m)/(S^n\times S^m)$}{(D\string^{n+1} x S\string^m)/(S\string^n x S\string^m)}]
Let
\[
X=\frac{D^{n+1}\times S^m}{S^n\times S^m},
\]
where \(S^n=\partial D^{n+1}\). Find a cell decomposition of \(X\) and compute its homology using the cellular chain complex.
\end{problem}

\begin{solution}
We collapse the whole subspace
\[
S^n\times S^m \subset D^{n+1}\times S^m
\]
to a single point.

We use the standard CW structure on \(S^m\), namely one \(0\)-cell \(e^0\) and one \(m\)-cell \(e^m\). On the other hand, the quotient
\[
D^{n+1}/S^n \cong S^{n+1}
\]
shows that, relative to the collapsed boundary, \(D^{n+1}\) contributes one \((n+1)\)-cell.

Therefore the quotient
\[
X=(D^{n+1}\times S^m)/(S^n\times S^m)
\]
inherits a CW structure whose cells come from the product of this relative \((n+1)\)-cell with the cells of \(S^m\). Hence \(X\) has exactly:
\[
e^0,\qquad e^{n+1},\qquad e^{n+m+1}.
\]
So a CW decomposition of \(X\) consists of
\begin{itemize}
	\item one \(0\)-cell,
	\item one \((n+1)\)-cell,
	\item one \((n+m+1)\)-cell.
\end{itemize}

Thus the cellular chain groups are
\[
C_q^{\mathrm{cell}}(X)=
\begin{cases}
	\mathbb Z, & q=0,n+1,n+m+1,\\[4pt]
	0, & \text{otherwise}.
\end{cases}
\]

Since there are no cells in dimensions \(n\) and \(n+m\), every cellular boundary map is zero. Therefore the cellular chain complex is
\[
0 \longrightarrow \mathbb Z
\overset{0}{\longrightarrow}
\mathbb Z
\overset{0}{\longrightarrow}
\mathbb Z
\longrightarrow 0,
\]
where the three copies of \(\mathbb Z\) sit in dimensions \(n+m+1\), \(n+1\), and \(0\), respectively.

It follows at once that
\[
H_q(X)=
\begin{cases}
	\mathbb Z, & q=0,\\[4pt]
	\mathbb Z, & q=n+1,\\[4pt]
	\mathbb Z, & q=n+m+1,\\[4pt]
	0, & \text{otherwise}.
\end{cases}
\]

Hence the homology of \(X\) is
\[
H_*(X)\cong \mathbb Z \text{ in degrees }0,n+1,n+m+1,
\]
and vanishes in all other degrees.

\par
\end{solution}

\begin{problem}[CP-VIII-0062 --- Cellular maps and cellular chain complexes]
Let \(X\) and \(Y\) be finite cell complexes. A map
\[
f:X\to Y
\]
is called \emph{cellular} if \(f(X_n)\subseteq Y_n\) for all \(n\), where \(X_n\) and \(Y_n\) denote the \(n\)-skeleta. Show that such a map induces a chain map
\[
f_*^{\mathrm{cell}}:C_*^{\mathrm{cell}}(X)\to C_*^{\mathrm{cell}}(Y),
\]
and that the induced map on homology agrees with the usual map
\[
f_*:H_*(X)\to H_*(Y).
\]

\par
\end{problem}

\begin{solution}
Recall that the cellular chain groups are defined by
\[
C_n^{\mathrm{cell}}(X)=H_n(X_n,X_{n-1}),
\qquad
C_n^{\mathrm{cell}}(Y)=H_n(Y_n,Y_{n-1}).
\]

Since \(f\) is cellular, for every \(n\) we have
\[
f(X_n)\subseteq Y_n
\qquad\text{and}\qquad
f(X_{n-1})\subseteq Y_{n-1}.
\]
Hence \(f\) induces a map of pairs
\[
f:(X_n,X_{n-1})\to (Y_n,Y_{n-1}),
\]
and therefore a homomorphism
\[
f_n^{\mathrm{cell}}:
H_n(X_n,X_{n-1})\longrightarrow H_n(Y_n,Y_{n-1}).
\]
So for each \(n\) we obtain a map
\[
f_n^{\mathrm{cell}}:C_n^{\mathrm{cell}}(X)\to C_n^{\mathrm{cell}}(Y).
\]

It remains to verify that these maps commute with the cellular boundary operators. The cellular differential
\[
d_n:C_n^{\mathrm{cell}}(X)\to C_{n-1}^{\mathrm{cell}}(X)
\]
is defined using the connecting homomorphisms in the long exact sequence of the triple
\[
(X_n,X_{n-1},X_{n-2}),
\]
and similarly for \(Y\). Since \(f\) preserves skeleta, it induces a map of triples
\[
(X_n,X_{n-1},X_{n-2})\to (Y_n,Y_{n-1},Y_{n-2}).
\]
By naturality of the connecting homomorphism, we obtain
\[
d_n\circ f_n^{\mathrm{cell}}
=
f_{n-1}^{\mathrm{cell}}\circ d_n.
\]
Thus
\[
f_*^{\mathrm{cell}}:C_*^{\mathrm{cell}}(X)\to C_*^{\mathrm{cell}}(Y)
\]
is a chain map.

We now compare the induced map on homology with the usual singular-homology map. For CW complexes, cellular homology is naturally isomorphic to singular homology:
\[
H_*^{\mathrm{cell}}(X)\cong H_*(X),
\qquad
H_*^{\mathrm{cell}}(Y)\cong H_*(Y).
\]
These isomorphisms are natural with respect to cellular maps. Consequently, under the canonical identifications
\[
H_*^{\mathrm{cell}}(X)\cong H_*(X),
\qquad
H_*^{\mathrm{cell}}(Y)\cong H_*(Y),
\]
the homomorphism induced by the cellular chain map \(f_*^{\mathrm{cell}}\) agrees exactly with the usual map
\[
f_*:H_*(X)\to H_*(Y).
\]

\par

Therefore a cellular map induces a cellular chain map, and its induced map on homology is the same as the standard homology map.

\par
\end{solution}

\begin{problem}[CP-VIII-0063 --- The quotient map for \(X=S^n\cup_f D^{n+1}\)]
Let
\[
f:S^n\to S^n
\]
have degree \(k\), and let
\[
X=S^n\cup_f D^{n+1}.
\]
Let
\[
\pi:X\to X/S^n\cong S^{n+1}
\]
be the quotient map. Compute
\[
\pi^*:H^*(S^{n+1})\to H^*(X)
\qquad\text{and}\qquad
\pi_*:H_*(X)\to H_*(S^{n+1}).
\]

\par
\end{problem}

\begin{solution}
The space \(X\) has a CW structure with one \(0\)-cell, one \(n\)-cell, and one \((n+1)\)-cell. The attaching map of the \((n+1)\)-cell has degree \(k\), so the cellular chain complex of \(X\) is
\[
0\longrightarrow C_{n+1}^{\mathrm{cell}}(X)\cong \mathbb Z
\xrightarrow{\ \times k\ }
C_n^{\mathrm{cell}}(X)\cong \mathbb Z
\longrightarrow 0,
\]
together with \(C_0^{\mathrm{cell}}(X)\cong \mathbb Z\).

Hence
\[
H_{n+1}(X)=\ker(\times k:\mathbb Z\to\mathbb Z),
\qquad
H_n(X)=\operatorname{coker}(\times k:\mathbb Z\to\mathbb Z)\cong \mathbb Z/k\mathbb Z.
\]
So, if \(k\neq 0\), then
\[
H_{n+1}(X)=0,
\qquad
H_n(X)\cong \mathbb Z/k\mathbb Z,
\qquad
H_0(X)\cong \mathbb Z,
\]
and all other homology groups vanish. If \(k=0\), then
\[
H_{n+1}(X)\cong \mathbb Z,
\qquad
H_n(X)\cong \mathbb Z,
\]
so in that case \(X\simeq S^n\vee S^{n+1}\).

Dually, the cellular cochain complex gives
\[
H^{n+1}(X)\cong \operatorname{coker}(\times k:\mathbb Z\to\mathbb Z)\cong \mathbb Z/k\mathbb Z,
\qquad
H^n(X)\cong \ker(\times k:\mathbb Z\to\mathbb Z),
\]
and therefore, for \(k\neq 0\),
\[
H^{n+1}(X)\cong \mathbb Z/k\mathbb Z,
\qquad
H^n(X)=0,
\qquad
H^0(X)\cong \mathbb Z.
\]

Now collapse the subspace \(S^n\subset X\). Since \(X\) is obtained from \(S^n\) by attaching one \((n+1)\)-cell, we have
\[
X/S^n\cong S^{n+1}.
\]
Moreover,
\[
\widetilde H_*(X/S^n)\cong H_*(X,S^n),
\qquad
\widetilde H^*(X/S^n)\cong H^*(X,S^n),
\]
and under these identifications the quotient map \(\pi\) corresponds to the canonical passage from absolute groups to relative groups.

\par

\textbf{Computation of \(\pi_*\).}
Consider the long exact sequence of the pair \((X,S^n)\):
\[
0\to H_{n+1}(X)\to H_{n+1}(X,S^n)\xrightarrow{\partial} H_n(S^n)\to H_n(X)\to 0.
\]
Since
\[
H_{n+1}(X,S^n)\cong \widetilde H_{n+1}(X/S^n)\cong H_{n+1}(S^{n+1})\cong \mathbb Z,
\]
and the connecting homomorphism \(\partial\) is multiplication by \(k\), it follows that
\[
H_{n+1}(X)=\ker(\times k:\mathbb Z\to\mathbb Z).
\]
The map
\[
\pi_*:H_{n+1}(X)\to H_{n+1}(S^{n+1})\cong \mathbb Z
\]
is precisely the inclusion of this kernel into \(\mathbb Z\). Therefore:
\[
\pi_*=
\begin{cases}
	\mathrm{id}_{\mathbb Z}, & *=0,\\[1mm]
	0, & *=n+1 \text{ and } k\neq 0,\\[1mm]
	\text{an isomorphism }\mathbb Z\to\mathbb Z, & *=n+1 \text{ and } k=0,\\[1mm]
	0, & \text{in all other positive degrees.}
\end{cases}
\]

\par

\textbf{Computation of \(\pi^*\).}
Now use the long exact sequence in cohomology:
\[
0\to H^n(X)\to H^n(S^n)\xrightarrow{\delta} H^{n+1}(X,S^n)\to H^{n+1}(X)\to 0.
\]
Again,
\[
H^{n+1}(X,S^n)\cong \widetilde H^{n+1}(X/S^n)\cong H^{n+1}(S^{n+1})\cong \mathbb Z,
\]
and the connecting map \(\delta\) is multiplication by \(k\). Hence
\[
H^{n+1}(X)\cong \mathbb Z/k\mathbb Z.
\]
Under the identification \(H^{n+1}(X,S^n)\cong H^{n+1}(S^{n+1})\), the map
\[
\pi^*:H^{n+1}(S^{n+1})\to H^{n+1}(X)
\]
is exactly the quotient map
\[
\mathbb Z\to \mathbb Z/k\mathbb Z,
\qquad
m\mapsto \overline{m}.
\]
Thus
\[
\pi^*=
\begin{cases}
	\mathrm{id}_{\mathbb Z}, & *=0,\\[1mm]
	\mathbb Z\to \mathbb Z/k\mathbb Z,\quad m\mapsto \overline{m}, & *=n+1 \text{ and } k\neq 0,\\[1mm]
	\text{an isomorphism }\mathbb Z\to\mathbb Z, & *=n+1 \text{ and } k=0,\\[1mm]
	0, & \text{in all other degrees.}
\end{cases}
\]

\par

In particular, for \(k\neq 0\), the quotient map \(\pi\) is trivial on positive-dimensional homology, but in cohomology it sends a generator of \(H^{n+1}(S^{n+1})\cong \mathbb Z\) to the generator of
\[
H^{n+1}(X)\cong \mathbb Z/k\mathbb Z
\]
modulo \(k\).

\par

\noindent
\textbf{Summary.}
These three problems illustrate three central themes in algebraic topology:

\begin{itemize}
	\item every homology class is represented by a map from a finite cell complex;
	\item cellular maps act directly on cellular chain complexes and recover the usual maps on homology;
	\item attaching maps of spheres control the algebraic structure of the resulting space, and quotient maps detect this structure naturally in homology and cohomology.
\end{itemize}
\end{solution}

\begin{problem}[CP-VIII-0064 --- Compatibility of boundary maps with the Kronecker pairing]
Let
\[
\partial : H_*(X,A)\longrightarrow H_{*-1}(A)
\]
and
\[
\delta : H^*(A)\longrightarrow H^{*+1}(X,A)
\]
be the boundary maps in the long exact homology and cohomology sequences of the pair \((X,A)\). If
\[
a\in H^k(A)
\qquad\text{and}\qquad
x\in H_{k+1}(X,A),
\]
show that
\[
\langle a,\partial x\rangle=\langle \delta a,x\rangle .
\]

\par
\end{problem}

\begin{solution}
We prove the identity directly from the chain-level definitions of the connecting homomorphisms.

Let
\[
\alpha\in Z^k(A)
\]
be a cocycle representing the cohomology class \(a\in H^k(A)\). Also let
\[
c\in C_{k+1}(X)
\]
be a singular chain representing the relative homology class
\[
x\in H_{k+1}(X,A).
\]
Since \(c\) represents a relative cycle, its boundary lies in \(A\), that is,
\[
\partial c\in C_k(A).
\]
Therefore the homology boundary \(\partial x\in H_k(A)\) is represented by the cycle \(\partial c\).

Hence
\[
\langle a,\partial x\rangle
=
\alpha(\partial c).
\]

Now we describe the cohomology connecting homomorphism. Choose a cochain
\[
\widetilde{\alpha}\in C^k(X)
\]
extending \(\alpha\), meaning that
\[
\widetilde{\alpha}|_{C_k(A)}=\alpha.
\]
Because \(\alpha\) is a cocycle on \(A\), we have
\[
d\alpha=0.
\]
Thus \(d\widetilde{\alpha}\) vanishes on \(C_{k+1}(A)\), so it defines a relative cocycle
\[
d\widetilde{\alpha}\in Z^{k+1}(X,A).
\]
By definition of the connecting homomorphism in cohomology,
\[
\delta a=[d\widetilde{\alpha}]\in H^{k+1}(X,A).
\]

Therefore
\[
\langle \delta a,x\rangle
=
(d\widetilde{\alpha})(c).
\]
Using the definition of the coboundary,
\[
(d\widetilde{\alpha})(c)
=
\widetilde{\alpha}(\partial c).
\]
But \(\partial c\in C_k(A)\), and \(\widetilde{\alpha}\) restricts to \(\alpha\) on \(A\). Hence
\[
\widetilde{\alpha}(\partial c)
=
\alpha(\partial c).
\]
Thus
\[
\langle \delta a,x\rangle
=
\alpha(\partial c)
=
\langle a,\partial x\rangle.
\]

Therefore
\[
\boxed{
	\langle a,\partial x\rangle=\langle \delta a,x\rangle .
}
\]

\par

\textbf{Remark.}
The proof says that the following square is compatible with the natural pairing:
\[
\begin{array}{ccc}
	H^k(A)\times H_{k+1}(X,A) & \longrightarrow & \mathbb Z \\[1mm]
	\delta\times \mathrm{id}\downarrow\phantom{\delta\times \mathrm{id}} 
	&& 
	\phantom{\mathrm{id}\times \partial}\downarrow \mathrm{id} \\[1mm]
	H^{k+1}(X,A)\times H_{k+1}(X,A) & \longrightarrow & \mathbb Z ,
\end{array}
\]
or more concretely,
\[
\boxed{
	\text{evaluate first and then take a boundary}
	=
	\text{take a coboundary first and then evaluate.}
}
\]

\par
\end{solution}

\begin{problem}[CP-VIII-0065 --- Coefficient short exact sequences and the Bockstein homomorphism]
Suppose
\[
0\longrightarrow K\longrightarrow G\longrightarrow H\longrightarrow 0
\]
is a short exact sequence of abelian groups. If \(X\) is a space, show that there is a short exact sequence
\[
0\longrightarrow C_*(X;K)\longrightarrow C_*(X;G)\longrightarrow C_*(X;H)\longrightarrow 0.
\]

Let
\[
\beta_X:H_*(X;H)\longrightarrow H_{*-1}(X;K)
\]
be the boundary map in the associated long exact sequence. This map is called the \emph{Bockstein homomorphism}. If
\[
f:X\longrightarrow Y
\]
is a map, show that
\[
f_*\beta_X=\beta_Y f_*.
\]
If \(X=Y=\mathbb{RP}^n\), deduce that
\[
f_*:H_*(\mathbb{RP}^n;\mathbb Z/2)
\longrightarrow
H_*(\mathbb{RP}^n;\mathbb Z/2)
\]
is determined by
\[
f_*:H_*(\mathbb{RP}^n;\mathbb Z)
\longrightarrow
H_*(\mathbb{RP}^n;\mathbb Z).
\]
Show that the converse holds if \(n\) is even, but not if \(n\) is odd.

\par
\end{problem}

\begin{solution}

\medskip
\noindent\textbf{Step 1. The short exact sequence of chain complexes}\par
\smallskip


For singular chains with coefficients in an abelian group \(G\), we have
\[
C_n(X;G)=C_n(X)\otimes G.
\]
The group \(C_n(X)\) is a free abelian group, generated by all singular \(n\)-simplices in \(X\). Since tensoring with a free abelian group preserves short exact sequences, the short exact sequence
\[
0\longrightarrow K\longrightarrow G\longrightarrow H\longrightarrow 0
\]
gives, for every \(n\),
\[
0\longrightarrow C_n(X)\otimes K
\longrightarrow C_n(X)\otimes G
\longrightarrow C_n(X)\otimes H
\longrightarrow 0.
\]
Equivalently,
\[
0\longrightarrow C_n(X;K)
\longrightarrow C_n(X;G)
\longrightarrow C_n(X;H)
\longrightarrow 0.
\]

The maps commute with the boundary operators because the boundary with coefficients is given by
\[
\partial\otimes \mathrm{id}.
\]
Therefore the degreewise short exact sequences assemble into a short exact sequence of chain complexes:
\[
\boxed{
	0\longrightarrow C_*(X;K)
	\longrightarrow C_*(X;G)
	\longrightarrow C_*(X;H)
	\longrightarrow 0.
}
\]


\medskip
\noindent\textbf{Step 2. The Bockstein homomorphism}\par
\smallskip


From the short exact sequence of chain complexes
\[
0\longrightarrow C_*(X;K)
\longrightarrow C_*(X;G)
\longrightarrow C_*(X;H)
\longrightarrow 0,
\]
we obtain a long exact sequence in homology:
\[
\cdots
\longrightarrow
H_q(X;K)
\longrightarrow
H_q(X;G)
\longrightarrow
H_q(X;H)
\xrightarrow{\beta_X}
H_{q-1}(X;K)
\longrightarrow
H_{q-1}(X;G)
\longrightarrow
\cdots .
\]
The connecting homomorphism
\[
\beta_X:H_q(X;H)\longrightarrow H_{q-1}(X;K)
\]
is the Bockstein homomorphism associated with the coefficient sequence
\[
0\longrightarrow K\longrightarrow G\longrightarrow H\longrightarrow 0.
\]


\medskip
\noindent\textbf{Step 3. Naturality of the Bockstein homomorphism}\par
\smallskip


Let
\[
f:X\longrightarrow Y
\]
be a continuous map. It induces chain maps
\[
f_\#:C_*(X;K)\longrightarrow C_*(Y;K),
\]
\[
f_\#:C_*(X;G)\longrightarrow C_*(Y;G),
\]
and
\[
f_\#:C_*(X;H)\longrightarrow C_*(Y;H).
\]
These maps fit into a commutative diagram of short exact sequences of chain complexes:
\[
\begin{array}{ccccccccc}
	0 &\longrightarrow& C_*(X;K) &\longrightarrow& C_*(X;G) &\longrightarrow& C_*(X;H) &\longrightarrow& 0 \\[1mm]
	&& \downarrow f_\# && \downarrow f_\# && \downarrow f_\# && \\[1mm]
	0 &\longrightarrow& C_*(Y;K) &\longrightarrow& C_*(Y;G) &\longrightarrow& C_*(Y;H) &\longrightarrow& 0 .
\end{array}
\]
Passing to homology gives a commutative diagram of long exact sequences. In particular, the connecting homomorphisms commute with the maps induced by \(f\). Hence
\[
\boxed{
	f_*\beta_X=\beta_Y f_*.
}
\]

This is the naturality of the Bockstein homomorphism.


\medskip
\noindent\textbf{Step 4. The coefficient sequence for \(\mathbb{Z}/2\)}\par
\smallskip


Now take the standard short exact sequence
\[
0\longrightarrow \mathbb Z
\xrightarrow{\times 2}
\mathbb Z
\longrightarrow
\mathbb Z/2
\longrightarrow 0.
\]
The associated Bockstein homomorphism is
\[
\beta:
H_q(X;\mathbb Z/2)
\longrightarrow
H_{q-1}(X;\mathbb Z).
\]

For \(X=\mathbb{RP}^n\), the integral homology groups are
\[
H_q(\mathbb{RP}^n;\mathbb Z)
\cong
\begin{cases}
	\mathbb Z, & q=0,\\[1mm]
	\mathbb Z/2, & q \text{ odd and } 0<q<n,\\[1mm]
	\mathbb Z, & q=n \text{ and } n \text{ odd},\\[1mm]
	0, & \text{otherwise}.
\end{cases}
\]
The mod \(2\) homology groups are
\[
H_q(\mathbb{RP}^n;\mathbb Z/2)
\cong
\mathbb Z/2
\qquad
\text{for }0\le q\le n.
\]

The long exact coefficient sequence is
\[
\begin{aligned}
\cdots&\longrightarrow H_q(\mathbb{RP}^n;\mathbb Z)
\xrightarrow{\times 2}H_q(\mathbb{RP}^n;\mathbb Z)\\
&\longrightarrow H_q(\mathbb{RP}^n;\mathbb Z/2)
\xrightarrow{\beta}H_{q-1}(\mathbb{RP}^n;\mathbb Z)\\
&\xrightarrow{\times 2}H_{q-1}(\mathbb{RP}^n;\mathbb Z)
\longrightarrow\cdots .
\end{aligned}
\]


\medskip
\noindent\textbf{Step 5. Why the integral map determines the mod \(2\) map}\par
\smallskip


We analyze degree by degree.

\par

\textbf{Case 1: \(q=0\).}
Since \(\mathbb{RP}^n\) is connected, every self-map
\[
f:\mathbb{RP}^n\longrightarrow \mathbb{RP}^n
\]
induces the identity map on \(H_0\), both with integral and mod \(2\) coefficients. Thus the mod \(2\) map in degree \(0\) is determined.

\par

\textbf{Case 2: \(q\) odd and \(0<q<n\).}
Then
\[
H_q(\mathbb{RP}^n;\mathbb Z)\cong \mathbb Z/2
\]
and
\[
H_{q-1}(\mathbb{RP}^n;\mathbb Z)=0.
\]
Therefore the reduction map
\[
H_q(\mathbb{RP}^n;\mathbb Z)
\longrightarrow
H_q(\mathbb{RP}^n;\mathbb Z/2)
\]
is an isomorphism. Hence the action of \(f_*\) on
\[
H_q(\mathbb{RP}^n;\mathbb Z/2)
\]
is determined by its action on
\[
H_q(\mathbb{RP}^n;\mathbb Z).
\]

\par

\textbf{Case 3: \(q\) even and \(0<q\le n\).}
Then
\[
H_q(\mathbb{RP}^n;\mathbb Z)=0
\]
and
\[
H_{q-1}(\mathbb{RP}^n;\mathbb Z)\cong \mathbb Z/2.
\]
The Bockstein map
\[
\beta:
H_q(\mathbb{RP}^n;\mathbb Z/2)
\longrightarrow
H_{q-1}(\mathbb{RP}^n;\mathbb Z)
\]
is an isomorphism. Since the Bockstein is natural, we have
\[
f_*\beta=\beta f_*.
\]
Therefore, because \(\beta\) is an isomorphism, the action of \(f_*\) on
\[
H_q(\mathbb{RP}^n;\mathbb Z/2)
\]
is determined by the action of \(f_*\) on
\[
H_{q-1}(\mathbb{RP}^n;\mathbb Z).
\]

\par

\textbf{Case 4: \(q=n\) and \(n\) is odd.}
Then
\[
H_n(\mathbb{RP}^n;\mathbb Z)\cong \mathbb Z
\]
and
\[
H_n(\mathbb{RP}^n;\mathbb Z/2)\cong \mathbb Z/2.
\]
The reduction map
\[
H_n(\mathbb{RP}^n;\mathbb Z)
\longrightarrow
H_n(\mathbb{RP}^n;\mathbb Z/2)
\]
sends an integer degree to its value modulo \(2\). Therefore the integral map on top homology determines the mod \(2\) map on top homology.

Combining all cases, we conclude that
\[
\boxed{
	f_*:H_*(\mathbb{RP}^n;\mathbb Z/2)
	\longrightarrow
	H_*(\mathbb{RP}^n;\mathbb Z/2)
}
\]
is determined by
\[
\boxed{
	f_*:H_*(\mathbb{RP}^n;\mathbb Z)
	\longrightarrow
	H_*(\mathbb{RP}^n;\mathbb Z).
}
\]


\medskip
\noindent\textbf{Step 6. The converse when \(n\) is even}\par
\smallskip


Assume now that \(n\) is even. Then
\[
H_q(\mathbb{RP}^n;\mathbb Z)
\cong
\begin{cases}
	\mathbb Z, & q=0,\\[1mm]
	\mathbb Z/2, & q \text{ odd and }0<q<n,\\[1mm]
	0, & \text{otherwise}.
\end{cases}
\]
For every odd \(q\) with \(0<q<n\), the reduction map
\[
H_q(\mathbb{RP}^n;\mathbb Z)
\longrightarrow
H_q(\mathbb{RP}^n;\mathbb Z/2)
\]
is an isomorphism. Therefore the mod \(2\) map determines the integral map in all nonzero positive degrees.

In degree \(0\), since \(\mathbb{RP}^n\) is connected, every self-map induces the identity map on
\[
H_0(\mathbb{RP}^n;\mathbb Z)\cong \mathbb Z.
\]
Thus degree \(0\) is also determined.

Therefore, when \(n\) is even, the mod \(2\) homology action determines the integral homology action. Hence
\[
\boxed{
	\text{if }n\text{ is even, the converse holds.}
}
\]


\medskip
\noindent\textbf{Step 7. Failure of the converse when \(n\) is odd}\par
\smallskip


Now assume \(n\) is odd. Then \(\mathbb{RP}^n\) is orientable and
\[
H_n(\mathbb{RP}^n;\mathbb Z)\cong \mathbb Z.
\]
The mod \(2\) homology group in top degree is only
\[
H_n(\mathbb{RP}^n;\mathbb Z/2)\cong \mathbb Z/2.
\]
Thus mod \(2\) homology cannot distinguish multiplication by \(+1\) from multiplication by \(-1\) on the top integral homology group.

Indeed, compare the identity map
\[
\mathrm{id}:\mathbb{RP}^n\longrightarrow \mathbb{RP}^n
\]
with the projective linear homeomorphism induced by the reflection
\[
(x_0,x_1,\dots,x_n)\longmapsto (-x_0,x_1,\dots,x_n).
\]
This reflection has determinant \(-1\). Since \(n\) is odd, \(\mathbb{RP}^n\) is orientable, and the induced map on top homology is multiplication by \(-1\):
\[
r_*=-\mathrm{id}
\qquad
\text{on }
H_n(\mathbb{RP}^n;\mathbb Z)\cong \mathbb Z.
\]
By contrast, the identity map induces multiplication by \(+1\) on
\[
H_n(\mathbb{RP}^n;\mathbb Z).
\]
Therefore the two maps induce different maps on integral homology.

However, with \(\mathbb Z/2\)-coefficients, we have
\[
+1=-1
\qquad
\text{in }
\mathbb Z/2.
\]
So these two maps induce the same map on
\[
H_n(\mathbb{RP}^n;\mathbb Z/2).
\]
In fact, since each group
\[
H_q(\mathbb{RP}^n;\mathbb Z/2)
\]
is either \(0\) or \(\mathbb Z/2\), and since a homeomorphism induces an isomorphism in homology, it must induce the identity on each nonzero mod \(2\) homology group.

Thus the identity map and the orientation-reversing projective homeomorphism have the same action on mod \(2\) homology but different actions on integral homology.

Therefore
\[
\boxed{
	\text{if }n\text{ is odd, the converse does not hold.}
}
\]


\medskip
\noindent\textbf{Summary}\par
\smallskip


For the coefficient sequence
\[
0\longrightarrow \mathbb Z
\xrightarrow{\times 2}
\mathbb Z
\longrightarrow
\mathbb Z/2
\longrightarrow 0,
\]
the associated Bockstein homomorphism is natural:
\[
\boxed{
	f_*\beta_X=\beta_Y f_*.
}
\]
For \(X=Y=\mathbb{RP}^n\), this naturality implies that the integral homology action determines the mod \(2\) homology action:
\[
\boxed{
	f_* \text{ on } H_*(\mathbb{RP}^n;\mathbb Z)
	\quad\Longrightarrow\quad
	f_* \text{ on } H_*(\mathbb{RP}^n;\mathbb Z/2).
}
\]
The converse holds when \(n\) is even, because there is no top-dimensional infinite cyclic homology group. It fails when \(n\) is odd, because the top homology group is
\[
H_n(\mathbb{RP}^n;\mathbb Z)\cong \mathbb Z,
\]
and mod \(2\) homology cannot distinguish the maps \(+1\) and \(-1\) on this group.
\end{solution}

\begin{problem}[CP-VIII-0066 --- Mod-\(p\) homology test and chain homotopy equivalences]
Let \(C_*\) be a free finitely generated chain complex defined over \(\mathbb Z\). Show that
\[
H_*(C)=0
\]
if and only if
\[
H_*(C\otimes \mathbb Z/p)=0
\]
for every prime \(p\).

Now suppose
\[
f:C_* \longrightarrow C'_*
\]
is a chain map. Using Problem 5 from Example Sheet 1, show that \(f\) is a chain homotopy equivalence if and only if the induced map
\[
f_*:H_*(C\otimes \mathbb Z/p)\longrightarrow H_*(C'\otimes \mathbb Z/p)
\]
is an isomorphism for every prime \(p\).
\end{problem}

\begin{solution}
Write
\[
\mathbb F_p=\mathbb Z/p.
\]
Since \(C_*\) is a chain complex of free abelian groups, we may apply the universal coefficient theorem for homology. For every \(n\), there is a short exact sequence
\[
0
\longrightarrow
H_n(C)\otimes \mathbb F_p
\longrightarrow
H_n(C\otimes \mathbb F_p)
\longrightarrow
\operatorname{Tor}_1^{\mathbb Z}(H_{n-1}(C),\mathbb F_p)
\longrightarrow
0.
\]

First suppose that
\[
H_*(C)=0.
\]
Then \(H_n(C)=0\) for every \(n\), and so both outer terms in the universal coefficient sequence vanish. Therefore
\[
H_n(C\otimes \mathbb F_p)=0
\]
for every \(n\) and every prime \(p\).

Conversely, suppose that
\[
H_*(C\otimes \mathbb F_p)=0
\]
for every prime \(p\). Then the universal coefficient sequence implies in particular that
\[
H_n(C)\otimes \mathbb F_p=0
\]
for every \(n\) and every prime \(p\).

Because \(C_*\) is a free finitely generated chain complex over \(\mathbb Z\), each homology group \(H_n(C)\) is a finitely generated abelian group. By the structure theorem for finitely generated abelian groups,
\[
H_n(C)\cong \mathbb Z^r\oplus
\bigoplus_{j=1}^m \mathbb Z/q_j^{a_j}.
\]
If \(r>0\), then
\[
H_n(C)\otimes \mathbb F_p
\]
contains a copy of \(\mathbb F_p\) for every prime \(p\), which is impossible. Hence \(r=0\).

If a torsion summand \(\mathbb Z/q^a\) occurs, then taking \(p=q\) gives
\[
\mathbb Z/q^a\otimes \mathbb F_q \cong \mathbb F_q\neq 0,
\]
again impossible. Therefore there is no torsion either. Hence
\[
H_n(C)=0
\]
for every \(n\). Thus
\[
H_*(C)=0.
\]

Now let
\[
f:C_*\longrightarrow C'_*
\]
be a chain map. We use the mapping cone
\[
\operatorname{Cone}(f).
\]
Recall that \(f\) is a quasi-isomorphism if and only if
\[
H_*(\operatorname{Cone}(f))=0.
\]
Moreover, for finite free chain complexes over \(\mathbb Z\), an acyclic cone is contractible, and by the standard cone criterion this is equivalent to \(f\) being a chain homotopy equivalence.

Suppose first that \(f\) is a chain homotopy equivalence. Then tensoring with \(\mathbb F_p\) preserves chain homotopies, so
\[
f\otimes \mathbb F_p:C_*\otimes \mathbb F_p
\longrightarrow
C'_*\otimes \mathbb F_p
\]
is also a chain homotopy equivalence. Therefore it induces an isomorphism on homology:
\[
f_*:H_*(C\otimes \mathbb F_p)
\longrightarrow
H_*(C'\otimes \mathbb F_p).
\]

Conversely, suppose that for every prime \(p\), the map
\[
f_*:H_*(C\otimes \mathbb F_p)
\longrightarrow
H_*(C'\otimes \mathbb F_p)
\]
is an isomorphism. Then the mapping cone of
\[
f\otimes \mathbb F_p
\]
is acyclic:
\[
H_*(\operatorname{Cone}(f\otimes \mathbb F_p))=0.
\]
But tensoring commutes with taking the mapping cone, so
\[
\operatorname{Cone}(f\otimes \mathbb F_p)
\cong
\operatorname{Cone}(f)\otimes \mathbb F_p.
\]
Hence
\[
H_*(\operatorname{Cone}(f)\otimes \mathbb F_p)=0
\]
for every prime \(p\).

By the first part of the problem, applied to the finite free chain complex
\[
\operatorname{Cone}(f),
\]
we get
\[
H_*(\operatorname{Cone}(f))=0.
\]
Thus \(\operatorname{Cone}(f)\) is acyclic. Since it is a finite free chain complex over \(\mathbb Z\), Problem 5 from Example Sheet 1 implies that \(\operatorname{Cone}(f)\) is contractible. Therefore, by the mapping cone criterion, \(f\) is a chain homotopy equivalence.

So \(f\) is a chain homotopy equivalence if and only if it induces an isomorphism on homology after tensoring with \(\mathbb Z/p\) for every prime \(p\).

\[
\begin{aligned}
f\text{ is a chain homotopy equivalence}
\quad\Longleftrightarrow\quad&
\\[-2mm]
f_*:H_*(C\otimes\mathbb Z/p)&\xrightarrow{\;\cong\;}
H_*(C'\otimes\mathbb Z/p)
\quad\text{for every prime }p .
\end{aligned}
\]
\end{solution}

\begin{problem}[CP-VIII-0067 --- Free resolution over \(\mathbb C\lbrack x\rbrack/(x^3)\) and Tor groups]
Let
\[
R=\mathbb C[x]/(x^3),
\]
and for \(i=1,2\), let
\[
M_i=\mathbb C[x]/(x^i)
\]
be regarded as an \(R\)-module. Find a free resolution of \(M_1\), and use it to compute
\[
\operatorname{Tor}^R_*(M_1,M_1)
\]
and
\[
\operatorname{Tor}^R_*(M_1,M_2).
\]
\end{problem}

\begin{solution}
We have
\[
R=\mathbb C[x]/(x^3).
\]
The module \(M_1\) is
\[
M_1=\mathbb C[x]/(x)\cong R/(x).
\]
The module \(M_2\) is
\[
M_2=\mathbb C[x]/(x^2)\cong R/(x^2).
\]

We first construct a free resolution of \(M_1\). Since
\[
x^3=0 \quad \text{in } R,
\]
multiplication by \(x\) and multiplication by \(x^2\) compose to zero:
\[
x\cdot x^2=x^3=0.
\]
Thus we get a \(2\)-periodic chain complex
\[
\cdots
\longrightarrow
R
\xrightarrow{\cdot x^2}
R
\xrightarrow{\cdot x}
R
\xrightarrow{\cdot x^2}
R
\xrightarrow{\cdot x}
R
\longrightarrow
M_1
\longrightarrow
0.
\]

We claim that this is exact. The final map
\[
R\longrightarrow M_1=R/(x)
\]
is the quotient map. Its kernel is
\[
(x).
\]
But \((x)\) is exactly the image of multiplication by \(x\):
\[
R\xrightarrow{\cdot x} R.
\]

Next,
\[
\ker(\cdot x:R\to R)
=
\{r\in R: xr=0\}.
\]
Since \(R=\mathbb C[x]/(x^3)\), every element has the form
\[
r=a+bx+cx^2.
\]
Then
\[
xr=ax+bx^2.
\]
This is zero in \(R\) if and only if
\[
a=0,\qquad b=0.
\]
Hence
\[
r=cx^2,
\]
so
\[
\ker(\cdot x)=(x^2).
\]
But this is exactly the image of multiplication by \(x^2\).

Similarly,
\[
\ker(\cdot x^2:R\to R)
=
\{r\in R:x^2r=0\}.
\]
For \(r=a+bx+cx^2\), we get
\[
x^2r=ax^2.
\]
Thus \(x^2r=0\) if and only if \(a=0\), so
\[
r=bx+cx^2.
\]
Hence
\[
\ker(\cdot x^2)=(x),
\]
which is the image of multiplication by \(x\). Therefore the resolution is exact.

So a free resolution of \(M_1\) is
\[
\boxed{
	\cdots
	\longrightarrow
	R
	\xrightarrow{\cdot x^2}
	R
	\xrightarrow{\cdot x}
	R
	\xrightarrow{\cdot x^2}
	R
	\xrightarrow{\cdot x}
	R
	\longrightarrow
	M_1
	\longrightarrow
	0.
}
\]

Now tensor this free resolution with \(M_1\). Since \(x=0\) on \(M_1=R/(x)\), both multiplication by \(x\) and multiplication by \(x^2\) become zero maps. Thus the tensor complex is
\[
\cdots
\longrightarrow
M_1
\xrightarrow{0}
M_1
\xrightarrow{0}
M_1
\xrightarrow{0}
M_1
\xrightarrow{0}
M_1
\longrightarrow
0.
\]
Therefore every homology group is just \(M_1\). Hence
\[
\boxed{
	\operatorname{Tor}^R_n(M_1,M_1)\cong M_1
	\cong \mathbb C
	\quad \text{for every } n\geq 0.
}
\]

Now tensor the same free resolution with
\[
M_2=R/(x^2).
\]
On \(M_2\), multiplication by \(x^2\) is zero, but multiplication by \(x\) is nonzero. Indeed, as a \(\mathbb C\)-vector space,
\[
M_2=\operatorname{span}_{\mathbb C}\{1,x\},
\]
with
\[
x\cdot 1=x,\qquad x\cdot x=x^2=0.
\]
Therefore the tensor complex becomes
\[
\cdots
\longrightarrow
M_2
\xrightarrow{\cdot x}
M_2
\xrightarrow{0}
M_2
\xrightarrow{\cdot x}
M_2
\xrightarrow{0}
M_2
\xrightarrow{\cdot x}
M_2
\longrightarrow
0.
\]

We compute the homology. First,
\[
\operatorname{im}(\cdot x)=xM_2=(x),
\]
which is the one-dimensional vector space spanned by \(x\). Also,
\[
\ker(\cdot x)=\{m\in M_2:xm=0\}.
\]
If
\[
m=a+bx,
\]
then
\[
xm=ax.
\]
Hence \(xm=0\) if and only if \(a=0\), so
\[
\ker(\cdot x)=(x).
\]

Thus
\[
\ker(\cdot x)=(x),
\qquad
\operatorname{im}(\cdot x)=(x),
\qquad
\ker(0)=M_2,
\qquad
\operatorname{im}(0)=0.
\]

For degree \(0\), we get
\[
\operatorname{Tor}^R_0(M_1,M_2)
\cong
M_1\otimes_R M_2
\cong
M_2/xM_2
\cong
M_2/(x)
\cong
\mathbb C.
\]

For degree \(1\),
\[
\operatorname{Tor}^R_1(M_1,M_2)
\cong
\ker(\cdot x)/\operatorname{im}(0)
=
(x)/0
\cong
\mathbb C.
\]

For degree \(2\),
\[
\operatorname{Tor}^R_2(M_1,M_2)
\cong
\ker(0)/\operatorname{im}(\cdot x)
=
M_2/(x)
\cong
\mathbb C.
\]

The same pattern repeats periodically. Therefore
\[
\boxed{
	\operatorname{Tor}^R_n(M_1,M_2)\cong \mathbb C
	\quad \text{for every } n\geq 0.
}
\]

Equivalently, as \(R\)-modules,
\[
\boxed{
	\operatorname{Tor}^R_n(M_1,M_2)\cong M_1
	\quad \text{for every } n\geq 0.
}
\]
\end{solution}

\begin{problem}[CP-VIII-0068 --- Extension from \(X\times\{0\}\cup A\times I\)]
Let \(X\) be a finite cell complex, and let \(A\) be a subcomplex of \(X\). Let
\[
Y=X\times\{0\}\cup A\times I
\subset X\times I.
\]
Show that any map
\[
Y\longrightarrow Z
\]
extends to a map
\[
X\times I\longrightarrow Z.
\]
Hint: start with the case
\[
X=D^n,\qquad A=S^{n-1}.
\]
\end{problem}

\begin{solution}
The space
\[
Y=X\times\{0\}\cup A\times I
\]
consists of two parts:

\[
X\times\{0\},
\]
which is the whole space \(X\) at time \(0\), and
\[
A\times I,
\]
which is the cylinder on the subcomplex \(A\).

So \(Y\) contains the initial copy of \(X\), together with the whole homotopy cylinder over \(A\). The problem asks us to prove that any map already defined on this part extends over the whole cylinder \(X\times I\).

The key point is that the inclusion
\[
Y\subset X\times I
\]
has a retraction. In other words, we will show that there exists a continuous map
\[
r:X\times I\longrightarrow Y
\]
such that
\[
r(y)=y
\]
for every \(y\in Y\). Then, given any map
\[
g:Y\longrightarrow Z,
\]
we can define its extension by
\[
\widetilde g=g\circ r:X\times I\longrightarrow Z.
\]
This clearly restricts to \(g\) on \(Y\), because \(r\) is the identity on \(Y\).

We first prove the basic cell case:
\[
X=D^n,\qquad A=S^{n-1}.
\]
Then
\[
Y=D^n\times\{0\}\cup S^{n-1}\times I.
\]
Geometrically, this is the bottom disk of the cylinder \(D^n\times I\), together with the side wall of the cylinder.

For example, when \(n=1\), the space \(D^1\times I\) is a square, and
\[
D^1\times\{0\}\cup S^0\times I
\]
is the union of the bottom edge and the two vertical side edges. This is a deformation retract of the square.

When \(n=2\), the space \(D^2\times I\) is a solid cylinder, and
\[
D^2\times\{0\}\cup S^1\times I
\]
is the union of the bottom disk and the cylindrical side wall. Again, this is a deformation retract of the solid cylinder.

In general, there is a retraction
\[
r_n:D^n\times I\longrightarrow D^n\times\{0\}\cup S^{n-1}\times I.
\]
One way to visualize it is to push every point of the cylinder \(D^n\times I\) outward and downward until it lands either on the bottom \(D^n\times\{0\}\) or on the side \(S^{n-1}\times I\). Points already on the bottom or on the side are fixed.

Therefore, for the pair
\[
(D^n,S^{n-1}),
\]
any map
\[
D^n\times\{0\}\cup S^{n-1}\times I
\longrightarrow Z
\]
extends to a map
\[
D^n\times I\longrightarrow Z
\]
by composing with this retraction.

Now let \(X\) be a finite cell complex and \(A\subset X\) a subcomplex. Since \(A\) is a subcomplex, the cells of \(X\setminus A\) are attached to \(A\) skeleton by skeleton. We extend the given map cell by cell.

Assume that the map has already been extended over the cylinder on the part of \(X\) constructed so far. Let \(e^n\) be an \(n\)-cell of \(X\setminus A\). Its characteristic map has the form
\[
\varphi:D^n\longrightarrow X,
\]
with boundary
\[
\varphi(S^{n-1})
\]
lying in the previously constructed part of \(X\).

Over this cell, the new part to be filled is
\[
D^n\times I.
\]
The map is already defined on
\[
D^n\times\{0\}\cup S^{n-1}\times I,
\]
because \(D^n\times\{0\}\) belongs to \(X\times\{0\}\), and \(S^{n-1}\times I\) lies over the previously constructed skeleton.

By the basic cell case, the map extends over
\[
D^n\times I.
\]
Doing this for every cell of \(X\setminus A\), and using the fact that \(X\) is finite, we obtain an extension over all of
\[
X\times I.
\]

Thus every map
\[
g:Y\longrightarrow Z
\]
extends to a map
\[
\widetilde g:X\times I\longrightarrow Z.
\]

\[
\boxed{
	\text{Every map }Y=X\times\{0\}\cup A\times I\to Z
	\text{ extends to }X\times I\to Z.
}
\]

This is exactly the cellular form of the homotopy extension property for the pair
\[
(A\subset X).
\]
\end{solution}

\begin{problem}[CP-VIII-0069 --- Acyclic non-contractible complexes and homology-equivalent maps]
This question assumes some knowledge of \(\pi_1\).

\begin{enumerate}
	\item[(a)] Construct a connected finite cell complex \(X\) with
	\[
	\widetilde H_*(X)=0
	\]
	but which is not contractible.

	\item[(b)] Find a map
	\[
	f:S^1\vee S^2\to S^1\vee S^2
	\]
	which induces the identity map on
	\[
	H_*(S^1\vee S^2),
	\]
	but is not homotopic to the identity.
\end{enumerate}
\end{problem}

\begin{solution}

\medskip
\noindent\textbf{(a) A finite acyclic complex which is not contractible}\par
\smallskip


We construct a finite \(2\)-dimensional cell complex from a group presentation.

Let
\[
G=
\left\langle
x,y,z
\ \middle|\
x^2=xyz,\quad
y^3=xyz,\quad
z^5=xyz
\right\rangle.
\]

This is a standard presentation of the binary icosahedral group, a non-trivial finite perfect group.

Let \(X\) be the presentation complex of this presentation. Thus \(X\) has:

\begin{itemize}
	\item one \(0\)-cell,
	\item three \(1\)-cells corresponding to \(x,y,z\),
	\item three \(2\)-cells corresponding to the relators
	\[
	r_1=x^2(xyz)^{-1},
	\]
	\[
	r_2=y^3(xyz)^{-1},
	\]
	\[
	r_3=z^5(xyz)^{-1}.
	\]
\end{itemize}

The cellular chain complex of \(X\) has the form
\[
0
\longrightarrow
C_2(X)
\xrightarrow{d_2}
C_1(X)
\xrightarrow{d_1}
C_0(X)
\longrightarrow
0.
\]

Since there is one \(0\)-cell, and all \(1\)-cells are loops based at it,
\[
d_1=0.
\]

Also,
\[
C_2(X)\cong \mathbb Z^3,
\qquad
C_1(X)\cong \mathbb Z^3.
\]

The boundary map \(d_2\) is given by the exponent-sum matrix of the relators. The exponent-sum vectors are
\[
r_1:\quad (1,-1,-1),
\]
\[
r_2:\quad (-1,2,-1),
\]
\[
r_3:\quad (-1,-1,4).
\]

Therefore
\[
d_2
=
\begin{pmatrix}
	1 & -1 & -1\\
	-1 & 2 & -1\\
	-1 & -1 & 4
\end{pmatrix}.
\]

Its determinant is
\[
\det
\begin{pmatrix}
	1 & -1 & -1\\
	-1 & 2 & -1\\
	-1 & -1 & 4
\end{pmatrix}
=
-1.
\]

Hence \(d_2\) is an isomorphism
\[
\mathbb Z^3\to \mathbb Z^3.
\]

Therefore
\[
H_2(X)=\ker d_2=0,
\]
and
\[
H_1(X)=\ker d_1/\operatorname{im}d_2
=
\mathbb Z^3/\mathbb Z^3
=
0.
\]

Since \(X\) is connected,
\[
H_0(X)\cong \mathbb Z.
\]

Thus
\[
\widetilde H_*(X)=0.
\]

So \(X\) is acyclic.

However,
\[
\pi_1(X)\cong G,
\]
and \(G\) is non-trivial. Therefore \(X\) is not simply connected.

If \(X\) were contractible, then it would be simply connected. Hence \(X\) is not contractible.

Thus we have constructed a connected finite cell complex such that
\[
\boxed{
	\widetilde H_*(X)=0
	\quad
	\text{but}
	\quad
	X\text{ is not contractible.}
}
\]


\medskip
\noindent\textbf{(b) A homology identity map which is not homotopic to the identity}\par
\smallskip


Let
\[
X=S^1\vee S^2.
\]

Let
\[
a\in \pi_1(X)
\]
be the generator coming from the \(S^1\)-summand.

Let
\[
u\in \pi_2(X)
\]
be the class represented by the inclusion
\[
S^2\hookrightarrow S^1\vee S^2.
\]

The universal cover of \(X=S^1\vee S^2\) is an infinite line, with a copy of \(S^2\) attached at each integer point. Therefore
\[
\pi_2(X)\cong \mathbb Z[t,t^{-1}]u,
\]
where \(t\) records the action of the generator \(a\in \pi_1(X)\).

The ordinary homology group
\[
H_2(X)\cong \mathbb Z
\]
is obtained from \(\pi_2(X)\) by forgetting the \(t\)-position. Algebraically, this corresponds to the augmentation map
\[
\varepsilon:\mathbb Z[t,t^{-1}]\to \mathbb Z,
\]
defined by
\[
\varepsilon\left(\sum_k m_k t^k\right)
=
\sum_k m_k.
\]

Now define a map
\[
f:S^1\vee S^2\to S^1\vee S^2
\]
as follows.

On the \(S^1\)-summand, let
\[
f|_{S^1}=\operatorname{id}_{S^1}.
\]

On the \(S^2\)-summand, choose a based map representing the element
\[
(2-t)u\in \pi_2(X)\cong \mathbb Z[t,t^{-1}]u.
\]

Thus
\[
f_*(u)=(2-t)u.
\]

We now check the induced map on homology.

In degree \(1\), \(f\) is the identity on \(S^1\), so
\[
f_*:H_1(X)\to H_1(X)
\]
is the identity.

In degree \(2\), the induced map is multiplication by
\[
\varepsilon(2-t).
\]
But
\[
\varepsilon(2-t)=2-1=1.
\]
Therefore
\[
f_*:H_2(X)\to H_2(X)
\]
is also the identity.

Hence \(f\) induces the identity on all homology groups:
\[
\boxed{
	f_*=\operatorname{id}
	\text{ on }H_*(S^1\vee S^2).
}
\]

However, \(f\) is not homotopic to the identity.

Indeed, the identity map sends
\[
u\mapsto u.
\]
But \(f\) sends
\[
u\mapsto (2-t)u.
\]

Since
\[
2-t\neq 1
\]
in
\[
\mathbb Z[t,t^{-1}],
\]
the induced maps on \(\pi_2(X)\) are different. Therefore \(f\) is not based-homotopic to the identity.

Even if one allows free homotopy, a homotopy from \(f\) to the identity could only change the induced action on \(\pi_2\) by multiplication by a monomial \(t^k\). But
\[
2-t
\]
is not a monomial. Hence \(f\) is not freely homotopic to the identity either.

Thus
\[
\boxed{
	f:S^1\vee S^2\to S^1\vee S^2
	\text{ induces the identity on homology but is not homotopic to }
	\operatorname{id}.
}
\]
\end{solution}

\begin{problem}[CP-VIII-0070 --- Products, Künneth Formula, Universal Coefficient Theorem, and Euler Characteristic]
Let
\[
X=L_4^3\times \mathbb{RP}^3.
\]
\begin{enumerate}
	\item[(a)] Write out \(C_*^{\mathrm{cell}}(X)\) and use it to compute \(H_*(X)\).
	\item[(b)] Compute \(H_*(X)\) using the Künneth formula and verify that it agrees with the answer in part (a).
	\item[(c)] Use the universal coefficient theorem and the answer to part (b) to compute \(H^*(X)\) and \(H_*(X;\mathbb Z/2)\).
	\item[(d)] Use the universal coefficient theorem to compute \(H_*(L_4^3;\mathbb Z/2)\) and \(H_*(\mathbb{RP}^3;\mathbb Z/2)\). Use this answer to compute \(H_*(X;\mathbb Z/2)\) and verify that it agrees with the answer in part (c).
\end{enumerate}
\end{problem}

\begin{solution}

\medskip
\noindent\textbf{The cellular chain complexes of \(L_4^3\) and \(\mathbb{RP}^3\)}\par
\smallskip


The lens space \(L_4^3\) has a CW-structure with one cell in each dimension
\[
0,1,2,3.
\]
Its cellular chain complex is
\[
0
\longrightarrow C_3(L_4^3)\cong \mathbb Z
\xrightarrow{\partial_3=0}
C_2(L_4^3)\cong \mathbb Z
\xrightarrow{\partial_2=4}
C_1(L_4^3)\cong \mathbb Z
\xrightarrow{\partial_1=0}
C_0(L_4^3)\cong \mathbb Z
\longrightarrow 0.
\]
Therefore
\[
H_0(L_4^3)\cong \mathbb Z,
\qquad
H_1(L_4^3)\cong \mathbb Z/4,
\qquad
H_2(L_4^3)=0,
\qquad
H_3(L_4^3)\cong \mathbb Z.
\]

Similarly, \(\mathbb{RP}^3\) has a CW-structure with one cell in each dimension
\[
0,1,2,3.
\]
Its cellular chain complex is
\[
0
\longrightarrow C_3(\mathbb{RP}^3)\cong \mathbb Z
\xrightarrow{\partial_3=0}
C_2(\mathbb{RP}^3)\cong \mathbb Z
\xrightarrow{\partial_2=2}
C_1(\mathbb{RP}^3)\cong \mathbb Z
\xrightarrow{\partial_1=0}
C_0(\mathbb{RP}^3)\cong \mathbb Z
\longrightarrow 0.
\]
Hence
\[
H_0(\mathbb{RP}^3)\cong \mathbb Z,
\qquad
H_1(\mathbb{RP}^3)\cong \mathbb Z/2,
\qquad
H_2(\mathbb{RP}^3)=0,
\qquad
H_3(\mathbb{RP}^3)\cong \mathbb Z.
\]


\medskip
\noindent\textbf{Part (a): The cellular chain complex of \(X\)}\par
\smallskip


The product CW-structure on
\[
X=L_4^3\times \mathbb{RP}^3
\]
has one cell \(e_{i,j}\) in dimension \(i+j\) for each pair
\[
0\leq i\leq 3,
\qquad
0\leq j\leq 3.
\]
Here \(e_{i,j}\) denotes the product of the \(i\)-cell of \(L_4^3\) with the \(j\)-cell of \(\mathbb{RP}^3\).

Thus
\[
C_n^{\mathrm{cell}}(X)
\cong
\bigoplus_{i+j=n}
C_i^{\mathrm{cell}}(L_4^3)\otimes C_j^{\mathrm{cell}}(\mathbb{RP}^3).
\]
Therefore the chain groups are
\[
C_0(X)\cong \mathbb Z,
\]
\[
C_1(X)\cong \mathbb Z^2,
\]
\[
C_2(X)\cong \mathbb Z^3,
\]
\[
C_3(X)\cong \mathbb Z^4,
\]
\[
C_4(X)\cong \mathbb Z^3,
\]
\[
C_5(X)\cong \mathbb Z^2,
\]
\[
C_6(X)\cong \mathbb Z.
\]

The cellular differential on the product is given by
\[
\partial(a\otimes b)
=
\partial a\otimes b
+
(-1)^{|a|}a\otimes \partial b.
\]

We use the following ordered bases:
\[
C_0=\langle e_{0,0}\rangle,
\]
\[
C_1=\langle e_{1,0},e_{0,1}\rangle,
\]
\[
C_2=\langle e_{2,0},e_{1,1},e_{0,2}\rangle,
\]
\[
C_3=\langle e_{3,0},e_{2,1},e_{1,2},e_{0,3}\rangle,
\]
\[
C_4=\langle e_{3,1},e_{2,2},e_{1,3}\rangle,
\]
\[
C_5=\langle e_{3,2},e_{2,3}\rangle,
\]
\[
C_6=\langle e_{3,3}\rangle.
\]

With respect to these bases, the boundary maps are as follows.

First,
\[
\partial_1=0.
\]
Next,
\[
\partial_2(e_{2,0})=4e_{1,0},
\qquad
\partial_2(e_{1,1})=0,
\qquad
\partial_2(e_{0,2})=2e_{0,1}.
\]
Hence
\[
[\partial_2]
=
\begin{pmatrix}
	4 & 0 & 0\\
	0 & 0 & 2
\end{pmatrix}.
\]

For \(\partial_3\), we get
\[
\partial_3(e_{3,0})=0,
\]
\[
\partial_3(e_{2,1})=4e_{1,1},
\]
\[
\partial_3(e_{1,2})=-2e_{1,1},
\]
\[
\partial_3(e_{0,3})=0.
\]
Thus
\[
[\partial_3]
=
\begin{pmatrix}
	0 & 0 & 0 & 0\\
	0 & 4 & -2 & 0\\
	0 & 0 & 0 & 0
\end{pmatrix}.
\]

For \(\partial_4\), we have
\[
\partial_4(e_{3,1})=0,
\]
\[
\partial_4(e_{2,2})
=
2e_{2,1}+4e_{1,2},
\]
\[
\partial_4(e_{1,3})=0.
\]
So
\[
[\partial_4]
=
\begin{pmatrix}
	0 & 0 & 0\\
	0 & 2 & 0\\
	0 & 4 & 0\\
	0 & 0 & 0
\end{pmatrix}.
\]

For \(\partial_5\), we obtain
\[
\partial_5(e_{3,2})=-2e_{3,1},
\qquad
\partial_5(e_{2,3})=4e_{1,3}.
\]
Hence
\[
[\partial_5]
=
\begin{pmatrix}
	-2 & 0\\
	0 & 0\\
	0 & 4
\end{pmatrix}.
\]

Finally,
\[
\partial_6(e_{3,3})=0.
\]
Therefore
\[
\partial_6=0.
\]

So the full cellular chain complex is
\[
0
\longrightarrow
\mathbb Z
\xrightarrow{\partial_6=0}
\mathbb Z^2
\xrightarrow{\partial_5}
\mathbb Z^3
\xrightarrow{\partial_4}
\mathbb Z^4
\xrightarrow{\partial_3}
\mathbb Z^3
\xrightarrow{\partial_2}
\mathbb Z^2
\xrightarrow{\partial_1=0}
\mathbb Z
\longrightarrow 0.
\]

We now compute homology degree by degree.

In degree \(0\),
\[
H_0(X)\cong \mathbb Z.
\]

In degree \(1\),
\[
H_1(X)
=
\ker \partial_1/\operatorname{im}\partial_2
=
\mathbb Z^2/\langle 4e_{1,0},2e_{0,1}\rangle.
\]
Hence
\[
H_1(X)\cong \mathbb Z/4\oplus \mathbb Z/2.
\]

In degree \(2\), the map \(\partial_2\) sends
\[
a e_{2,0}+b e_{1,1}+c e_{0,2}
\]
to
\[
4a e_{1,0}+2c e_{0,1}.
\]
Thus
\[
\ker \partial_2=\langle e_{1,1}\rangle\cong \mathbb Z.
\]
Also,
\[
\operatorname{im}\partial_3
=
\langle 4e_{1,1},-2e_{1,1}\rangle
=
\langle 2e_{1,1}\rangle.
\]
Therefore
\[
H_2(X)
\cong
\mathbb Z/\langle 2\rangle
\cong
\mathbb Z/2.
\]

In degree \(3\), an element of \(C_3\) has the form
\[
a e_{3,0}+b e_{2,1}+c e_{1,2}+d e_{0,3}.
\]
The condition that it lies in \(\ker \partial_3\) is
\[
4b-2c=0.
\]
Equivalently,
\[
c=2b.
\]
Thus
\[
\ker \partial_3
=
\langle e_{3,0},\, e_{2,1}+2e_{1,2},\, e_{0,3}\rangle
\cong \mathbb Z^3.
\]
Meanwhile
\[
\operatorname{im}\partial_4
=
\langle 2e_{2,1}+4e_{1,2}\rangle
=
\langle 2(e_{2,1}+2e_{1,2})\rangle.
\]
Therefore
\[
H_3(X)
\cong
\mathbb Z\oplus \mathbb Z/2\oplus \mathbb Z
\cong
\mathbb Z^2\oplus \mathbb Z/2.
\]

In degree \(4\), we have
\[
\partial_4(ae_{3,1}+be_{2,2}+ce_{1,3})
=
b(2e_{2,1}+4e_{1,2}).
\]
Therefore
\[
\ker \partial_4
=
\langle e_{3,1},e_{1,3}\rangle
\cong \mathbb Z^2.
\]
Also,
\[
\operatorname{im}\partial_5
=
\langle -2e_{3,1},4e_{1,3}\rangle.
\]
Hence
\[
H_4(X)
\cong
\mathbb Z/2\oplus \mathbb Z/4.
\]
Equivalently,
\[
H_4(X)\cong \mathbb Z/4\oplus \mathbb Z/2.
\]

In degree \(5\),
\[
\partial_5(ae_{3,2}+be_{2,3})
=
-2a e_{3,1}+4b e_{1,3}.
\]
This vanishes only if
\[
a=0,
\qquad
b=0.
\]
Therefore
\[
\ker \partial_5=0.
\]
Since \(\operatorname{im}\partial_6=0\), we get
\[
H_5(X)=0.
\]

Finally, in degree \(6\),
\[
\partial_6=0,
\]
and there are no \(7\)-cells. Therefore
\[
H_6(X)\cong \mathbb Z.
\]

Thus
\[
\boxed{
	H_k(X)\cong
	\begin{cases}
		\mathbb Z, & k=0,6,\\[1mm]
		\mathbb Z/4\oplus \mathbb Z/2, & k=1,4,\\[1mm]
		\mathbb Z/2, & k=2,\\[1mm]
		\mathbb Z^2\oplus \mathbb Z/2, & k=3,\\[1mm]
		0, & k=5.
\end{cases}}
\]


\medskip
\noindent\textbf{Part (b): Computation using the Künneth formula}\par
\smallskip


The integral Künneth theorem gives a short exact sequence
\[
0
\longrightarrow
\bigoplus_{i+j=n}
H_i(L_4^3)\otimes H_j(\mathbb{RP}^3)
\longrightarrow
H_n(X)
\longrightarrow
\bigoplus_{i+j=n-1}
\operatorname{Tor}\bigl(H_i(L_4^3),H_j(\mathbb{RP}^3)\bigr)
\longrightarrow
0.
\]

We use the standard facts
\[
\mathbb Z\otimes G\cong G,
\]
\[
\mathbb Z/m\otimes \mathbb Z/n\cong \mathbb Z/\gcd(m,n),
\]
and
\[
\operatorname{Tor}(\mathbb Z/m,\mathbb Z/n)
\cong
\mathbb Z/\gcd(m,n).
\]

Since
\[
H_*(L_4^3)
=
(\mathbb Z,\mathbb Z/4,0,\mathbb Z),
\]
and
\[
H_*(\mathbb{RP}^3)
=
(\mathbb Z,\mathbb Z/2,0,\mathbb Z),
\]
we compute degree by degree.

For \(n=0\),
\[
H_0(X)
\cong
H_0(L_4^3)\otimes H_0(\mathbb{RP}^3)
\cong
\mathbb Z.
\]

For \(n=1\),
\[
H_1(X)
\cong
(H_1(L_4^3)\otimes H_0(\mathbb{RP}^3))
\oplus
(H_0(L_4^3)\otimes H_1(\mathbb{RP}^3)).
\]
Thus
\[
H_1(X)
\cong
\mathbb Z/4\oplus \mathbb Z/2.
\]

For \(n=2\), the only nonzero tensor contribution is
\[
H_1(L_4^3)\otimes H_1(\mathbb{RP}^3)
\cong
\mathbb Z/4\otimes \mathbb Z/2
\cong
\mathbb Z/2.
\]
There is no nonzero Tor contribution in this degree. Therefore
\[
H_2(X)\cong \mathbb Z/2.
\]

For \(n=3\), the tensor contributions are
\[
H_3(L_4^3)\otimes H_0(\mathbb{RP}^3)
\cong
\mathbb Z,
\]
and
\[
H_0(L_4^3)\otimes H_3(\mathbb{RP}^3)
\cong
\mathbb Z.
\]
The Tor contribution is
\[
\operatorname{Tor}(H_1(L_4^3),H_1(\mathbb{RP}^3))
=
\operatorname{Tor}(\mathbb Z/4,\mathbb Z/2)
\cong
\mathbb Z/2.
\]
Hence
\[
H_3(X)
\cong
\mathbb Z^2\oplus \mathbb Z/2.
\]

For \(n=4\), the tensor contributions are
\[
H_3(L_4^3)\otimes H_1(\mathbb{RP}^3)
\cong
\mathbb Z/2,
\]
and
\[
H_1(L_4^3)\otimes H_3(\mathbb{RP}^3)
\cong
\mathbb Z/4.
\]
There is no nonzero Tor contribution. Thus
\[
H_4(X)
\cong
\mathbb Z/4\oplus \mathbb Z/2.
\]

For \(n=5\), all tensor and Tor terms vanish. Therefore
\[
H_5(X)=0.
\]

For \(n=6\),
\[
H_6(X)
\cong
H_3(L_4^3)\otimes H_3(\mathbb{RP}^3)
\cong
\mathbb Z.
\]

Thus the Künneth formula gives
\[
\boxed{
	H_k(X)\cong
	\begin{cases}
		\mathbb Z, & k=0,6,\\[1mm]
		\mathbb Z/4\oplus \mathbb Z/2, & k=1,4,\\[1mm]
		\mathbb Z/2, & k=2,\\[1mm]
		\mathbb Z^2\oplus \mathbb Z/2, & k=3,\\[1mm]
		0, & k=5.
\end{cases}}
\]
This agrees exactly with the computation using cellular chains.


\medskip
\noindent\textbf{Part (c): Integral cohomology and mod-\(2\) homology of \(X\)}\par
\smallskip


The universal coefficient theorem for cohomology says that there is a short exact sequence
\[
0
\longrightarrow
\operatorname{Ext}(H_{n-1}(X),\mathbb Z)
\longrightarrow
H^n(X)
\longrightarrow
\operatorname{Hom}(H_n(X),\mathbb Z)
\longrightarrow
0.
\]

We use
\[
\operatorname{Hom}(\mathbb Z,\mathbb Z)\cong \mathbb Z,
\qquad
\operatorname{Hom}(\mathbb Z/m,\mathbb Z)=0,
\]
and
\[
\operatorname{Ext}(\mathbb Z,\mathbb Z)=0,
\qquad
\operatorname{Ext}(\mathbb Z/m,\mathbb Z)\cong \mathbb Z/m.
\]

Since
\[
H_0(X)\cong \mathbb Z,
\]
\[
H_1(X)\cong \mathbb Z/4\oplus \mathbb Z/2,
\]
\[
H_2(X)\cong \mathbb Z/2,
\]
\[
H_3(X)\cong \mathbb Z^2\oplus \mathbb Z/2,
\]
\[
H_4(X)\cong \mathbb Z/4\oplus \mathbb Z/2,
\]
\[
H_5(X)=0,
\]
\[
H_6(X)\cong \mathbb Z,
\]
we obtain
\[
H^0(X)\cong \mathbb Z,
\]
\[
H^1(X)=0,
\]
\[
H^2(X)\cong \mathbb Z/4\oplus \mathbb Z/2,
\]
\[
H^3(X)\cong \mathbb Z^2\oplus \mathbb Z/2,
\]
\[
H^4(X)\cong \mathbb Z/2,
\]
\[
H^5(X)\cong \mathbb Z/4\oplus \mathbb Z/2,
\]
\[
H^6(X)\cong \mathbb Z.
\]

Therefore
\[
\boxed{
	H^k(X)\cong
	\begin{cases}
		\mathbb Z, & k=0,6,\\[1mm]
		0, & k=1,\\[1mm]
		\mathbb Z/4\oplus \mathbb Z/2, & k=2,5,\\[1mm]
		\mathbb Z^2\oplus \mathbb Z/2, & k=3,\\[1mm]
		\mathbb Z/2, & k=4.
\end{cases}}
\]

Now we compute \(H_*(X;\mathbb Z/2)\). The universal coefficient theorem for homology says that
\[
0
\longrightarrow
H_n(X)\otimes \mathbb Z/2
\longrightarrow
H_n(X;\mathbb Z/2)
\longrightarrow
\operatorname{Tor}(H_{n-1}(X),\mathbb Z/2)
\longrightarrow
0.
\]

Since the coefficient group is a field, the sequence splits, although not canonically. We use
\[
\mathbb Z\otimes \mathbb Z/2\cong \mathbb Z/2,
\]
\[
\mathbb Z/m\otimes \mathbb Z/2
\cong
\mathbb Z/\gcd(m,2),
\]
and
\[
\operatorname{Tor}(\mathbb Z/m,\mathbb Z/2)
\cong
\mathbb Z/\gcd(m,2).
\]

For \(n=0\),
\[
H_0(X;\mathbb Z/2)\cong \mathbb Z/2.
\]

For \(n=1\),
\[
H_1(X)\otimes \mathbb Z/2
\cong
(\mathbb Z/4\oplus \mathbb Z/2)\otimes \mathbb Z/2
\cong
(\mathbb Z/2)^2.
\]
Also,
\[
\operatorname{Tor}(H_0(X),\mathbb Z/2)=0.
\]
Hence
\[
H_1(X;\mathbb Z/2)\cong (\mathbb Z/2)^2.
\]

For \(n=2\),
\[
H_2(X)\otimes \mathbb Z/2
\cong
\mathbb Z/2,
\]
and
\[
\operatorname{Tor}(H_1(X),\mathbb Z/2)
\cong
\operatorname{Tor}(\mathbb Z/4,\mathbb Z/2)
\oplus
\operatorname{Tor}(\mathbb Z/2,\mathbb Z/2)
\cong
(\mathbb Z/2)^2.
\]
Therefore
\[
H_2(X;\mathbb Z/2)\cong (\mathbb Z/2)^3.
\]

For \(n=3\),
\[
H_3(X)\otimes \mathbb Z/2
\cong
(\mathbb Z^2\oplus \mathbb Z/2)\otimes \mathbb Z/2
\cong
(\mathbb Z/2)^3.
\]
Also,
\[
\operatorname{Tor}(H_2(X),\mathbb Z/2)
\cong
\mathbb Z/2.
\]
Thus
\[
H_3(X;\mathbb Z/2)\cong (\mathbb Z/2)^4.
\]

For \(n=4\),
\[
H_4(X)\otimes \mathbb Z/2
\cong
(\mathbb Z/4\oplus \mathbb Z/2)\otimes \mathbb Z/2
\cong
(\mathbb Z/2)^2.
\]
Also,
\[
\operatorname{Tor}(H_3(X),\mathbb Z/2)
\cong
\mathbb Z/2.
\]
Therefore
\[
H_4(X;\mathbb Z/2)\cong (\mathbb Z/2)^3.
\]

For \(n=5\),
\[
H_5(X)\otimes \mathbb Z/2=0,
\]
and
\[
\operatorname{Tor}(H_4(X),\mathbb Z/2)
\cong
\operatorname{Tor}(\mathbb Z/4,\mathbb Z/2)
\oplus
\operatorname{Tor}(\mathbb Z/2,\mathbb Z/2)
\cong
(\mathbb Z/2)^2.
\]
Hence
\[
H_5(X;\mathbb Z/2)\cong (\mathbb Z/2)^2.
\]

For \(n=6\),
\[
H_6(X)\otimes \mathbb Z/2
\cong
\mathbb Z/2,
\]
and
\[
\operatorname{Tor}(H_5(X),\mathbb Z/2)=0.
\]
Thus
\[
H_6(X;\mathbb Z/2)\cong \mathbb Z/2.
\]

Therefore
\[
\boxed{
	H_k(X;\mathbb Z/2)\cong
	\begin{cases}
		\mathbb Z/2, & k=0,6,\\[1mm]
		(\mathbb Z/2)^2, & k=1,5,\\[1mm]
		(\mathbb Z/2)^3, & k=2,4,\\[1mm]
		(\mathbb Z/2)^4, & k=3.
\end{cases}}
\]


\medskip
\noindent\textbf{Part (d): Mod-\(2\) homology of the factors and verification}\par
\smallskip


First consider \(L_4^3\). We know
\[
H_0(L_4^3)\cong \mathbb Z,
\qquad
H_1(L_4^3)\cong \mathbb Z/4,
\qquad
H_2(L_4^3)=0,
\qquad
H_3(L_4^3)\cong \mathbb Z.
\]

Using the universal coefficient theorem for homology,
\[
0
\longrightarrow
H_n(L_4^3)\otimes \mathbb Z/2
\longrightarrow
H_n(L_4^3;\mathbb Z/2)
\longrightarrow
\operatorname{Tor}(H_{n-1}(L_4^3),\mathbb Z/2)
\longrightarrow
0,
\]
we get
\[
H_0(L_4^3;\mathbb Z/2)\cong \mathbb Z/2,
\]
\[
H_1(L_4^3;\mathbb Z/2)\cong \mathbb Z/2,
\]
\[
H_2(L_4^3;\mathbb Z/2)\cong \mathbb Z/2,
\]
\[
H_3(L_4^3;\mathbb Z/2)\cong \mathbb Z/2.
\]
Thus
\[
\boxed{
	H_k(L_4^3;\mathbb Z/2)\cong \mathbb Z/2
	\quad\text{for }k=0,1,2,3.}
\]

Now consider \(\mathbb{RP}^3\). We know
\[
H_0(\mathbb{RP}^3)\cong \mathbb Z,
\qquad
H_1(\mathbb{RP}^3)\cong \mathbb Z/2,
\qquad
H_2(\mathbb{RP}^3)=0,
\qquad
H_3(\mathbb{RP}^3)\cong \mathbb Z.
\]
Again by the universal coefficient theorem,
\[
H_0(\mathbb{RP}^3;\mathbb Z/2)\cong \mathbb Z/2,
\]
\[
H_1(\mathbb{RP}^3;\mathbb Z/2)\cong \mathbb Z/2,
\]
\[
H_2(\mathbb{RP}^3;\mathbb Z/2)\cong \mathbb Z/2,
\]
\[
H_3(\mathbb{RP}^3;\mathbb Z/2)\cong \mathbb Z/2.
\]
Therefore
\[
\boxed{
	H_k(\mathbb{RP}^3;\mathbb Z/2)\cong \mathbb Z/2
	\quad\text{for }k=0,1,2,3.}
\]

Now apply the Künneth formula over the field \(\mathbb Z/2\). Since \(\mathbb Z/2\) is a field, there are no Tor terms:
\[
H_n(X;\mathbb Z/2)
\cong
\bigoplus_{i+j=n}
H_i(L_4^3;\mathbb Z/2)
\otimes_{\mathbb Z/2}
H_j(\mathbb{RP}^3;\mathbb Z/2).
\]

Each factor has one copy of \(\mathbb Z/2\) in degrees \(0,1,2,3\). Therefore the dimension of \(H_n(X;\mathbb Z/2)\) is the number of pairs
\[
(i,j)
\]
such that
\[
i+j=n,
\qquad
0\leq i,j\leq 3.
\]

These numbers are
\[
1,2,3,4,3,2,1
\]
in degrees
\[
0,1,2,3,4,5,6.
\]
Hence
\[
\boxed{
	H_k(X;\mathbb Z/2)\cong
	\begin{cases}
		\mathbb Z/2, & k=0,6,\\[1mm]
		(\mathbb Z/2)^2, & k=1,5,\\[1mm]
		(\mathbb Z/2)^3, & k=2,4,\\[1mm]
		(\mathbb Z/2)^4, & k=3.
\end{cases}}
\]
This agrees exactly with the computation in part (c).
\end{solution}

\begin{problem}[CP-VIII-0071 --- Naturality of the Boundary Map for Cross Products]
Let
\[
\Phi_{(X,A)}:
H^*(X,A)\otimes H^*(Y)
\longrightarrow
H^*(X\times Y,A\times Y)
\]
be defined by
\[
\Phi_{(X,A)}(\alpha\otimes \beta)=\alpha\times \beta.
\]
Similarly, define
\[
\Phi_A:
H^*(A)\otimes H^*(Y)
\longrightarrow
H^*(A\times Y)
\]
by
\[
\Phi_A(\alpha\otimes \beta)=\alpha\times \beta.
\]
Show that the cross product commutes with the connecting homomorphism in the long exact sequence of a pair. In the usual convention, this means
\[
\delta_{(X\times Y,A\times Y)}\circ \Phi_A
=
\Phi_{(X,A)}\circ(\delta_{(X,A)}\otimes 1).
\]
Equivalently, with the boundary notation used in the problem, this is written as
\[
\delta\Phi_{(X,A)}=\Phi_A\delta .
\]
\end{problem}

\begin{solution}
Let
\[
\alpha\in H^p(A),
\qquad
\beta\in H^q(Y).
\]
We want to show that
\[
\delta(\alpha\times \beta)
=
(\delta\alpha)\times \beta.
\]

Choose cocycle representatives
\[
a\in C^p(A),
\qquad
b\in C^q(Y),
\]
with
\[
da=0,
\qquad
db=0.
\]
To compute the connecting homomorphism
\[
\delta_{(X,A)}:H^p(A)\to H^{p+1}(X,A),
\]
extend \(a\) to a cochain
\[
\widetilde a\in C^p(X).
\]
Then the relative cochain
\[
d\widetilde a
\]
vanishes on \(A\), because \(a\) was a cocycle on \(A\). Hence
\[
[d\widetilde a]\in H^{p+1}(X,A),
\]
and by definition
\[
\delta[\alpha]=[d\widetilde a].
\]

Now consider the product pair
\[
(X\times Y,A\times Y).
\]
The cochain \(a\times b\) on \(A\times Y\) is extended by
\[
\widetilde a\times b
\]
on \(X\times Y\). Therefore
\[
\delta(\alpha\times \beta)
=
[d(\widetilde a\times b)].
\]

Using the Leibniz rule for the coboundary of a cross product,
\[
d(\widetilde a\times b)
=
d\widetilde a\times b
+
(-1)^p\widetilde a\times db.
\]
But \(b\) is a cocycle, so \(db=0\). Therefore
\[
d(\widetilde a\times b)
=
d\widetilde a\times b.
\]
Hence
\[
\delta(\alpha\times \beta)
=
[d\widetilde a\times b]
=
[d\widetilde a]\times [b]
=
(\delta\alpha)\times \beta.
\]

Thus the diagram



\[
\begin{array}{ccc}
H^p(A)\otimes H^q(Y)
&\xrightarrow{\ \Phi_A\ }&
H^{p+q}(A\times Y)\\
\big\downarrow{\scriptstyle \delta\otimes 1}
&&
\big\downarrow{\scriptstyle \delta}\\
H^{p+1}(X,A)\otimes H^q(Y)
&\xrightarrow{\ \Phi_{(X,A)}\ }&
H^{p+q+1}(X\times Y,A\times Y)
\end{array}
\]



commutes.

Therefore
\[
\boxed{
	\delta(\alpha\times\beta)
	=
	(\delta\alpha)\times\beta
}
\]
and so
\[
\boxed{
	\delta\Phi_{(X,A)}
	=
	\Phi_A\delta
}
\]
in the notation of the problem.
\end{solution}

\section{Cohomology, Bundles and Manifolds}

\begin{problem}[CP-VIII-0072 --- Basic properties of the Euler characteristic]
If \(X\) is a finite cell complex, define the Euler characteristic by
\[
\chi(X)=\sum_k (-1)^k n_k,
\]
where \(n_k\) is the number of \(k\)-cells in \(X\). Prove the following properties of \(\chi\):
\begin{enumerate}
	\item[(a)] \(\chi(X)=P_{\mathbb F}(X)\big|_{t=-1}\), where \(\mathbb F\) is any field.
	\item[(b)] \(\chi(X\times Y)=\chi(X)\chi(Y)\).
	\item[(c)] If \(A\) and \(B\) are subcomplexes of \(X\), then
	\[
	\chi(A\cup B)=\chi(A)+\chi(B)-\chi(A\cap B).
	\]
\end{enumerate}
\end{problem}

\begin{solution}

\medskip
\noindent\textbf{Part (a): Euler characteristic and the Poincaré polynomial}\par
\smallskip


Let \(\mathbb F\) be a field. The Poincaré polynomial of \(X\) over \(\mathbb F\) is
\[
P_{\mathbb F}(X)
=
\sum_k \beta_k^{\mathbb F}(X)t^k,
\]
where
\[
\beta_k^{\mathbb F}(X)
=
\dim_{\mathbb F} H_k(X;\mathbb F)
\]
is the \(k\)-th Betti number of \(X\) with coefficients in \(\mathbb F\).

We want to prove that
\[
\chi(X)=P_{\mathbb F}(X)\big|_{t=-1}.
\]

The cellular chain complex of \(X\) with coefficients in \(\mathbb F\) is
\[
\cdots
\longrightarrow
C_k(X;\mathbb F)
\xrightarrow{\partial_k}
C_{k-1}(X;\mathbb F)
\longrightarrow
\cdots.
\]
Because \(X\) has \(n_k\) many \(k\)-cells, we have
\[
C_k(X;\mathbb F)\cong \mathbb F^{n_k}.
\]
Therefore
\[
\dim_{\mathbb F} C_k(X;\mathbb F)=n_k.
\]

For a finite-dimensional chain complex over a field, the alternating sum of the dimensions of the chain groups equals the alternating sum of the dimensions of the homology groups:
\[
\sum_k (-1)^k \dim_{\mathbb F} C_k(X;\mathbb F)
=
\sum_k (-1)^k \dim_{\mathbb F} H_k(X;\mathbb F).
\]

For completeness, we recall why this is true. Let
\[
Z_k=\ker \partial_k
\]
be the group of \(k\)-cycles, and let
\[
B_k=\operatorname{im}\partial_{k+1}
\]
be the group of \(k\)-boundaries. Then
\[
H_k(X;\mathbb F)=Z_k/B_k.
\]
Since all groups are finite-dimensional vector spaces over \(\mathbb F\),
\[
\dim Z_k=\dim H_k+\dim B_k.
\]
Also, by rank-nullity applied to
\[
\partial_k:C_k\to C_{k-1},
\]
we have
\[
\dim C_k=\dim Z_k+\dim B_{k-1}.
\]
Therefore
\[
\dim C_k
=
\dim H_k+\dim B_k+\dim B_{k-1}.
\]
Taking alternating sums gives
\[
\sum_k (-1)^k \dim C_k
=
\sum_k (-1)^k \dim H_k
+
\sum_k (-1)^k \dim B_k
+
\sum_k (-1)^k \dim B_{k-1}.
\]
The last two sums cancel because
\[
\sum_k (-1)^k \dim B_{k-1}
=
-\sum_k (-1)^k \dim B_k.
\]
Hence
\[
\sum_k (-1)^k \dim C_k
=
\sum_k (-1)^k \dim H_k.
\]

Now
\[
\chi(X)
=
\sum_k (-1)^k n_k
=
\sum_k (-1)^k \dim_{\mathbb F} C_k(X;\mathbb F).
\]
Therefore
\[
\chi(X)
=
\sum_k (-1)^k \dim_{\mathbb F} H_k(X;\mathbb F)
=
\sum_k (-1)^k \beta_k^{\mathbb F}(X).
\]
But
\[
P_{\mathbb F}(X)\big|_{t=-1}
=
\sum_k \beta_k^{\mathbb F}(X)(-1)^k.
\]
Hence
\[
\boxed{
	\chi(X)=P_{\mathbb F}(X)\big|_{t=-1}.}
\]


\medskip
\noindent\textbf{Part (b): Multiplicativity under products}\par
\smallskip


We prove that
\[
\chi(X\times Y)=\chi(X)\chi(Y).
\]

Let \(X\) and \(Y\) be finite cell complexes. Suppose \(X\) has \(n_i\) cells of dimension \(i\), and \(Y\) has \(m_j\) cells of dimension \(j\).

The product CW-structure on \(X\times Y\) has one \((i+j)\)-cell for each pair consisting of an \(i\)-cell of \(X\) and a \(j\)-cell of \(Y\). Therefore, the number of \(k\)-cells of \(X\times Y\) is
\[
\sum_{i+j=k} n_i m_j.
\]
Thus
\[
\chi(X\times Y)
=
\sum_k (-1)^k
\left(
\sum_{i+j=k} n_i m_j
\right).
\]
Changing the order of summation gives
\[
\chi(X\times Y)
=
\sum_{i,j} (-1)^{i+j}n_i m_j.
\]
Since
\[
(-1)^{i+j}=(-1)^i(-1)^j,
\]
we get
\[
\chi(X\times Y)
=
\left(\sum_i (-1)^i n_i\right)
\left(\sum_j (-1)^j m_j\right).
\]
Therefore
\[
\chi(X\times Y)=\chi(X)\chi(Y).
\]

Hence
\[
\boxed{\chi(X\times Y)=\chi(X)\chi(Y).}
\]

Equivalently, one can prove this using part (a) and the Künneth formula over a field. Indeed, over a field \(\mathbb F\),
\[
H_n(X\times Y;\mathbb F)
\cong
\bigoplus_{i+j=n}
H_i(X;\mathbb F)\otimes_{\mathbb F}H_j(Y;\mathbb F).
\]
Taking dimensions gives
\[
\beta_n^{\mathbb F}(X\times Y)
=
\sum_{i+j=n}
\beta_i^{\mathbb F}(X)\beta_j^{\mathbb F}(Y).
\]
Hence
\[
P_{\mathbb F}(X\times Y)
=
P_{\mathbb F}(X)P_{\mathbb F}(Y).
\]
Evaluating at \(t=-1\), we obtain
\[
\chi(X\times Y)
=
P_{\mathbb F}(X\times Y)\big|_{t=-1}
=
P_{\mathbb F}(X)\big|_{t=-1}
P_{\mathbb F}(Y)\big|_{t=-1}.
\]
Using part (a), this becomes
\[
\chi(X\times Y)=\chi(X)\chi(Y).
\]


\medskip
\noindent\textbf{Part (c): Inclusion-exclusion for subcomplexes}\par
\smallskip


Let \(A\) and \(B\) be subcomplexes of \(X\). We want to prove that
\[
\chi(A\cup B)=\chi(A)+\chi(B)-\chi(A\cap B).
\]

Because \(A\) and \(B\) are subcomplexes, their union \(A\cup B\) and their intersection \(A\cap B\) are also cell complexes built from cells of \(X\).

For each \(k\), let
\[
n_k(A)
\]
denote the number of \(k\)-cells of \(A\), and similarly define
\[
n_k(B),
\qquad
n_k(A\cup B),
\qquad
n_k(A\cap B).
\]

For each fixed dimension \(k\), the ordinary inclusion-exclusion principle for finite sets gives
\[
n_k(A\cup B)
=
n_k(A)+n_k(B)-n_k(A\cap B).
\]

Now take alternating sums over all \(k\):
\[
\chi(A\cup B)
=
\sum_k (-1)^k n_k(A\cup B).
\]
Substituting the formula above gives
\[
\chi(A\cup B)
=
\sum_k (-1)^k
\left(
n_k(A)+n_k(B)-n_k(A\cap B)
\right).
\]
Therefore
\[
\chi(A\cup B)
=
\sum_k (-1)^k n_k(A)
+
\sum_k (-1)^k n_k(B)
-
\sum_k (-1)^k n_k(A\cap B).
\]
Hence
\[
\chi(A\cup B)
=
\chi(A)+\chi(B)-\chi(A\cap B).
\]

Thus
\[
\boxed{
	\chi(A\cup B)=\chi(A)+\chi(B)-\chi(A\cap B).}
\]


\medskip
\noindent\textbf{Summary}\par
\smallskip


For
\[
X=L_4^3\times \mathbb{RP}^3,
\]
we have
\[
\boxed{
	H_k(X)\cong
	\begin{cases}
		\mathbb Z, & k=0,6,\\[1mm]
		\mathbb Z/4\oplus \mathbb Z/2, & k=1,4,\\[1mm]
		\mathbb Z/2, & k=2,\\[1mm]
		\mathbb Z^2\oplus \mathbb Z/2, & k=3,\\[1mm]
		0, & k=5.
\end{cases}}
\]

Also,
\[
\boxed{
	H^k(X)\cong
	\begin{cases}
		\mathbb Z, & k=0,6,\\[1mm]
		0, & k=1,\\[1mm]
		\mathbb Z/4\oplus \mathbb Z/2, & k=2,5,\\[1mm]
		\mathbb Z^2\oplus \mathbb Z/2, & k=3,\\[1mm]
		\mathbb Z/2, & k=4.
\end{cases}}
\]

Finally,
\[
\boxed{
	H_k(X;\mathbb Z/2)\cong
	\begin{cases}
		\mathbb Z/2, & k=0,6,\\[1mm]
		(\mathbb Z/2)^2, & k=1,5,\\[1mm]
		(\mathbb Z/2)^3, & k=2,4,\\[1mm]
		(\mathbb Z/2)^4, & k=3.
\end{cases}}
\]

For every finite cell complex \(X\),
\[
\boxed{
	\chi(X)=P_{\mathbb F}(X)\big|_{t=-1}}
\]
for any field \(\mathbb F\). Moreover,
\[
\boxed{
	\chi(X\times Y)=\chi(X)\chi(Y)}
\]
and, for subcomplexes \(A,B\subseteq X\),
\[
\boxed{
	\chi(A\cup B)=\chi(A)+\chi(B)-\chi(A\cap B).}
\]
\end{solution}

\begin{problem}[CP-VIII-0073 --- Diagonal Maps and Cup Products]
If \(X\) is a space, let
\[
\Delta_X:X\to X\times X
\]
be the diagonal map
\[
\Delta_X(x)=(x,x).
\]
Compute
\[
\Delta_{S^2}^*:H^*(S^2\times S^2)\to H^*(S^2)
\]
and
\[
\Delta_{T^2}^*:H^*(T^2\times T^2)\to H^*(T^2).
\]
\end{problem}

\begin{solution}

\medskip
\noindent\textbf{Main idea}\par
\smallskip


The diagonal map turns external products into cup products. Namely, if
\[
\alpha\in H^p(X),
\qquad
\beta\in H^q(X),
\]
then
\[
\Delta_X^*(\alpha\times \beta)
=
\alpha\smile \beta.
\]

Indeed, the cup product is defined by
\[
\alpha\smile \beta
=
\Delta_X^*(\alpha\times \beta).
\]


\medskip
\noindent\textbf{Part I. The case of \(S^2\)}\par
\smallskip


Let
\[
u\in H^2(S^2;\mathbb Z)
\]
be a generator. Then
\[
H^*(S^2;\mathbb Z)
\cong
\mathbb Z[u]/(u^2),
\qquad
|u|=2.
\]
The relation \(u^2=0\) holds because
\[
H^4(S^2;\mathbb Z)=0.
\]

By the Künneth formula,
\[
H^*(S^2\times S^2;\mathbb Z)
\cong
H^*(S^2;\mathbb Z)\otimes H^*(S^2;\mathbb Z).
\]
Let
\[
u_1=\pi_1^*(u),
\qquad
u_2=\pi_2^*(u),
\]
where
\[
\pi_1,\pi_2:S^2\times S^2\to S^2
\]
are the two projections.

Then
\[
H^*(S^2\times S^2;\mathbb Z)
\cong
\mathbb Z[u_1,u_2]/(u_1^2,u_2^2),
\qquad
|u_1|=|u_2|=2.
\]

The diagonal map satisfies
\[
\pi_1\circ \Delta_{S^2}=\operatorname{id}_{S^2},
\qquad
\pi_2\circ \Delta_{S^2}=\operatorname{id}_{S^2}.
\]
Therefore
\[
\Delta_{S^2}^*(u_1)=u,
\qquad
\Delta_{S^2}^*(u_2)=u.
\]

Since \(\Delta_{S^2}^*\) is a ring homomorphism,
\[
\Delta_{S^2}^*(u_1u_2)
=
\Delta_{S^2}^*(u_1)\Delta_{S^2}^*(u_2)
=
u^2
=
0.
\]

Hence
\[
\boxed{
	\Delta_{S^2}^*(1)=1
}
\]
\[
\boxed{
	\Delta_{S^2}^*(u_1)=u,
	\qquad
	\Delta_{S^2}^*(u_2)=u
}
\]
and
\[
\boxed{
	\Delta_{S^2}^*(u_1u_2)=0.
}
\]

Equivalently, in degrees:
\[
H^0(S^2\times S^2)\cong \mathbb Z
\longrightarrow
H^0(S^2)\cong \mathbb Z
\]
is the identity,
\[
H^2(S^2\times S^2)\cong \mathbb Z u_1\oplus \mathbb Z u_2
\longrightarrow
H^2(S^2)\cong \mathbb Z u
\]
is given by
\[
au_1+bu_2\longmapsto (a+b)u,
\]
and
\[
H^4(S^2\times S^2)\cong \mathbb Z(u_1u_2)
\longrightarrow
H^4(S^2)=0
\]
is the zero map.


\medskip
\noindent\textbf{Part II. The case of \(T^2\)}\par
\smallskip


Let
\[
T^2=S^1\times S^1.
\]
The cohomology ring of the torus is
\[
H^*(T^2;\mathbb Z)
\cong
\Lambda_{\mathbb Z}(\alpha,\beta),
\]
where
\[
|\alpha|=|\beta|=1.
\]
Thus
\[
H^0(T^2)\cong \mathbb Z,
\]
\[
H^1(T^2)\cong \mathbb Z\alpha\oplus \mathbb Z\beta,
\]
and
\[
H^2(T^2)\cong \mathbb Z(\alpha\beta).
\]
Write
\[
\omega=\alpha\beta.
\]
Because the cup product is graded-commutative,
\[
\beta\alpha=-\alpha\beta=-\omega.
\]

Now consider
\[
T^2\times T^2.
\]
Let
\[
\alpha_1=\pi_1^*(\alpha),
\qquad
\beta_1=\pi_1^*(\beta),
\]
and
\[
\alpha_2=\pi_2^*(\alpha),
\qquad
\beta_2=\pi_2^*(\beta),
\]
where
\[
\pi_1,\pi_2:T^2\times T^2\to T^2
\]
are the two projections.

Since
\[
\pi_1\circ \Delta_{T^2}
=
\operatorname{id}_{T^2},
\qquad
\pi_2\circ \Delta_{T^2}
=
\operatorname{id}_{T^2},
\]
we get
\[
\Delta_{T^2}^*(\alpha_1)=\alpha,
\qquad
\Delta_{T^2}^*(\beta_1)=\beta,
\]
and
\[
\Delta_{T^2}^*(\alpha_2)=\alpha,
\qquad
\Delta_{T^2}^*(\beta_2)=\beta.
\]

Therefore \(\Delta_{T^2}^*\) is determined by
\[
\boxed{
	\alpha_1\mapsto \alpha,
	\qquad
	\beta_1\mapsto \beta,
	\qquad
	\alpha_2\mapsto \alpha,
	\qquad
	\beta_2\mapsto \beta.
}
\]

Since \(\Delta_{T^2}^*\) is a ring homomorphism, we can compute all products.

For degree \(2\), we obtain:
\[
\Delta_{T^2}^*(\alpha_1\beta_1)
=
\alpha\beta
=
\omega,
\]
\[
\Delta_{T^2}^*(\alpha_2\beta_2)
=
\alpha\beta
=
\omega,
\]
\[
\Delta_{T^2}^*(\alpha_1\alpha_2)
=
\alpha\alpha
=
0,
\]
\[
\Delta_{T^2}^*(\beta_1\beta_2)
=
\beta\beta
=
0,
\]
\[
\Delta_{T^2}^*(\alpha_1\beta_2)
=
\alpha\beta
=
\omega,
\]
and
\[
\Delta_{T^2}^*(\beta_1\alpha_2)
=
\beta\alpha
=
-\omega.
\]

Thus
\[
\boxed{
	\Delta_{T^2}^*(\alpha_1\beta_1)=\omega
}
\]
\[
\boxed{
	\Delta_{T^2}^*(\alpha_2\beta_2)=\omega
}
\]
\[
\boxed{
	\Delta_{T^2}^*(\alpha_1\alpha_2)=0
}
\]
\[
\boxed{
	\Delta_{T^2}^*(\beta_1\beta_2)=0
}
\]
\[
\boxed{
	\Delta_{T^2}^*(\alpha_1\beta_2)=\omega
}
\]
and
\[
\boxed{
	\Delta_{T^2}^*(\beta_1\alpha_2)=-\omega.
}
\]

All classes of degree \(3\) and \(4\) in \(H^*(T^2\times T^2)\) map to zero because
\[
H^3(T^2)=0
\]
and
\[
H^4(T^2)=0.
\]

Therefore the diagonal map sends the generators from both factors to the same generators on \(T^2\), and then multiplies them using the ordinary cup product on \(T^2\).
\end{solution}

\begin{problem}[CP-VIII-0074 --- Relative Cup Products and Open Covers by Contractible Sets]
Let \(U,V\subset X\) be open sets. If
\[
x\in H^*(X,U)
\]
and
\[
y\in H^*(X,V),
\]
show that
\[
x\smile y\in H^*(X,U\cup V).
\]

Using this, show that if \(X\) has a covering by \(n\) contractible open subsets, then
\[
a_1\smile a_2\smile \cdots \smile a_n=0
\]
whenever
\[
a_1,\dots,a_n\in H^*(X)
\]
have positive degree.

Deduce that if \(T^k\) can be covered by \(n\) contractible open subsets, then
\[
n>k.
\]
Finally, find a covering of \(T^2\) by three contractible open subsets.
\end{problem}

\begin{solution}

\medskip
\noindent\textbf{Part I. The relative cup product}\par
\smallskip


Let
\[
x\in H^p(X,U),
\qquad
y\in H^q(X,V).
\]
The external product gives
\[
x\times y
\in
H^{p+q}(X\times X, U\times X\cup X\times V).
\]

Now consider the diagonal map
\[
\Delta:X\to X\times X,
\qquad
\Delta(z)=(z,z).
\]
If
\[
z\in U\cup V,
\]
then either \(z\in U\) or \(z\in V\). Hence
\[
(z,z)\in U\times X
\]
or
\[
(z,z)\in X\times V.
\]
Therefore
\[
\Delta(U\cup V)
\subset
U\times X\cup X\times V.
\]

Thus \(\Delta\) is a map of pairs
\[
\Delta:
(X,U\cup V)
\longrightarrow
(X\times X, U\times X\cup X\times V).
\]

Pulling back the external product by \(\Delta\), we obtain
\[
\Delta^*(x\times y)
\in
H^{p+q}(X,U\cup V).
\]

By definition, this is the relative cup product:
\[
x\smile y
=
\Delta^*(x\times y).
\]

Therefore
\[
\boxed{
	x\smile y\in H^{p+q}(X,U\cup V).
}
\]


\medskip
\noindent\textbf{Part II. Vanishing of products from a contractible open cover}\par
\smallskip


Assume that
\[
X=U_1\cup U_2\cup \cdots \cup U_n,
\]
where each \(U_i\) is contractible and open.

Let
\[
a_i\in H^{d_i}(X)
\]
with
\[
d_i>0.
\]
Since \(U_i\) is contractible,
\[
H^{d_i}(U_i)=0
\]
for \(d_i>0\). Therefore the restriction of \(a_i\) to \(U_i\) is zero:
\[
a_i|_{U_i}=0.
\]

From the long exact sequence of the pair \((X,U_i)\),
\[
\cdots
\to
H^{d_i}(X,U_i)
\to
H^{d_i}(X)
\to
H^{d_i}(U_i)
\to
\cdots,
\]
we conclude that \(a_i\) comes from a relative cohomology class. Thus there exists
\[
\widetilde a_i\in H^{d_i}(X,U_i)
\]
such that its image in \(H^{d_i}(X)\) is \(a_i\).

Now use the relative cup product repeatedly:
\[
\widetilde a_1\smile \widetilde a_2
\in
H^*(X,U_1\cup U_2),
\]
then
\[
\widetilde a_1\smile \widetilde a_2\smile \widetilde a_3
\in
H^*(X,U_1\cup U_2\cup U_3),
\]
and so on. Finally,
\[
\widetilde a_1\smile \cdots \smile \widetilde a_n
\in
H^*(X,U_1\cup\cdots\cup U_n).
\]
But
\[
U_1\cup\cdots\cup U_n=X.
\]
Hence
\[
H^*(X,U_1\cup\cdots\cup U_n)
=
H^*(X,X)
=
0.
\]
Therefore
\[
\widetilde a_1\smile\cdots\smile \widetilde a_n=0.
\]

Passing to absolute cohomology gives
\[
a_1\smile a_2\smile \cdots \smile a_n=0.
\]

Thus
\[
\boxed{
	a_1\smile a_2\smile\cdots\smile a_n=0
}
\]
for all positive-degree classes \(a_1,\dots,a_n\in H^*(X)\).


\medskip
\noindent\textbf{Part III. Application to the torus \(T^k\)}\par
\smallskip


The cohomology ring of the \(k\)-torus is
\[
H^*(T^k;\mathbb Z)
\cong
\Lambda_{\mathbb Z}(\alpha_1,\dots,\alpha_k),
\]
where
\[
|\alpha_i|=1.
\]
The product
\[
\alpha_1\smile \alpha_2\smile\cdots\smile \alpha_k
\]
is the orientation class of \(T^k\). In particular,
\[
\alpha_1\smile \cdots\smile \alpha_k\neq 0.
\]

Moreover, for every \(r\leq k\),
\[
\alpha_1\smile\cdots\smile\alpha_r\neq 0.
\]

Suppose \(T^k\) could be covered by \(n\) contractible open subsets with
\[
n\leq k.
\]
Then by the result above, every product of \(n\) positive-degree cohomology classes would have to vanish.

But
\[
\alpha_1\smile\cdots\smile\alpha_n\neq 0.
\]
This is a contradiction.

Therefore
\[
\boxed{
	n>k.
}
\]

So \(T^k\) cannot be covered by \(k\) or fewer contractible open subsets.


\medskip
\noindent\textbf{Part IV. A covering of \(T^2\) by three contractible open subsets}\par
\smallskip


We represent \(T^2\) as the quotient of a square by identifying opposite sides:

\[
T^2
=
[0,1]\times[0,1]/\sim,
\]
where
\[
(0,y)\sim (1,y)
\]
and
\[
(x,0)\sim (x,1).
\]

The square model has one \(0\)-cell, two \(1\)-cells, and one \(2\)-cell. We construct three open subsets.

\begin{enumerate}
	\item \(U_0\): a small open neighborhood of the unique vertex of the quotient. In the square picture, this appears as four small neighborhoods of the four corners. After identifying the corners, these four pieces become one open disk. Hence \(U_0\) is contractible.

	\item \(U_2\): a large open disk contained in the interior of the square. This covers the middle of the \(2\)-cell. Since it is an open disk, \(U_2\) is contractible.

	\item \(U_1\): an open regular neighborhood of a tree joining the middle parts of the two \(1\)-cells. In the square picture, this appears as thin strips near the middle portions of the boundary edges, together with a small bridge through the interior. Since a regular neighborhood of a tree in a surface is an open disk, \(U_1\) is contractible.
\end{enumerate}

Choosing the neighborhoods large enough, these three open sets cover the whole torus:
\[
T^2=U_0\cup U_1\cup U_2.
\]

The following picture is schematic. Opposite edges of the square are identified.



\[
\boxed{
\begin{array}{c}
\text{Square model of }T^2:\ \text{opposite }a\text{-edges and }b\text{-edges are identified}\\[1mm]
U_0:\ \text{small neighborhoods of the four corners}\\
U_1:\ \text{edge strips joined by a thin interior bridge}\\
U_2:\ \text{a large open disk in the interior}\\[1mm]
T^2=U_0\cup U_1\cup U_2
\end{array}}
\]



After quotienting the square to form the torus:

\[
\boxed{
	U_0\simeq \text{open disk},
	\qquad
	U_1\simeq \text{open disk},
	\qquad
	U_2\simeq \text{open disk}.
}
\]

Hence each \(U_i\) is contractible, and
\[
\boxed{
	T^2=U_0\cup U_1\cup U_2.
}
\]

This agrees with the lower bound from cup products. Since
\[
H^*(T^2)\cong \Lambda(\alpha,\beta)
\]
and
\[
\alpha\smile\beta\neq 0,
\]
the torus cannot be covered by two contractible open sets. But the construction above shows that it can be covered by three. Therefore the minimal number is exactly three:
\[
\boxed{
	T^2 \text{ can be covered by } 3 \text{ contractible open subsets, but not by } 2.
}
\]
\end{solution}

\begin{problem}[CP-VIII-0075 --- Maps between closed orientable surfaces of different genus]
Let \(\Sigma_g\) be the closed orientable surface of genus \(g\). Suppose

\[
\varphi:\Sigma_g\longrightarrow \Sigma_h
\]

is a continuous map, where

\[
g<h.
\]

Show that the induced map

\[
\varphi_*:H_2(\Sigma_g)\longrightarrow H_2(\Sigma_h)
\]

is the zero map.

The hint suggests considering

\[
\varphi^*:H^1(\Sigma_h)\longrightarrow H^1(\Sigma_g).
\]
\end{problem}

\begin{solution}

\medskip
\noindent\textbf{Geometric idea}\par
\smallskip


The group

\[
H_2(\Sigma_g;\mathbb Z)\cong \mathbb Z
\]

is generated by the fundamental class

\[
[\Sigma_g].
\]

Therefore the map \(\varphi_*\) is multiplication by an integer \(d\):

\[
\varphi_*[\Sigma_g]=d[\Sigma_h].
\]

This integer \(d\) is the degree of \(\varphi\), written

\[
d=\deg(\varphi).
\]

Thus the problem is equivalent to proving

\[
\deg(\varphi)=0.
\]

The key idea is:

\[
\boxed{
	\text{If }\deg(\varphi)\neq 0,\text{ then }\varphi^*:H^1(\Sigma_h;\mathbb Q)\to H^1(\Sigma_g;\mathbb Q)
	\text{ must be injective.}
}
\]

But

\[
\dim_{\mathbb Q}H^1(\Sigma_h;\mathbb Q)=2h
\]

and

\[
\dim_{\mathbb Q}H^1(\Sigma_g;\mathbb Q)=2g.
\]

Since \(g<h\), we have

\[
2g<2h.
\]

So an injective linear map

\[
\mathbb Q^{2h}\longrightarrow \mathbb Q^{2g}
\]

cannot exist. Therefore the degree must be zero.


\medskip
\noindent\textbf{Dimension diagram}\par
\smallskip


The contradiction can be visualized as follows:



\[
\boxed{
H^1(\Sigma_h;\mathbb Q)\cong\mathbb Q^{2h}
\ \xrightarrow[\text{would be injective if }\deg\varphi\ne0]{\ \varphi^*\ }\
H^1(\Sigma_g;\mathbb Q)\cong\mathbb Q^{2g}}
\]
\[
2h>2g,
\qquad\text{so such an injection cannot exist.}
\]



\medskip
\noindent\textbf{Solution}\par
\smallskip


Choose orientations on \(\Sigma_g\) and \(\Sigma_h\). Since both are closed connected orientable surfaces,

\[
H_2(\Sigma_g;\mathbb Z)\cong \mathbb Z,
\qquad
H_2(\Sigma_h;\mathbb Z)\cong \mathbb Z.
\]

Let

\[
[\Sigma_g]\in H_2(\Sigma_g;\mathbb Z),
\qquad
[\Sigma_h]\in H_2(\Sigma_h;\mathbb Z)
\]

be the fundamental classes.

Then there exists an integer \(d\in \mathbb Z\) such that

\[
\varphi_*[\Sigma_g]=d[\Sigma_h].
\]

We need to show that

\[
d=0.
\]

Suppose, for contradiction, that

\[
d\neq 0.
\]

We work with rational coefficients. Consider the induced map

\[
\varphi^*:H^1(\Sigma_h;\mathbb Q)\longrightarrow H^1(\Sigma_g;\mathbb Q).
\]

We claim that \(\varphi^*\) is injective.

Let

\[
0\neq \alpha\in H^1(\Sigma_h;\mathbb Q).
\]

By Poincaré duality on the closed orientable surface \(\Sigma_h\), the pairing

\[
H^1(\Sigma_h;\mathbb Q)\times H^1(\Sigma_h;\mathbb Q)
\longrightarrow
\mathbb Q
\]

given by

\[
(\alpha,\beta)
\longmapsto
\langle \alpha\smile \beta,[\Sigma_h]\rangle
\]

is nondegenerate. Therefore there exists

\[
\beta\in H^1(\Sigma_h;\mathbb Q)
\]

such that

\[
\langle \alpha\smile \beta,[\Sigma_h]\rangle\neq 0.
\]

Now use naturality of the cup product and the definition of degree:

\[
\begin{aligned}
	\left\langle \varphi^*\alpha\smile \varphi^*\beta,[\Sigma_g]\right\rangle
	&=
	\left\langle \varphi^*(\alpha\smile \beta),[\Sigma_g]\right\rangle \\
	&=
	\left\langle \alpha\smile \beta,\varphi_*[\Sigma_g]\right\rangle \\
	&=
	\left\langle \alpha\smile \beta,d[\Sigma_h]\right\rangle \\
	&=
	d\left\langle \alpha\smile \beta,[\Sigma_h]\right\rangle.
\end{aligned}
\]

Since

\[
d\neq 0
\]

and

\[
\left\langle \alpha\smile \beta,[\Sigma_h]\right\rangle\neq 0,
\]

we obtain

\[
\left\langle \varphi^*\alpha\smile \varphi^*\beta,[\Sigma_g]\right\rangle\neq 0.
\]

In particular,

\[
\varphi^*\alpha\neq 0.
\]

Thus every nonzero class in \(H^1(\Sigma_h;\mathbb Q)\) has nonzero pullback. Therefore

\[
\varphi^*:H^1(\Sigma_h;\mathbb Q)\to H^1(\Sigma_g;\mathbb Q)
\]

is injective.

But

\[
\dim_{\mathbb Q}H^1(\Sigma_h;\mathbb Q)=2h
\]

and

\[
\dim_{\mathbb Q}H^1(\Sigma_g;\mathbb Q)=2g.
\]

Since \(g<h\),

\[
2g<2h.
\]

So an injective linear map

\[
\mathbb Q^{2h}\longrightarrow \mathbb Q^{2g}
\]

cannot exist.

This contradiction shows that our assumption was false. Hence

\[
d=0.
\]

Therefore

\[
\varphi_*[\Sigma_g]=0.
\]

Since \(H_2(\Sigma_g)\) is generated by \([\Sigma_g]\), it follows that

\[
\boxed{
	\varphi_*:H_2(\Sigma_g)\longrightarrow H_2(\Sigma_h)
	\text{ is the zero map.}
}
\]


\medskip
\noindent\textbf{Summary chart}\par
\smallskip


\[
\begin{array}{c|c|c}
	\text{Surface} & H^1(-;\mathbb Q) & \dim_{\mathbb Q}H^1(-;\mathbb Q) \\
	\hline
	\Sigma_g & \mathbb Q^{2g} & 2g \\
	\Sigma_h & \mathbb Q^{2h} & 2h
\end{array}
\]

Since \(g<h\), one has

\[
2g<2h.
\]

Thus a nonzero-degree map would force an impossible injection

\[
H^1(\Sigma_h;\mathbb Q)\hookrightarrow H^1(\Sigma_g;\mathbb Q).
\]

Therefore the degree is zero.
\end{solution}

\begin{problem}[CP-VIII-0076 --- The Hopf invariant]
Given a map

\[
\phi:S^{2n-1}\longrightarrow S^n,
\]

let

\[
X_\phi=S^n\cup_\phi D^{2n}.
\]

The cohomology of \(X_\phi\) has the form

\[
H^*(X_\phi)=\langle 1,x_n,x_{2n}\rangle,
\]

where

\[
x_i\in H^i(X_\phi).
\]

Thus

\[
x_n^2=H(\phi)x_{2n}
\]

for some integer

\[
H(\phi)\in \mathbb Z.
\]

\begin{enumerate}
	\item[(a)] Show that the map
	\[
	[\phi]\longmapsto H(\phi)
	\]
	defines a homomorphism
	\[
	H:\pi_{2n-1}(S^n)\longrightarrow \mathbb Z.
	\]
	The integer \(H(\phi)\) is called the Hopf invariant.

	\item[(b)] By considering a cell decomposition of \(S^n\times S^n\), show that \(H\) is nontrivial for every even \(n\). Deduce that
	\[
	\pi_{4m-1}(S^{2m})
	\]
	is infinite for all \(m>0\).
\end{enumerate}
\end{problem}

\begin{solution}

\medskip
\noindent\textbf{The mapping cone \(X_\phi\)}\par
\smallskip


The space

\[
X_\phi=S^n\cup_\phi D^{2n}
\]

is obtained by attaching a \(2n\)-cell to \(S^n\) using the attaching map

\[
\phi:S^{2n-1}\to S^n.
\]

Thus \(X_\phi\) has one cell in dimension \(0\), one cell in dimension \(n\), and one cell in dimension \(2n\).

Schematically:



\[
\boxed{
S^n
\ \xleftarrow{\ \phi\ }\ 
S^{2n-1}=\partial D^{2n}
\ \hookrightarrow\
D^{2n}
\quad\leadsto\quad
X_\phi=S^n\cup_\phi D^{2n}}
\]
\[
\text{cells of }X_\phi:\qquad e^0,\ e^n,\ e^{2n}.
\]



The cellular cochain groups are

\[
C^0(X_\phi)\cong \mathbb Z,
\qquad
C^n(X_\phi)\cong \mathbb Z,
\qquad
C^{2n}(X_\phi)\cong \mathbb Z,
\]

and there are no cells in the dimensions between \(n\) and \(2n\). Therefore

\[
H^0(X_\phi)\cong \mathbb Z,
\qquad
H^n(X_\phi)\cong \mathbb Z,
\qquad
H^{2n}(X_\phi)\cong \mathbb Z.
\]

Let

\[
x_n\in H^n(X_\phi),
\qquad
x_{2n}\in H^{2n}(X_\phi)
\]

be generators. Then the product

\[
x_n^2=x_n\smile x_n
\]

lies in

\[
H^{2n}(X_\phi)\cong \mathbb Z.
\]

Hence there is a unique integer \(H(\phi)\) such that

\[
\boxed{
	x_n^2=H(\phi)x_{2n}.
}
\]

This integer is the Hopf invariant of \(\phi\).


\medskip
\noindent\textbf{Part (a). The Hopf invariant is a homomorphism}\par
\smallskip


\medskip
\noindent\textbf{Step 1. Dependence only on the homotopy class}\par
\smallskip


If

\[
\phi\simeq \psi:S^{2n-1}\to S^n,
\]

then the corresponding cell complexes

\[
X_\phi=S^n\cup_\phi D^{2n}
\]

and

\[
X_\psi=S^n\cup_\psi D^{2n}
\]

are homotopy equivalent. The equivalence preserves the \(n\)-cell and the \(2n\)-cell up to orientation choices. Therefore the coefficient of \(x_{2n}\) in the square \(x_n^2\) is unchanged up to the corresponding orientation convention.

Thus

\[
H(\phi)=H(\psi).
\]

So the Hopf invariant depends only on the homotopy class

\[
[\phi]\in \pi_{2n-1}(S^n).
\]


\medskip
\noindent\textbf{Step 2. Odd-dimensional case}\par
\smallskip


If \(n\) is odd, then graded commutativity gives

\[
x_n^2
=
(-1)^{n^2}x_n^2
=
(-1)^n x_n^2
=
-x_n^2.
\]

Hence

\[
2x_n^2=0.
\]

But

\[
H^{2n}(X_\phi)\cong \mathbb Z
\]

has no torsion. Therefore

\[
x_n^2=0.
\]

Thus

\[
H(\phi)=0
\]

for every \(\phi\) when \(n\) is odd.

So in odd dimensions, the Hopf invariant is the zero homomorphism:

\[
\boxed{
	H:\pi_{2n-1}(S^n)\to \mathbb Z
	\text{ is identically zero if }n\text{ is odd.}
}
\]


\medskip
\noindent\textbf{Step 3. Additivity}\par
\smallskip


Now suppose \(n\) is even.

Let

\[
\phi,\psi:S^{2n-1}\to S^n
\]

represent elements of

\[
\pi_{2n-1}(S^n).
\]

Their sum

\[
\phi+\psi
\]

is represented by first pinching the domain sphere to a wedge and then applying \(\phi\) on one summand and \(\psi\) on the other:

\[
S^{2n-1}
\xrightarrow{\text{pinch}}
S^{2n-1}\vee S^{2n-1}
\xrightarrow{\phi\vee \psi}
S^n.
\]

Diagrammatically:



\[
S^{2n-1}
\xrightarrow{\ \mathrm{pinch}\ }
S^{2n-1}\vee S^{2n-1}
\xrightarrow{\ \phi\vee\psi\ }
S^n .
\]



The Hopf invariant can also be described cochain-theoretically.

Let

\[
u\in C^n(S^n)
\]

be a cocycle representing the generator of

\[
H^n(S^n;\mathbb Z).
\]

Since

\[
H^n(S^{2n-1};\mathbb Z)=0,
\]

the pullback

\[
\phi^*u
\]

is exact. Hence there exists a cochain

\[
v_\phi\in C^{n-1}(S^{2n-1})
\]

such that

\[
dv_\phi=\phi^*u.
\]

One cochain formula for the Hopf invariant is

\[
H(\phi)
=
\left\langle v_\phi\smile \phi^*u,\,[S^{2n-1}]\right\rangle.
\]

Equivalently,

\[
H(\phi)
=
\left\langle v_\phi\smile dv_\phi,\,[S^{2n-1}]\right\rangle.
\]

For the sum \(\phi+\psi\), the domain sphere may be viewed as two large pieces joined along a thin region. On one piece the map is \(\phi\), and on the other piece the map is \(\psi\). A primitive for \((\phi+\psi)^*u\) can be chosen to agree with \(v_\phi\) on the first piece and with \(v_\psi\) on the second piece.

Since the two pieces contribute independently, the integral splits:

\[
H(\phi+\psi)=H(\phi)+H(\psi).
\]

Therefore

\[
\boxed{
	H([\phi]+[\psi])=H([\phi])+H([\psi]).
}
\]

So

\[
\boxed{
	H:\pi_{2n-1}(S^n)\longrightarrow \mathbb Z
}
\]

is a group homomorphism.


\medskip
\noindent\textbf{Part (b). Nontriviality for even \(n\)}\par
\smallskip


We now show that for every even \(n\), there exists a map

\[
\phi:S^{2n-1}\to S^n
\]

with

\[
H(\phi)\neq 0.
\]

The map will arise from the top cell of

\[
S^n\times S^n.
\]


\medskip
\noindent\textbf{Cell decomposition of \(S^n\times S^n\)}\par
\smallskip


The sphere \(S^n\) has one \(0\)-cell and one \(n\)-cell:

\[
S^n=e^0\cup e^n.
\]

Therefore

\[
S^n\times S^n
\]

has cells in dimensions

\[
0,\quad n,\quad n,\quad 2n.
\]

More precisely,

\[
S^n\times S^n
=
e^0
\cup e^n_1
\cup e^n_2
\cup e^{2n}.
\]

The \(n\)-skeleton is

\[
S^n\vee S^n.
\]

The top cell is attached by a map

\[
S^{2n-1}\longrightarrow S^n\vee S^n.
\]

This attaching map is the Whitehead product

\[
[\iota_1,\iota_2],
\]

where

\[
\iota_1:S^n\to S^n\vee S^n
\]

and

\[
\iota_2:S^n\to S^n\vee S^n
\]

are the inclusions of the two wedge summands.

Thus

\[
S^n\times S^n
\simeq
(S^n\vee S^n)\cup_{[\iota_1,\iota_2]}e^{2n}.
\]


\medskip
\noindent\textbf{Cell-structure diagram}\par
\smallskip



\[
\boxed{
S^n\vee S^n
\ \xleftarrow{\ [\iota_1,\iota_2]\ }\ 
S^{2n-1}=\partial e^{2n}
\quad\leadsto\quad
S^n\times S^n}
\]
\[
S^n\times S^n
\simeq
(S^n\vee S^n)\cup_{[\iota_1,\iota_2]}e^{2n}.
\]



Now let

\[
\nabla:S^n\vee S^n\longrightarrow S^n
\]

be the fold map. It is defined by taking each wedge summand identically onto \(S^n\):

\[
\nabla|_{S^n_1}=\operatorname{id}_{S^n},
\qquad
\nabla|_{S^n_2}=\operatorname{id}_{S^n}.
\]

Compose the attaching map with the fold map:

\[
S^{2n-1}
\xrightarrow{[\iota_1,\iota_2]}
S^n\vee S^n
\xrightarrow{\nabla}
S^n.
\]

Define

\[
\phi=\nabla\circ[\iota_1,\iota_2].
\]

This map is the Whitehead square

\[
\phi=[\iota,\iota]\in \pi_{2n-1}(S^n).
\]


\medskip
\noindent\textbf{A map from \(S^n\times S^n\) to \(X_\phi\)}\par
\smallskip


Since the top cell of \(S^n\times S^n\) is attached by

\[
[\iota_1,\iota_2],
\]

and since \(X_\phi\) is obtained by attaching a \(2n\)-cell to \(S^n\) by

\[
\phi=\nabla\circ[\iota_1,\iota_2],
\]

the fold map extends to a cellular map

\[
q:S^n\times S^n\longrightarrow X_\phi.
\]

The diagram is:



\[
\begin{array}{ccc}
S^{2n-1}&\xrightarrow{\ [\iota_1,\iota_2]\ }&S^n\vee S^n\\
\big\Vert&&\big\downarrow{\scriptstyle \nabla}\\
S^{2n-1}&\xrightarrow{\ \phi\ }&S^n
\end{array}
\qquad
\bigl(\phi=\nabla\circ[\iota_1,\iota_2]\bigr).
\]



After attaching the \(2n\)-cells, this gives

\[
q:S^n\times S^n\to X_\phi.
\]


\medskip
\noindent\textbf{Cohomology computation}\par
\smallskip


Let

\[
a=\pi_1^*(u)\in H^n(S^n\times S^n),
\]

and

\[
b=\pi_2^*(u)\in H^n(S^n\times S^n),
\]

where

\[
u\in H^n(S^n)
\]

is a generator.

Then

\[
H^*(S^n\times S^n)
\]

has generators

\[
1,\quad a,\quad b,\quad ab.
\]

Here

\[
ab\in H^{2n}(S^n\times S^n)
\]

is the top cohomology class.

Since \(a\) and \(b\) come from the two sphere factors, we have

\[
a^2=0,
\qquad
b^2=0.
\]

Also,

\[
ba=(-1)^n ab.
\]

Now let

\[
x=x_n\in H^n(X_\phi),
\qquad
y=x_{2n}\in H^{2n}(X_\phi).
\]

By construction, the fold map sends the \(n\)-dimensional generator of \(X_\phi\) to the sum of the two generators in \(S^n\times S^n\). Therefore

\[
q^*(x)=a+b.
\]

Also, choosing orientations consistently,

\[
q^*(y)=ab.
\]

Now compute:

\[
q^*(x^2)=q^*(x)^2.
\]

Since

\[
q^*(x)=a+b,
\]

we get

\[
q^*(x^2)=(a+b)^2.
\]

Expand:

\[
(a+b)^2=a^2+ab+ba+b^2.
\]

Using

\[
a^2=0,
\qquad
b^2=0,
\qquad
ba=(-1)^n ab,
\]

we obtain

\[
(a+b)^2
=
ab+(-1)^n ab.
\]

Hence

\[
(a+b)^2
=
\bigl(1+(-1)^n\bigr)ab.
\]

On the other hand, by definition of the Hopf invariant,

\[
x^2=H(\phi)y.
\]

Applying \(q^*\), we get

\[
q^*(x^2)
=
H(\phi)q^*(y)
=
H(\phi)ab.
\]

Therefore

\[
H(\phi)ab
=
\bigl(1+(-1)^n\bigr)ab.
\]

Since \(ab\neq 0\), it follows that

\[
\boxed{
	H(\phi)=1+(-1)^n.
}
\]

Thus

\[
H(\phi)=
\begin{cases}
	2, & n\text{ even},\\
	0, & n\text{ odd}.
\end{cases}
\]

In particular, when \(n\) is even,

\[
\boxed{
	H(\phi)=2\neq 0.
}
\]

Therefore \(H\) is nontrivial for every even \(n\).


\medskip
\noindent\textbf{Summary chart for \(S^n\times S^n\)}\par
\smallskip


\[
\begin{array}{c|c|c}
	\text{Object} & \text{Description} & \text{Role}\\ \hline
	S^n\vee S^n & n\text{-skeleton of }S^n\times S^n & \text{two }n\text{-cells}\\
	e^{2n} & \text{top cell} & \text{creates }H^{2n}\\
	\left[\iota_1,\iota_2\right] & S^{2n-1}\to S^n\vee S^n & \text{attaching map of top cell}\\
	\nabla & S^n\vee S^n\to S^n & \text{fold map}\\
	\phi=\nabla\circ\left[\iota_1,\iota_2\right] & S^{2n-1}\to S^n & \text{Whitehead square}\\
	H(\phi) & \text{Hopf invariant} & 1+(-1)^n
\end{array}
\]


\medskip
\noindent\textbf{Deduction: \(\pi_{4m-1}(S^{2m})\) is infinite}\par
\smallskip


Now take

\[
n=2m.
\]

Then \(n\) is even. From the previous subsection, there exists

\[
\phi:S^{2n-1}\to S^n
\]

such that

\[
H(\phi)=2\neq 0.
\]

Therefore the homomorphism

\[
H:\pi_{2n-1}(S^n)\to \mathbb Z
\]

has nonzero image.

A nonzero subgroup of \(\mathbb Z\) is infinite. Hence

\[
\pi_{2n-1}(S^n)
\]

must be infinite.

Since

\[
n=2m,
\]

we have

\[
2n-1=4m-1.
\]

Therefore

\[
\pi_{2n-1}(S^n)
=
\pi_{4m-1}(S^{2m}).
\]

Thus

\[
\boxed{
	\pi_{4m-1}(S^{2m})\text{ is infinite for every }m>0.
}
\]


\medskip
\noindent\textbf{Conceptual summary}\par
\smallskip


Problem 6 uses the interaction between degree, cup products, and Poincaré duality.

If

\[
\deg(\varphi)\neq 0,
\]

then the pullback

\[
\varphi^*:H^1(\Sigma_h;\mathbb Q)\to H^1(\Sigma_g;\mathbb Q)
\]

would preserve enough of the intersection pairing to be injective. But this is impossible when \(g<h\), because

\[
\dim H^1(\Sigma_h;\mathbb Q)=2h>2g=\dim H^1(\Sigma_g;\mathbb Q).
\]

Therefore

\[
\varphi_*:H_2(\Sigma_g)\to H_2(\Sigma_h)
\]

is zero.

Problem 7 uses cup products in a two-cell complex

\[
X_\phi=S^n\cup_\phi D^{2n}.
\]

The Hopf invariant \(H(\phi)\) measures the square of the middle-dimensional cohomology class:

\[
x_n^2=H(\phi)x_{2n}.
\]

The cell decomposition

\[
S^n\times S^n=(S^n\vee S^n)\cup_{[\iota_1,\iota_2]}e^{2n}
\]

shows that the Whitehead square

\[
[\iota,\iota]\in \pi_{2n-1}(S^n)
\]

has Hopf invariant

\[
H([\iota,\iota])=1+(-1)^n.
\]

Hence for even \(n\),

\[
H([\iota,\iota])=2\neq 0.
\]

Therefore the Hopf invariant homomorphism is nontrivial for even \(n\), and consequently

\[
\pi_{4m-1}(S^{2m})
\]

is infinite for all \(m>0\).
\end{solution}

\begin{problem}[CP-VIII-0077 --- The M\"obius bundle and its sums/products]

\medskip
\noindent\textbf{Exact statement.}\par
\smallskip

	Let \(M\) be the M\"obius bundle over \(S^{1}\). Show that \(M\oplus M\) is a trivial bundle. Is \(M\times M\), considered as a bundle over \(S^{1}\times S^{1}\), trivial?
\end{problem}

\begin{solution}
The M\"obius bundle
	\[
	M\to S^{1}
	\]
	is the nontrivial real line bundle over \(S^{1}\). Its transition function changes sign after going once around the circle. Equivalently, its first Stiefel--Whitney class is nonzero:
	\[
	w_{1}(M)\neq 0
	\]
	in
	\[
	H^{1}(S^{1};\mathbb Z_{2})\cong \mathbb Z_{2}.
	\]
	Let
	\[
	w_{1}(M)=\alpha,
	\]
	where \(\alpha\) is the nonzero generator of \(H^{1}(S^{1};\mathbb Z_{2})\).

	Now consider the direct sum bundle
	\[
	M\oplus M\to S^{1}.
	\]
	The first Stiefel--Whitney class is additive under direct sums:
	\[
	w_{1}(M\oplus M)
	=
	w_{1}(M)+w_{1}(M).
	\]
	Therefore
	\[
	w_{1}(M\oplus M)
	=
	\alpha+\alpha
	=
	0,
	\]
	because the coefficients are in \(\mathbb Z_{2}\).

	Thus \(M\oplus M\) is orientable. Since real rank \(2\) vector bundles over \(S^{1}\) are classified by whether their monodromy is orientation-preserving or orientation-reversing, orientability implies triviality in this case. Hence
	\[
	M\oplus M\cong S^{1}\times \mathbb R^{2}.
	\]

	There is also a concrete transition-function explanation. The M\"obius bundle has transition function given by multiplication by \(-1\). Therefore the transition function of \(M\oplus M\) is
	\[
	\begin{pmatrix}
		-1&0\\
		0&-1
	\end{pmatrix}
	=
	-I_{2}.
	\]
	But \(-I_{2}\) is a rotation by \(\pi\), so
	\[
	-I_{2}\in SO(2).
	\]
	Since \(SO(2)\) is connected, \(-I_{2}\) is homotopic to \(I_{2}\). Hence the clutching function is null-homotopic, and therefore the bundle is trivial.

	Thus
	\[
	\boxed{M\oplus M\cong S^{1}\times \mathbb R^{2}.}
	\]

	Now consider
	\[
	M\times M\to S^{1}\times S^{1}.
	\]
	Here \(M\times M\) means the external product bundle. Its fiber over a point
	\[
	(x,y)\in S^{1}\times S^{1}
	\]
	is
	\[
	M_{x}\times M_{y}\cong \mathbb R^{2}.
	\]
	Equivalently,
	\[
	M\times M
	=
	\pi_{1}^{*}M\oplus \pi_{2}^{*}M,
	\]
	where
	\[
	\pi_{1},\pi_{2}:S^{1}\times S^{1}\to S^{1}
	\]
	are the two projection maps.

	Using additivity and naturality of the first Stiefel--Whitney class, we get
	\[
	w_{1}(M\times M)
	=
	w_{1}(\pi_{1}^{*}M\oplus \pi_{2}^{*}M),
	\]
	so
	\[
	w_{1}(M\times M)
	=
	w_{1}(\pi_{1}^{*}M)+w_{1}(\pi_{2}^{*}M).
	\]
	By naturality,
	\[
	w_{1}(\pi_{1}^{*}M)=\pi_{1}^{*}w_{1}(M),
	\]
	and
	\[
	w_{1}(\pi_{2}^{*}M)=\pi_{2}^{*}w_{1}(M).
	\]
	Therefore
	\[
	w_{1}(M\times M)
	=
	\pi_{1}^{*}\alpha+\pi_{2}^{*}\alpha.
	\]

	Now
	\[
	H^{1}(S^{1}\times S^{1};\mathbb Z_{2})
	\cong
	H^{1}(S^{1};\mathbb Z_{2})\oplus H^{1}(S^{1};\mathbb Z_{2})
	\cong
	\mathbb Z_{2}\oplus \mathbb Z_{2}.
	\]
	The class
	\[
	\pi_{1}^{*}\alpha+\pi_{2}^{*}\alpha
	\]
	is nonzero. Hence
	\[
	w_{1}(M\times M)\neq 0.
	\]
	Therefore \(M\times M\) is not orientable. Since every trivial real vector bundle is orientable, \(M\times M\) cannot be trivial.

	Thus
	\[
	\boxed{M\times M\to S^{1}\times S^{1}\text{ is not trivial}.}
	\]
\end{solution}

\begin{problem}[CP-VIII-0078 --- The unit tangent bundle of \(S^{2}\)]

\medskip
\noindent\textbf{Exact statement.}\par
\smallskip

	Let
	\[
	E=TS^{2}
	\]
	be the tangent bundle of \(S^{2}\). Show that the unit sphere bundle \(S(E)\) is homeomorphic to \(SO(3)\), which is also homeomorphic to \(\mathbb{RP}^{3}\). Deduce that
	\[
	e(E)=\pm 2a,
	\]
	where
	\[
	H^{2}(S^{2})=\langle a\rangle.
	\]
\end{problem}

\begin{solution}
The unit sphere bundle of \(TS^{2}\) is the unit tangent bundle
	\[
	S(TS^{2})
	=
	\{(x,v)\in S^{2}\times \mathbb R^{3}:v\in T_{x}S^{2},\ \|v\|=1\}.
	\]
	Since
	\[
	T_{x}S^{2}
	=
	\{v\in \mathbb R^{3}:x\cdot v=0\},
	\]
	we can write
	\[
	S(TS^{2})
	=
	\{(x,v):\|x\|=1,\ \|v\|=1,\ x\cdot v=0\}.
	\]
	Thus a point of \(S(TS^{2})\) is an ordered pair of perpendicular unit vectors:
	\[
	(x,v).
	\]
	From such a pair, we form the ordered triple
	\[
	(x,v,x\times v).
	\]
	This is a positively oriented orthonormal basis of \(\mathbb R^{3}\).

	Define
	\[
	\Phi:S(TS^{2})\to SO(3)
	\]
	by
	\[
	\Phi(x,v)
	=
	\begin{pmatrix}
		|&|&|\\
		x&v&x\times v\\
		|&|&|
	\end{pmatrix}.
	\]
	The columns of this matrix are orthonormal, so the matrix is orthogonal. Moreover, because the basis
	\[
	(x,v,x\times v)
	\]
	is positively oriented, the determinant is \(1\). Therefore
	\[
	\Phi(x,v)\in SO(3).
	\]

	We now construct the inverse map. Let
	\[
	A\in SO(3).
	\]
	Write the columns of \(A\) as
	\[
	A=(u_{1},u_{2},u_{3}).
	\]
	Since \(A\in SO(3)\), the vectors
	\[
	u_{1},u_{2},u_{3}
	\]
	form a positively oriented orthonormal basis of \(\mathbb R^{3}\). In particular,
	\[
	u_{1}\in S^{2}.
	\]
	Also,
	\[
	u_{1}\cdot u_{2}=0,
	\]
	so
	\[
	u_{2}\in T_{u_{1}}S^{2}.
	\]
	Since
	\[
	\|u_{2}\|=1,
	\]
	we have
	\[
	(u_{1},u_{2})\in S(TS^{2}).
	\]
	Thus the inverse map is
	\[
	A\mapsto (u_{1},u_{2}).
	\]
	Therefore
	\[
	\boxed{S(TS^{2})\cong SO(3).}
	\]

	Next we recall the standard homeomorphism
	\[
	SO(3)\cong \mathbb{RP}^{3}.
	\]
	One way to see this is through the double covering
	\[
	S^{3}\to SO(3).
	\]
	The unit quaternions \(S^{3}\) act on \(\mathbb R^{3}\) by rotations, and the two unit quaternions \(q\) and \(-q\) determine the same rotation. Therefore
	\[
	SO(3)\cong S^{3}/\{\pm 1\}.
	\]
	But
	\[
	S^{3}/\{\pm 1\}=\mathbb{RP}^{3}.
	\]
	Hence
	\[
	\boxed{S(TS^{2})\cong SO(3)\cong \mathbb{RP}^{3}.}
	\]

	Now we deduce the Euler class of \(E=TS^{2}\). Since \(TS^{2}\) is an oriented real rank \(2\) vector bundle, its unit sphere bundle
	\[
	S(E)\to S^{2}
	\]
	is an oriented circle bundle. Oriented circle bundles over \(S^{2}\) are classified by their Euler class
	\[
	e(E)\in H^{2}(S^{2};\mathbb Z).
	\]
	Since
	\[
	H^{2}(S^{2};\mathbb Z)\cong \mathbb Z,
	\]
	and since
	\[
	H^{2}(S^{2};\mathbb Z)=\langle a\rangle,
	\]
	we may write
	\[
	e(E)=ka
	\]
	for some integer \(k\).

	We use the Gysin sequence for the oriented circle bundle
	\[
	S(E)\to S^{2}.
	\]
	The relevant part of the Gysin sequence is
	\[
	H^{0}(S^{2};\mathbb Z)
	\xrightarrow{\ \smile e(E)\ }
	H^{2}(S^{2};\mathbb Z)
	\longrightarrow
	H^{2}(S(E);\mathbb Z)
	\longrightarrow
	H^{1}(S^{2};\mathbb Z).
	\]
	Now
	\[
	H^{0}(S^{2};\mathbb Z)\cong \mathbb Z,
	\]
	\[
	H^{2}(S^{2};\mathbb Z)\cong \mathbb Z,
	\]
	and
	\[
	H^{1}(S^{2};\mathbb Z)=0.
	\]
	The map
	\[
	H^{0}(S^{2};\mathbb Z)\to H^{2}(S^{2};\mathbb Z)
	\]
	is multiplication by the Euler class \(e(E)=ka\). Therefore the cokernel is
	\[
	H^{2}(S(E);\mathbb Z)\cong \mathbb Z/k\mathbb Z.
	\]

	But we have shown that
	\[
	S(E)\cong \mathbb{RP}^{3}.
	\]
	The integral cohomology of \(\mathbb{RP}^{3}\) gives
	\[
	H^{2}(\mathbb{RP}^{3};\mathbb Z)\cong \mathbb Z_{2}.
	\]
	Hence
	\[
	H^{2}(S(E);\mathbb Z)\cong \mathbb Z_{2}.
	\]
	Therefore
	\[
	\mathbb Z/k\mathbb Z\cong \mathbb Z_{2}.
	\]
	It follows that
	\[
	|k|=2.
	\]
	Thus
	\[
	k=\pm 2.
	\]
	Therefore
	\[
	\boxed{e(E)=\pm 2a.}
	\]
	Since \(E=TS^{2}\), we also write
	\[
	\boxed{e(TS^{2})=\pm 2a.}
	\]
	The sign depends on the chosen orientation convention for the generator \(a\).
\end{solution}

\begin{problem}[CP-VIII-0079 --- Orthogonal complements of subbundles]

\medskip
\noindent\textbf{Exact statement.}\par
\smallskip

	Let
	\[
	E\to B
	\]
	be a real vector bundle equipped with a Riemannian metric, and let
	\[
	F\subset E
	\]
	be a subbundle. Show that \(F^{\perp}\) is a vector bundle, and that
	\[
	F\oplus F^{\perp}\cong E.
	\]
\end{problem}

\begin{solution}
Since \(E\to B\) is equipped with a Riemannian metric, every fiber \(E_{b}\) is a finite-dimensional real inner product space. Since \(F\subset E\) is a subbundle, for each \(b\in B\), the fiber
	\[
	F_{b}
	\]
	is a linear subspace of
	\[
	E_{b}.
	\]

	Define the orthogonal complement fiberwise by
	\[
	F_{b}^{\perp}
	=
	\{v\in E_{b}:\langle v,w\rangle_{b}=0
	\text{ for all }w\in F_{b}\}.
	\]
	Then define
	\[
	F^{\perp}
	=
	\coprod_{b\in B}F_{b}^{\perp}.
	\]
	We want to show that
	\[
	F^{\perp}\to B
	\]
	is itself a vector bundle.

	Let \(U\subset B\) be an open set over which \(F\) is trivial. Choose a local frame for \(F\) over \(U\):
	\[
	s_{1},\dots,s_{k}.
	\]
	Since \(F\subset E\) is a subbundle, we may extend this to a local frame for \(E\) over \(U\):
	\[
	s_{1},\dots,s_{k},s_{k+1},\dots,s_{n}.
	\]

	Now apply the Gram--Schmidt process fiber by fiber, using the Riemannian metric on \(E\). Because the metric varies continuously, respectively smoothly, over \(B\), the resulting orthonormal frame also varies continuously, respectively smoothly.

	Thus we obtain a local orthonormal frame
	\[
	e_{1},\dots,e_{n}
	\]
	for \(E|_{U}\), such that
	\[
	e_{1},\dots,e_{k}
	\]
	span \(F|_{U}\). Since the frame is orthonormal, the remaining vectors
	\[
	e_{k+1},\dots,e_{n}
	\]
	span the orthogonal complement \(F^{\perp}|_{U}\).

	Therefore
	\[
	F^{\perp}|_{U}\cong U\times \mathbb R^{n-k}.
	\]
	This proves that \(F^{\perp}\to B\) is locally trivial, and hence \(F^{\perp}\) is a vector bundle.

	Finally, for every point \(b\in B\), ordinary linear algebra gives the orthogonal decomposition
	\[
	E_{b}=F_{b}\oplus F_{b}^{\perp}.
	\]
	Thus we define a bundle map
	\[
	\Psi:F\oplus F^{\perp}\to E
	\]
	by
	\[
	\Psi(u,v)=u+v.
	\]
	On each fiber,
	\[
	\Psi_{b}:F_{b}\oplus F_{b}^{\perp}\to E_{b}
	\]
	is a linear isomorphism. Since this map is expressed locally in continuous, respectively smooth, frames, it is a vector bundle isomorphism.

	Hence
	\[
	\boxed{F^{\perp}\text{ is a vector bundle}}
	\]
	and
	\[
	\boxed{F\oplus F^{\perp}\cong E.}
	\]

	
\medskip
\noindent\textbf{Final conclusions}\par
\smallskip


	For Problem 8,
	\[
	\boxed{M\oplus M\text{ is trivial, but }M\times M\to S^{1}\times S^{1}\text{ is not trivial}.}
	\]

	For Problem 9,
	\[
	\boxed{S(TS^{2})\cong SO(3)\cong \mathbb{RP}^{3},\qquad e(TS^{2})=\pm 2a.}
	\]

	For Problem 10,
	\[
	\boxed{F^{\perp}\text{ is a vector bundle and }F\oplus F^{\perp}\cong E.}
	\]
\end{solution}

\begin{problem}[CP-VIII-0080 --- Vector bundles over spheres via clutching functions]
Given a continuous map
\[
\gamma:S^{n-1}\longrightarrow O(k),
\]
let
\[
E_{\gamma}
=
D_N^n\times \mathbb{R}^k
\;\sqcup\;
D_S^n\times \mathbb{R}^k
\big/ \sim,
\]
where the equivalence relation identifies
\[
(x,v)\in S_N^{n-1}\times \mathbb{R}^k
\]
with
\[
(x,\gamma(x)v)\in S_S^{n-1}\times \mathbb{R}^k.
\]
Show that \(E_{\gamma}\) is a vector bundle over \(S^n\), and that every vector bundle over \(S^n\) is isomorphic to some \(E_{\gamma}\). Show further that
\[
E_{\gamma}\cong E_{\gamma'}
\quad\Longleftrightarrow\quad
\gamma\simeq \gamma'.
\]
\end{problem}

\begin{solution}
We view \(S^n\) as the union of two closed hemispheres
\[
S^n=D_N^n\cup D_S^n,
\]
where
\[
D_N^n\cap D_S^n=S^{n-1}.
\]
The construction glues two trivial vector bundles
\[
D_N^n\times \mathbb{R}^k
\quad\text{and}\quad
D_S^n\times \mathbb{R}^k
\]
along their common boundary by the map
\[
\gamma:S^{n-1}\to O(k).
\]

Define a projection
\[
\pi:E_{\gamma}\longrightarrow S^n
\]
by
\[
\pi([x,v])=x.
\]
This is well-defined because the gluing relation identifies only points lying over the same base point \(x\in S^{n-1}\). Hence the fiber over \(x\in S^n\) is naturally a copy of \(\mathbb{R}^k\).

Over the northern hemisphere \(D_N^n\), the bundle is trivial:
\[
\pi^{-1}(D_N^n)\cong D_N^n\times \mathbb{R}^k.
\]
Similarly, over the southern hemisphere,
\[
\pi^{-1}(D_S^n)\cong D_S^n\times \mathbb{R}^k.
\]
On the overlap \(S^{n-1}\), the transition function is precisely
\[
\gamma(x)\in O(k)\subset GL(k,\mathbb{R}).
\]
Since each \(\gamma(x)\) is a linear isomorphism of \(\mathbb{R}^k\), the vector space structures on the fibers are compatible with the gluing. Therefore \(E_{\gamma}\to S^n\) is a rank \(k\) real vector bundle.

Now let
\[
p:E\to S^n
\]
be any rank \(k\) real vector bundle. Since both hemispheres \(D_N^n\) and \(D_S^n\) are contractible, the restrictions
\[
E|_{D_N^n}
\quad\text{and}\quad
E|_{D_S^n}
\]
are trivial. Choose trivializations
\[
\varphi_N:E|_{D_N^n}\longrightarrow D_N^n\times \mathbb{R}^k
\]
and
\[
\varphi_S:E|_{D_S^n}\longrightarrow D_S^n\times \mathbb{R}^k.
\]
On the overlap \(S^{n-1}\), these two trivializations differ by a transition function
\[
g:S^{n-1}\longrightarrow GL(k,\mathbb{R}).
\]
Using a bundle metric and applying Gram--Schmidt, we may deform this transition function into \(O(k)\). Thus we obtain a map
\[
\gamma:S^{n-1}\to O(k)
\]
such that
\[
E\cong E_{\gamma}.
\]
Therefore every rank \(k\) real vector bundle over \(S^n\) is obtained by a clutching construction.

It remains to discuss when two such bundles are isomorphic. Suppose first that
\[
\gamma\simeq \gamma'.
\]
Let
\[
H:S^{n-1}\times I\longrightarrow O(k)
\]
be a homotopy from \(\gamma\) to \(\gamma'\), so that
\[
H(x,0)=\gamma(x),
\qquad
H(x,1)=\gamma'(x).
\]
Using \(H\), one constructs a vector bundle over
\[
S^n\times I
\]
whose restriction to \(S^n\times\{0\}\) is \(E_{\gamma}\), and whose restriction to \(S^n\times\{1\}\) is \(E_{\gamma'}\). Since \(I\) is contractible, these endpoint restrictions are isomorphic. Hence
\[
E_{\gamma}\cong E_{\gamma'}.
\]

Conversely, suppose that
\[
E_{\gamma}\cong E_{\gamma'}.
\]
An isomorphism between these bundles, when restricted to the two hemispheres, is represented by maps
\[
A_N:D_N^n\to O(k),
\qquad
A_S:D_S^n\to O(k).
\]
On the equator \(S^{n-1}\), compatibility with the clutching functions gives
\[
\gamma'(x)A_N(x)=A_S(x)\gamma(x).
\]
Thus
\[
\gamma'(x)
=
A_S(x)\gamma(x)A_N(x)^{-1}.
\]
Since \(A_N\) and \(A_S\) extend over disks, their restrictions to \(S^{n-1}\) are null-homotopic. Hence \(\gamma'\) is homotopic to \(\gamma\). Therefore
\[
E_{\gamma}\cong E_{\gamma'}
\quad\Longleftrightarrow\quad
\gamma\simeq \gamma'.
\]

Thus isomorphism classes of rank \(k\) real vector bundles over \(S^n\) are classified by homotopy classes of clutching maps
\[
[S^{n-1},O(k)].
\]
\end{solution}

\begin{problem}[CP-VIII-0081 --- The unique nontrivial real $k$-plane bundle over $S^1$]
Show that up to isomorphism there is a unique nontrivial \(k\)-dimensional real vector bundle \(E_k\) over \(S^1\). Use Problem 11 to compute
\[
H^*(S(E_k)).
\]
\end{problem}

\begin{solution}
For \(S^1\), the clutching construction says that rank \(k\) real vector bundles over \(S^1\) are classified by maps into \(O(k)\). Equivalently, one may describe such a bundle as a mapping torus
\[
E_A
=
[0,1]\times \mathbb{R}^k
\big/
(1,v)\sim (0,Av),
\]
where
\[
A\in O(k).
\]

The orthogonal group \(O(k)\) has two path components:
\[
O(k)=SO(k)\sqcup \{A\in O(k):\det A=-1\}.
\]
Thus the determinant of \(A\) gives the only invariant:
\[
\det A=+1
\quad\text{or}\quad
\det A=-1.
\]

If
\[
\det A=+1,
\]
then \(A\) lies in the identity component \(SO(k)\), so it is homotopic to the identity. The corresponding bundle is therefore trivial:
\[
E_A\cong S^1\times \mathbb{R}^k.
\]

If
\[
\det A=-1,
\]
then \(A\) lies in the other component of \(O(k)\), and it is homotopic to a reflection, for example
\[
R=\operatorname{diag}(-1,1,\dots,1).
\]
All such matrices give isomorphic bundles. Hence there is exactly one nontrivial rank \(k\) real vector bundle over \(S^1\). It is
\[
E_k\cong \mu\oplus \varepsilon^{k-1},
\]
where \(\mu\) is the Möbius line bundle and \(\varepsilon^{k-1}\) is the trivial bundle of rank \(k-1\).

Now consider the unit sphere bundle
\[
S(E_k)\to S^1.
\]
Since \(E_k\) is obtained from the reflection \(R\), its sphere bundle is the mapping torus
\[
S(E_k)
\cong
S^{k-1}\times [0,1]
\big/
(u,1)\sim (Ru,0).
\]
Thus \(S(E_k)\) is the mapping torus of the reflection
\[
R:S^{k-1}\to S^{k-1}.
\]

The map \(R\) has degree
\[
\deg(R)=-1.
\]
Therefore the induced map on cohomology satisfies
\[
R^*:H^0(S^{k-1};\mathbb{Z})\to H^0(S^{k-1};\mathbb{Z})
\]
is the identity, while
\[
R^*:H^{k-1}(S^{k-1};\mathbb{Z})
\to
H^{k-1}(S^{k-1};\mathbb{Z})
\]
is multiplication by \(-1\).

Using the Wang sequence for the mapping torus, we obtain
\[
\cdots
\to
H^{q-1}(S^{k-1})
\xrightarrow{R^*-\operatorname{id}}
H^{q-1}(S^{k-1})
\to
H^q(S(E_k))
\to
H^q(S^{k-1})
\xrightarrow{R^*-\operatorname{id}}
H^q(S^{k-1})
\to
\cdots .
\]

Because
\[
H^q(S^{k-1};\mathbb{Z})
=
\begin{cases}
	\mathbb{Z}, & q=0,\\
	\mathbb{Z}, & q=k-1,\\
	0, & \text{otherwise},
\end{cases}
\]
the only nontrivial maps are the identity difference in degree \(0\), and multiplication by \(-2\) in degree \(k-1\):
\[
R^*-\operatorname{id}
=
-1-1
=
-2.
\]

Hence, for \(k\geq 2\),
\[
H^q(S(E_k);\mathbb{Z})
\cong
\begin{cases}
	\mathbb{Z}, & q=0,\\[2mm]
	\mathbb{Z}, & q=1,\\[2mm]
	\mathbb{Z}/2, & q=k,\\[2mm]
	0, & \text{otherwise}.
\end{cases}
\]

Thus, additively,
\[
H^*(S(E_k);\mathbb{Z})
\cong
\mathbb{Z}
\oplus
\mathbb{Z}\alpha
\oplus
(\mathbb{Z}/2)\tau,
\]
where
\[
|\alpha|=1,
\qquad
|\tau|=k.
\]
The class \(\alpha\) comes from the base circle \(S^1\), while the class \(\tau\) is the top torsion class.

As a ring, for \(k\geq 2\),
\[
H^*(S(E_k);\mathbb{Z})
\cong
\mathbb{Z}[\alpha,\tau]
\big/
(\alpha^2,\;2\tau,\;\alpha\tau,\;\tau^2),
\]
with
\[
|\alpha|=1,
\qquad
|\tau|=k.
\]

For \(k=1\), the nontrivial line bundle over \(S^1\) is the Möbius line bundle. Its sphere bundle has fiber
\[
S^0=\{-1,1\},
\]
and the reflection interchanges the two points. The total space is again a circle:
\[
S(E_1)\cong S^1.
\]
Therefore
\[
H^*(S(E_1);\mathbb{Z})
\cong
H^*(S^1;\mathbb{Z}).
\]


\medskip
\noindent\textbf{Geometric interpretation}\par
\smallskip


The nontrivial bundle \(E_k\) over \(S^1\) is obtained by taking a cylinder
\[
[0,1]\times \mathbb{R}^k
\]
and gluing the two ends with a reflection:
\[
(1,v)\sim (0,Rv).
\]
For \(k=1\), this gives the Möbius line bundle.

For \(k=2\), the sphere bundle \(S(E_2)\) has fiber \(S^1\). Since the monodromy is a reflection of the circle, the total space is the Klein bottle:
\[
S(E_2)\cong K.
\]
Hence
\[
H^0(K;\mathbb{Z})\cong \mathbb{Z},
\qquad
H^1(K;\mathbb{Z})\cong \mathbb{Z},
\qquad
H^2(K;\mathbb{Z})\cong \mathbb{Z}/2.
\]

For \(k\geq 3\), \(S(E_k)\) is a higher-dimensional analogue of the Klein bottle: it is the mapping torus of a reflection of \(S^{k-1}\).


\medskip
\noindent\textbf{Summary chart}\par
\smallskip


\begin{itemize}
\item If \(\det A=+1\), the bundle over \(S^1\) is trivial, with total
space \(S^1\times\mathbb R^k\).
\item If \(\det A=-1\), then \(E_k\cong\mu\oplus\varepsilon^{k-1}\), the
unique nontrivial rank-\(k\) bundle over \(S^1\).
\item Its sphere bundle \(S(E_k)\) is the mapping torus of a map
\(R:S^{k-1}\to S^{k-1}\) of degree \(-1\), and
\[
H^0\cong\mathbb Z,\qquad H^1\cong\mathbb Z,\qquad H^k\cong\mathbb Z/2.
\]
\end{itemize}
\end{solution}

\begin{problem}[CP-VIII-0082 --- Orientation reversal on $\mathbb{CP}^2$ and parity of degree]
\noindent\textbf{(a)} Show that no homeomorphism of \(\mathbb{CP}^{2}\) reverses
orientation. Equivalently, prove that a homeomorphism
\[
f:\mathbb{CP}^{2}\longrightarrow \mathbb{CP}^{2}
\]
cannot have degree \(-1\).

\begin{enumerate}
	\item[(b)] If
	\[
	f:S^{2}\times S^{2}\longrightarrow \mathbb{CP}^{2},
	\]
	show that
	\[
	f_{*}:H_{4}(S^{2}\times S^{2})\longrightarrow H_{4}(\mathbb{CP}^{2})
	\]
	has even degree.
\end{enumerate}
\end{problem}

\begin{solution}
\subsubsection*{Part (a): No orientation-reversing homeomorphism of \texorpdfstring{$\mathbb{CP}^{2}$}{CP2}}

We use the cohomology ring of complex projective space:
\[
H^{*}(\mathbb{CP}^{2};\mathbb{Z})
\cong
\mathbb{Z}[\alpha]/(\alpha^{3}),
\]
where
\[
\alpha\in H^{2}(\mathbb{CP}^{2};\mathbb{Z})
\]
is a generator, and
\[
\alpha^{2}\in H^{4}(\mathbb{CP}^{2};\mathbb{Z})
\]
is a generator of the top cohomology group.

Let
\[
f:\mathbb{CP}^{2}\longrightarrow \mathbb{CP}^{2}
\]
be a homeomorphism. Since
\[
H^{2}(\mathbb{CP}^{2};\mathbb{Z})\cong \mathbb{Z},
\]
the induced map
\[
f^{*}:H^{2}(\mathbb{CP}^{2};\mathbb{Z})
\longrightarrow
H^{2}(\mathbb{CP}^{2};\mathbb{Z})
\]
must be multiplication by either \(1\) or \(-1\). Hence
\[
f^{*}(\alpha)=\pm \alpha.
\]

Now consider the induced map on the top-dimensional cohomology class
\[
\alpha^{2}\in H^{4}(\mathbb{CP}^{2};\mathbb{Z}).
\]
Because \(f^{*}\) preserves cup products, we have
\[
f^{*}(\alpha^{2})
=
f^{*}(\alpha\smile \alpha)
=
f^{*}(\alpha)\smile f^{*}(\alpha).
\]
Therefore
\[
f^{*}(\alpha^{2})
=
(\pm \alpha)\smile(\pm \alpha)
=
\alpha^{2}.
\]

So every homeomorphism of \(\mathbb{CP}^{2}\) acts trivially on
\[
H^{4}(\mathbb{CP}^{2};\mathbb{Z}).
\]

However, if \(f\) were orientation-reversing, then \(f^{*}\) would act by multiplication by \(-1\) on the top cohomology group. Thus we would have
\[
f^{*}(\alpha^{2})=-\alpha^{2}.
\]

This contradicts the previous computation
\[
f^{*}(\alpha^{2})=\alpha^{2}.
\]

Hence no orientation-reversing homeomorphism
\[
f:\mathbb{CP}^{2}\longrightarrow \mathbb{CP}^{2}
\]
exists.

Therefore,
\[
\boxed{\mathbb{CP}^{2}\text{ admits no orientation-reversing homeomorphism.}}
\]

\subsubsection*{Part (b): A map \texorpdfstring{$S^{2}\times S^{2}\to \mathbb{CP}^{2}$}{S2 x S2 to CP2} has even degree}

Let
\[
u\in H^{2}(\mathbb{CP}^{2};\mathbb{Z})
\]
be a generator. Then
\[
u^{2}\in H^{4}(\mathbb{CP}^{2};\mathbb{Z})
\]
is a generator of the top cohomology group.

Now consider the cohomology ring of \(S^{2}\times S^{2}\). Let
\[
a,b\in H^{2}(S^{2}\times S^{2};\mathbb{Z})
\]
be the pullbacks of the fundamental cohomology classes from the two sphere factors. Then
\[
H^{*}(S^{2}\times S^{2};\mathbb{Z})
\cong
\mathbb{Z}[a,b]/(a^{2},b^{2}),
\]
where
\[
ab\in H^{4}(S^{2}\times S^{2};\mathbb{Z})
\]
is a generator.

Since
\[
H^{2}(S^{2}\times S^{2};\mathbb{Z})
\cong
\mathbb{Z}\oplus \mathbb{Z},
\]
we may write
\[
f^{*}(u)=ma+nb
\]
for some integers \(m,n\in\mathbb{Z}\).

Now compute the pullback of the top-dimensional generator \(u^{2}\). Since \(f^{*}\) preserves cup products,
\[
f^{*}(u^{2})
=
f^{*}(u)^2.
\]
Thus
\[
f^{*}(u^{2})
=
(ma+nb)^2.
\]
Expanding, we obtain
\[
(ma+nb)^2
=
m^{2}a^{2}+2mnab+n^{2}b^{2}.
\]
But in the cohomology ring of \(S^{2}\times S^{2}\), we have
\[
a^{2}=0,
\qquad
b^{2}=0.
\]
Therefore
\[
f^{*}(u^{2})=2mnab.
\]

The degree of \(f\) is the integer \(d\) such that
\[
f^{*}(u^{2})=d\,ab.
\]
Comparing this with
\[
f^{*}(u^{2})=2mnab,
\]
we get
\[
d=2mn.
\]

Hence \(d\) is even.

Therefore,
\[
\boxed{\deg(f)\text{ is even.}}
\]

Equivalently,
\[
f_{*}[S^{2}\times S^{2}]
=
2mn[\mathbb{CP}^{2}],
\]
so the induced homomorphism
\[
f_{*}:H_{4}(S^{2}\times S^{2})
\longrightarrow
H_{4}(\mathbb{CP}^{2})
\]
has even degree.
\end{solution}

\begin{problem}[CP-VIII-0083 --- Connected sums and the fundamental class]
Suppose \(M_{1}\) and \(M_{2}\) are closed, connected, oriented
\(n\)-manifolds and \(x_i\in M_i\). Choose charts
\[
\phi_i:U_i\longrightarrow \mathbb{R}^{n},
\]
where
\[
x_i\in U_i\subset M_i,
\qquad
\phi_i(x_i)=0,
\]
with the property that
\[
\phi_{1*}([M_1]|_{x_1})
=
-\phi_{2*}([M_2]|_{x_2}).
\]
Let
\[
M_i'
=
M_i-\phi_i^{-1}(B_1(0)),
\]
where \(B_1(0)\) is the open ball of radius \(1\). Define the connected sum
\[
M_1\# M_2
:=
M_1'\cup M_2'/\sim,
\]
where
\[
\phi_1^{-1}(x)\sim \phi_2^{-1}(x)
\qquad
\text{for }x\in S^{n-1}.
\]

\begin{enumerate}
	\item[(a)] Explain why \(M_1\# M_2\) is a manifold.

	\item[(b)] Show that
	\[
	H^i(M_1\#M_2)
	\cong
	H^i(M_1)\oplus H^i(M_2)
	\]
	for \(0<i<n\).

	\item[(c)] Show that there is an orientation \([M_1\#M_2]\) on
	\(M_1\#M_2\) such that
	\[
	[M_1\#M_2]|_{M_i'}
	=
	[M_i]|_{M_i'}
	\]
	for \(i=1,2\).

	\item[(d)] Show that the cup product pairing on \(H^*(M_1\#M_2)\)
	is given by
	\[
	((a_1,a_2),(b_1,b_2))_{\#}
	=
	(a_1,b_1)_1+(a_2,b_2)_2,
	\]
	for
	\[
	0<|a_1|=|a_2|<n,
	\qquad
	0<|b_1|=|b_2|<n,
	\]
	where \((\cdot,\cdot)_i\) is the cup product pairing on \(M_i\).
\end{enumerate}
\end{problem}

\begin{solution}

\medskip
\noindent\textbf{}\par
\smallskip


\medskip
\noindent\textbf{Part (a): Why \texorpdfstring{$M_1\#M_2$}{M1\#M2} is a manifold}\par
\smallskip


The connected sum is constructed by deleting an open \(n\)-ball from each
manifold and then gluing the resulting boundary spheres.

For each \(i=1,2\), the space
\[
M_i'
=
M_i-\phi_i^{-1}(B_1(0))
\]
is an \(n\)-manifold with boundary. Its boundary is
\[
\partial M_i'
=
\phi_i^{-1}(S^{n-1}).
\]

Away from the boundary, nothing changes. Every point in the interior of
\(M_1'\) or \(M_2'\) already has a neighborhood homeomorphic to an open
subset of \(\mathbb{R}^n\).

Thus the only point to check is what happens near the glued sphere. Near a
boundary point, each \(M_i'\) looks locally like a half-space:
\[
\mathbb{R}^{n-1}\times [0,\varepsilon).
\]

When two such half-spaces are glued along their common boundary
\[
\mathbb{R}^{n-1}\times \{0\},
\]
the result is locally
\[
\mathbb{R}^{n-1}\times (-\varepsilon,\varepsilon)
\cong
\mathbb{R}^n.
\]

Therefore every point of the glued space has a neighborhood homeomorphic to
\(\mathbb{R}^n\). Hence
\[
\boxed{M_1\#M_2\text{ is an }n\text{-manifold}.}
\]

Since \(M_1\) and \(M_2\) are closed and connected, the connected sum is also
closed and connected.


\medskip
\noindent\textbf{Part (b): Cohomology of the connected sum}\par
\smallskip


Let
\[
X=M_1\#M_2.
\]

We cover \(X\) by two open sets \(A\) and \(B\), where \(A\) is a small open
thickening of \(M_1'\), and \(B\) is a small open thickening of \(M_2'\).
Then
\[
A\simeq M_1',
\qquad
B\simeq M_2',
\]
and
\[
A\cap B\simeq S^{n-1}.
\]

Now use the Mayer--Vietoris sequence:
\[
\cdots
\longrightarrow
H^{i-1}(A\cap B)
\longrightarrow
H^i(X)
\longrightarrow
H^i(A)\oplus H^i(B)
\longrightarrow
H^i(A\cap B)
\longrightarrow
\cdots .
\]

Since
\[
A\cap B\simeq S^{n-1},
\]
we have
\[
H^i(A\cap B)=0
\]
for
\[
0<i<n-1.
\]

Therefore, for \(0<i<n-1\), the Mayer--Vietoris sequence gives
\[
H^i(X)
\cong
H^i(A)\oplus H^i(B).
\]

Because
\[
A\simeq M_1',
\qquad
B\simeq M_2',
\]
we obtain
\[
H^i(X)
\cong
H^i(M_1')\oplus H^i(M_2').
\]

Removing an open \(n\)-ball from a connected closed \(n\)-manifold does not
change cohomology in degrees \(0<i<n\). Hence
\[
H^i(M_j')
\cong
H^i(M_j)
\]
for \(j=1,2\). Thus
\[
H^i(M_1\#M_2)
\cong
H^i(M_1)\oplus H^i(M_2)
\]
for \(0<i<n-1\).

It remains to discuss the degree \(i=n-1\). The relevant part of the
Mayer--Vietoris sequence is
\[
0
\longrightarrow
H^{n-1}(X)
\longrightarrow
H^{n-1}(M_1')\oplus H^{n-1}(M_2')
\longrightarrow
H^{n-1}(S^{n-1})
\longrightarrow
H^n(X)
\longrightarrow
0.
\]

Since \(X\) is a closed connected oriented \(n\)-manifold,
\[
H^n(X)\cong \mathbb{Z}.
\]
The connecting map
\[
H^{n-1}(S^{n-1})
\longrightarrow
H^n(X)
\]
identifies the generator of the gluing sphere with the top orientation
class of \(X\). Hence this map is an isomorphism. Therefore the preceding
map
\[
H^{n-1}(M_1')\oplus H^{n-1}(M_2')
\longrightarrow
H^{n-1}(S^{n-1})
\]
has zero image, and we get
\[
H^{n-1}(X)
\cong
H^{n-1}(M_1')\oplus H^{n-1}(M_2').
\]

Again,
\[
H^{n-1}(M_j')
\cong
H^{n-1}(M_j),
\]
so
\[
H^{n-1}(M_1\#M_2)
\cong
H^{n-1}(M_1)\oplus H^{n-1}(M_2).
\]

Therefore, for every \(0<i<n\),
\[
\boxed{
	H^i(M_1\#M_2)
	\cong
	H^i(M_1)\oplus H^i(M_2).
}
\]


\medskip
\noindent\textbf{Part (c): Orientation of the connected sum}\par
\smallskip


Each \(M_i'\) inherits an orientation from \(M_i\). The only issue is whether
the two orientations agree along the glued boundary sphere.

The assumption
\[
\phi_{1*}([M_1]|_{x_1})
=
-\phi_{2*}([M_2]|_{x_2})
\]
means that the two coordinate balls are identified with opposite
orientations.

This is exactly the condition needed when forming the connected sum of
oriented manifolds. Indeed, when an open ball is removed from an oriented
manifold, the induced boundary orientation on the resulting boundary sphere
is determined by the outward normal convention. The boundary orientations on
\[
\partial M_1'
\qquad\text{and}\qquad
\partial M_2'
\]
are opposite under the gluing map. Hence the two orientations fit together
after gluing.

Therefore there exists an orientation class
\[
[M_1\#M_2]\in H_n(M_1\#M_2)
\]
such that
\[
\boxed{
	[M_1\#M_2]|_{M_i'}
	=
	[M_i]|_{M_i'}
}
\]
for \(i=1,2\).

Thus the connected sum orientation restricts to the original orientations on
both punctured pieces.


\medskip
\noindent\textbf{Part (d): Cup product pairing on the connected sum}\par
\smallskip


By part (b), for \(0<i<n\), we have
\[
H^i(M_1\#M_2)
\cong
H^i(M_1)\oplus H^i(M_2).
\]

Thus a cohomology class on \(M_1\#M_2\) may be written as a pair
\[
(a_1,a_2),
\]
where
\[
a_i\in H^*(M_i).
\]

Similarly, write another class as
\[
(b_1,b_2).
\]

The cup product pairing on \(M_1\#M_2\) is defined by
\[
((a_1,a_2),(b_1,b_2))_{\#}
=
\left\langle
(a_1,a_2)\smile (b_1,b_2),
[M_1\#M_2]
\right\rangle.
\]

Under the splitting of cohomology from part (b), the cup product is computed
componentwise:
\[
(a_1,a_2)\smile (b_1,b_2)
=
(a_1\smile b_1,\; a_2\smile b_2).
\]

There are no cross terms such as \(a_1\smile b_2\) or \(a_2\smile b_1\),
because these classes come from different sides of the connected sum neck.

Using the orientation from part (c), we have
\[
[M_1\#M_2]|_{M_i'}
=
[M_i]|_{M_i'}.
\]
Therefore evaluation on the fundamental class splits as the sum of the
evaluations on the two summands:
\[
\begin{aligned}
	((a_1,a_2),(b_1,b_2))_{\#}
	&=
	\left\langle
	a_1\smile b_1,[M_1]
	\right\rangle
	+
	\left\langle
	a_2\smile b_2,[M_2]
	\right\rangle.
\end{aligned}
\]

By definition,
\[
(a_i,b_i)_i
=
\left\langle
a_i\smile b_i,[M_i]
\right\rangle.
\]

Hence
\[
\boxed{
	((a_1,a_2),(b_1,b_2))_{\#}
	=
	(a_1,b_1)_1+(a_2,b_2)_2.
}
\]

Thus the cup product pairing on the connected sum is the orthogonal direct
sum of the cup product pairings on \(M_1\) and \(M_2\).
\end{solution}

\begin{problem}[CP-VIII-0084 --- Intersection forms on orientable $4$-manifolds]
Suppose \(M\) is a \(\mathbb Z\)-orientable \(4\)-manifold. If \(H_1(M)\) is free over
\(\mathbb Z\), show that \(H_*(M)\) and \(H^*(M)\) are free over \(\mathbb Z\).
Assume that this is the case, and choose a basis
\[
\langle a_1,\ldots,a_n\rangle
\]
for \(H^2(M)\). Let \(A\) be the matrix whose \(ij\)-th entry is
\[
(a_i,a_j).
\]
Show that \(A\) is symmetric, and that
\[
\det A=\pm 1.
\]

\par

\noindent
\end{problem}

\begin{solution}
Since \(M\) is a closed orientable \(4\)-manifold, Poincaré duality gives
\[
H^k(M)\cong H_{4-k}(M).
\]

We also use the universal coefficient theorem for cohomology:
\[
0\longrightarrow \operatorname{Ext}(H_{k-1}(M),\mathbb Z)
\longrightarrow H^k(M)
\longrightarrow \operatorname{Hom}(H_k(M),\mathbb Z)
\longrightarrow 0.
\]

Because \(H_1(M)\) is free, we have
\[
\operatorname{Ext}(H_1(M),\mathbb Z)=0.
\]
Therefore, for \(k=2\), the universal coefficient theorem gives
\[
H^2(M)\cong \operatorname{Hom}(H_2(M),\mathbb Z).
\]

On the other hand, Poincaré duality gives
\[
H^2(M)\cong H_2(M).
\]

Combining the two isomorphisms, we get
\[
H_2(M)\cong \operatorname{Hom}(H_2(M),\mathbb Z).
\]

Since \(H_2(M)\) is finitely generated, this implies that \(H_2(M)\) is free. Indeed,
if \(H_2(M)\) had a nonzero torsion subgroup, then
\(\operatorname{Hom}(H_2(M),\mathbb Z)\) would kill that torsion, so
\(H_2(M)\) could not be isomorphic to its dual.

Now
\[
H_0(M)\cong \mathbb Z,
\qquad
H_4(M)\cong \mathbb Z.
\]
By assumption, \(H_1(M)\) is free. Also,
\[
H_3(M)\cong H^1(M)
\]
by Poincaré duality. The universal coefficient theorem gives
\[
H^1(M)\cong \operatorname{Hom}(H_1(M),\mathbb Z),
\]
because \(\operatorname{Ext}(H_0(M),\mathbb Z)=0\). Since \(H_1(M)\) is free,
\(\operatorname{Hom}(H_1(M),\mathbb Z)\) is also free. Hence \(H_3(M)\) is free.

Thus all homology groups \(H_*(M)\) are free over \(\mathbb Z\). Since the homology
groups are free, the universal coefficient theorem gives
\[
H^k(M)\cong \operatorname{Hom}(H_k(M),\mathbb Z),
\]
so all cohomology groups \(H^*(M)\) are also free over \(\mathbb Z\).

Now define the intersection pairing on \(H^2(M)\) by
\[
(a,b)=\langle a\smile b,[M]\rangle.
\]

Let \(a,b\in H^2(M)\). Since cup product is graded-commutative,
\[
a\smile b=(-1)^{2\cdot 2}b\smile a=b\smile a.
\]
Therefore
\[
(a,b)=\langle a\smile b,[M]\rangle
=\langle b\smile a,[M]\rangle
=(b,a).
\]
Hence the intersection pairing is symmetric, and so the matrix \(A\) is symmetric.

Finally, by Poincaré duality, the map
\[
H^2(M)\longrightarrow \operatorname{Hom}(H^2(M),\mathbb Z),
\qquad
a\longmapsto \bigl(b\mapsto (a,b)\bigr),
\]
is an isomorphism. With respect to the basis
\[
\langle a_1,\ldots,a_n\rangle
\]
and its dual basis, this map is represented by the matrix \(A\). Therefore \(A\) is
an invertible integer matrix whose inverse is also an integer matrix. Hence \(A\) is
unimodular, and so
\[
\det A=\pm 1.
\]

Thus \(A\) is symmetric and unimodular.

\[
\boxed{A^T=A,\qquad \det A=\pm 1.}
\]
\end{solution}

\begin{problem}[CP-VIII-0085 --- Equal homology but distinct homotopy types of $4$-manifolds]
Orient \(\mathbb CP^2\) over \(\mathbb Z\) so that
\[
(a,a)=1,
\]
where
\[
\langle a\rangle=H^2(\mathbb CP^2),
\]
and let \(\overline{\mathbb CP}^{\,2}\) denote the same manifold \(\mathbb CP^2\) with
the opposite orientation. Define
\[
X_1=\mathbb CP^2\#\mathbb CP^2,
\qquad
X_2=\mathbb CP^2\#\overline{\mathbb CP}^{\,2},
\qquad
X_3=S^2\times S^2.
\]
Show that
\[
H_*(X_1)\cong H_*(X_2)\cong H_*(X_3),
\]
but that no two of the \(X_i\) are homotopy equivalent.

\par

\noindent
\end{problem}

\begin{solution}
First compute the homology groups. We know that
\[
H_k(\mathbb CP^2)\cong
\begin{cases}
	\mathbb Z, & k=0,2,4,\\
	0, & \text{otherwise}.
\end{cases}
\]

The connected sum of two closed connected orientable \(4\)-manifolds has top and
bottom homology
\[
H_0\cong \mathbb Z,
\qquad
H_4\cong \mathbb Z,
\]
and in dimension \(2\), the groups add. Hence
\[
H_2(X_1)\cong H_2(X_2)\cong \mathbb Z^2.
\]

Also,
\[
H_1(X_1)=H_1(X_2)=0,
\qquad
H_3(X_1)=H_3(X_2)=0.
\]

For \(S^2\times S^2\), by the Künneth formula,
\[
H_k(S^2\times S^2)\cong
\begin{cases}
	\mathbb Z, & k=0,4,\\
	\mathbb Z^2, & k=2,\\
	0, & \text{otherwise}.
\end{cases}
\]

Therefore
\[
H_k(X_i)\cong
\begin{cases}
	\mathbb Z, & k=0,4,\\
	\mathbb Z^2, & k=2,\\
	0, & k=1,3,
\end{cases}
\]
for each \(i=1,2,3\). Hence
\[
H_*(X_1)\cong H_*(X_2)\cong H_*(X_3).
\]

Now we distinguish the spaces using intersection forms.

For \(\mathbb CP^2\), choose the generator \(a\in H^2(\mathbb CP^2)\) such that
\[
(a,a)=1.
\]
Reversing orientation changes the sign of the fundamental class. Hence the
intersection form on \(\overline{\mathbb CP}^{\,2}\) is
\[
(a,a)=-1.
\]

Therefore, with respect to natural bases,
\[
Q_{X_1}\cong
\begin{pmatrix}
	1 & 0\\
	0 & 1
\end{pmatrix},
\]
and
\[
Q_{X_2}\cong
\begin{pmatrix}
	1 & 0\\
	0 & -1
\end{pmatrix}.
\]

For \(X_3=S^2\times S^2\), let
\[
u=[S^2\times \{\mathrm{pt}\}],
\qquad
v=[\{\mathrm{pt}\}\times S^2].
\]
Then
\[
u\cdot u=0,
\qquad
v\cdot v=0,
\qquad
u\cdot v=1.
\]
Thus
\[
Q_{X_3}\cong
\begin{pmatrix}
	0 & 1\\
	1 & 0
\end{pmatrix}.
\]

So the three intersection forms are
\[
Q_{X_1}=
\begin{pmatrix}
	1 & 0\\
	0 & 1
\end{pmatrix},
\qquad
Q_{X_2}=
\begin{pmatrix}
	1 & 0\\
	0 & -1
\end{pmatrix},
\qquad
Q_{X_3}=
\begin{pmatrix}
	0 & 1\\
	1 & 0
\end{pmatrix}.
\]

A homotopy equivalence between closed oriented \(4\)-manifolds preserves the
intersection form up to sign. Therefore, if two of the \(X_i\) were homotopy
equivalent, their intersection forms would be isomorphic up to multiplication by
\(-1\).

Now \(Q_{X_1}\) is definite. Indeed,
\[
Q_{X_1}(x,y)=x^2+y^2,
\]
so it is positive definite. Its signature is
\[
\sigma(X_1)=2.
\]

On the other hand, \(Q_{X_2}\) is indefinite:
\[
Q_{X_2}(x,y)=x^2-y^2,
\]
and \(Q_{X_3}\) is also indefinite:
\[
Q_{X_3}(x,y)=2xy.
\]
Both have signature \(0\). Therefore \(X_1\) cannot be homotopy equivalent to
\(X_2\) or \(X_3\).

It remains to distinguish \(X_2\) and \(X_3\). The form \(Q_{X_2}\) is odd because
there exists a class with odd self-intersection. For example,
\[
(1,0)\cdot (1,0)=1.
\]

However, \(Q_{X_3}\) is even. Indeed, for any class
\[
x=au+bv,
\]
we have
\[
x\cdot x=(au+bv)\cdot(au+bv)
=2ab,
\]
which is always even.

Parity of an integral intersection form is preserved under isomorphism, and it is
also preserved after multiplying the form by \(-1\). Hence
\[
Q_{X_2}\not\cong Q_{X_3}
\]
even up to sign. Therefore \(X_2\) and \(X_3\) are not homotopy equivalent.

Thus the three manifolds have isomorphic homology groups, but no two of them are
homotopy equivalent.

\[
\boxed{
	H_*(X_1)\cong H_*(X_2)\cong H_*(X_3),
	\quad
	\text{but } X_i\not\simeq X_j \text{ for } i\neq j.
}
\]
\end{solution}

\begin{problem}[CP-VIII-0086 --- Parity of the middle Betti number in dimension $4n+2$]
If \(M\) is a \(\mathbb Z\)-orientable manifold of dimension \(4n+2\), show that
the dimension of
\[
H_{2n+1}(M;\mathbb Q)
\]
is even.

\par

\noindent
\end{problem}

\begin{solution}
Assume
\[
\dim M=4n+2.
\]
The middle dimension is therefore
\[
2n+1.
\]

Consider the middle-dimensional cup-product pairing over \(\mathbb Q\):
\[
H^{2n+1}(M;\mathbb Q)\times H^{2n+1}(M;\mathbb Q)
\longrightarrow \mathbb Q,
\]
defined by
\[
(a,b)=\langle a\smile b,[M]\rangle.
\]

By Poincaré duality, this pairing is nondegenerate.

Now let
\[
a,b\in H^{2n+1}(M;\mathbb Q).
\]
Since \(2n+1\) is odd, graded commutativity gives
\[
a\smile b
=
(-1)^{(2n+1)(2n+1)}b\smile a.
\]
But
\[
(2n+1)(2n+1)
\]
is odd, so
\[
(-1)^{(2n+1)(2n+1)}=-1.
\]
Therefore
\[
a\smile b=-b\smile a.
\]

Hence
\[
(a,b)=-(b,a).
\]
Thus the middle-dimensional intersection pairing is skew-symmetric.

So \(H^{2n+1}(M;\mathbb Q)\) is a finite-dimensional vector space over
\(\mathbb Q\) equipped with a nondegenerate skew-symmetric bilinear form.

We now use the linear algebra fact that a nondegenerate skew-symmetric bilinear
form can exist only on an even-dimensional vector space.

Indeed, suppose the form has matrix \(A\). Since the form is skew-symmetric,
\[
A^T=-A.
\]
Taking determinants gives
\[
\det(A^T)=\det(-A).
\]
Since \(\det(A^T)=\det(A)\), we get
\[
\det(A)=(-1)^m\det(A),
\]
where \(m\) is the size of the matrix, i.e.
\[
m=\dim_{\mathbb Q}H^{2n+1}(M;\mathbb Q).
\]
If \(m\) were odd, then
\[
(-1)^m=-1,
\]
so
\[
\det(A)=-\det(A).
\]
Therefore
\[
2\det(A)=0.
\]
Over \(\mathbb Q\), this implies
\[
\det(A)=0,
\]
contradicting the nondegeneracy of the pairing. Hence \(m\) must be even.

Therefore
\[
\dim_{\mathbb Q}H^{2n+1}(M;\mathbb Q)
\]
is even.

Finally, over a field, the universal coefficient theorem gives
\[
H^{2n+1}(M;\mathbb Q)
\cong
\operatorname{Hom}_{\mathbb Q}(H_{2n+1}(M;\mathbb Q),\mathbb Q).
\]
Thus
\[
\dim_{\mathbb Q}H^{2n+1}(M;\mathbb Q)
=
\dim_{\mathbb Q}H_{2n+1}(M;\mathbb Q).
\]

Hence
\[
\boxed{
	\dim_{\mathbb Q}H_{2n+1}(M;\mathbb Q)
	\text{ is even.}
}
\]
\end{solution}

\begin{problem}[CP-VIII-0087 --- Thom classes under direct sums of vector bundles]
\noindent\textbf{(a)} For \(i=1,2\), choose a generator
\[
u_i\in H^{n_i}(\mathbb R^{n_i}\mid 0).
\]
Let
\[
p_i:\mathbb R^{n_1}\times \mathbb R^{n_2}\longrightarrow \mathbb R^{n_i}
\]
be the projection. Show that
\[
p_1^*(u_1)\smile p_2^*(u_2)
\]
generates
\[
H^{n_1+n_2}(\mathbb R^{n_1}\times \mathbb R^{n_2}\mid 0).
\]

\begin{enumerate}
	\item[(b)] Now suppose that
	\[
	\pi_i:E_i\to B
	\qquad (i=1,2)
	\]
	are vector bundles over \(B\), and that \(U_i\) is a Thom class for \(E_i\).
	If
	\[
	p_i:E_1\oplus E_2\to E_i
	\]
	is the projection, show that
	\[
	p_1^*(U_1)\smile p_2^*(U_2)
	\]
	is a Thom class for \(E_1\oplus E_2\).

	\item[(c)] Deduce that \(E_1\oplus E_2\) can be oriented so that
	\[
	e(E_1\oplus E_2)=e(E_1)\smile e(E_2).
	\]
\end{enumerate}

\par

\noindent
\end{problem}

\begin{solution}

\medskip
\noindent\textbf{Part (a)}\par
\smallskip


We use the notation
\[
H^k(\mathbb R^n\mid 0)
=
H^k(\mathbb R^n,\mathbb R^n\setminus\{0\}).
\]
It is well known that
\[
H^k(\mathbb R^n,\mathbb R^n\setminus\{0\})
\cong
\begin{cases}
	\mathbb Z, & k=n,\\
	0, & k\neq n.
\end{cases}
\]
Hence
\[
H^{n_i}(\mathbb R^{n_i}\mid 0)\cong \mathbb Z
\]
for \(i=1,2\), and
\[
H^{n_1+n_2}(\mathbb R^{n_1}\times \mathbb R^{n_2}\mid 0)\cong \mathbb Z.
\]

Let
\[
u_i\in H^{n_i}(\mathbb R^{n_i},\mathbb R^{n_i}\setminus\{0\})
\]
be a generator. The product
\[
p_1^*(u_1)\smile p_2^*(u_2)
\]
is the local cohomology version of the external product
\[
u_1\times u_2.
\]
Indeed, pulling back by the two projections and taking the cup product gives the
class supported at the common zero
\[
(0,0)\in \mathbb R^{n_1}\times \mathbb R^{n_2}.
\]

By the K\"unneth theorem for these local cohomology groups, the external product
induces an isomorphism
\[
H^{n_1}(\mathbb R^{n_1}\mid 0)
\otimes
H^{n_2}(\mathbb R^{n_2}\mid 0)
\cong
H^{n_1+n_2}(\mathbb R^{n_1}\times \mathbb R^{n_2}\mid 0).
\]
Since both \(u_1\) and \(u_2\) are generators of infinite cyclic groups, their
product is again a generator, up to sign.

Therefore
\[
p_1^*(u_1)\smile p_2^*(u_2)
\]
generates
\[
H^{n_1+n_2}(\mathbb R^{n_1}\times \mathbb R^{n_2}\mid 0).
\]

Thus
\[
\boxed{
	p_1^*(u_1)\smile p_2^*(u_2)
	\text{ is a generator of }
	H^{n_1+n_2}(\mathbb R^{n_1}\times \mathbb R^{n_2}\mid 0).
}
\]


\medskip
\noindent\textbf{Part (b)}\par
\smallskip


Let
\[
\pi_i:E_i\to B
\]
be vector bundles of ranks \(n_i\). A Thom class for \(E_i\) is a class
\[
U_i\in H^{n_i}(E_i,E_i\setminus B)
\]
whose restriction to every fiber
\[
(E_i)_b\cong \mathbb R^{n_i}
\]
is a generator of
\[
H^{n_i}((E_i)_b,(E_i)_b\setminus\{0\})
\cong \mathbb Z.
\]

Now consider the direct sum bundle
\[
E_1\oplus E_2\to B.
\]
Its fiber over \(b\in B\) is
\[
(E_1\oplus E_2)_b
\cong
(E_1)_b\oplus (E_2)_b
\cong
\mathbb R^{n_1}\oplus \mathbb R^{n_2}.
\]

Let
\[
p_i:E_1\oplus E_2\to E_i
\]
be the bundle projection. Consider the class
\[
U
=
p_1^*(U_1)\smile p_2^*(U_2).
\]
This class has degree
\[
n_1+n_2.
\]
Moreover, its support is along the zero section of \(E_1\oplus E_2\), so
\[
U\in H^{n_1+n_2}(E_1\oplus E_2,(E_1\oplus E_2)\setminus B).
\]

We must show that \(U\) restricts to a generator on each fiber.

Fix \(b\in B\). Restricting \(U\) to the fiber gives
\[
U|_{(E_1\oplus E_2)_b}
=
p_1^*(U_1|_{(E_1)_b})
\smile
p_2^*(U_2|_{(E_2)_b}).
\]
Since \(U_i\) is a Thom class, each restriction
\[
U_i|_{(E_i)_b}
\]
is a generator of the local cohomology group of the fiber:
\[
H^{n_i}((E_i)_b,(E_i)_b\setminus\{0\}).
\]

By part (a), the cup product of these two pulled-back generators is a generator of
\[
H^{n_1+n_2}
\bigl(
(E_1)_b\oplus (E_2)_b,
((E_1)_b\oplus (E_2)_b)\setminus\{0\}
\bigr).
\]

Therefore the restriction of
\[
p_1^*(U_1)\smile p_2^*(U_2)
\]
to every fiber is a generator. Hence this class is a Thom class for
\(E_1\oplus E_2\).

Thus
\[
\boxed{
	p_1^*(U_1)\smile p_2^*(U_2)
	\text{ is a Thom class for }E_1\oplus E_2.
}
\]


\medskip
\noindent\textbf{Part (c)}\par
\smallskip


The choice of a Thom class determines an orientation of a vector bundle. By part
(b), the class
\[
U
=
p_1^*(U_1)\smile p_2^*(U_2)
\]
is a Thom class for \(E_1\oplus E_2\). We orient \(E_1\oplus E_2\) using this Thom
class.

Let
\[
s:B\to E_1\oplus E_2
\]
be the zero section of \(E_1\oplus E_2\). By definition, the Euler class of an
oriented vector bundle is the pullback of its Thom class along the zero section.
Therefore
\[
e(E_1\oplus E_2)
=
s^*(U).
\]
Using the formula for \(U\), we get
\[
e(E_1\oplus E_2)
=
s^*\bigl(p_1^*(U_1)\smile p_2^*(U_2)\bigr).
\]
By naturality of cup products,
\[
e(E_1\oplus E_2)
=
s^*p_1^*(U_1)\smile s^*p_2^*(U_2).
\]

But
\[
p_i\circ s:B\to E_i
\]
is precisely the zero section of \(E_i\). Hence
\[
s^*p_i^*(U_i)
=
(p_i\circ s)^*(U_i)
=
e(E_i).
\]
Therefore
\[
e(E_1\oplus E_2)
=
e(E_1)\smile e(E_2).
\]

Thus
\[
\boxed{
	e(E_1\oplus E_2)=e(E_1)\smile e(E_2).
}
\]
\end{solution}

\begin{problem}[CP-VIII-0088 --- B\'ezout's theorem in $\mathbb{CP}^2$ via intersection theory]
If
\[
p\in \mathbb C[z_0,z_1,z_2]
\]
is a homogeneous polynomial, define
\[
V_p
=
\{[z_0:z_1:z_2]\in \mathbb{CP}^2
\mid
p(z_0,z_1,z_2)=0\}.
\]
If \(p\) and \(q\) are chosen so that \(V_p\) and \(V_q\) are embedded
submanifolds of \(\mathbb{CP}^2\) which intersect transversely, show that
\(V_p\) and \(V_q\) intersect in precisely
\[
(\deg p)(\deg q)
\]
points.

\par

\noindent
\end{problem}

\begin{solution}
Let
\[
d=\deg p,
\qquad
e=\deg q.
\]
The zero set \(V_p\) is a complex hypersurface in \(\mathbb{CP}^2\), hence a
complex curve. Since \(V_p\) is assumed to be an embedded submanifold, it defines
a homology class
\[
[V_p]\in H_2(\mathbb{CP}^2;\mathbb Z).
\]

Let
\[
a\in H^2(\mathbb{CP}^2;\mathbb Z)
\]
be the standard generator. Then
\[
H^*(\mathbb{CP}^2;\mathbb Z)
\cong
\mathbb Z[a]/(a^3),
\qquad
\deg a=2,
\]
and
\[
\langle a^2,[\mathbb{CP}^2]\rangle=1.
\]

A homogeneous polynomial of degree \(d\) defines a section of the complex line
bundle
\[
\mathcal O(d)\to \mathbb{CP}^2.
\]
Its zero set \(V_p\) has Poincar\'e dual equal to the first Chern class of
\(\mathcal O(d)\):
\[
[V_p]^\vee
=
c_1(\mathcal O(d)).
\]
Since
\[
c_1(\mathcal O(d))=d\,a,
\]
we obtain
\[
[V_p]^\vee=d\,a.
\]

Similarly,
\[
[V_q]^\vee=e\,a.
\]

The algebraic intersection number of \(V_p\) and \(V_q\) is therefore
\[
V_p\cdot V_q
=
\langle [V_p]^\vee\smile [V_q]^\vee,[\mathbb{CP}^2]\rangle.
\]
Substituting the Poincar\'e dual classes gives
\[
V_p\cdot V_q
=
\langle (d\,a)\smile(e\,a),[\mathbb{CP}^2]\rangle.
\]
Hence
\[
V_p\cdot V_q
=
de\langle a^2,[\mathbb{CP}^2]\rangle.
\]
Since
\[
\langle a^2,[\mathbb{CP}^2]\rangle=1,
\]
we get
\[
V_p\cdot V_q=de.
\]

Because \(V_p\) and \(V_q\) are complex submanifolds of the complex surface
\(\mathbb{CP}^2\), every transverse intersection point has local intersection
number \(+1\). Therefore the algebraic intersection number is equal to the actual
number of intersection points.

Thus
\[
|V_p\cap V_q|
=
de
=
(\deg p)(\deg q).
\]

Therefore
\[
\boxed{
	|V_p\cap V_q|=(\deg p)(\deg q).
}
\]

This is the smooth transverse case of B\'ezout's theorem for projective plane
curves.
\end{solution}

\begin{problem}[CP-VIII-0089 --- Homology of the sphere bundle of $T\mathbb{CP}^n$]
Let \(E\) be the tangent bundle of \(\mathbb{CP}^n\). Compute
\[
H_*(S(E)).
\]

\par

\noindent
\end{problem}

\begin{solution}
Let
\[
B=\mathbb{CP}^n,
\qquad
E=T\mathbb{CP}^n.
\]
Since \(\mathbb{CP}^n\) is a complex manifold of complex dimension \(n\), its
tangent bundle \(T\mathbb{CP}^n\) is a complex vector bundle of complex rank
\(n\). As a real vector bundle, it has rank
\[
2n.
\]

Therefore the sphere bundle
\[
S(E)=S(T\mathbb{CP}^n)
\]
has fiber
\[
S^{2n-1}.
\]
So we have an oriented sphere bundle
\[
S^{2n-1}
\longrightarrow
S(T\mathbb{CP}^n)
\xrightarrow{\ \pi\ }
\mathbb{CP}^n.
\]


\medskip
\noindent\textbf{Step 1: Cohomology of \(\mathbb{CP}^n\)}\par
\smallskip


The cohomology ring of complex projective space is
\[
H^*(\mathbb{CP}^n;\mathbb Z)
\cong
\mathbb Z[x]/(x^{n+1}),
\qquad
\deg x=2.
\]
Hence
\[
H^{2k}(\mathbb{CP}^n;\mathbb Z)\cong \mathbb Z
\qquad
(0\le k\le n),
\]
and
\[
H^{2k+1}(\mathbb{CP}^n;\mathbb Z)=0.
\]


\medskip
\noindent\textbf{Step 2: Euler class of \(T\mathbb{CP}^n\)}\par
\smallskip


For \(\mathbb{CP}^n\), the total Chern class of the tangent bundle is
\[
c(T\mathbb{CP}^n)=(1+x)^{n+1}.
\]
Therefore the top Chern class is
\[
c_n(T\mathbb{CP}^n)=(n+1)x^n.
\]

Since \(T\mathbb{CP}^n\) is a complex vector bundle, the Euler class of its
underlying oriented real bundle is the top Chern class:
\[
e(T\mathbb{CP}^n)=c_n(T\mathbb{CP}^n).
\]
Hence
\[
e(E)
=
e(T\mathbb{CP}^n)
=
(n+1)x^n.
\]


\medskip
\noindent\textbf{Step 3: The Gysin sequence}\par
\smallskip


For an oriented sphere bundle coming from an oriented real vector bundle of rank
\(2n\), the Gysin sequence has the form
\[
\begin{aligned}
\cdots&\longrightarrow H^{q-2n}(\mathbb{CP}^n)
\xrightarrow{\smile e(E)}H^q(\mathbb{CP}^n)\\
&\longrightarrow H^q(S(E))
\longrightarrow H^{q-2n+1}(\mathbb{CP}^n)\\
&\xrightarrow{\smile e(E)}H^{q+1}(\mathbb{CP}^n)
\longrightarrow\cdots .
\end{aligned}
\]

Since
\[
H^{\text{odd}}(\mathbb{CP}^n)=0,
\]
this sequence simplifies significantly.


\medskip
\noindent\textbf{Step 4: Degrees below \(2n-1\)}\par
\smallskip


For
\[
q<2n-1,
\]
we have
\[
H^{q-2n}(\mathbb{CP}^n)=0
\]
and
\[
H^{q-2n+1}(\mathbb{CP}^n)=0.
\]
Therefore the Gysin sequence gives
\[
H^q(S(E))\cong H^q(\mathbb{CP}^n).
\]

Hence
\[
H^q(S(E))\cong \mathbb Z
\]
for even \(q\) satisfying
\[
0\le q\le 2n-2,
\]
and
\[
H^q(S(E))=0
\]
for odd \(q<2n-1\).


\medskip
\noindent\textbf{Step 5: The middle part}\par
\smallskip


Now take
\[
q=2n-1
\]
and
\[
q=2n.
\]
The relevant portion of the Gysin sequence is
\[
0
\longrightarrow
H^{2n-1}(S(E))
\longrightarrow
H^0(\mathbb{CP}^n)
\xrightarrow{\ \smile e(E)\ }
H^{2n}(\mathbb{CP}^n)
\longrightarrow
H^{2n}(S(E))
\longrightarrow
0.
\]

Since
\[
H^0(\mathbb{CP}^n)\cong \mathbb Z,
\qquad
H^{2n}(\mathbb{CP}^n)\cong \mathbb Z,
\]
and
\[
e(E)=(n+1)x^n,
\]
the map
\[
H^0(\mathbb{CP}^n)\longrightarrow H^{2n}(\mathbb{CP}^n)
\]
is multiplication by \(n+1\). Therefore the exact sequence becomes
\[
0
\longrightarrow
H^{2n-1}(S(E))
\longrightarrow
\mathbb Z
\xrightarrow{\ \times(n+1)\ }
\mathbb Z
\longrightarrow
H^{2n}(S(E))
\longrightarrow
0.
\]

The map
\[
\mathbb Z\xrightarrow{\ \times(n+1)\ }\mathbb Z
\]
is injective, so
\[
H^{2n-1}(S(E))=0.
\]
Its cokernel is
\[
\mathbb Z/(n+1).
\]
Hence
\[
H^{2n}(S(E))\cong \mathbb Z/(n+1).
\]


\medskip
\noindent\textbf{Step 6: Degrees above \(2n\)}\par
\smallskip


For
\[
q>2n,
\]
we have
\[
H^q(\mathbb{CP}^n)=0.
\]
The Gysin sequence then gives
\[
H^q(S(E))\cong H^{q-2n+1}(\mathbb{CP}^n).
\]

Therefore
\[
H^q(S(E))\cong \mathbb Z
\]
when
\[
q-2n+1
\]
is even and lies between \(0\) and \(2n\). Equivalently,
\[
q=2n+1,2n+3,\ldots,4n-1.
\]
In all other degrees above \(2n\), the cohomology vanishes.

Thus the cohomology groups of \(S(E)\) are
\[
H^q(S(E);\mathbb Z)\cong
\begin{cases}
	\mathbb Z, & q=0,2,4,\ldots,2n-2,\\[1mm]
	0, & q=2n-1,\\[1mm]
	\mathbb Z/(n+1), & q=2n,\\[1mm]
	\mathbb Z, & q=2n+1,2n+3,\ldots,4n-1,\\[1mm]
	0, & \text{otherwise}.
\end{cases}
\]


\medskip
\noindent\textbf{Step 7: Convert cohomology to homology}\par
\smallskip


The total space \(S(E)\) is a closed orientable manifold. Its dimension is
\[
\dim S(E)
=
\dim \mathbb{CP}^n+\dim S^{2n-1}
=
2n+(2n-1)
=
4n-1.
\]

By Poincar\'e duality,
\[
H_k(S(E);\mathbb Z)
\cong
H^{4n-1-k}(S(E);\mathbb Z).
\]

Using the cohomology computation above, we obtain
\[
H_k(S(E);\mathbb Z)\cong
\begin{cases}
	\mathbb Z, & k=0,2,4,\ldots,2n-2,\\[1mm]
	\mathbb Z/(n+1), & k=2n-1,\\[1mm]
	0, & k=2n,\\[1mm]
	\mathbb Z, & k=2n+1,2n+3,\ldots,4n-1,\\[1mm]
	0, & \text{otherwise}.
\end{cases}
\]

Therefore
\[
\boxed{
	H_k(S(T\mathbb{CP}^n);\mathbb Z)\cong
	\begin{cases}
		\mathbb Z, & k=0,2,4,\ldots,2n-2,\\[1mm]
		\mathbb Z/(n+1), & k=2n-1,\\[1mm]
		0, & k=2n,\\[1mm]
		\mathbb Z, & k=2n+1,2n+3,\ldots,4n-1,\\[1mm]
		0, & \text{otherwise}.
	\end{cases}
}
\]


\medskip
\noindent\textbf{Check: The case \(n=1\)}\par
\smallskip


When
\[
n=1,
\]
we have
\[
\mathbb{CP}^1\cong S^2.
\]
The tangent bundle is \(TS^2\), and its unit sphere bundle is
\[
S(TS^2).
\]
This is the unit tangent bundle of \(S^2\), which is diffeomorphic to
\[
SO(3)\cong \mathbb{RP}^3.
\]

Our formula gives
\[
H_0(S(TS^2))\cong \mathbb Z,
\]
\[
H_1(S(TS^2))\cong \mathbb Z/2,
\]
\[
H_2(S(TS^2))=0,
\]
and
\[
H_3(S(TS^2))\cong \mathbb Z.
\]
These are exactly the homology groups of \(\mathbb{RP}^3\). Hence the formula is
consistent in this basic case.
\end{solution}

\begin{problem}[CP-VIII-0090 --- A rank-$3$ bundle over $S^4$ with no nonvanishing section]
Construct a \(3\)-dimensional real vector bundle over \(S^4\) which has no
nonvanishing section.

\par

\noindent
\end{problem}

\begin{solution}
We construct the bundle by the clutching construction.

Write
\[
S^4=D^4_+\cup D^4_-,
\]
where
\[
D^4_+\cap D^4_-=S^3.
\]
A real oriented rank-\(3\) vector bundle over \(S^4\) is obtained by taking two
trivial bundles
\[
D^4_+\times \mathbb R^3
\qquad\text{and}\qquad
D^4_-\times \mathbb R^3
\]
and gluing them along the equator \(S^3\) by a map
\[
g:S^3\to SO(3).
\]

Such maps are classified up to homotopy by
\[
\pi_3(SO(3)).
\]
It is known that
\[
\pi_3(SO(3))\cong \mathbb Z.
\]

Choose \(g:S^3\to SO(3)\) to represent a nonzero element of
\[
\pi_3(SO(3)).
\]
For example, one may take the universal covering map
\[
S^3\cong SU(2)\longrightarrow SO(3).
\]

Let
\[
E_g\to S^4
\]
be the rank-\(3\) real vector bundle obtained from this nontrivial clutching map.

We claim that \(E_g\) has no nonvanishing section.

Suppose, for contradiction, that \(E_g\) had a nowhere-zero section
\[
s:S^4\to E_g.
\]
After choosing a bundle metric, we can normalize \(s\) and assume that it is a
unit section. Then each fiber splits as
\[
(E_g)_x
=
\langle s(x)\rangle\oplus s(x)^\perp.
\]
Therefore the bundle splits as
\[
E_g\cong \varepsilon^1\oplus \xi,
\]
where \(\varepsilon^1\) is the trivial real line bundle and \(\xi\) is an
oriented real \(2\)-plane bundle.

Equivalently, the structure group of \(E_g\) reduces from \(SO(3)\) to the
stabilizer of a unit vector in \(\mathbb R^3\). This stabilizer is
\[
SO(2)\subset SO(3).
\]
Thus the clutching map \(g:S^3\to SO(3)\) would be homotopic to a map whose image
lies in \(SO(2)\). In other words, its homotopy class would lie in the image of
\[
\pi_3(SO(2))\longrightarrow \pi_3(SO(3)).
\]

But
\[
SO(2)\cong S^1,
\]
and hence
\[
\pi_3(SO(2))=\pi_3(S^1)=0.
\]
Therefore the image of
\[
\pi_3(SO(2))\to \pi_3(SO(3))
\]
is zero. This would force the clutching class of \(E_g\) in
\[
\pi_3(SO(3))\cong \mathbb Z
\]
to be zero, contradicting the choice of \(g\).

Hence \(E_g\) has no nonvanishing section.

Therefore a rank-\(3\) real vector bundle over \(S^4\) defined by a
nonzero clutching map \(S^3\to SO(3)\) has no nonvanishing section.
\end{solution}

\begin{problem}[CP-VIII-0091 --- Euler characteristic and boundaries of odd-dimensional manifolds]
If \(M\) is a closed odd-dimensional manifold, show that
\[
\chi(M)=0.
\]
Next, suppose \(M\) is a compact odd-dimensional manifold with boundary. Show
that
\[
\chi(M)=\frac{1}{2}\chi(\partial M).
\]
Conclude that \(\mathbb{RP}^2\) does not bound a compact \(3\)-manifold and that
\(\mathbb{CP}^2\) does not bound a compact \(5\)-manifold. Does
\(\mathbb{RP}^3\) bound a compact \(4\)-manifold? Does \(\mathbb{CP}^3\) bound a
compact \(7\)-manifold?

\par

\noindent
\end{problem}

\begin{solution}

\medskip
\noindent\textbf{Step 1: Closed odd-dimensional manifolds}\par
\smallskip


Let \(M\) be a closed \(n\)-dimensional manifold, where \(n\) is odd. We compute
the Euler characteristic using Betti numbers over a field \(F\):
\[
\chi(M)=\sum_{i=0}^n (-1)^i b_i(M),
\]
where
\[
b_i(M)=\dim_F H_i(M;F).
\]

By Poincar\'e duality,
\[
H_i(M;F)\cong H_{n-i}(M;F),
\]
so
\[
b_i(M)=b_{n-i}(M).
\]

Since \(n\) is odd, the terms in the Euler characteristic cancel in pairs:
\[
(-1)^i b_i + (-1)^{n-i}b_{n-i}
=
(-1)^i b_i + (-1)^{n-i}b_i.
\]
But \(n\) is odd, so
\[
(-1)^{n-i}=-(-1)^i.
\]
Therefore
\[
(-1)^i b_i + (-1)^{n-i}b_i=0.
\]

All terms cancel pairwise, and hence
\[
\chi(M)=0.
\]

Thus
\[
\boxed{
	\chi(M)=0
	\quad
	\text{for every closed odd-dimensional manifold }M.
}
\]


\medskip
\noindent\textbf{Step 2: Compact odd-dimensional manifolds with boundary}\par
\smallskip


Now suppose \(M\) is a compact odd-dimensional manifold with boundary
\[
\partial M.
\]

Consider the double of \(M\):
\[
DM=M\cup_{\partial M} M.
\]
The double \(DM\) is a closed odd-dimensional manifold. Therefore, by Step 1,
\[
\chi(DM)=0.
\]

Using additivity of Euler characteristic, we get
\[
\chi(DM)
=
\chi(M)+\chi(M)-\chi(\partial M).
\]
Hence
\[
0=2\chi(M)-\chi(\partial M).
\]
Therefore
\[
\chi(M)=\frac{1}{2}\chi(\partial M).
\]

Thus
\[
\boxed{
	\chi(M)=\frac{1}{2}\chi(\partial M)
}
\]
for every compact odd-dimensional manifold \(M\) with boundary.


\medskip
\noindent\textbf{Step 3: Consequence for boundaries}\par
\smallskip


If a closed manifold \(N\) bounds a compact odd-dimensional manifold \(M\), then
\[
N=\partial M.
\]
The formula above gives
\[
\chi(M)=\frac{1}{2}\chi(N).
\]
Since \(\chi(M)\) is an integer, \(\chi(N)\) must be even.

Thus a necessary condition for \(N\) to bound a compact odd-dimensional manifold
is:
\[
\chi(N)\equiv 0\pmod 2.
\]


\medskip
\noindent\textbf{Step 4: \(\mathbb{RP}^2\)}\par
\smallskip


The Euler characteristic of real projective space is
\[
\chi(\mathbb{RP}^{2})=1.
\]
This is odd. Therefore \(\mathbb{RP}^2\) cannot be the boundary of a compact
\(3\)-manifold.

Hence
\[
\boxed{
	\mathbb{RP}^2
	\text{ does not bound a compact }3\text{-manifold.}
}
\]


\medskip
\noindent\textbf{Step 5: \(\mathbb{CP}^2\)}\par
\smallskip


The Euler characteristic of complex projective space is
\[
\chi(\mathbb{CP}^n)=n+1.
\]
Thus
\[
\chi(\mathbb{CP}^2)=3.
\]
This is odd. Therefore \(\mathbb{CP}^2\) cannot be the boundary of a compact
\(5\)-manifold.

Hence
\[
\boxed{
	\mathbb{CP}^2
	\text{ does not bound a compact }5\text{-manifold.}
}
\]


\medskip
\noindent\textbf{Step 6: Does \(\mathbb{RP}^3\) bound a compact \(4\)-manifold?}\par
\smallskip


Yes.

A standard construction uses the disk bundle of an oriented real \(2\)-plane
bundle over \(S^2\) with Euler number \(2\). Let
\[
\xi\to S^2
\]
be an oriented real \(2\)-plane bundle with Euler class
\[
e(\xi)=2\in H^2(S^2;\mathbb Z)\cong \mathbb Z.
\]
Let
\[
D(\xi)
\]
be its disk bundle and
\[
S(\xi)
\]
its circle bundle. Then
\[
\partial D(\xi)=S(\xi).
\]

The oriented circle bundle over \(S^2\) with Euler number \(2\) is the lens space
\[
L(2,1),
\]
and
\[
L(2,1)\cong \mathbb{RP}^3.
\]
Hence
\[
\partial D(\xi)\cong \mathbb{RP}^3.
\]

Therefore
\[
\boxed{
	\mathbb{RP}^3
	\text{ bounds a compact }4\text{-manifold.}
}
\]


\medskip
\noindent\textbf{Step 7: Does \(\mathbb{CP}^3\) bound a compact \(7\)-manifold?}\par
\smallskip


Yes.

Recall that \(\mathbb{CP}^3\) may be realized as the twistor space of \(S^4\).
Equivalently,
\[
\mathbb{CP}^3
\cong
S(\Lambda^2_+(TS^4)),
\]
where
\[
\Lambda^2_+(TS^4)\to S^4
\]
is the real rank-\(3\) vector bundle of self-dual \(2\)-forms on \(S^4\), and
\(S(\Lambda^2_+(TS^4))\) denotes its sphere bundle.

Let
\[
D(\Lambda^2_+(TS^4))
\]
be the corresponding disk bundle. Since a sphere bundle is the boundary of the
disk bundle, we have
\[
\partial D(\Lambda^2_+(TS^4))
=
S(\Lambda^2_+(TS^4)).
\]
Therefore
\[
\partial D(\Lambda^2_+(TS^4))
\cong
\mathbb{CP}^3.
\]

The base \(S^4\) has dimension \(4\), and the disk fibers have dimension \(3\).
Hence the disk bundle has dimension
\[
4+3=7.
\]

Therefore
\[
\boxed{
	\mathbb{CP}^3
	\text{ bounds a compact }7\text{-manifold.}
}
\]
\end{solution}

\begin{problem}[CP-VIII-0092 --- The Lefschetz number as an intersection number]
Suppose \(F\) is a field, and \(M\) is a closed \(F\)-oriented manifold. Let
\[
\{a_i\}
\]
be a homogeneous basis of \(H^*(M;F)\), and let
\[
\{b_i\}
\]
be the dual basis with respect to the cup product pairing, meaning
\[
\langle a_i\smile b_j,[M]\rangle=\delta_{ij}.
\]
Given a map
\[
f:M\to M,
\]
let
\[
\Lambda_f=\{(p,f(p))\in M\times M\}
\]
be the graph of \(f\).

\begin{enumerate}
	\item[(a)] Express \(PD_{M\times M}(\Lambda_f)\) in terms of the \(a_i\) and \(b_i\).
	\item[(b)] Define the Lefschetz number of \(f\) by
	\[
	L(f)=\sum_{j=0}^{n}(-1)^j\operatorname{Tr} f_j^*,
	\]
	where
	\[
	f_j^*:H^j(M;F)\to H^j(M;F).
	\]
	Show that
	\[
	L(f)=\Lambda_f\cdot \Delta,
	\]
	where \(\Delta\subset M\times M\) is the diagonal.
	\item[(c)] Deduce that if \(L(f)\neq 0\), then \(f\) must have a fixed point.
	\item[(d)] Show that any map
	\[
	f:\mathbb{CP}^2\to \mathbb{CP}^2
	\]
	has a fixed point.
\end{enumerate}

\par

\noindent
\end{problem}

\begin{solution}

\medskip
\noindent\textbf{Part (a)}\par
\smallskip


Let
\[
\Delta=\{(p,p)\in M\times M\}
\]
be the diagonal. The standard formula for the Poincar\'e dual of the diagonal is
\[
PD_{M\times M}(\Delta)
=
\sum_i (-1)^{|a_i|}
p_1^*(a_i)\smile p_2^*(b_i),
\]
where
\[
p_1,p_2:M\times M\to M
\]
are the projections, and \(|a_i|\) denotes the degree of \(a_i\).

The graph of \(f\) is
\[
\Lambda_f=\{(p,f(p))\mid p\in M\}.
\]
It is the image of the embedding
\[
\Gamma_f:M\to M\times M,
\qquad
\Gamma_f(p)=(p,f(p)).
\]

Equivalently, \(\Lambda_f\) is obtained from the diagonal by applying the map
\[
\operatorname{id}\times f:M\times M\to M\times M
\]
to the first copy of \(M\) in the diagonal picture. The corresponding
Poincar\'e dual formula is obtained by applying \(f^*\) to the second factor in
the diagonal formula. Thus
\[
PD_{M\times M}(\Lambda_f)
=
\sum_i (-1)^{|a_i|}
p_1^*(a_i)\smile p_2^*(f^*b_i).
\]

Equivalently, depending on the convention for the graph
\[
p\longmapsto (f(p),p)
\]
rather than
\[
p\longmapsto (p,f(p)),
\]
one obtains the equivalent formula
\[
PD_{M\times M}(\Lambda_f)
=
\sum_i (-1)^{|a_i|}
p_1^*(f^*a_i)\smile p_2^*(b_i).
\]

With the convention in this problem,
\[
\Lambda_f=\{(p,f(p))\},
\]
we use
\[
\boxed{
	PD_{M\times M}(\Lambda_f)
	=
	\sum_i (-1)^{|a_i|}
	p_1^*(a_i)\smile p_2^*(f^*b_i).
}
\]


\medskip
\noindent\textbf{Part (b)}\par
\smallskip


We want to show that
\[
L(f)=\Lambda_f\cdot \Delta.
\]

The intersection number of \(\Lambda_f\) with the diagonal \(\Delta\) is
\[
\Lambda_f\cdot \Delta
=
\langle PD(\Lambda_f)\smile PD(\Delta),[M\times M]\rangle.
\]

There is a useful equivalent way to compute this intersection number. Pull back
\(PD(\Lambda_f)\) along the diagonal map
\[
\delta:M\to M\times M,
\qquad
\delta(p)=(p,p).
\]
Then
\[
\Lambda_f\cdot \Delta
=
\langle \delta^*PD(\Lambda_f),[M]\rangle.
\]

Using the formula from part (a),
\[
PD(\Lambda_f)
=
\sum_i (-1)^{|a_i|}
p_1^*(a_i)\smile p_2^*(f^*b_i).
\]
Since
\[
p_1\circ \delta=\operatorname{id}_M,
\qquad
p_2\circ \delta=\operatorname{id}_M,
\]
we get
\[
\delta^*PD(\Lambda_f)
=
\sum_i (-1)^{|a_i|}
a_i\smile f^*b_i.
\]

Therefore
\[
\Lambda_f\cdot \Delta
=
\sum_i (-1)^{|a_i|}
\langle a_i\smile f^*b_i,[M]\rangle.
\]

Now group the basis elements by degree. For each \(j\), let
\[
\{a_{j,\alpha}\}_\alpha
\]
be a basis of \(H^j(M;F)\), and let
\[
\{b_{j,\alpha}\}_\alpha
\]
be the dual basis in \(H^{n-j}(M;F)\), so that
\[
\langle a_{j,\alpha}\smile b_{j,\beta},[M]\rangle=\delta_{\alpha\beta}.
\]

The map
\[
f^*:H^{n-j}(M;F)\to H^{n-j}(M;F)
\]
has a matrix in the basis \(\{b_{j,\alpha}\}\). By duality, this trace agrees
with the trace of
\[
f^*:H^j(M;F)\to H^j(M;F).
\]
Therefore
\[
\sum_\alpha
\langle a_{j,\alpha}\smile f^*b_{j,\alpha},[M]\rangle
=
\operatorname{Tr}\left(f^*:H^j(M;F)\to H^j(M;F)\right).
\]

Since \(|a_{j,\alpha}|=j\), we obtain
\[
\Lambda_f\cdot \Delta
=
\sum_{j=0}^{n}
(-1)^j
\operatorname{Tr}\left(f^*:H^j(M;F)\to H^j(M;F)\right).
\]

By definition, the right-hand side is the Lefschetz number:
\[
L(f)
=
\sum_{j=0}^{n}
(-1)^j
\operatorname{Tr}(f_j^*).
\]

Hence
\[
\boxed{
	L(f)=\Lambda_f\cdot \Delta.
}
\]


\medskip
\noindent\textbf{Part (c)}\par
\smallskip


A point of the graph \(\Lambda_f\) has the form
\[
(p,f(p)).
\]
A point of the diagonal \(\Delta\) has the form
\[
(p,p).
\]

Therefore
\[
(p,f(p))\in \Delta
\]
if and only if
\[
f(p)=p.
\]

Thus
\[
\Lambda_f\cap \Delta
\]
is precisely the set of fixed points of \(f\).

If \(f\) had no fixed point, then
\[
\Lambda_f\cap \Delta=\varnothing.
\]
In that case their intersection number would be zero:
\[
\Lambda_f\cdot \Delta=0.
\]

But by part (b),
\[
\Lambda_f\cdot \Delta=L(f).
\]
Hence if \(f\) has no fixed point, then
\[
L(f)=0.
\]

Taking the contrapositive, if
\[
L(f)\neq 0,
\]
then \(f\) must have a fixed point.

Therefore
\[
\boxed{
	L(f)\neq 0
	\quad\Longrightarrow\quad
	f\text{ has a fixed point.}
}
\]

This is the Lefschetz fixed point theorem in this setting.


\medskip
\noindent\textbf{Part (d)}\par
\smallskip


Let
\[
f:\mathbb{CP}^2\to \mathbb{CP}^2
\]
be any continuous map. We will show that \(f\) has a fixed point.

Work with rational coefficients:
\[
F=\mathbb Q.
\]
The cohomology ring of \(\mathbb{CP}^2\) is
\[
H^*(\mathbb{CP}^2;\mathbb Q)
\cong
\mathbb Q[h]/(h^3),
\qquad
\deg h=2.
\]

Thus
\[
H^0(\mathbb{CP}^2;\mathbb Q)\cong \mathbb Q,
\]
\[
H^2(\mathbb{CP}^2;\mathbb Q)\cong \mathbb Q\langle h\rangle,
\]
and
\[
H^4(\mathbb{CP}^2;\mathbb Q)\cong \mathbb Q\langle h^2\rangle.
\]
All other cohomology groups vanish.

Since \(H^2(\mathbb{CP}^2;\mathbb Q)\cong \mathbb Q\), the map \(f^*\) acts by
multiplication by some integer \(d\):
\[
f^*(h)=d\,h.
\]
Because \(f^*\) preserves cup products,
\[
f^*(h^2)=f^*(h\smile h)=f^*(h)\smile f^*(h)=d^2h^2.
\]

Therefore the traces of \(f^*\) on cohomology are:
\[
\operatorname{Tr}(f^*:H^0\to H^0)=1,
\]
\[
\operatorname{Tr}(f^*:H^2\to H^2)=d,
\]
and
\[
\operatorname{Tr}(f^*:H^4\to H^4)=d^2.
\]

Since all nonzero cohomology groups occur in even degrees, the Lefschetz number
is
\[
L(f)=1+d+d^2.
\]

Now
\[
d^2+d+1
\]
is never zero for an integer \(d\). Indeed,
\[
d^2+d+1
=
\left(d+\frac12\right)^2+\frac34
>0.
\]
Therefore
\[
L(f)\neq 0.
\]

By part (c), \(f\) must have a fixed point.

Hence
\[
\boxed{
	\text{Every continuous map }
	f:\mathbb{CP}^2\to \mathbb{CP}^2
	\text{ has a fixed point.}
}
\]
\end{solution}
